% % % % % % % % % % % % % % % % % % % % % % % % % % %
%DIF LATEXDIFF DIFFERENCE FILE
%DIF DEL /tmp/1F6K40aZud/latexdiff-vc-main/thesis.tex   Sun Apr  5 20:34:45 2026
%DIF ADD thesis.tex                                     Mon Aug 24 17:48:37 2026
%         ____          _                           %
%        |    \ _ _ ___| |_ ___ _____               %
%        |  |  | | |  _|   | .'|     |              %
%        |____/|___|_| |_|_|__,|_|_|_|              %
%                                                   %
% ------------------------------------------------- %
%                                                   %
%               This is thesis.tex.                 %
%   This should be placed on the main thesis dir.   %
%           Written by Samet Akcay, 2019            %
% % % % % % % % % % % % % % % % % % % % % % % % % % %
                                                    %
\documentclass[asymmetricmargins]{library/duthesis}%
% \documentclass[symmetricmargins]{library/duthesis}%
                                                    %
% Explicitly declare bib resources in the root file so TeX Workshop can index citations.
\addbibresource{references/bib1.bib}
\addbibresource{references/bib2.bib}

\documenttype{Thesis}  % <n-Month Report | Thesis>  %                 
\title{A Step towards the Origin of Fracture and Yielding in Amorphous Solids}       %
\author{Samuel Walker}                              %
\dept{Department of Physics}                        %
\degree{Doctor of Philosophy}                       %
\degreedate{March 2026}                             %
\copyrightyear{2026}                                %
                                                    %    
%% File to be included while running latex.         
%\includeonly{chapters/1-Introduction/chapter,       
%             chapters/2-Background/chapter,          
%             references/ref,                        
%             appendices/a/appendix}                 
                                                    %
%DIF PREAMBLE EXTENSION ADDED BY LATEXDIFF
%DIF UNDERLINE PREAMBLE %DIF PREAMBLE
\RequirePackage[normalem]{ulem} %DIF PREAMBLE
\RequirePackage{color}\definecolor{RED}{rgb}{1,0,0}\definecolor{BLUE}{rgb}{0,0,1} %DIF PREAMBLE
\providecommand{\DIFadd}[1]{{\protect\color{blue}\uwave{#1}}} %DIF PREAMBLE
\providecommand{\DIFdel}[1]{{\protect\color{red}\sout{#1}}}                      %DIF PREAMBLE
%DIF SAFE PREAMBLE %DIF PREAMBLE
\providecommand{\DIFaddbegin}{} %DIF PREAMBLE
\providecommand{\DIFaddend}{} %DIF PREAMBLE
\providecommand{\DIFdelbegin}{} %DIF PREAMBLE
\providecommand{\DIFdelend}{} %DIF PREAMBLE
\providecommand{\DIFmodbegin}{} %DIF PREAMBLE
\providecommand{\DIFmodend}{} %DIF PREAMBLE
%DIF FLOATSAFE PREAMBLE %DIF PREAMBLE
\providecommand{\DIFaddFL}[1]{\DIFadd{#1}} %DIF PREAMBLE
\providecommand{\DIFdelFL}[1]{\DIFdel{#1}} %DIF PREAMBLE
\providecommand{\DIFaddbeginFL}{} %DIF PREAMBLE
\providecommand{\DIFaddendFL}{} %DIF PREAMBLE
\providecommand{\DIFdelbeginFL}{} %DIF PREAMBLE
\providecommand{\DIFdelendFL}{} %DIF PREAMBLE
%DIF COLORLISTINGS PREAMBLE %DIF PREAMBLE
\RequirePackage{listings} %DIF PREAMBLE
\RequirePackage{color} %DIF PREAMBLE
\lstdefinelanguage{DIFcode}{ %DIF PREAMBLE
%DIF DIFCODE_UNDERLINE %DIF PREAMBLE
  moredelim=[il][\color{red}\sout]{\%DIF\ <\ }, %DIF PREAMBLE
  moredelim=[il][\color{blue}\uwave]{\%DIF\ >\ } %DIF PREAMBLE
} %DIF PREAMBLE
\lstdefinestyle{DIFverbatimstyle}{ %DIF PREAMBLE
	language=DIFcode, %DIF PREAMBLE
	basicstyle=\ttfamily, %DIF PREAMBLE
	columns=fullflexible, %DIF PREAMBLE
	keepspaces=true %DIF PREAMBLE
} %DIF PREAMBLE
\lstnewenvironment{DIFverbatim}{\lstset{style=DIFverbatimstyle}}{} %DIF PREAMBLE
\lstnewenvironment{DIFverbatim*}{\lstset{style=DIFverbatimstyle,showspaces=true}}{} %DIF PREAMBLE
%DIF END PREAMBLE EXTENSION ADDED BY LATEXDIFF

\begin{document}                                    %
% Frontmatter.                                      %
\frontmatter                                        %
\titlepage                                          %
\abstract{frontmatter/abstract}                     %
\declaration                                        %
\acknowledgement{frontmatter/acknowledgement.tex}   %
\tableofcontents                                    %
\listoffigures                                      %
%\listoftables                                       %
\notations{frontmatter/notations.tex}          %
% \dedication{frontmatter/dedication}{Dedication}     %
% \newpage
% \listoftodos[Notes]
% Mainmatter                                        %
\mainmatter                                         %   
 \newpage % !TeX root = ../../thesis.tex

\setcounter{equation}{0}
\chapter{Introduction}

Material failure is a part of everyday life, occurring whenever materials are subjected to stresses that exceed their ability to withstand them. These processes span a wide range of scales, from the rupture of tectonic plates during earthquakes \cite{Cocco2023} and the shearing of ice sheets \cite{Jennings2021}, to the tearing of collagen fibre networks that provide structural support in connective tissues \cite{Burla2020}. Understanding the mechanics of failure is crucial for predicting and preventing catastrophic collapse in engineering structures, mitigating the effects of natural disasters, and developing new materials with improved resistance to failure.

In this thesis, we focus on two broad modes of failure: yielding, in which a material undergoes irreversible plastic deformation while retaining overall integrity and eventually attaining a state of plastic flow, and fracture, which involves the formation and propagation of macroscopic cracks. The central aim of this thesis is to understand how structural disorder, stress redistribution, and microscopic failure processes together determine the macroscopic mechanical response of soft amorphous materials, and in particular when they yield, fracture, or reach a quiescent steady-state under an applied deformation.

At the scales relevant to their mechanical response, the systems listed above are effectively amorphous: rather than exhibiting long-range structural order, their constituent parts (grains, fibres, droplets, etc.) are arranged in a disordered manner.
%In this sense, what matters is not the detailed microscopic structure of the underlying material, but the absence of a regular organisation at the scales over which stress is transmitted and failure develops. 
For example, the heterogeneous arrangement of grains in geological materials often gives rise to the complex fracture patterns observed in practice \cite{Cocco2023}. Similarly, at large scales, glaciers and ice sheets behave as heterogeneous disordered solids, in which structural irregularity leads to the formation of crevasses and cracks \cite{Jennings2021}. At smaller scales, the random network of collagen fibres in connective tissues likewise gives rise to complex failure behaviour that is central to tissue mechanics \cite{Burla2020,ruiz-franco_force_2022}. The study of how such amorphous materials yield and fracture is therefore broad in scope and remains the subject of extensive debate \cite{falk_deformation_2011,nicolas_deformation_2018,larson_structure_1999,bonn_yield_2017}.

The mechanical response of these materials is often complex and difficult to predict. Shear experiments on a range of soft disordered materials, including foams (e.g. aqueous foams and bubble raft systems) \cite{durian_foam_1995,Khan1988,WEAIRE2008}, emulsions (e.g. mayonnaise and oil-in-water emulsions) \cite{Mason1996,Datta2011,Derkach2009,AbuJdayil2003}, and colloidal suspensions (e.g. paint and toothpaste) \cite{Mewis2011,Koumakis2011,Petekidis2004,Cloitre2003}, have shown that the onset of irreversible flow in these materials is often associated with a critical point \cite{nicolas_deformation_2018,Divoux2024,lin_scaling_2014,Ozawa2018,Dahmen2011} beyond which small differences in loading, material properties, preparation, and the spatial arrangement of the underlying microstructure can produce vastly different outcomes. Depending on both the material and the imposed loading or deformation conditions, amorphous solids can display behaviour that combines features traditionally associated with both solids and liquids. These transitions between solid-like and liquid-like states, which are often associated with the onset of yielding or fracture \cite{nicolas_deformation_2018,Divoux2024,lin_scaling_2014,Ozawa2018,Dahmen2011}, are an essential part of understanding the behaviour of the material.

The physical origin of yielding in amorphous materials has been associated with a variety of microscopic phenomena \cite{bonn_yield_2017}. In earthquakes and landslides, the accumulation of stress along geological faults or shear zones can eventually overcome frictional resistance, leading to sudden slip and the rapid release of stored elastic energy \cite{Cocco2023}. In other systems, such as foams, emulsions, and colloidal suspensions, solid-like behaviour can arise through the jamming of the constituent particles, which become so densely packed that they are unable to rearrange and flow \cite{Mason1996,Datta2011,Derkach2009,AbuJdayil2003,Mewis2011,Koumakis2011,Petekidis2004,Cloitre2003,Torquato2010,Paredes2013,durian_foam_1995,WEAIRE2008}. At small deformations or loads, these materials respond elastically because the jammed packing of particles resists the imposed deformation. At larger loads or deformations, however, local particle rearrangements occur, allowing the material to accommodate the applied strain and relieve the stored stress through yielding. When such rearrangements become sufficiently widespread and persist over long enough timescales, fluid-like behaviour can emerge. This has been directly observed in colloidal suspensions and polymer melts, where the material typically possesses a characteristic relaxation timescale over which local stresses can dissipate \cite{Koumakis2011,Petekidis2004,Nordstrom2010}. For deformations applied on timescales shorter than this, the material can behave like a solid and fracture in a brittle manner, whereas for deformations applied on longer timescales, it can flow like a liquid \cite{nicolas_deformation_2018,bonn_yield_2017}.

A final class of amorphous materials that is particularly relevant to the work considered in this thesis is soft network materials. This term is commonly used to describe a broad range of materials that share a common structural motif: a disordered network of interconnected fibres \cite{tauber_microscopic_2020,gong_why_2010,tanaka_determination_2005,biotTheoryFiniteDeformations1972,mackintosh_elasticity_1995,Arevalo_Stress_Heterogeneities_2015}. Such materials arise in many different contexts, from everyday objects such as elastomers, which make up the rubber of car and bicycle tyres, to biological systems such as collagen fibre networks that provide structural support in connective tissues. They also appear in specialised materials such as polymer gels and engineered fibre networks used in manufacturing and medical applications \cite{yasudaNovelDoubleNetworkHydrogel2009,yinSelfHealingHydrogelsSynthesis2023,hoHydrogelsPropertiesApplications2022}. Despite occupying only a small fraction of the total material volume, the fibres in these systems form a load-bearing network that largely determines the mechanical response of the material \cite{Bonamy2011,sharma_strain-controlled_2016,Arevalo2015}. As a result, properties such as stiffness, yielding, and fracture are governed primarily by the structure and failure of this underlying network rather than by the surrounding medium. These ideas are central to the work developed in this thesis, which examines how disorder, stress redistribution, and microscopic failure processes govern macroscopic mechanical response.

\newpage
\section{Layout of thesis}

The work presented in this thesis consists of three main investigations, which can be grouped into two broader themes. The first theme focuses on yielding and fatigue in amorphous materials under cyclic deformation. The second concerns the fracture of soft network materials, first in a single disordered network and then in composite double-network systems. 

The overall structure of the thesis reflects this progression. \Cref{sec:Background} introduces the basic concepts and theoretical frameworks that are used throughout\DIFdelbegin \DIFdel{this thesis}\DIFdelend . This is followed by \Cref{sec:Fatigue}, which studies absorption and failure under oscillatory shear in an elastoplastic model. The next two chapters, \cref{sec:Ultradelayed,sec:DoubleNetworks}, turn to fracture in soft network materials, first under step strain and then in double-networks under extension. The thesis concludes with a final chapter that synthesises the main results and discusses their broader implications.

\subsection*{Theoretical Framework and Methods}
In \cref{sec:Background}, we introduce the concepts and theoretical frameworks used throughout this thesis. The chapter provides a common language for subsequent work, allowing the research chapters to focus on their specific physical problems without repeatedly reintroducing definitions and assumptions. In particular, it establishes the continuum and mesoscopic ideas needed to discuss deformation, yielding, and fracture in disordered soft materials.

The first part of the chapter is devoted to rheology and continuum mechanics. We introduce the basic concepts of stress and strain, distinguish between elastic and plastic deformation, and outline the rheological protocols used in later chapters. These include oscillatory shear, step strain, and shear startup, each employed to probe a different aspect of material response. A central component of this section is the force-balance equations in the Stokes limit, which provide a framework for understanding stress redistribution and underpin the arguments developed in later chapters.

The second part of the chapter turns from continuum descriptions to the second theme of this thesis: the fracture of soft network materials. This section connects the continuum framework used to describe bulk rheology with the discrete fibre networks that form simple models of many materials of interest. We introduce in more detail the types of soft network materials that are the focus of this section and consider the phenomena of rigidity, strain stiffening, and fracture, which are important for understanding the mechanical behaviour of these materials. Finally, we introduce the explicit mesoscale network model studied in the later chapters, along with the numerical methods used to simulate it. Taken together, these two parts provide the theoretical and methodological foundation on which the rest of the thesis is built.

\subsection*{Absorption and Failure in Oscillatory Shear}
\Cref{sec:Fatigue} is the first research chapter and addresses the first of our two broader themes. We examine the response of an amorphous material under large-amplitude oscillatory shear using a mesoscopic elastoplastic model, with the aim of understanding how repeated deformation drives the system either towards an absorbing state, in which plastic activity eventually ceases, or towards a yielded state, characterised by persistent plastic activity.

Amorphous materials subjected to cyclic loading can display strikingly different long-time responses. For strain amplitudes below a threshold, repeated oscillations over many cycles may organise the system into an absorbing state in which plastic activity ceases. At larger amplitudes, the same system can instead sustain irreversible rearrangements, promote shear localisation, and ultimately trigger macroscopic failure \cite{pineChaosThresholdIrreversibility2005,Cort2008,Fiocco2013,Regev_onset_2013,kawasakiMacroscopicYieldingJammed2016,Leishangthem2017,JocteurYielding2024}. Understanding this competition between absorption and yielding is central to fatigue in disordered solids.

In both particle-based simulations and coarse-grained descriptions, well-annealed systems tend to yield more abruptly and localise deformation more readily. In contrast, poorly annealed systems exhibit more distributed plastic activity and require a larger number of cycles before they reach an absorbed or yielded state \cite{Leishangthem2017,Popovi2018,barlow_ductile_2020,liuFateShearoscillatedAmorphous2022}. Elastoplastic models provide a useful framework for studying this behaviour because they retain the essential ingredients for yielding in amorphous materials, namely local plastic events, long-ranged stress redistribution, and structural disorder, while remaining simple enough to probe delayed dynamics over many deformation cycles \cite{nicolas_deformation_2018,pollard_yielding_2022}. Recent work has shown that such models can reproduce the rapid divergence of failure times near cyclic yielding and can exhibit a coexistence region in which absorbing and yielding trajectories remain possible over a finite interval of strain amplitudes \cite{khandareElastoplasticModellingCyclic2026,cochran_slow_2024}.

In this chapter, we use an athermal quasi-static elastoplastic model to examine this transition region in detail. We focus on how the imposed strain amplitude and the degree of annealing control both the probability of reaching an absorbing or yielded state and the number of cycles required to do so. We then show that, for sufficiently low annealing levels, the remaining plastic activity becomes progressively more dilute as the dynamics evolve, so that the approach to absorption can be interpreted as a first-passage process for weakly interacting mesoscopic elements. In this way, the chapter presents numerical results and develops a mechanistic interpretation, culminating in a stochastic picture of how delayed absorption emerges in the oscillatory response of the elastoplastic model.

\subsection*{Ultradelayed Fracture After Step Strain in Networks}
The work presented in \cref{sec:Ultradelayed} introduces the second, and largest, theme of this thesis. At this point, the focus shifts from yielding in amorphous solids to the study of fracture in soft network materials. 

The chapter studies the response of a disordered network subjected to an imposed step strain. In this protocol, the material is rapidly deformed to a fixed strain and then held at that deformation. Within a simple athermal picture, the subsequent response appears relatively straightforward: the stress either relaxes gradually towards a residual plateau, or the material fails shortly after loading. Recent constitutive modelling has challenged this picture by showing that elastoplastic materials subjected to a step strain can exhibit ultra-delayed instabilities, in which catastrophic failure occurs only after a long period of apparently stable behaviour extending over many decades in time \cite{lockwood_ultradelayed_2025}. Related delayed-yielding and delayed-fracture phenomenology has also been reported in soft gels, elastomers, and polymer networks, indicating that long-lived post-strain failure is observed in real materials \cite{sprakel_stress_2011,blindstrom_structures_2012,van_der_kooij_laser_2018,wangDelayedFractureGels2012,Ju2023}.

Inspired by these experimental works, we consider a simple mesoscopic model of a disordered network of fibres with thermally activated bond-breaking. This model provides a natural setting in which to study the problem, because its response to deformation is governed by non-affine motion, rigidity, and strain stiffening, all of which strongly influence how bond failures accumulate after loading \cite{shivers_strain-controlled_2023}. These features emerge from the intrinsic structure of the network and are not captured by traditional continuum models. By considering this more explicit model, we show that the phenomenon persists in a direct simulation of a disordered network. To incorporate the thermal effects necessary to capture the experimentally and theoretically observed phenomenology, we introduce a simple stochastic bond-breaking rate that models the thermal fluctuations that can cause a bond to break.

Using this minimal model, the chapter shows how ultra-delayed fracture can emerge in an explicit simulation of a disordered network. At low imposed strains, the network first relaxes towards a long-lived quasi-stable state that is almost indistinguishable from its counterpart at low temperatures, before the eventual rupture of a rare critical bond triggers a brittle system-spanning crack. At higher strains, substantial diffuse failure can accumulate while the macroscopic stress remains nearly unchanged, yet the network can persist for extremely long times before ultimately failing. In this way, \cref{sec:Ultradelayed} connects recent theoretical predictions of ultra-delayed post-strain instability to a concrete microscopic picture in which thermally activated bond failure, non-affine relaxation, and stress redistribution together control when and how catastrophic fracture finally occurs.

\subsection*{Toughness of Double Network Hydrogels}
\Cref{sec:DoubleNetworks}, the final research chapter of this thesis, continues the second of our two broader themes by turning from fracture in single disordered networks to composite double-network materials. Double-network hydrogels \cite{gong_why_2010, nakajima_synthesis_2013, sun_highly_2012, xin_effect_2013}, elastomers \cite{ducrot_toughening_2014, millereauMechanicsElastomericMolecular2018}, and macroscopic composites \cite{kingMacroscaleDoubleNetworks2019,takahashiCreatingStiffTough2018} combine a stiff, brittle sacrificial network with a soft, extensible matrix, allowing them to retain a high initial stiffness while failing at strains far beyond those of the stiff component alone \cite{gong_double-network_2003,gong_why_2010,ducrot_toughening_2014,millereauMechanicsElastomericMolecular2018}. Understanding how the interplay between these two networks produces this combination of stiffness and toughness is a central question in the study of soft composite network materials \cite{Creton50th2017,gong_why_2010,kingMacroscaleDoubleNetworks2019,tauber_stretchy_2022}.

Experiments on double-network materials suggest that their enhanced toughness is associated with stress delocalisation and diffuse failure within the sacrificial network, which suppress the early nucleation of a single catastrophic crack \cite{ducrot_structure_2015,huang_importance_2007,zheng_how_2021,fukao_effect_2020,ju_role_nodate}: scattering measurements indicate that deformation becomes more spatially delocalised in double networks \cite{ducrot_structure_2015,tominaga_molecular_2007,Orr2026}, while birefringence imaging shows that materials with better mechanical performance exhibit more homogeneous internal stress fields \cite{gong_double-network_2003,zheng_how_2021}. Compositional studies likewise reveal a transition from brittle macroscopic cracking to diffuse microcracking as the matrix becomes increasingly dominant \cite{fukao_effect_2020,ju_role_nodate,tauber_sharing_2021}. Together, these results suggest that stress delocalisation and diffuse failure are central to the development of toughness, although their structural origin remains unclear.

In this chapter, we introduce a minimal double-network model composed of a sparse, stiff sacrificial network and a denser, ductile matrix. By simulating its response under quasi-static extension, we show how the ratio of the two network length scales and relative stiffness control the transition from sacrificial-dominated to matrix-dominated behaviour. Increasing the capacity of the matrix to absorb and redistribute stress suppresses the cascade of neighbouring sacrificial-bond failures that would otherwise drive brittle fracture, enabling diffuse failure to accumulate in the sacrificial component while preserving the high stiffness of the sacrificial network and the large failure strain of the matrix. 

\subsection*{Conclusion}
The final chapter draws together the main results of the three investigations just outlined. Its purpose is to identify the broader physical ideas that connect the two main themes of the thesis, despite the differences in protocol and modelling between the chapters. In particular, it returns to the common roles of disorder, stress redistribution, and microscopic failure processes in controlling the macroscopic yielding and fracture behaviour of soft amorphous materials.
 \newpage       
 \newpage % !TeX root = ../../thesis.tex

\setcounter{equation}{0}
\chapter{Theoretical Framework and Methods} \label{sec:Background}

This thesis examines how yielding and fracture emerge in simplified theoretical models of soft disordered solids, with particular emphasis on disordered networks. The central problem is to understand how bulk mechanical response is connected to local failure, rearrangement, and the eventual loss of load-bearing capacity.

A key difficulty is the separation of scales involved in failure. Stress and strain are imposed and measured at the macroscopic scale, but yielding and fracture originate through local rearrangements and microscopic failure. Continuum variables are therefore needed to describe loading, force balance, and macroscopic response \cite{chandrasekharaiahContinuumMechanics2014,lubliner2008plasticity,macoskoRheologyPrinciplesMeasurements1994}, but they do not by themselves identify where damage begins or how it redistributes stress through a heterogeneous structure.

For this reason, the chapter starts with a discussion of continuum mechanics, and then moves to consider mesoscopic and discrete descriptions. We first introduce the rheological variables of stress and strain and the overdamped force-balance framework used throughout this thesis. We then turn to models in which local rearrangements, stress relaxation, and bond failure can be resolved explicitly \cite{falk_dynamics_1998,maloney_amorphous_2006,tanguy_plastic_2006,gartner_nonlinear_2016,picard_elastic_2004}. We consider in particular two-dimensional lattice elastoplastic models and disordered network models, which provide a minimal setting for studying damage accumulation, stress redistribution, and fracture in heterogeneous solids \cite{mackintosh_elasticity_1995,storm_nonlinear_2005,broedersz_modeling_2014,broedersz_criticality_2011,feng_nonlinear_2016,tauber_stretchy_2022}.

\section{Introduction to Rheology} \label{sec:IntroRheology}

Rheology is the study of how materials respond to applied stress, relating stress to the resulting deformation or rate of deformation \cite{larson_structure_1999,chenRheologySoftMaterials2010,tannerEngineeringRheology2000,barnes_introduction_1989} (or vice versa). The term originates from the early 20th century and comes from the Greek words \quotes{rheo}, meaning \quotes{to flow}, and \quotes{logos}, meaning \quotes{study} \cite{barnes_introduction_1989}. Early work by Hooke \cite{hooke_lectures_1678} and Newton \cite{newton_philosophiae_1687} established the simplest linear relations between stress and strain, and between stress and strain rate. Modern rheology, however, is primarily concerned with materials for which these simple linear relations break down and non-linear behaviour governs the response. Because the later chapters focus on such regimes, here we introduce only those concepts of stress and strain needed for the remainder of this thesis. The aim is not to provide a complete treatment of rheology, but rather a concise framework for interpreting the results that follow. For a fuller account, the reader is referred to the standard texts cited throughout this chapter.

\subsection{Stress and Strain} \label{sec:StressAndStrain}

The concepts of stress and strain form the basis of rheology. Strain quantifies the deformation of a material, while stress measures the internal forces that arise in response to that deformation.

The simplest measure of extensional strain is the relative change in material length,
\begin{equation}
    \epsilon = \frac{\Delta x}{L_x},
    \label{eq:simpleStrain}
\end{equation}
where $\Delta x$ is the change in length and $L_x$ is the original length measured along the same axis \cite{lubliner2008plasticity,macoskoRheologyPrinciplesMeasurements1994,barnes_introduction_1989}. This definition corresponds to extensional (or normal) deformation, in which a material element changes its length along a given direction. In contrast, materials may also deform through shear, in which layers of material slide past one another. In simple shear, if the upper boundary of a material is displaced by $\Delta x$ over a separation $L_y$, the shear strain is defined as
\begin{equation}
    \gamma_{xy} = \frac{\Delta x}{L_y}.
\end{equation}
\cite{lubliner2008plasticity,macoskoRheologyPrinciplesMeasurements1994,barnes_introduction_1989}. These two fundamental modes of deformation are illustrated in \cref{fig:typesOfStrain}.

Extension and shear represent two key modes of deformation, with more general deformations arising as combinations of the two. Because both $\epsilon$ and $\gamma_{xy}$ are defined relative to the undeformed configuration, they correspond to engineering strains, and the term strain will be used in this sense unless otherwise specified.

\begin{figure}
    \centering
    \includegraphics[width=.95\linewidth]{chapters/2-Background/Figures/StrainModes.pdf}
    \caption[Illustration of the two fundamental strain modes in 2D]{Illustration of the two fundamental strain modes in 2D: extensional strain $\epsilon$, where the element changes length, and shear strain \DIFdelbeginFL \DIFdelFL{$\gamma$}\DIFdelendFL \DIFaddbeginFL \DIFaddFL{$\gamma\equiv\gamma_{xy}$}\DIFaddendFL , where the element changes shape through angular distortion. The dashed lines denote the material's undeformed configuration, while the solid lines indicate its deformed state.}
    \label{fig:typesOfStrain}
\end{figure}

A simple measure of stress is the force per unit area transmitted across a surface,
\begin{equation}
    \sigma = \frac{F}{A},
    \label{eq:simpleStress}
\end{equation}
where $F$ is the force acting on an area $A$ \cite{macoskoRheologyPrinciplesMeasurements1994,barnes_introduction_1989}. However, both stress and strain are inherently directional quantities. Strain depends on the direction along which deformation is measured, while stress depends on both the orientation of the surface and the direction of the force acting across it. Consequently, they are most naturally described by second-rank tensors, the stress tensor $\bm{\sigma}$ and the strain tensor $\bm{\epsilon}$. Different components of these tensors reflect the two different deformation modes introduced above. Diagonal components correspond to normal stresses and extensional or compressive deformations, while off-diagonal components correspond to shear deformation and shear stresses. In \cref{sec:StrainTensor,sec:StressTensor}, we construct these tensors more formally, but we refer the reader to Refs.~\cite{tannerEngineeringRheology2000, chandrasekharaiahContinuumMechanics2014} for a more detailed treatment.

\subsection{Strain Tensor} \label{sec:StrainTensor}

The engineering strains introduced above are useful for simple extension and shear, but in general, a material element may stretch, compress, and shear in different directions simultaneously. To describe such local deformations, it is necessary to work in terms of the displacement field and its gradients.

We start by considering a material point initially at position $\bm{x}$, which is displaced to a new position $\bm{X}(\bm{x})$. The resulting displacement field is defined as
\[
    \bm{u}(\bm{x}) = \bm{X}(\bm{x}) - \bm{x}.
\]
While $\bm{u}(\bm{x})$ describes the displacement of an individual material point, deformation is determined by spatial variations of this displacement field. To quantify this, we consider two nearby material points initially separated in the reference configuration by $d\bm{x}$. Because $\bm{X}(\bm{x}) = \bm{x} + \bm{u}(\bm{x})$, the separation of the two points after the applied deformation is
\[
    d\bm{X} = \frac{\partial \bm{X}}{\partial \bm{x}} d\bm{x} = \bm{F} d\bm{x} = \left( \bm{I} + \frac{\partial \bm{u}}{\partial \bm{x}} \right) d\bm{x},
\]
where $\bm{F}$ is the deformation gradient \DIFaddbegin \DIFadd{and $\bm{I}$ is the identity tensor}\DIFaddend . The change in separation is therefore
\[
    d\bm{u} = d\bm{X} - d\bm{x} = \frac{\partial \bm{u}}{\partial \bm{x}} d\bm{x}.
\]
Thus, the local deformation is characterised by the displacement gradient $\nabla \bm{u}$, or in components, $\partial u_i/\partial x_j$.

Because the displacement gradient contains contributions from both deformation and rigid-body rotation, it may be decomposed as 
\[
    \nabla \bm{u} = \bm{\epsilon} + \bm{\omega},
\]
where $\bm{\epsilon}$ is symmetric and $\bm{\omega}$ is antisymmetric. The antisymmetric part $\bm{\omega} = \frac{1}{2}\left( \nabla \bm{u} - (\nabla \bm{u})^T \right)$ corresponds to bulk rigid-body rotations and does not contribute to the measured local deformation. The strain tensor, which quantifies this local deformation, is defined as the symmetric part of the displacement gradient,
\begin{equation}
\bm{\epsilon} = \frac{1}{2}\left( \nabla \bm{u} + (\nabla \bm{u})^T \right),
\label{eq:linearStrain}
\end{equation}
where we have neglected higher-order terms in the displacement gradient, because these are small in the limit of small deformations. This is known as the linearised strain tensor or infinitesimal strain tensor \cite{chandrasekharaiahContinuumMechanics2014}.

In Cartesian coordinates, the corresponding strain tensor in three dimensions may therefore be written as
\[
    \bm{\epsilon} =
    \begin{bmatrix}
        \dfrac{\partial u_x}{\partial x} & \dfrac{1}{2}\left( \dfrac{\partial u_x}{\partial y} + \dfrac{\partial u_y}{\partial x} \right) & \dfrac{1}{2}\left( \dfrac{\partial u_x}{\partial z} + \dfrac{\partial u_z}{\partial x} \right) \\[6pt]
        \dfrac{1}{2}\left( \dfrac{\partial u_y}{\partial x} + \dfrac{\partial u_x}{\partial y} \right) & \dfrac{\partial u_y}{\partial y} & \dfrac{1}{2}\left( \dfrac{\partial u_y}{\partial z} + \dfrac{\partial u_z}{\partial y} \right) \\[6pt]
        \dfrac{1}{2}\left( \dfrac{\partial u_z}{\partial x} + \dfrac{\partial u_x}{\partial z} \right) & \dfrac{1}{2}\left( \dfrac{\partial u_z}{\partial y} + \dfrac{\partial u_y}{\partial z} \right) & \dfrac{\partial u_z}{\partial z}
    \end{bmatrix}.
\]
The diagonal components of this tensor represent extensional strains, corresponding to fractional changes in length along the $x$, $y$, and $z$ directions, while the off-diagonal components represent shear strains, corresponding to changes in the angle between material line elements. These off-diagonal components are related to the engineering shear strains introduced in \cref{sec:StressAndStrain} through
\begin{equation}
    \gamma_{ij} = 2\epsilon_{ij}, \qquad i \neq j,
    \label{eq:engineeringShearStrain}
\end{equation}
so that, for example, $\gamma_{xy} = 2\epsilon_{xy}$. Because the strain tensor is symmetric, there exists a basis, known as the principal axes, in which the tensor is diagonal and the off-diagonal components vanish.

A further decomposition of the symmetric strain tensor $\bm{\epsilon}$ separates it into deviatoric and isotropic parts,
\[
    \bm{\epsilon} = \bm{d} + \frac{1}{D} \mathrm{Tr}(\bm{\epsilon}) \bm{I},
\]
where $\bm{d}$ is the deviatoric (traceless) part, $\bm{I}$ is the identity tensor, and $D$ is the dimensionality of the space. For incompressible materials, the volumetric strain vanishes, so that $\mathrm{Tr}(\bm{\epsilon}) = 0$ and the strain is purely deviatoric, $\bm{\epsilon} = \bm{d}$. This decomposition will be useful in the continuum descriptions developed later.

The form presented above is only valid in the limit of small deformations, where the displacement gradients are sufficiently small that higher-order terms can be neglected. For larger deformations, finite-strain measures such as the Green–Lagrange strain tensor must be used \cite{TruesdellNoll1992}. In this thesis, the linearised strain tensor serves as the common language used to discuss local elastic response, force balance, and stress redistribution.

\subsection{Stress Tensor} \label{sec:StressTensor}
We now introduce the stress tensor, which describes how internal stresses are distributed through a material.

If we consider a surface element of area $A$ within a material, the force $\bm{T}$ acting across this surface defines the traction,
\begin{equation}
    \bm{s} = \lim_{A \to 0} \frac{\bm{T}}{A},
\end{equation}
which represents the force per unit area acting on a surface with a given orientation.

To characterise the internal state of stress at a point, we may imagine making an arbitrary cut through the material passing through that point. The traction $\bm{s}$ then represents the force exerted across this cut. However, the form of the traction vector $\bm{s}$ depends on the orientation of the cut. Different cuts through the same point, with different normals, will in general yield different traction vectors $\bm{s}$.

A single traction vector $\bm{s}$ therefore provides only partial information about the internal state of stress. Instead, the full state of stress at a point is determined by the set of all traction vectors acting on all possible planes through that point. To relate these traction vectors, we consider an infinitesimal tetrahedral element with three faces aligned with the Cartesian coordinate planes, with normals $\bm{n}^{(x)}$, $\bm{n}^{(y)}$, and $\bm{n}^{(z)}$, and a fourth inclined face with normal $\bm{n}$. Requiring force balance on this element shows that the traction $\bm{s}(\bm{n})$ acting on the inclined face can be written as a linear combination of the tractions $\bm{s}(\bm{n}^{(x)})$, $\bm{s}(\bm{n}^{(y)})$, and $\bm{s}(\bm{n}^{(z)})$ acting on the coordinate planes. In the limit as the tetrahedron shrinks to a point, the traction therefore depends linearly on the surface normal $\bm{n}$. This construction is illustrated in \cref{fig:CauchyTetrahedron}. More generally, the stress tensor may be expressed in any orthonormal basis. As described for strain, there exists a set of principal axes for which the traction vectors on the corresponding coordinate planes are parallel to the plane normals, so that the stress tensor is diagonal in this basis.

\begin{figure}[t!]
    \centering
    \includegraphics[width=0.5\linewidth]{chapters/2-Background/Figures/cauchy_tetrahedron.pdf}
    \caption[Force equilibrium on a Cauchy tetrahedron for the derivation of the traction vector.]{Cauchy tetrahedron used to derive the traction vector on an arbitrary plane from force equilibrium, relating stress components on coordinate planes to the stress acting on an inclined plane.}
    \label{fig:CauchyTetrahedron}
\end{figure}

This linear relation can be written as
\begin{equation}
    \bm{s}(\bm{n}) = \bm{\sigma}\cdot\bm{n} \qquad (\textrm{or in components, } s_i = \sigma_{ij} n_j).
\end{equation}
This relation defines the Cauchy stress tensor $\bm{\sigma}$. Because we defined the three faces of the tetrahedron to be aligned with the coordinate planes, the stress tensor in three-dimensional Cartesian coordinates can be constructed as
\[
    \bm{\sigma} = 
    \begin{bmatrix}
    | & | & | \\
    \bm{s}^{(x)} & \bm{s}^{(y)} & \bm{s}^{(z)} \\
    | & | & |
    \end{bmatrix}
    = 
    \begin{bmatrix}
        \sigma_{xx} & \sigma_{xy} & \sigma_{xz} \\
        \sigma_{yx} & \sigma_{yy} & \sigma_{yz} \\
        \sigma_{zx} & \sigma_{zy} & \sigma_{zz}
    \end{bmatrix}.
\]

As with the strain tensor $\bm{\epsilon}$, the components of the stress tensor $\bm{\sigma}$ have a clear physical interpretation. The diagonal components correspond to normal stresses, acting perpendicular to a surface, while the off-diagonal components correspond to shear stresses, acting parallel to the surface. In the absence of body torques, local torque balance implies that the stress tensor is symmetric, i.e., $\sigma_{ij} = \sigma_{ji}$.
In Cartesian coordinates, it may therefore be written as
\[
    \bm{\sigma} =
    \begin{bmatrix}
        \sigma_{xx} & \sigma_{xy} & \sigma_{xz} \\
        \sigma_{xy} & \sigma_{yy} & \sigma_{yz} \\
        \sigma_{xz} & \sigma_{yz} & \sigma_{zz}
    \end{bmatrix}.
\]

\DIFaddbegin \DIFadd{Here, $\bm{x}$ and $\bm{X}$ distinguish positions in the reference and deformed configurations, respectively. In the subsequent small-strain calculations, this distinction may be neglected to leading order, and we use $\bm{r}=(x,y,z)$ to denote a generic spatial position.
}

\DIFaddend For notational clarity, throughout the remainder of this thesis we distinguish between local and bulk measures of stress. In continuum descriptions, the bulk stress tensor $\bm{\Sigma}$ is defined as the spatial average of the local stress,
\begin{equation}
    \Sigma_{ij}=\frac{1}{V}\int_V \sigma_{ij}(\bm{r})\, dV.
\end{equation}
This may differ from the local stress in the presence of spatial heterogeneity, such as that arising from plastic deformation. Throughout this thesis, $\bm{\Sigma}$ denotes a bulk or system-averaged stress. The precise method used to compute the stress differs across the models considered here, and is introduced where needed in the relevant chapter. In \cref{sec:Fatigue}, stress is defined explicitly as part of the construction of the elastoplastic model, while in \cref{sec:Ultradelayed,sec:DoubleNetworks} we use the virial stress tensor defined in \cref{sec:virialStress}.

\subsection{Linear Stress--Strain Relation}

To connect stress and strain, we require a constitutive relation. For a linearly elastic material under small deformation, the stress is a linear function of the strain, as given by Hooke’s law,
\[
    \sigma_{ij} = C_{ijkm}\epsilon_{km},
\]
where $C_{ijkm}$ is a fourth-rank tensor, often called the stiffness or elasticity tensor.

For an isotropic material, the stiffness tensor can be expressed in terms of two independent elastic moduli, $\mu$ and $\lambda$ (the Lam\'e parameters). The stress can therefore be written as
\[
    \sigma_{ij} = \lambda \epsilon_{kk} \delta_{ij} + 2\mu \epsilon_{ij}
\DIFdelbegin \DIFdel{.
}\DIFdelend \]
\DIFaddbegin \DIFadd{where $\delta_{ij}$ is the Kronecker delta.
}\DIFaddend 

For incompressible materials, the first term is replaced by a hydrostatic pressure $p$ that enforces the incompressibility constraint. The stress therefore reduces to
\begin{equation}
    \sigma_{ij} = -p\delta_{ij} + 2\mu\epsilon_{ij},
    \label{eq:incompressibleStress}
\end{equation}
which is the form used later when incompressibility is assumed.

\subsection{Plasticity} \label{sec:Plasticity}
The constitutive relation introduced above assumes a purely elastic response, in which deformation is fully reversible and the material returns to its original configuration upon unloading. However, many materials deviate from this behaviour once the applied stress exceeds a characteristic threshold. Beyond this point, the response is no longer purely elastic and irreversible deformation occurs. Here, yielding denotes the onset of this irreversible response, whether through plastic rearrangements, bond rupture, or other local irreversible events. Such processes introduce history dependence in the system's response and, in some cases, may progress to fracture, by which we mean the nucleation and propagation of cracks that lead to loss of the material's capacity to bear load.

In contrast to elastic deformation, in which all strain is recovered upon unloading and the material returns to its initial configuration, plasticity refers to the ability of a material to undergo permanent deformation under applied loading. Upon unloading, the original state is not recovered and the material instead retains a memory of its deformation history in its new configuration \cite{lubliner2008plasticity,Simo2000-ComputationalInelasticity,Wu2004}. Microscopically, these plastic deformations may correspond to particle rearrangements \cite{falk_dynamics_1998,maloney_amorphous_2006,tanguy_plastic_2006,gartner_nonlinear_2016}, dislocations \cite{shiba_plastic_2010, ispanovity_avalanches_2014,ben-zion_earthquake_1993}, or bond rupture \cite{tauber_stretchy_2022,ducrot_toughening_2014,Creton50th2017,vernereyStatisticalDamageMechanics2018,gong_why_2010}, all of which reduce the total stored elastic energy. Rheologically, the sum of these microscopic deformations may manifest in the material's macroscopic response as stress relaxation, yielding, and flow under sustained loading.

\subsection{Shear Banding} \label{sec:ShearBanding}

In the discussion so far, we have implicitly assumed that deformation remains homogeneous across the material. However, this assumption can break down, and the deformation may become spatially heterogeneous. An important example is shear banding, in which deformation localises into one or more bands that accommodate a large fraction of the total strain, while the material outside these bands remains only weakly deformed by comparison. Such shear bands are observed both experimentally \cite{manneville_recent_2008,divouxShearBandingComplex2016,berretTransientRheologyWormlike1997,berretNonlinearRheologyTelechelic2001,tapadiaBandingEntangledPolymer2006,goyonSpatialCooperativitySoft2008,besselingShearBandingFlowConcentration2010,divouxTransientShearBanding2010,rogersAgingYieldingShear2008} and theoretically \cite{fielding_complex_2007,fielding_shear_2014,mfielding_shear_2009,fielding_triggers_2016,moorcroft_age-dependent_2011,carterShearBandingPolymeric,OlmstedShearBanding2008,doi_theory_1988} across a wide range of materials, including metallic glasses \cite{greer_shear_2013,SCHUH2007}, colloidal glasses \cite{aimeMicroscopicDynamicsFailure2018,besselingShearBandingFlowConcentration2010}, granular materials \cite{Shang2024,nessAbsorbingStateTransitionsGranular2020}, and polymeric materials \cite{carterShearBandingPolymeric,tapadiaBandingEntangledPolymer2006,moorcroft_shear_2014}.

\begin{figure}[b!]
    \centering
        \begin{subfigure}{0.45\textwidth}
        \centering
        \includegraphics[width=\textwidth]{chapters/2-Background/Figures/ShearBanding_Homogeneous.pdf}
        \caption{Homogeneous deformation field}
        \label{fig:HomogeneousShear}
    \end{subfigure}
    \hfill
    \begin{subfigure}{0.45\textwidth}
        \centering
        \includegraphics[width=\textwidth]{chapters/2-Background/Figures/ShearBanding_Localized.pdf}
        \caption{Shear banded deformation field}
        \label{fig:ShearBanded}
    \end{subfigure}
    \caption[Illustration of homogeneous shear and shear banding.]{Schematic displacement profiles for {\bf a)} a homogeneous deformation field and {\bf b)} a shear-banded deformation field $u_x(y)$. Arrows indicate the direction and magnitude of the displacement. The system consists of two parallel plates; the top plate is moved in the $x$ direction by a distance $\Delta x$, while the bottom plate is fixed. Spatial invariance in the $x$ and $z$ directions is assumed.
    }
    \label{fig:shearBandSketch}
\end{figure}

In simple shear, with displacement in the $x$ direction and gradient along $y$, this localisation corresponds to a spatially varying shear strain,
\begin{equation}
    \gamma(y) = \frac{\partial u_x}{\partial y},
\end{equation}
which becomes strongly heterogeneous across the sample. A schematic illustration contrasting homogeneous deformation and shear banding is shown in \cref{fig:shearBandSketch}. Under homogeneous simple shear, the displacement field varies smoothly across the gap, as illustrated in \cref{fig:HomogeneousShear}. In the presence of shear banding, by contrast, the displacement profile becomes sharply non-uniform, as shown in \cref{fig:ShearBanded}: a narrow region accommodates much of the imposed deformation, while neighbouring regions deform much more weakly.

Shear banding should be distinguished from fracture. A shear band is a localised zone of intense deformation, typically associated with plastic flow or local weakening, whereas fracture involves the nucleation and propagation of a crack that leads to macroscopic separation or loss of load-bearing capacity. The two are nevertheless closely related. In some materials, a shear band remains a mode of heterogeneous flow, while in others it acts as a precursor to fracture by concentrating deformation and damage into a narrow region. This distinction is particularly important in amorphous solids, where shear bands often control plasticity and can also govern the route to final failure \cite{greer_shear_2013,hufnagel_deformation_2016,nicolas_deformation_2018}.

In many disordered solids, such as metallic glasses \cite{greer_shear_2013}, shear banding can be understood as the result of a mechanical instability of the homogeneous state. In particular, if the material exhibits strain softening, such that the stress decreases with increasing strain, $\partial_{\gamma}\sigma < 0$, then a uniformly deformed state becomes unstable \cite{barlow_ductile_2020}. Small spatial variations in strain are then amplified under continued loading, because regions that deform more undergo local softening, requiring less additional stress for further deformation and therefore promoting continued localisation. In experimental systems, the perturbations that trigger this instability may arise from structural heterogeneity or thermal fluctuations. This positive feedback leads to the formation of a shear band, which is sustained under continued deformation.

Depending on the applied deformation protocol and material properties, these localised regions may either be transient or persist in time. Of particular interest are transient shear bands, which form during the early stages of deformation but subsequently weaken, leading to a more spatially distributed deformation field \cite{moorcroft_age-dependent_2011,barlow_ductile_2020,divouxTransientShearBanding2010,adams_transient_2011}. These bands typically arise from an instability of the initially homogeneous response. However, as the material restructures, relaxes, or otherwise evolves, the conditions that favoured localisation may no longer be sustained, allowing deformation to become more evenly distributed again.

Alternatively, deformation may become permanently concentrated within well-defined bands that accommodate most of the imposed strain and act as precursors to macroscopic failure \cite{greer_shear_2013,hufnagel_deformation_2016,nicolas_deformation_2018,Falk2003,maloney_amorphous_2006,Shi2007,falk_deformation_2011,Barrat2011}. In such systems, once a dominant band has formed, subsequent deformation remains largely confined to this region, leading to progressive local weakening of the material along the band. This localisation is then self-reinforcing and may ultimately result either in catastrophic fracture or sustained, highly heterogeneous flow, depending on the material and loading conditions.

In general, shear banding reflects a competition between mechanisms that promote localisation, such as strain softening or yielding-induced weakening, and processes that redistribute deformation, including structural relaxation and long-range elastic stress redistribution, which can either suppress or enhance localisation depending on the material and geometry \cite{fielding_shear_2014,divouxShearBandingComplex2016,martens_spontaneous_2012,moorcroft_criteria_2013,moorcroft_age-dependent_2011,nicolas_deformation_2018}. Whether deformation remains homogeneous, develops transient bands, or evolves towards persistent localisation is therefore determined by the balance of these effects together with the imposed loading conditions \cite{fielding_shear_2014,divouxShearBandingComplex2016}. Understanding this competition is central to predicting failure and flow in disordered materials. In the following chapters, we will examine how these mechanisms emerge within the models considered here, and how they control the onset, evolution, and stability of strain localisation.

\section{Rheological protocol} \label{sec:Protocols}

\begin{figure}[p]
    \centering
    \includegraphics[width=0.8\linewidth]{chapters/2-Background/Figures/ShearProtocols.pdf}
    \caption[Schematics of deformation protocols]{Schematics of the three deformation protocols studied in this thesis. The imposed strain is shown by the solid black line, and a representative stress response is shown by the dashed red line. The strain is defined according to the relevant deformation mode in each chapter: $\gamma$ denotes the shear strain in \Cref{sec:Fatigue,sec:Ultradelayed}, and $\epsilon$ denotes the extensional strain in \Cref{sec:DoubleNetworks}. The protocols are as follows: \newline
    {\bf a)} Large-amplitude oscillatory shear (LAOS), in which the strain varies periodically between $-\gamma_0$ and $\gamma_0$. \newline
    {\bf b)} Step strain, in which a strain of amplitude $\gamma_0$ is applied instantaneously at $t=0$ and then held constant. The material subsequently relaxes towards a new equilibrium state. \newline
    {\bf c)} Strain startup, in which deformation is imposed at a constant rate $\dot{\epsilon}$ for $t>0$. A yielding response is illustrated, where the stress initially increases with strain before dropping beyond the critical strain $\epsilon^*$.
    }
    \label{fig:ShearProtocols}
\end{figure}

There are many deformation protocols by which the mechanical response of a material can be probed, each emphasising different aspects of that response. In this thesis, we focus on strain-controlled protocols\DIFdelbegin \DIFdel{, in which the bulk deformation is prescribed throughout the simulation and the resulting stress response is measured }\DIFdelend \DIFaddbegin \DIFadd{. In such protocols, a prescribed history of a macroscopic strain component, $\gamma(t)$ in shear or $\epsilon(t)$ in extension, is imposed through the boundaries or simulation cell, and the conjugate bulk stress $\Sigma(t)$ is measured as the material response}\DIFaddend . By contrast, in a stress-controlled protocol, the \DIFdelbegin \DIFdel{applied stress }\DIFdelend \DIFaddbegin \DIFadd{corresponding macroscopic stress history $\Sigma(t)$ }\DIFaddend is prescribed and the \DIFdelbegin \DIFdel{material is allowed to deform accordingly \mbox{%DIFAUXCMD
\cite{macoskoRheologyPrinciplesMeasurements1994}}\hskip0pt%DIFAUXCMD
. Stress-controlled }\DIFdelend \DIFaddbegin \DIFadd{strain evolves as the response \mbox{%DIFAUXCMD
\cite{macoskoRheologyPrinciplesMeasurements1994}}\hskip0pt%DIFAUXCMD
. Constant-stress }\DIFaddend tests are therefore naturally suited to studying \DIFdelbegin \DIFdel{the phenomena of creep and to the distinction }\DIFdelend \DIFaddbegin \DIFadd{creep and distinguishing }\DIFaddend between arrested, weakly creeping, and continuously deforming states, while strain-controlled tests \DIFdelbegin \DIFdel{instead resolve how the }\DIFdelend \DIFaddbegin \DIFadd{resolve how }\DIFaddend stress evolves under a prescribed deformation history, including through yielding and subsequent stress redistribution.

Within the linear elastic regime, stress- and strain-controlled protocols give equivalent information, because stress and strain are uniquely related by the elastic moduli. In the nonlinear regime, however, this equivalence no longer holds, and the two classes of protocol can produce qualitatively different responses. Since the focus of this thesis is on the evolution of stress under imposed deformation, including yielding, relaxation, and post-yield stress redistribution, strain-controlled protocols provide the most natural framework.

We consider two modes of deformation in this thesis: simple shear in \Cref{sec:Fatigue,sec:Ultradelayed} and elongational stretching in \Cref{sec:DoubleNetworks}. In addition, each research chapter employs a different strain-controlled protocol, as illustrated in \cref{fig:ShearProtocols}.

% There are many deformation protocols by which flow can be induced in a material, each probing different features of its mechanical response. In this work, we adopt strain-controlled deformation protocols, in which a prescribed bulk deformation is imposed at all times during the simulation, and the resulting stress response is measured. By contrast, in a stress-controlled protocol, a fixed (though potentially time-dependent) stress is imposed, and the material is allowed to deform accordingly \cite{macoskoRheologyPrinciplesMeasurements1994}. 

% Within the linear elastic regime, stress- and strain-controlled protocols are equivalent because stress and strain are uniquely related by the elastic moduli. In the nonlinear regime, however, the two protocols can yield different responses. In stress-controlled experiments, the imposed stress can produce time-dependent deformation (creep) even below yielding, and, if sufficiently large, lead to sustained flow with effectively unbounded deformation. Such protocols are therefore well suited to probing both creep and the onset of yielding through the transition from arrested or weakly creeping states to continuous flow. By contrast, in strain-controlled experiments, the deformation history is prescribed and the stress is measured as the material response. These protocols are particularly well suited to resolving the evolution of stress through yielding, including post-yield stress redistribution, which is the focus of this thesis.

% We will consider two forms of deformation in this thesis: simple shear is used in \Cref{sec:Fatigue,sec:Ultradelayed} and elongational stretching in \Cref{sec:DoubleNetworks}. In addition, each chapter considers a different strain-controlled protocol, as illustrated in \cref{fig:ShearProtocols}.

\subsection{Oscillatory \DIFdelbegin \DIFdel{Strain}\DIFdelend \DIFaddbegin \DIFadd{Shear}\DIFaddend } \label{sec:lAOS}

Oscillatory shear is a widely used protocol for characterising the rheology of soft materials, with large-amplitude oscillatory shear (LAOS) being particularly well studied \cite{rouyerLargeAmplitudeOscillatory2008a,yoshimuraResponseElasticBingham1987a,viasnoffHowAreColloidal2003,carterShearBandingPolymeric,ewoldtLargeAmplitudeOscillatory2010,renouYieldingProcessesColloidal2010}. In a conventional experiment, the strain is prescribed as a sinusoidal function of time,
\[
    \gamma(t) = \gamma_0 \sin(\omega t),
\]
where $\gamma_0$ is the strain amplitude and the frequency $\omega$ sets the inverse timescale over which the deformation is applied. For sufficiently small amplitudes, $\gamma_0 \ll 1$, the material responds in the linear viscoelastic regime, in which the stress is proportional to the applied strain. In this regime, it can be expressed in terms of the storage and loss moduli, which quantify the elastic and viscous contributions to the stress. At larger amplitudes, the response becomes nonlinear, leading to waveform distortion and the appearance of higher harmonics in the stress signal. These provide a sensitive probe of yielding and dissipative processes within the material \cite{ewoldtLargeAmplitudeOscillatory2010}.

In addition to resolving the intra-cycle response, oscillatory shear provides access to the evolution of a material over successive cycles. In many amorphous systems, repeated driving leads to the emergence of either periodic, reversible responses in which the system returns to the same configuration stroboscopically in each cycle, or to progressively evolving dynamics associated with irreversible rearrangements. The former corresponds to absorbing or limit-cycle states, while the latter often reflects fatigue, in which plastic deformation accumulates over many cycles and can ultimately lead to yielding or failure \cite{liuFateShearoscillatedAmorphous2022,Leishangthem2017,Fiocco_memory_2014,Regev_onset_2013}. Oscillatory shear is therefore particularly well suited to probing the transition between reversible and irreversible behaviour in driven amorphous materials. \DIFaddbegin \DIFadd{Beyond its use as a rheological probe, oscillatory deformation also has practical relevance in materials science: cyclic loading forms the basis of many accelerated fatigue tests used to investigate damage accumulation and estimate material lifetimes \mbox{%DIFAUXCMD
\cite{Suresh1998,DowlingMechanical2013}}\hskip0pt%DIFAUXCMD
.
}

\DIFadd{Experimentally, oscillatory shear is commonly realised using a rotational rheometer, with the sample confined between a stationary surface and a moving plate or cone. The instrument imposes a sinusoidal angular displacement and measures the resulting torque, which is converted into the macroscopic shear stress using the dimensions of the measuring geometry \mbox{%DIFAUXCMD
\cite{macoskoRheologyPrinciplesMeasurements1994}}\hskip0pt%DIFAUXCMD
.
}\DIFaddend 

In \cref{sec:Fatigue}, we consider the quasi-static limit of this oscillatory shear protocol, in which the driving frequency is taken to be small compared to the intrinsic inverse relaxation timescales of the material, $\omega \to 0$. In this limit, inertial and rate-dependent effects are negligible and the system evolves through a sequence of mechanically equilibrated states, even as plastic events occur. In practice in our computational simulations, this corresponds to cycling the strain between $-\gamma_0$ and $\gamma_0$ over successive cycles, pausing at strains where any local yielding events occur within the material, and evolving the system to a new equilibrium state before proceeding to the next strain in the cycle at which any further yielding events occur. This protocol has become a common tool for probing the stability and memory of amorphous solids under repeated loading, because it emphasises the role of accumulated plastic activity while minimising rate-dependent effects \cite{liuFateShearoscillatedAmorphous2022,Leishangthem2017,Fiocco_memory_2014,Regev_onset_2013}.

\subsection{Step Strain} \label{sec:StepStrain}
In \Cref{sec:Ultradelayed} we focus on the step strain protocol. In this protocol, after the sample is initialised, a strain of amplitude $\gamma_0$ is applied instantaneously at $t=0$ and is held constant for all subsequent times. The response is then characterised by the evolution of the stress $\Sigma(t)$, which typically exhibits a rapid initial increase as the material is strained, followed by a time-dependent decay as the sample relaxes towards a new equilibrium state. This relaxation reflects the progressive redistribution of internal stresses and provides a direct probe of the relaxation dynamics of the material.

Mathematically, the imposed strain may be modelled as
\[
    \gamma(t) = \gamma_0 \Theta(t),
\]
where $\Theta(t)$ is the Heaviside function. Formally, this corresponds to a Dirac delta function in the strain rate, $\dot{\gamma}(t) = \gamma_0 \delta(t)$. In \DIFdelbegin \DIFdel{practice, an }\DIFdelend \DIFaddbegin \DIFadd{a physical experiment, the sample is placed in a strain-controlled rheometer, whose moving surface is rapidly displaced by an amount corresponding to $\gamma_0$ and then held fixed while the resulting torque is measured. An }\DIFaddend idealised step strain cannot be realised experimentally due to the finite inertia of rheological instruments. Nevertheless, a strain ramp applied on a timescale which is short compared with the intrinsic material relaxation time can be modelled effectively as an instantaneous step \cite{Fang_inhomogeneity_2011}.

The step strain protocol is particularly useful because it isolates the intrinsic relaxation processes of the material. In linear viscoelastic systems, the stress typically decays smoothly over time, reflecting a spectrum of relaxation times associated with the underlying microstructure \cite{macoskoRheologyPrinciplesMeasurements1994,larson_structure_1999}. In disordered and glassy materials, however, the relaxation may be highly nontrivial, exhibiting multiple timescales, slow decay, ageing, or delayed yielding following the initial deformation \cite{fielding_complex_2007,moorcroft_age-dependent_2011,lockwood_ultradelayed_2025}.

\subsection{Strain Startup} \label{sec:ShearStartup}

Another common protocol for characterising a material’s rheological response prior to steady flow or failure is strain startup, in which deformation is imposed continuously at a constant rate $\dot{\epsilon}$ for all times $t>0$ in a sample that was previously undeformed for time $t<0$ \cite{larson_structure_1999,tannerEngineeringRheology2000,barnes_introduction_1989}. This protocol probes how the material evolves from its initial undeformed state under sustained driving and provides a direct means of accessing the transient dynamics leading to yielding and steady flow. In particular, it allows one to resolve the progression from elastic loading, through yielding, to a flowing or plastically deforming state, or to catastrophic material failure, and it is widely used to identify elastic regimes, and the onset of irreversible deformation \cite{divouxShearBandingComplex2016,moorcroft_age-dependent_2011}.

\DIFaddbegin \DIFadd{In a shear experiment, strain startup is realised by accelerating the moving surface of a rheometer to a prescribed constant velocity and measuring the resulting torque and forces. In an extensional experiment, as considered in }\Cref{sec:DoubleNetworks}\DIFadd{, the boundaries of the sample are separated at a prescribed rate while the resulting tensile force is measured. The latter is the basis of a conventional tensile test.
}

\DIFaddend For an elastoplastic material with a yield threshold, the stress initially increases linearly with strain, $\Sigma \propto \epsilon = \dot{\epsilon} t$, until a characteristic yield strain is reached, beyond which the response deviates from linear elasticity. By contrast, a Newtonian fluid exhibits a constant stress response, $\Sigma = \eta \dot{\epsilon}$, determined entirely by the imposed rate\DIFaddbegin \DIFadd{,where $\eta$ is the fluid viscosity}\DIFaddend . Yield stress materials display behaviour intermediate between these limits: the stress initially increases approximately linearly with strain, reflecting an elastic regime, before approaching a steady-state in which the value of stress is set by the flow curve $\Sigma(\dot{\epsilon})$. In many such materials, particularly in shear, the stress reaches a maximum (a stress overshoot) and subsequently relaxes towards a steady-state value $\Sigma(t \to \infty) = \Sigma_{\mathrm{steady}}(\dot{\epsilon})$ consistent with the underlying flow curve \cite{divouxShearBandingComplex2016,larson_structure_1999}. However, not all materials evolve towards steady flow under continued deformation. In some cases, the material is unable to sustain any applied load beyond a critical strain $\epsilon^*$, and instead undergoes fracture or catastrophic failure.

In \Cref{sec:DoubleNetworks}, we consider a quasi-static version of this protocol, in which the imposed deformation is applied sufficiently slowly that the system remains close to mechanical equilibrium at all times, corresponding to $\dot{\epsilon} \to 0$. In this limit, inertial and finite-rate effects are negligible, and the response can be interpreted as a sequence of mechanically equilibrated states, allowing the onset of plasticity and failure to be resolved in terms of the underlying structural evolution.

% \subsection{Experiments} \label{sec:Experements}

% In order to understand these nonlinear effects, which are often controlled by the dynamics of internal microstructure, experimental rheologists use a number of experimental techniques which often involve shearing a material between two plates.

\section{Force balance in the Stokes limit} \label{sec:Hydrodynamics}

The constitutive descriptions introduced above specify how local stress responds to deformation. Across this thesis, we also require a general framework for how stress relaxes and redistributes after local perturbations, including plastic rearrangements in mesoscopic models and bond-level damage in networks. In the regime of interest in this thesis, for the disordered soft materials that we consider, this redistribution is governed by incompressible hydrodynamics in the low-Reynolds-number limit. We therefore summarise the force-balance equations in a form relevant to this work.

% \subsection{Navier-Stokes}
% From conservation of mass and momentum, the Navier--Stokes equations for a generalised incompressible flow read
% \begin{subequations}
% \label{eq:NavierStokes}
% \begin{align}
%     \rho \left(\partial_t+(\bm{u} \cdot \nabla)\right)\bm{u} &= \nabla \cdot \bm{\sigma} \label{eq:incompressibleNavier1} \\
%     \nabla\cdot\bm{v} &= 0 \label{eq:incompressibleNavier2}
% \end{align}
% \end{subequations}

% where $\rho$ is the fluid's density, $\bm{v}$ the fluid's flow field, and $\bm{\sigma}$ is a tensor which characterises the total stress applied to the fluid. The total stress may be decomposed as $\bm{\sigma} = \bm{\tau} - p\bm{I}$, where $\bm{\tau}$ is the deviatoric stress tensor and $p$ is the hydrostatic pressure.

% The physics we wish to consider in this thesis is the slow flow of pasty materials where the internal viscous forces dominate over the inertial forces. We therefore consider the low-Reynolds-number limit of \cref{eq:incompressibleNavier1} as 
% \[
%     Re=\rho U L / \mu \to 0
% \]
% where $U$ and $L$ are characteristic displacements and length scales. In this limit, the inertial term on the left-hand side of \cref{eq:incompressibleNavier1} becomes negligible, and the equations reduce to
% \begin{subequations}
% \label{eq:StokesLimit}
% \begin{align}
%     \nabla \cdot (\bm{\tau} - p\bm{I}) &= 0 \\
%     \nabla \cdot \bm{u} &= 0.
% \end{align}
% \end{subequations}
% This result may also be obtained directly from the mechanical equilibrium condition in \cref{eq:incompressibleStress}, which requires the net force in the system to vanish:
% \[
% \nabla \cdot (\mu \bm{\epsilon}) = \nabla \cdot \bm{\sigma} = 0.
% \]

\subsection{\DIFaddbegin \DIFadd{The }\DIFaddend Navier-Stokes \DIFaddbegin \DIFadd{Equations}\DIFaddend }
From conservation of mass and momentum, the Navier--Stokes equations for a generalised incompressible flow read
\begin{subequations}
\label{eq:NavierStokes}
\begin{align}
    \rho \left(\partial_t+(\bm{v} \cdot \nabla)\right)\bm{v} &= \nabla \cdot \bm{\sigma} \label{eq:incompressibleNavier1} \\
    \nabla\cdot\bm{v} &= 0 \label{eq:incompressibleNavier2}
\end{align}
\end{subequations}
where $\rho$ is the fluid density, $\bm{v}$ the velocity field, and $\bm{\sigma}$ the total Cauchy stress tensor. The stress may be decomposed as
\begin{equation*}
    \bm{\sigma} = \bm{\tau} - p\bm{I},
\end{equation*}
where $\bm{\tau}$ is the deviatoric stress tensor and $p$ the hydrostatic pressure enforcing incompressibility.

The physics considered in this thesis concerns the slow deformation of soft disordered materials, in which viscous and elastic stresses dominate over inertia. We therefore consider the low-Reynolds-number limit of \cref{eq:incompressibleNavier1},
\[
    \text{Re} = \frac{v L}{\nu} \to 0,
\]
where $v$ and $L$ are characteristic velocity and length scales and $\nu$ is the kinematic viscosity. In this limit, inertial terms become negligible and the equations reduce to the Stokes (force-balance) equations
\begin{subequations}
\label{eq:StokesLimit}
\begin{align}
    \nabla \cdot (\bm{\tau} - p\bm{I}) &= 0, \\
    \nabla \cdot \bm{v} &= 0.
\end{align}
\end{subequations}

In what follows, we adopt a displacement-based formulation. Introducing a displacement field $\bm{u}(\bm{r},t)$ with $\bm{v}=\partial_t\bm{u}$, the quasi-static limit corresponds to determining $\bm{u}$ at each time from the mechanical equilibrium condition \cref{eq:StokesLimit}.

\DIFaddbegin \DIFadd{Although these equations are conventionally introduced for incompressible fluids, the force-balance and incompressibility conditions are not restricted to fluids. Whether the material responds as a fluid or a solid is determined by the constitutive relation for the deviatoric stress. For a Newtonian fluid, $\bm{\tau}$ is proportional to the strain-rate tensor, whereas for the solid-like materials considered here, it is proportional to the elastic strain. In the quasi-static limit, the equilibrium equations for an incompressible, linearly elastic solid have the same mathematical form as the Stokes equations, with the displacement $\bm{u}$ replacing the velocity $bm{v}$, and the shear modulus $\mu$ replacing the viscosity. We therefore adopt a displacement-based formulation, with $\bm{v}=\partial_t\bm{u}$, and determine $\bm{u}$ at each time from mechanical equilibrium.
}

\DIFaddend To connect this framework to models with local irreversible rearrangements (notably the elastoplastic model in \cref{sec:Fatigue}), we introduce a plastic eigenstrain field following Ref.~\cite{picard_elastic_2004}. Here, the eigenstrain $\bm{\epsilon}^{\mathrm{pl}}$ denotes a locally imposed, irreversible strain representing only the deformation generated by the rearrangement events themselves, excluding any elastic strain induced in the surrounding material. Assuming a homogeneous, isotropic, linearly elastic medium, the total strain may be written as
\begin{subequations}    \label{eq:decompositionAndDeformation} \label{eq:EPStrain}
\begin{gather}
    \bm{\epsilon} = \bm{\epsilon}^{\mathrm{el}} + \bm{\epsilon}^{\mathrm{pl}}
    \label{eq:strainDecomp} \\
    \bm{\epsilon} =
    \frac{1}{2}\left(\nabla \bm{u} + (\nabla \bm{u})^{T}\right)
\end{gather}
\end{subequations}
where $\bm{u}(\bm{r})$ is the total displacement vector at a position $\bm{r}$. For a homogeneous, isotropic, linearly elastic and incompressible medium ($\nabla\cdot\bm{u}=0$), the deviatoric stress becomes
\begin{equation}
    \bm{\tau} = 2\mu \bm{\epsilon}^{\mathrm{el}} = 2\mu\bm{\epsilon} - 2\mu\bm{\epsilon}^{\mathrm{pl}}.
    \label{eq:elasticStress}
\end{equation}

% Now, considering a single localised plastic event induced by elastic loading, the strain field $\bm{\epsilon}$ can be decomposed into a superposition of $\bm{\epsilon}^{\mathrm{el}0}$, $\bm{\epsilon}^{\mathrm{el}1}$ and $\bm{\epsilon}^{\mathrm{pl}1}$ where the superscript $0$ denotes the contribution from the elastic loading and $1$ the contribution from the localised plastic event. The materials deformation field can also be decomposed in the same way into a pair of deformation fields $\bm{u}^0$ and $\bm{u}^1$, which are the material's elastic and plastic responses respectively, so that $\bm{u} = \bm{u}^0 + \bm{u}^1$. As the elastic contribution must obey \cref{eq:StokesLimit} and the incompressibility condition $\nabla \cdot \bm{u}^0 = 0$, the plastic deformation $\bm{u}^1$ resulting from the localised plastic event must satisfy, 
% \begin{subequations}
%     \label{eq:plasticStokes}
%     \begin{align}
%         \nabla \cdot \left( 2\mu \bm{\epsilon}^{\mathrm{el}1} - p^1 \bm{I} \right) &= 2\mu \nabla \cdot \bm{\epsilon}^{\mathrm{pl}1}, \label{eq:plasticStokes1}\\
%         \nabla \cdot \bm{u}^1 &= 0.
%     \end{align}
% \end{subequations}

% Under the assumption that the deformation is isotropic and incompressible ($\nabla \cdot \bm{u} = \bm{0}$) for any deformation at rest, we can express the left-hand side of \cref{eq:plasticStokes1} as 
% \[\nabla \cdot \left( 2\mu \bm{\epsilon}^{\mathrm{el}1} - p^1 \bm{I} \right) = -\nabla p^1 + \mu \nabla^2 \bm{u}^1.\]

\subsection{\DIFaddbegin \DIFadd{The }\DIFaddend Oseen Tensor} \label{sec:Oseen}
\DIFdelbegin \DIFdel{The solution }\DIFdelend \DIFaddbegin \DIFadd{Solutions }\DIFaddend to \cref{eq:StokesLimit} \DIFdelbegin \DIFdel{is }\DIFdelend \DIFaddbegin \DIFadd{are }\DIFaddend well known and \DIFdelbegin \DIFdel{has }\DIFdelend \DIFaddbegin \DIFadd{have }\DIFaddend been studied in a wide range of systems. For a local plastic eigenstrain, the displacement field in an unbounded medium obeys:
\begin{subequations}
\label{eq:Stokes}
\begin{align}
- \nabla p(\bm{r}) + \mu \nabla^2 \bm{u}(\bm{r}) &= 2\mu \nabla \cdot \bm{\epsilon}^{\mathrm{pl}}(\bm{r}), \label{eq:Stokes1} \\
\nabla \cdot \bm{u} &= 0. \label{eq:Stokes2}
\end{align}
\end{subequations}
where $(\nabla \cdot \bm{\epsilon}^{\mathrm{pl}})_i = \partial_j \epsilon^{\mathrm{pl}}_{ij}$. The displacement field can then be expressed as a convolution of this eigenstrain source with the Green's function $\mathbb{O}_{ij}(\bm{r},\bm{r}^\prime)$,
\begin{equation}
    u_i(\bm{r}) = -2\mu \int \mathbb{O}_{ij}(\bm{r}-\bm{r}^\prime)\,\partial_l^\prime \epsilon_{jl}^{\mathrm{pl}}(\bm{r}^\prime)\, d\bm{r}^\prime
\DIFdelbegin \DIFdel{.
}\DIFdelend \end{equation}
\DIFaddbegin \DIFadd{where $\partial_l^\prime$ denotes differentiation with respect to the primed coordinates.
}\DIFaddend The tensor $\bm{\mathbb{O}}$ was first derived by Oseen in 1927 \cite{oseen_neuere_1927} and refined by Burgers in \DIFdelbegin \DIFdel{1939 }\DIFdelend \DIFaddbegin \DIFadd{1938 }\DIFaddend \cite{Burgers1939}. In the present context, it provides the projection kernel that maps the eigenstrain source onto an incompressible displacement field. This Green-function structure is closely related to Eshelby's treatment of an inclusion with a prescribed internal transformation strain \cite{eshelby_determination_1997}.

The \DIFaddbegin \DIFadd{Fourier transform of the }\DIFaddend Oseen tensor is obtained as
\begin{equation}
    \hat{\mathbb{O}}_{ij} = \frac{1}{\mu k^2}\left( \delta_{ij} - \frac{k_ik_j}{k^2} \right).
    \label{eq:Oseen}
\end{equation}
This yields the displacement field in Fourier space,
\begin{equation}
    \hat{u}_i = -2\mu \hat{\mathbb{O}}_{ij} i k_l \hat{\epsilon}_{jl}^{\mathrm{pl}}
    = -\frac{2i}{k^2}\left( \delta_{ij} - \frac{k_ik_j}{k^2} \right) k_l \hat{\epsilon}_{jl}^{\mathrm{pl}}
    \label{eq:fourierDeformation}
\end{equation}
and the corresponding deformation gradients are
\begin{equation}
    ik_m\hat{u}_i
    = 2\mu k_m \hat{\mathbb{O}}_{ij} k_l \hat{\epsilon}_{jl}^{\mathrm{pl}}
    = \frac{2 k_m}{k^2}\left( \delta_{ij} - \frac{k_ik_j}{k^2} \right) k_l \hat{\epsilon}_{jl}^{\mathrm{pl}}.
    \label{eq:fourierDeformationGradient}
\end{equation}
A full derivation of the results given in \cref{eq:Oseen,eq:fourierDeformation,eq:fourierDeformationGradient} is provided in \Cref{app:OseenDerivation}. The expressions in \cref{eq:fourierDeformation,eq:fourierDeformationGradient} describe the deformation and deformation gradient fields for an infinite system only up to an arbitrary constant displacement $\bm{u}_0$, reflecting the singular behaviour at $\bm{k}=\bm{0}$. This constant will be fixed below.

\subsection{Eshelby Stress Propagator} \label{sec:EshelbyPropogator}

We now use the deformation gradients obtained above to construct the corresponding strain and stress fields. This provides one concrete application of the mechanical-equilibrium framework: the Eshelby stress propagator describing elastic redistribution after a localised plastic event.

From \cref{eq:EPStrain}, the resulting strain and stress fields are
\begin{equation}
        \epsilon_{ij} = \epsilon_{ij}^{\mathrm{el}} + \epsilon_{ij}^{\mathrm{pl}} \implies \epsilon^{\mathrm{el}}_{ij} = \frac{1}{2}\left( \partial_i u_j + \partial_j u_i \right) - \epsilon_{ij}^{\mathrm{pl}}
        \label{eq:eshelbyStrain}
\end{equation}
and
\begin{equation}
    \tau_{ij} = 2\mu \epsilon^{\mathrm{el}}_{ij}
    \label{eq:eshelbyStress}
\end{equation}
These can then be expressed in terms of the deformation gradients found in \cref{eq:fourierDeformationGradient}.
\DIFaddbegin 

\DIFaddend For simplicity, we consider a two-dimensional system and focus on the elastic stress redistribution \DIFaddbegin \DIFadd{$\tau_{xy}$ }\DIFaddend associated with a \DIFdelbegin \DIFdel{local plastic shear event $\epsilon_{xy}^{\mathrm{pl}}$, characterised by the corresponding stress component $\tau_{xy}$}\DIFdelend \DIFaddbegin \DIFadd{localised plastic event of simple-shear symmetry. Following Picard et al.~\mbox{%DIFAUXCMD
\cite{picard_elastic_2004}}\hskip0pt%DIFAUXCMD
, we represent a point-like event centred at $\bm{r}_0$ as
}\begin{equation}
    \DIFadd{\epsilon_{xy}^{\mathrm{pl}}(\bm{r})
    =
    a^2\epsilon_0\delta^{(2)}(\bm{r}-\bm{r}_0),
}\end{equation}
\DIFadd{where $\epsilon_0$ is the dimensionless plastic-strain amplitude, $a^2$ is the microscopic area of the event and $\delta^{(2)}$ is the two-dimensional dirac delta function. The coefficient $a^2\epsilon_0$ therefore has dimensions of area. Its Fourier transform is
}\[
    \DIFadd{\hat{\epsilon}_{xy}^{\mathrm{pl}}(\bm{k})
    =
    a^2\epsilon_0e^{-i\bm{k}\cdot\bm{r}_0}.
}\]
\DIFadd{For an event at the origin, $\bm{r}_0=\bm{0}$, this reduces to $\hat{\epsilon}_{xy}^{\mathrm{pl}}=a^2\epsilon_0$, equal to the coefficient multiplying the real-space delta function}\DIFaddend .
This choice is made for convenience, because the same analysis can be carried out for any component of the resulting stress and eigenstrain components. While the explicit form of the response differs between components and in three dimensions, the underlying mathematical structure of the formalism remains unchanged and extends straightforwardly to more general cases.

In two dimensions, the Fourier representations of the displacement components $u_x$ and $u_y$ take the form,
\begin{subequations}
\begin{align}
    \hat{u}_x &= -2\mu\left(\hat{\mathbb{O}}_{xx}ik_y + \hat{\mathbb{O}}_{xy} ik_x \right)\hat{\epsilon}_{xy}^{\mathrm{pl}}, \label{eq:ux}\\
    \hat{u}_y &= -2\mu\left(\hat{\mathbb{O}}_{yx}ik_y + \hat{\mathbb{O}}_{yy} ik_x \right)\hat{\epsilon}_{xy}^{\mathrm{pl}}. \label{eq:uy}
\end{align}
\end{subequations}
Substituting $\hat{u}_x$ and $\hat{u}_y$ into \cref{eq:eshelbyStrain,eq:eshelbyStress} and resolving for $\hat{\tau}_{xy}$, we obtain
\begin{align*}
    \hat{\tau}_{xy} = 2\mu \hat{\epsilon}^{\mathrm{el}}_{xy} &= \mu\left( i k_x \hat{u}_y + i k_y \hat{u}_x \right) - 2\mu\hat{\epsilon}_{xy}^{\mathrm{pl}}\\
    & = 2\mu^2 \left( k_y^2 \hat{\mathbb{O}}_{xx} + k_x^2 \hat{\mathbb{O}}_{yy} + 2 k_x k_y\hat{\mathbb{O}}_{xy} - \frac{1}{\mu} \right)\hat{\epsilon}_{xy}^{\mathrm{pl}}
\end{align*}
where we use the symmetry of the Oseen tensor $\hat{\mathbb{O}}_{xy}=\hat{\mathbb{O}}_{yx}$.

By collecting terms, we obtain the shear stress propagator $\mathbb{G}_{xy}$, which describes the stress redistribution produced by a single plastic event in an infinite medium:
\begin{equation}
\tau_{xy}(\bm{r}) = \DIFaddbegin \DIFadd{2}\DIFaddend \mu \int \mathbb{G}_{xy}(\bm{r}-\bm{r}') \epsilon_{xy}^{\mathrm{pl}}(\bm{r}') \, d\DIFaddbegin \DIFadd{^2}\DIFaddend \bm{r}',
\end{equation}
where the Fourier-space kernel is
\begin{equation}
\hat{\mathbb{G}}_{xy}(\bm{k}) = -4\frac{k_x^2 k_y^2}{k^4}.
\label{eq:StressPropogator}
\end{equation}
We illustrate in \cref{fig:Propogator} the \DIFdelbegin \DIFdel{stress redistribution resulting from a single plastic event with $\epsilon_{xy}^{\mathrm{pl}} = -1$ }\DIFdelend \DIFaddbegin \DIFadd{response to a point plastic event centred }\DIFaddend at the origin\DIFaddbegin \DIFadd{, $\bm{r}_0=\bm{0}$. For this normalised plot, we use grid units in which the area of one element is unity and choose $a^2\epsilon_0=-1$. This unit value normalises the linear Green-function response and does not represent a physical strain of magnitude one; the response to a physical event with $|\epsilon_0| \ll 1$ is obtained by linear rescaling. Numerically, the plastic eigenstrain is assigned to the central grid element, Fourier transformed, multiplied by the propogator $2\mu\hat{\mathbb{G}}_{xy}(\bm{k})$, and transformed back to real space using the FFTW3 library \mbox{%DIFAUXCMD
\cite{Frigo2005}}\hskip0pt%DIFAUXCMD
}\DIFaddend .

\DIFdelbegin %DIFDELCMD < \begin{figure}
%DIFDELCMD <     %%%
\DIFdelendFL \DIFaddbeginFL \begin{figure}[t]
    \DIFaddendFL \centering
    \includegraphics[width=0.75\linewidth]{chapters/2-Background/Figures/Prop_512_inset17.pdf}
    %\missingfigure{Plot of Propagator}
    \caption[Quadrupolar elastoplastic stress propagator]{Quadrupolar redistribution of shear stress $\tau_{xy}$ in two dimensions \DIFdelbeginFL \DIFdelFL{, after }\DIFdelendFL \DIFaddbeginFL \DIFaddFL{following }\DIFaddendFL a \DIFdelbeginFL \DIFdelFL{single local }\DIFdelendFL \DIFaddbeginFL \DIFaddFL{normalised point }\DIFaddendFL plastic event in a system with $\mu=1$. The system is \DIFdelbeginFL \DIFdelFL{initialised with $\tau_{xy}=0$ everywhere; one central element is then yielded }\DIFdelendFL \DIFaddbeginFL \DIFaddFL{initially unstressed, }\DIFaddendFL and \DIFaddbeginFL \DIFaddFL{the discrete plastic eigenstrain is }\DIFaddendFL set to \DIFdelbeginFL \DIFdelFL{$\tau_{xy}=-1$}\DIFdelendFL \DIFaddbeginFL \DIFaddFL{$-1$ at the central grid element and zero elsewhere. The plastic strain, rather than the stress}\DIFaddendFL , \DIFdelbeginFL \DIFdelFL{after which mechanical equilibrium }\DIFdelendFL is \DIFdelbeginFL \DIFdelFL{enforced}\DIFdelendFL \DIFaddbeginFL \DIFaddFL{prescribed at the central element; the resulting equilibrated stress field is calculated using the Fourier-space propagator}\DIFaddendFL . The main panel shows the full \DIFdelbeginFL \DIFdelFL{$N_x\times N_y = 512\times512$ propagator}\DIFdelendFL \DIFaddbeginFL \DIFaddFL{$N_x\times N_y=512\times512$ field}\DIFaddendFL , and the inset shows the central $17\times17$ region\DIFdelbeginFL \DIFdelFL{around the propagator}\DIFdelendFL .}
    \label{fig:Propogator}
\end{figure}

Equivalent expressions can be derived for the normal stress components $\tau_{xx}$ and $\tau_{yy}$, as well as for plastic events associated with normal stresses $\mu \epsilon_{xx}^{\mathrm{pl}}$ and $\mu \epsilon_{yy}^{\mathrm{pl}}$. Alternatively, the propagator may be generalised to three dimensions, for which
\begin{equation}
\hat{\mathbb{G}}_{xy}(\bm{k}) = -4 \frac{k_x^2 k_y^2}{k^4} - \frac{k_z^2}{k^2}.
\end{equation}
Expressing both the two- and three-dimensional kernels in the $xy$ plane, the real-space form asymptotically becomes
\begin{equation}
\mathbb{G}_{xy}(\bm{r}) \sim \frac{1}{r^{d}} \cos(4\theta),
\label{eq:eshelbyReal}
\end{equation}
consistent with the quadrupolar Eshelby field for an elliptical inclusion \cite{eshelby_determination_1997}. Consequently, in a system driven in simple shear, the stress perturbation generated by a local plastic event exhibits an eight-lobed structure: four lobes in which the shear stress increases and four in which it decreases (see \cref{fig:Propogator}).

The Eshelby stress propagator defined in \cref{eq:StressPropogator} is specified only up to an additive constant because of the singular mode at $\bm{k}=\bm{0}$. In practice, this constant is fixed by the boundary conditions. In a strain-controlled protocol, the total elastic strain is conserved. The zero mode is therefore chosen such that $\hat{\mathbb{G}}_{xy}(\bm{0}) = -1$, ensuring that the total stress relaxed by plastic events equals
\[
    \DIFaddbegin \DIFadd{2}\DIFaddend \mu\int \epsilon_{xy}^{\mathrm{pl}}(\bm{r}^\prime) \, d\DIFaddbegin \DIFadd{^2}\DIFaddend \bm{r}^\prime.
\]

By contrast, under stress-controlled conditions, we set $\hat{\mathbb{G}}_{xy}(\bm{0}) = 0$. In this case, no net stress is removed from the system; instead, the stress released by plastic deformation is compensated by elastic loading of the surrounding material \cite{nicolas_deformation_2018,martens_spontaneous_2012,jocteur_protocol_2025,lin_scaling_2014}.

The form of the stress propagator given in \cref{eq:StressPropogator} is defined for an infinite medium. However, it remains a good approximation for events sufficiently far from a solid wall or inclusion and for a large enough periodic system, due to its rapid $1/r^{D}$ decay. For a more rigorous treatment, the method of images may be used to simulate solid walls with a non-slip boundary condition to derive an equivalent Green function to the Oseen tensor (\cref{eq:Oseen}) \cite{blake_note_1971}, and also the stress propagator $\mathbb{G}_{xy}$ \cite{picard_elastic_2004}. However, for a two-dimensional system, it is often easier numerically to use an equivalent stream-function formalism as described in Ref.~\cite{acheson_elementary_1990}.

Although the explicit example above is written for plastic eigenstrains, the same Green-function viewpoint underpins stress relaxation more broadly throughout this thesis, including redistribution driven by bond-level rearrangements and failure in network models.

\section{Soft Networks} \label{sec:Networks}

The continuum framework outlined above provides a basic description of stress transmission and flow in soft materials, and forms the foundation of the work considered in \cref{sec:Fatigue}. However, many of the systems of interest in this thesis derive their mechanical response from an underlying filamentous architecture that is not fully captured at this level. We will therefore adopt a mesoscale description in terms of soft networks, which retains the essential microscopic detail of how filaments are organised and connected, while remaining sufficiently coarse-grained to describe macroscopic behaviour. In such networks, the connectivity and arrangement of the constituent filaments play a central role in determining the emergent mechanics. The failure mechanisms of these networks are examined in detail in \cref{sec:Ultradelayed,sec:DoubleNetworks}.

\subsection{Types of Network} \label{sec:TypesOfNetworks}
With the term \quotes{soft networks}, we refer to materials composed of interconnected filaments. While this shared structure suggests some network-level behaviour may be universal, key features remain material-specific. To make clear where our models are expected to apply\DIFaddbegin \DIFadd{, }\DIFaddend we therefore distinguish between generic precursors to failure and material-dependent mechanisms. Below, we provide a brief description of common classes of soft networks.

\paragraph{Elastomers:}
Composed of long, flexible polymer chains cross-linked into a network, elastomers are perhaps the best-known class of soft networks and are commonly found in thermosetting rubbers and plastics. Existing above the glass transition temperature $T_g$ of polymer networks, their mechanical behaviour is largely governed by entropic elasticity: deformation reduces the configurational entropy of polymer chains, generating a restoring force \cite{treloarElasticityRelatedProperties1973}. As a result, elastomers typically exhibit large reversible deformations until bond or cross-linker (i.e., molecules that connect distinct polymer chains together) failure occurs, often with pronounced strain stiffening before failure \cite{kimMeasurementNonlinearMechanical2011}.

\paragraph{Hydrogels:}
As the name implies, hydrogels are crosslinked polymer gels swollen with water, which can constitute up to around $90\%$ of the material volume \cite{gong_why_2010,tanaka_determination_2005,sun_highly_2012}. This high solvent content leads to a lower polymer density than in elastomers, resulting in materials that are typically softer and mechanically weaker. As a result of the low polymer volume fraction, interactions below the network scale are diminished, and the macroscopic properties become strongly dependent on the network structure. The presence of water also introduces additional physical phenomena, including swelling of the polymer chains and a coupling between network deformation and fluid flow, known as poroelasticity \cite{biotTheoryFiniteDeformations1972}. Owing to their high solvent content and mechanical compliance, hydrogels are widely explored for biomedical applications, particularly as artificial tendons, ligaments, and cartilage \cite{yasudaNovelDoubleNetworkHydrogel2009,yinSelfHealingHydrogelsSynthesis2023,hoHydrogelsPropertiesApplications2022}.

\paragraph{Fibre networks:}
Biological fibre networks are typically composed of semi-flexible filaments such as actin and collagen. These networks possess a characteristic length scale much larger than that of elastomers and hydrogels, with mesh sizes on the order of $1 \mu m$ to $10 \mu m$ \cite{mackintosh_elasticity_1995,Arevalo_Stress_Heterogeneities_2015}.
Much like hydrogels, these networks are highly porous, with a large fraction of their volume being occupied by solvent \cite{Jansen_Role_Collagen_2018}; however, due to their larger mesh size, they do not exhibit significant poroelastic effects \cite{Arevalo_Stress_Heterogeneities_2015}. The nature of crosslinking in such networks is also diverse, ranging from flexible linker proteins to branching and interwoven fibre bundles, leading to highly heterogeneous and mechanically complex interactions between filaments. Owing to their low fibre density, these materials typically exhibit a very small linear stiffness, making conventional extensional testing challenging in practice. As a result, their elastic and fracture behaviour is most commonly characterised under shear rather than through standard extensional protocols. Such fibre networks arise predominantly in biological contexts, where they provide structural support both within and outside of cells \cite{burlaMechanicalResilienceActive2019}.

\paragraph{Colloidal networks:}
Colloidal networks consist of suspended particles that aggregate via short-range attractive interactions to form a space-spanning, percolated structure \cite{Mewis2011,nabizadeh_network_2024,Chen2020,zhang_correlated_2019}. These systems are typically out of equilibrium and evolve over time, exhibiting aging and restructuring as particle bonds rearrange \cite{viasnoffHowAreColloidal2003,renouYieldingProcessesColloidal2010,manley_gravitational_2005}. Mechanically, they behave as soft solids that can sustain stress at rest but yield and flow once interparticle bonds are disrupted, with properties strongly dependent on particle concentration and interaction strength \cite{Mewis2011}.

\paragraph{Macroscopic networks:}
Macroscopic network materials are systems in which the network structure can be directly observed with the naked eye. Common examples include textiles, fishing nets, and architected metamaterials. Such well-defined networks have also been used to study fracture and mechanical response \cite{driscollRoleRigidityControlling2016}. Their properties can vary widely depending on the material from which the network is made; however, the network architecture typically leads to behaviour that is softer and more extensible than that of the corresponding material in its bulk form. For example, a fishing net made of nylon fibres can stretch to several times its original length before breaking, even though nylon itself is relatively stiff and brittle.

To describe their mechanical response more precisely, it is necessary to consider the properties of the individual filaments that make up the network.

\subsection{Semi-Flexible Filaments} \label{sec:SemiFlexible}

Filaments can be broadly classified according to the ratio of their persistence length \(l_p\) to a relevant contour length \(l_c\), which in networks is typically taken as the mean distance between crosslinks. The persistence length quantifies filament stiffness and sets the length scale over which the filament orientation remains correlated. Intuitively, $l_p$ can be considered as the length scale at which thermal fluctuations of order $k_BT$ have a large enough effect to perturb the filament direction. More formally, it is defined through the decay of tangent--tangent correlations, $\langle \bm{t}(s) \cdot \bm{t} (0) \rangle = e^{-s/l_p}$, where $\bm{t}(s)$ is the tangent vector at a position $s$ along the filament backbone \cite{doi_theory_1988}. 

For short or stiff segments where $l_p \gg l_c$, the filaments act like a rigid rod with its mechanical response governed by enthalpic stretching of molecular bonds. For segments where persistence length is much smaller than the contour length $l_p \ll l_c$, thermal fluctuations are sufficient to induce significant bending, and the filament adopts a naturally coiled configuration. In this flexible \DIFdelbegin \DIFdel{region}\DIFdelend \DIFaddbegin \DIFadd{regime}\DIFaddend , the mechanical response is primarily entropic\DIFdelbegin \DIFdel{: small extensions reduce }\DIFdelend \DIFaddbegin \DIFadd{. Stretching straightens the coiled filament and reduces }\DIFaddend the number of \DIFdelbegin \DIFdel{accessible conformations, generating an effective }\DIFdelend \DIFaddbegin \DIFadd{configurations it can adopt. The filament therefore tends to return to its more disordered, coiled state, producing a }\DIFaddend restoring force. At larger extensions, the response again becomes governed by the enthalpic stretching of molecular bonds, as in the stiff case. 

Semi-flexible filaments, therefore, occupy an intermediate regime between these two limits, with a persistence length $l_p$ comparable to their contour length $l_c$. In this case, the filaments are stiff enough to avoid collapsing into a random coil but large enough that thermally induced bending remains relevant. 

The standard model for semi-flexible filaments is the extensible worm-like chain (EWLC), in which each filament is treated as a homogeneous cylindrical elastic rod of radius $r$ and Young's modulus $E$ \cite{doi_theory_1988}. The stretching rigidity of the rod is given by the product of its Young's modulus and cross-sectional area,
\[
    \mu = E \pi r^2,
\]
while the bending rigidity is determined by its Young's modulus and moment of inertia
\[
    \kappa = E I = E \pi r^4 / 4.
\]
The Hamiltonian of a single filament can therefore be written as the sum of two integrals, the stretching and bending contributions,
\begin{equation*}
    \mathcal{H} = \frac{\mu}{2} \int u(s)^2 \, ds + \frac{\kappa}{2} \int \left| \frac{\partial \vec{t}}{\partial s} \right|^2 \, ds
\end{equation*}
where \(u(s)\) denotes the local extensional strain and \(\bm{t}(s)\) the tangent vector at a position $s$ along the filament backbone. The second term penalises curvature and captures the bending energy of the filament, while torsional contributions are neglected.

\subsubsection*{Numerical Discretisation of Filaments}

In this thesis, we numerically discretise each filament into a series of segments, following Ref.~\cite{mackintosh_elasticity_1995}. The Hamiltonian of each filament can then be expressed as
\begin{equation}
    \mathcal{H} = \frac{\mu}{2} \sum_{ij} \frac{( l_{ij} - l_{ij,0} )^2}{l_{ij,0}} + \frac{\kappa}{2} \sum_{ijk} \frac{( \theta_{ijk} - \theta_{ijk,0} )^2}{l_{ijk,0}}.
    \label{eq:NetworkHamiltonian}
\end{equation}
The first sum runs over all pairs $ij$ along the filament backbone, while the second runs over all triplets $ijk$. Here $l_{ij}$ denotes the instantaneous segment length between nodes $i$ and $j$, and $\theta_{ijk}$ the instantaneous bond angle formed by nodes $i$, $j$, and $k$. Quantities with subscript $0$ denote their natural (not necessarily stress-free) values in the rest configuration. We further define $l_{ijk,0} = (l_{ij,0} + l_{jk,0})/2$. Multiple filaments can then be combined into a single network structure by connecting filaments at shared nodes, with the set of triplets $ijk$ updated accordingly. 

In this work, we restrict attention to central-force networks and therefore neglect bending rigidity, setting $\kappa = 0$. This approximation is justified in regimes where the mechanical response is dominated by the uniaxial stretching of filament segments rather than bending deformations. At the level of individual filaments, bending modes are generically much softer than stretching modes due to geometric scaling, with bending stiffness scaling as $EI \sim r^4$ and stretching stiffness as $EA \sim r^2$, such that slender filaments are comparatively easy to bend but resist extension. At the network level, in sufficiently connected or highly strained networks, these soft bending modes are progressively suppressed and the deformation becomes stretching-dominated, in which case the network elasticity is well captured by central-force interactions alone \cite{mackintosh_elasticity_1995,broedersz_criticality_2011,shivers_scaling_2019}. In this limit, the dominant contribution to the energy arises from changes in filament length, while angular distortions play a subleading role and may be neglected without qualitatively altering the mechanical response. Accordingly, we adopt a central-force description to isolate the role of network connectivity and axial elasticity in governing the behaviour studied in this thesis.

\subsection{Dynamics} \label{sec:NetworkDynamics}
To resolve the network dynamics, we assume overdamped motion in which elastic forces within the network are balanced at all times by a dissipative drag force arising from the surrounding solvent. The elastic force acting on node $i$, arising from filament stretching, is obtained from the Hamiltonian as
\begin{equation}
    \bm{F}_{n,i}=-\frac{\partial \mathcal{H}}{\partial \bm{r}_i}.
\end{equation}

This is balanced by a drag force from the solvent, $\bm{F}_{d,i}$, such that
\begin{equation}
    \bm{F}_{d,i} + \bm{F}_{n,i}=\bm{0}.
    \label{eq:ForceBalance}
\end{equation}

We neglect long-range hydrodynamic interactions in the drag term due to hydrodynamic screening arising from the network structure. In such networks, the fluid momentum is dissipated through friction with the network, leading to screening of hydrodynamic interactions beyond the mesh size \cite{shankar_theory_2002,visscher_dynamic_1994}. Instead, we consider a background Newtonian fluid with velocity field $\bm{v}_f(\bm{r}_i)$, which exerts a drag force on node $i$,
\begin{equation}
    \bm{F}_{d,i} = -\zeta\left(\dot{\bm{r}}_i - \bm{v}_f(\bm{r}_i)\right),
    \label{eq:Drag}
\end{equation}
where $\zeta$ is the drag coefficient. This form of drag is not Galilean invariant, because it depends on the absolute velocity of each node rather than relative velocities. However, it is widely used in the literature \cite{gittes_dynamic_1998,broedersz_modeling_2014,ikeda_unified_2012,puosi_time-dependent_2014,durian_foam_1995,salerno_avalanches_2012,liuDrivingRateDependence2016}, and yields results comparable to the more accurate form used in dissipative particle dynamics, in which dissipation is formulated in terms of pairwise relative velocities \cite{groot_dissipative_1997}. 

Deformations are imposed affinely at the level of the background flow, as in the affine solvent model \cite{durian_foam_1995,lerner_simulations_2013}, such that the fluid velocity field $\bm{v}_f$ follows the imposed macroscopic strain. The network nodes, however, are free to move non-affinely in response to the balance of elastic and drag forces.

\subsubsection{Integration Method} \label{sec:EulerHeun}
Rearranging \cref{eq:ForceBalance} and \cref{eq:Drag}, we obtain the equation of motion for the velocity of the $i$-th node,
\begin{equation*}
    \dot{\bm{r}}_i = \frac{1}{\zeta}\bm{F}_{n,i} + \bm{v}_f(\bm{r}_i).
\end{equation*}

To evolve the system numerically, we employ a splitting scheme in which the dynamics are decomposed into distinct contributions over each time step $\delta t$. Specifically, the displacement of each node is updated sequentially by (i) advection due to the background solvent velocity $\bm{v}_f$, and (ii) motion driven by elastic forces $\bm{F}_n$.

First, each node position, together with the system unit cell (see \cref{sec:Lees-Edwards} for implementation details), is advected affinely to account for the imposed flow field $\bm{v}_f$. The resulting advected position, $\bm{r}_i'$, is given by
\begin{equation*}
    \bm{r}_i' = \bm{r}_i(t) + \delta t \bm{v}_f(\bm{r}_i(t)).
\end{equation*}

In this thesis, deformations are either imposed instantaneously, as in \cref{sec:Ultradelayed}, or applied quasi-statically, as in \cref{sec:DoubleNetworks}. \DIFdelbegin \DIFdel{Under these conditions, the dynamics are effectively overdamped, and }\DIFdelend \DIFaddbegin \DIFadd{The subsequent dynamics are assumed to be overdamped, so that the node motion is governed by }\DIFaddend a first-order \DIFdelbegin \DIFdel{integration scheme is sufficient to capture the solvent drag and resulting node motion}\DIFdelend \DIFaddbegin \DIFadd{equation in time and no inertial dynamics need to be resolved. The numerical scheme must nevertheless resolve the transient relaxation and the evolution of fracture events accurately. We therefore employ an adaptive Euler--Heun method, in which a first-order predictor and second-order corrected solution provide an estimate of the local discretisation error and allow the timestep to be reduced when the dynamics vary rapidly. A maximum timestep is additionally imposed to maintain sufficient temporal resolution of fracture events}\DIFaddend .

In the second substep, we employ an adaptive Euler--Heun method (also known as a Runge–\DIFdelbegin \DIFdel{Kutta }\DIFdelend \DIFaddbegin \DIFadd{-Kutta }\DIFaddend 1(2) scheme) inspired by Ref.~\cite{sammullerAdaptiveBrownianDynamics2021}. This method uses an embedded Runge–\DIFdelbegin \DIFdel{Kutta }\DIFdelend \DIFaddbegin \DIFadd{-Kutta }\DIFaddend pair, in which two solutions of different order are computed from the same intermediate evaluations to estimate the local discretisation error and adapt the time step $\delta t$. 

Starting from the advected position $\bm{r}^\prime_i$, we first compute the predictor estimate $\tilde{\bm{r}}_i$,
\begin{equation*}
    \tilde{\bm{r}}_i = \bm{r}^\prime_i + \frac{\delta t}{\zeta} \bm{F}_{n,i}(\bm{r}^\prime).
\end{equation*}
The forces are then re-evaluated at the predicted positions to obtain the corrected position at time $t+\delta t$:
\begin{equation*}
    \bm{r}_i(t + \delta t) = \bm{r}^\prime_i + \frac{\delta t}{2\zeta} \left[ \bm{F}_{n,i}( \bm{r}^\prime_i ) + \bm{F}_{n,i}(\tilde{\bm{r}_i}) \right].
\end{equation*}
The particle-wise error is defined as the difference between the first- and second-order updates,
\begin{equation*}
    E_i = \left\| \Delta\tilde{\bm{r}}_i - \Delta\bm{r}_i \right\| 
    = \frac{\delta t}{2\zeta}\left\| \bm{F}_{n,i}(\tilde{\bm{r}}_i) - \bm{F}_{n,i}(\bm{r}_i') \right\|_2,
\end{equation*}
where $\Delta\tilde{\bm{r}}_i = \tilde{\bm{r}}_i - \bm{r}_i'$ and $\Delta\bm{r}_i = \bm{r}_i(t+\delta t) - \bm{r}_i'$. We introduce an error tolerance $\tau$, which constrains the maximum allowed absolute error. The global error estimate is then
\begin{equation*}
    \mathcal{E} = \left\| \frac{E_i}{\tau} \right\|_\infty,
\end{equation*}
where $\|\cdot\|_\infty$ denotes the maximum norm.

Following Ref.~\cite{sammullerAdaptiveBrownianDynamics2021}, we \DIFdelbegin \DIFdel{define the timestep scaling factor
}\DIFdelend \DIFaddbegin \DIFadd{adopt the conservative timestep scaling
}\DIFaddend \begin{equation*}
    q = \left( \frac{1}{2\mathcal{E}} \right)^2\DIFdelbegin \DIFdel{,
}\DIFdelend \DIFaddbegin \DIFadd{.
}\DIFaddend \end{equation*}
\DIFdelbegin \DIFdel{which is restricted }\DIFdelend \DIFaddbegin \DIFadd{This choice includes a safety factor of two and causes the timestep to decrease rapidly when the estimated error exceeds the prescribed tolerance, while allowing only gradual increases when the error is small. We additionally restrict $q$ }\DIFaddend to the interval $[0.001, 1.2]$\DIFaddbegin \DIFadd{, following the same reference, to prevent excessively large changes in timestep between successive updates}\DIFaddend . Two cases are then distinguished:
\begin{enumerate}
    \item $q \geq 1$: accept $\bm{r}_i(t+\delta t)$.
    \item $q < 1$: reject the step and restore the configuration to time $t$\DIFdelbegin \DIFdel{before retrying with a smaller timestep}\DIFdelend .
\end{enumerate}
In both cases, the time step is updated via $\delta t \gets q \delta t$. \DIFaddbegin \DIFadd{For a rejected step, $q < 1$, so the integration is retried using the resulting smaller timestep.
}

\DIFaddend We additionally impose a maximum time step \DIFdelbegin \DIFdel{$\delta t_{\max}$ to ensure adequate temporal resolution of fracture events}\DIFdelend \DIFaddbegin \DIFadd{$\delta t_{\max} = 0.1$ in all simulations, which limits the timestep when the estimated local error is small. The numerical accuracy of the integration is controlled primarily through the error tolerance. The values of $\tau$ used throughout this thesis was selected from convergence tests, presented in }\cref{app:stepStrainConvergence} \DIFadd{, in which the simulated dynamics were compared for progressively smaller tolerances and first order fixed timestep simulation and verified for each scenario separately }\DIFaddend .  

\subsection{Rigidity} \label{sec:Rigidity}

One of the fundamental properties of soft network materials, absent from conventional continuum descriptions, is the origin of rigidity. In continuum elasticity, rigidity is assumed a priori through constitutive relations, and materials are taken to support stress under infinitesimal deformation. By contrast, in soft networks, mechanical stability emerges from a balance between constraints and degrees of freedom at the filament level. Consequently, a network's connectivity, topology, and filament mechanics collectively determine whether it can support stress. In the section that follows, we distinguish between two complementary origins of rigidity in soft networks: constraint-based rigidity, governed primarily by network connectivity, and strain-induced rigidity, which emerges under applied deformation through non-linear dynamics in the network.

We begin by considering rigidity arising from the number of constraints in the system. To do so, we examine idealised frameworks in which each connection has a fixed length $l_{ij,0}$ and can freely hinge at its nodes, with no bending penalty. This corresponds to the limit of \cref{eq:NetworkHamiltonian} with $\mu \to \infty$, enforcing $l_{ij} = l_{ij,0}$, and $\kappa = 0$.

Defining $\bm{r}_i(t)$ as the position of the $i$th node of the system, the network configuration is specified by the set of constraints
\[|\bm{r}_i(t) - \bm{r}_j(t)|^2 = l_{ij,0}^2\] 
imposed for all connected pairs $ij$. By differentiating with respect to $t$, this yields a set of compatibility conditions on the nodal velocities $\dot{\bm{r}}$,
\begin{equation} \label{eq:FrameworkRigidity}
    (\bm{\dot{r}}_i(t) - \bm{\dot{r}}_j(t))\cdot(\bm{r}_i(t) - \bm{r}_j(t))=0
\end{equation}
for every connected pair $ij$.

If there exists a non-trivial set of velocities $\dot{\bm{r}}_i$ satisfying \cref{eq:FrameworkRigidity}, the framework admits an infinitesimal deformation (a floppy mode), because at least one node can be moved freely, and the structure is not first-order rigid. Conversely, the framework is defined to be first-order rigid if the only solutions to \cref{eq:FrameworkRigidity} are the trivial rigid-body motions. In simple terms, the network is first-order rigid if the constraints allow for a single unique configuration (up to rigid-body translations and rotations).

\DIFaddbegin \begin{figure}[t]
    \centering
    \includegraphics[width=0.9\linewidth]{chapters/2-Background/Figures/maxwell_counting_diagram.pdf}
    \caption[Illustration of Maxwell counting for a connected central-force network.]{\DIFaddFL{Illustration of Maxwell counting for a connected central-force network with $N=4$ nodes in two dimensions. Removing two uniform translations and one rigid-body rotation leaves $2N-3=5$ non-trivial degrees of freedom, while each independent central-force bond supplies one constraint. (}\textbf{\DIFaddFL{a}}\DIFaddFL{) With $N_{\mathrm{c}}=4$, the framework is sub-isostatic; the grey dashed configuration and blue arrows show its floppy shear mode. (}\textbf{\DIFaddFL{b}}\DIFaddFL{) Adding one diagonal gives $N_{\mathrm{c}}=5$, producing an isostatic, marginally rigid framework. (}\textbf{\DIFaddFL{c}}\DIFaddFL{) Adding the second diagonal gives $N_{\mathrm{c}}=6$, producing a hyperstatic framework. The red bond is redundant: it does not remove another motion but instead produces a state of self-stress. Apart from such redundancies, the count assumes independent constraints.}}
    \label{fig:MaxwellCounting}
\end{figure}

\DIFaddend This idea is most simply understood through a counting argument due to Maxwell \cite{maxwell_calculation_2011}, in which rigidity is determined by comparing the number of constraints to the number of degrees of freedom. For a $d$-dimensional system of $N$ nodes, there are $dN$ translational degrees of freedom, corresponding to the independent displacements of each node. However, not all of these contribute to internal deformation: $d$ degrees of freedom correspond to uniform translations of the entire system, and $d(d-1)/2$ correspond to rigid-body rotations, which leave all distances between nodes unchanged. Subtracting these trivial modes, the number of non-deforming degrees of freedom is $d(d+1)/2$, leaving $dN - d(d+1)/2$ non-trivial degrees of freedom. The total number of constraints $N_c$ needed for the system to be first-order rigid must therefore be at least equal to the number of non-trivial degrees of freedom
\[
    N_c \geq d N - d ( d+1 ) / 2.
\]
Assuming the network is homogeneous, we define the average connectivity as $z=2N_c/N$, where a bond between two nodes contributes one constraint. In the thermodynamic limit $N \to \infty$, this yields the condition $z \ge 2d$ for rigidity. The critical point $z_c = 2d$ is referred to as the isostatic point. More generally, networks are classified according to their connectivity relative to this point: sub-isostatic networks possess fewer constraints than degrees of freedom and admit floppy modes, isostatic networks are marginally rigid, and hyperstatic networks contain redundant constraints and can support states of self-stress \cite{LernerCoordination2024,Mao2018,ellenbroek_rigidity_2015,wyart_elasticity_2008}. \DIFaddbegin \DIFadd{We illustrate this counting argument in }\cref{fig:MaxwellCounting} \DIFadd{for a simple two-dimensional network of four nodes.
}\DIFaddend 

% This idea is more simply understood through the work of Maxwell \cite{maxwell_calculation_2011}, who initially pointed out that this definition of rigidity is directly related to the number of constraints in the system. For a $d$-dimensional system of $N$ nodes, the condition for the number of constraints $N_c$ for the system to be rigid is 
% \[
%     N_c \geq d N - d ( d+1 ) / 2.
% \]
% Assuming the network is homogenous, we can define an average connectivity of the network $z=2N_c/N=2N_B/N$, where $N_B$ is the number of bonds in the system. By letting $N\to\infty$, then for a homogenous network to be rigid, $z\geq 2d$. The critical point $z_c=2d$ is often referred to as the isostatic point. More generally, networks are classified according to their connectivity relative to the isostatic point: sub-isostatic networks possess fewer constraints than degrees of freedom and admit floppy modes, isostatic networks are marginally rigid, and hyperstatic (over-constrained) networks contain redundant constraints and can support states of self-stress \cite{LernerCoordination2024}.

This first-order rigidity criterion holds only when the network forms a single rigid cluster and the bond constraints are independent. In general, however, these conditions need not be satisfied, particularly in over-constrained systems. In such networks, states of self-stress can arise, defined as assignments of internal forces that satisfy force balance in the absence of external loading. The presence of these redundant constraints means that the number of independent constraints is smaller than the total number of bonds, and consequently the relationship between degrees of freedom and constraints cannot, in general, be inferred from simple counting arguments alone \cite{calladine_buckminster_1978}.

Often, this definition of first-order rigidity is overly restrictive and does not fully capture the behaviour of heterogeneous networks, for which the average connectivity $z$ alone is insufficient to determine rigidity. This limitation is particularly relevant for the disordered networks considered in later chapters.

As a simple example, consider attaching a single freely hanging node to an otherwise rigid network. This addition formally violates the strict definition of first-order rigidity, because the new node introduces a floppy mode. However, it does not alter the global load-bearing capacity of the original rigid cluster, which remains able to support stress. For such systems, it is necessary to consider not only the number of constraints but also their spatial organisation to determine global rigidity. In particular, sparsely connected networks below the isostatic point often consist of small rigid clusters embedded within floppy, under-constrained regions. In strongly heterogeneous networks, even systems with average connectivity exceeding the Maxwell threshold $z_c = 2d$ can reorganise such that global deformation occurs without requiring the extension of any bonds. Consequently, the ability of a network to sustain stress is governed by the existence of a system-spanning rigid cluster capable of supporting states of self-stress. For this reason, the onset of global rigidity is commonly described as a rigidity percolation transition \cite{zhang_correlated_2019}. This should be distinguished from geometrical percolation, which identifies the formation of a purely connected path between opposite boundaries of the network and does not, in general, imply mechanical rigidity \cite{kendall_percolation_2009,broadbent_percolation_1957}. In two dimensions, this problem has been well studied \cite{ellenbroek_rigidity_2015} through algorithms such as the pebble game \cite{jacobs_algorithm_1997, jacobs_generic_1995}. Extending these definitions to show rigidity in dimensions higher than two has been proven to be computationally intractable using current methods \cite{abbott_generalizations_2008}.

While the definitions above provide a framework for characterising structural, or first-order, rigidity, they do not account for the fact that the networks considered here consist of deformable springs rather than fixed-length linkages. It is well known that under-constrained (sub-isostatic) spring networks can undergo a rigidity transition at coordination numbers less than $z_c=2d$ \cite{feng_percolation_1984, sharma_strain-controlled_2016, wyart_elasticity_2008, feng_nonlinear_2016, vermeulen_geometry_2017, rens_micromechanical_2018}.

The simplest illustration of such a transition is a string fixed at both ends. When the separation between the fixed ends is smaller than the contour length of the string, the string carries no tension and admits a floppy mode, offering no resistance to transverse deformation. However, once the fixed points are separated sufficiently that the string comes under tension, the floppy mode is removed, and the string can resist transverse deflection, thereby appearing mechanically rigid \cite{biot_mechanics_1965}. Although the connectivity remains fixed at $z = 2$, rigidity emerges through the development of internal tension rather than changes in network topology. The disordered networks studied in this thesis also undergo this transition. We therefore define a critical shear strain $\gamma_c$ as the strain at which a sub-isostatic network first develops a finite shear modulus, marking the onset of strain-induced stiffening.

This transition is often referred to as second-order rigidity \cite{connelly_second-order_1996}, where non-structural changes can cause the system to act as if it is in a rigid configuration. Mathematically, second-order rigidity is obtained by differentiating \cref{eq:FrameworkRigidity} with respect to time, yielding
\begin{equation*}
(\bm{r}_i(t) - \bm{r}_j(t))\cdot(\bm{\ddot{r}}_i(t) - \bm{\ddot{r}}_j(t)) + (\bm{\dot{r}}_i(t) - \bm{\dot{r}}_j(t))\cdot(\bm{\dot{r}}_i(t) - \bm{\dot{r}}_j(t)) = 0,
\end{equation*}
where $\ddot{\bm{r}}_i(t)$ is the acceleration of a node $i$. This criterion allows rigidity to arise even when non-trivial solutions to \cref{eq:FrameworkRigidity} exist, through the stabilising effect of internal force balance in the material. It therefore represents a weaker notion of rigidity, but one that still permits the network to support states of self-stress. Such configurations are \DIFdelbegin \DIFdel{significantly more difficult to predict \mbox{%DIFAUXCMD
\cite{holmes-cerfon_almost-rigidity_2021}}\hskip0pt%DIFAUXCMD
, because they emerge from nonlinear geometric and mechanical effects that have been extensively studied in the context of frameworks andtensegrities \mbox{%DIFAUXCMD
\cite{connelly_frameworks_2022,connelly_second-order_1996,calladine_first-order_1991}}\hskip0pt%DIFAUXCMD
}\DIFdelend \DIFaddbegin \DIFadd{difficult to identify because their stability depends on higher-order geometric effects and, where present, internal stresses \mbox{%DIFAUXCMD
\cite{holmes-cerfon_almost-rigidity_2021,connelly_frameworks_2022,connelly_second-order_1996,calladine_first-order_1991}}\hskip0pt%DIFAUXCMD
}\DIFaddend . Indeed, determining whether a general planar spring network is rigid is NP-hard \cite{abbott_generalizations_2008}, and consequently, no simple, universally applicable criterion exists for establishing rigidity.

For systems composed of semi-flexible filaments with extensible bonds, both first- and second-order rigidity provide useful proxies for mechanical stability \cite{van_hecke_jamming_2010}. However, these kinematic criteria do not fully capture the complexity of the energy landscape explored by such networks. A more recent and intuitive perspective is provided by energetic rigidity \cite{damavandi_energetic_2022}, which characterises rigidity in terms of the local structure of the elastic energy landscape. In this framework, a configuration is defined to be rigid if any infinitesimal deformation increases the elastic energy, such that the network resides at a local minimum of the energy.

This viewpoint has yielded important insights into the mechanics of spring networks, vertex models, and jammed packings of soft particles \cite{damavandi_energetic_2022-1}, particularly in the presence of pre-stress, i.e. internal stresses that exist in the absence of external loading. Nevertheless, even the energetic criterion does not provide a fully predictive framework for determining rigidity in highly disordered networks.

\subsection{Fracture in networks}\label{sec:Fracture}
In \cref{sec:Plasticity}, the idea of plasticity and yielding in elastoplastic materials is discussed and is directly used in \cref{sec:Fatigue}. Throughout this thesis, yielding and fracture are treated as related but distinct mechanisms of irreversible failure. Yielding refers to the onset of irreversible deformation and loss of purely elastic response, whereas fracture additionally implies material separation or a substantial loss of structural connectivity and load-bearing capacity. In network materials, this distinction is particularly important because macroscopic failure typically emerges from the progressive scission of individual connections within a highly heterogeneous structure.

Generally, material failure is a complex problem, stemming from the large separation in length scales between the failure event (e.g. the rupture of a single bond) and the global failure of the material through the evolution of macroscopic cracks. At the microscopic level, fracture begins with the scission of individual connections between nodes that collectively grow into macroscopic failures \cite{Creton50th2017,ju_role_nodate,sanner_less_2025,higuchi_fracture_2018,bonnDelayedFractureInhomogeneous1998,teece_gels_2014}. However, there is generally no one-to-one mapping between the failure of a single bond and the resulting macroscopic fracture. Moreover, the range of time scales involved is also large: damage may accumulate slowly as the material is deformed, whereas macroscopic cracks can propagate rapidly across the system.

This challenge is further compounded by the inherent disorder present in real-world soft networks \cite{broedersz_modeling_2014,broedersz_criticality_2011,feng_nonlinear_2016}. Unlike ordered materials such as crystalline solids, which possess an inherent repeating unit cell and well-defined defect states that act as nucleation sites for fracture, disordered networks lack any obvious structural signature by which such critical defects can be identified. Furthermore, this inherent disorder and the large non-affine rearrangements which occur in the networks often lead to heterogeneity in the stress at the level of individual filaments (i.e. the load is not distributed evenly across the network) \cite{Liang2016,tauber_microscopic_2020,tauber_sharing_2021,broedersz_modeling_2014,storm_nonlinear_2005,bouzidNetworkTopologySoft2018}. Here networks redistribute stress through non-affine rearrangements, often delaying catastrophic failure \cite{shivers_scaling_2019,tauber_sharing_2021,ducrot_structure_2015, tominaga_molecular_2007}. Consequently, predicting when and where any individual bond may fail is challenging.

Experimentally, it has been shown that soft network materials, including elastomers, hydrogels, and fibre networks, can undergo very large deformations before visible macroscopic damage occurs \cite{gong_why_2010,Kim2021,ducrot_toughening_2014,Creton50th2017,aimeMicroscopicDynamicsFailure2018,storm_nonlinear_2005}. Experimental methods such as small-angle neutron scattering \cite{ducrot_structure_2015,tominaga_molecular_2007}, birefringence imaging \cite{gong_double-network_2003,zheng_how_2021}, and light scattering \cite{porr_probing_2025} have helped to reveal how stress is distributed across the network. 

Recent work on mesoscopic computational models, which provide a balance between structural details and computational cost, has characterised this stress distribution \cite{shivers_normal_2019,shivers_scaling_2019,LernerCoordination2024,storm_nonlinear_2005} as well as how fracture evolves \cite{tauber_microscopic_2020,tauber_sharing_2021,tauber_role_2020,dussi_athermal_2020}. In these models, it is common to assign a local rupture threshold to each bond. For central-force network models, this is typically implemented by rupturing bonds whose strain exceeds a critical threshold $\lambda$ \cite{tauber_role_2020,dussi_athermal_2020}. This scale, therefore, plays a central role in determining the system dynamics \cite{tauber_microscopic_2020}.

At higher breakage thresholds $\lambda$, bonds rupture only at higher extensions, typically beyond the critical strain $\gamma_c$, defined as the strain at which a sub-isostatic network first develops a finite shear modulus and begins to stiffen. At these larger deformations, stress concentrates along heterogeneous load-bearing paths within the network, often referred to as force chains, which become increasingly apparent at higher strains \cite{ruiz-franco_force_2022,Heussinger2007}. In this regime, the fracture response is governed less by the immediate proximity to the rigidity transition and more by the mature load-bearing backbone that develops at large strain \cite{Sarkar2022}. Moreover, because stress is already localised along the evolving force chains, increasing system size can promote a crossover towards more abrupt, crack-like failure, even in nominally sub-isostatic networks \cite{zhang_fiber_2017,driscollRoleRigidityControlling2016,dussi_athermal_2020,tauber_stretchy_2022}. By contrast, in the low-$\lambda$ limit, bond rupture occurs closer to $\gamma_c$, before these large-strain backbones are fully developed. Fracture is therefore more sensitive to the onset of strain-induced stiffening and to the network's proximity to the rigidity transition. 

In \cref{sec:Ultradelayed,sec:DoubleNetworks} we consider the failure of soft network materials in two different situations. In \cref{sec:Ultradelayed}, we consider delayed timescales that can be associated with the thermal failure of bonds \cite{tauber_role_2020,Bullerjahn2014,arruda_three-dimensional_1993,james_theory_1943,Skrzeszewska2010}. Here we explore the brittle-to-ductile transition seen extensively in networks under both increasing temperature and imposed shear deformation \cite{dussi_athermal_2020,ahmedBrittleDuctileTransition2014,ducrot_toughening_2014,tauber_role_2020}, characterising a material's response in both the high- and low-failure-threshold regimes. In \cref{sec:DoubleNetworks} the failure of bonds in networks composed of multiple components is considered. Here, we resolve how load sharing and stress redistribution between interpenetrating networks modify damage accumulation and can dramatically enhance toughness compared to single-network systems.

\subsection{Disordered Network Models} \label{sec:NetworkStructure}
Throughout the literature, a wide range of network geometries has been considered. The simplest are affine isotropic models, which approximate the network response by considering a single filament subjected to an imposed deformation and averaging over all initial filament orientations \cite{mackintosh_elasticity_1995, storm_nonlinear_2005}. While such models provide a useful baseline in the limits of high connectivity $z$ or large bending stiffness $\kappa$, they fail to capture key physics arising from non-affine rearrangements and orientational disorder in real-world networks. To incorporate these effects, it is necessary to explicitly specify a network topology that describes how bonds connect the nodes of the network. 

In this work, we consider two classes of networks: those derived from lattice structures and those constructed from the packing of bidisperse soft particles. These topologies are idealised and are not intended to represent a specific experimental system directly, but rather to provide controlled model networks that capture generic features of disordered filamentous materials.

\subsubsection{Lattice-based networks} \label{sec:LatticeBasedNetworks}

Lattice-based networks are among the simplest network topologies commonly used to study mechanical properties \cite{jacobs_generic_1995,feng_percolation_1984,shivers_scaling_2019,sharma_strain-controlled_2016}. Frequently studied examples include two-dimensional square \cite{mao_soft_2010}, triangular \cite{broedersz_criticality_2011,feng_percolation_1984}, and honeycomb \cite{rens_nonlinear_2016} lattices. Owing to their straightforward construction relative to more disordered network types, and their numerical robustness arising from the uniform lattice spacing $l_0$, these networks provide a convenient and well-controlled framework for analysing the mechanical response of fibre networks.

Such networks are constructed by placing nodes on a regular grid defined by a set of lattice vectors and subsequently connecting these nodes according to a prescribed set of rules. For example, in a two-dimensional square lattice, the node positions are generated by the lattice vectors $(1,0)$ and $(0,1)$, with nearest-neighbour nodes connected along the horizontal and vertical directions, subject to the chosen boundary conditions. Other simple lattices, such as the triangular lattice, can be generated analogously, while more complex lattices may be derived from these fundamental cases. For instance, a hexagonal lattice can be obtained from a triangular lattice by imposing the additional constraint that each node has connectivity $z = 3$. Owing to their regular structure, such lattices exhibit highly predictable behaviour, characterised by discrete values of connectivity and bond length.

\begin{figure}[t]
    \centering
    % \missingfigure[]{Lattice based network structure}
    \includegraphics[width=0.95\linewidth]{chapters/2-Background/Figures/lattice_full_diluted_generic.pdf}
    \caption[Periodic lattices: full, diluted, and generic diluted networks.]{
    Examples of periodic lattice structures. We begin with regular lattices (\textbf{left column}), from which disordered diluted networks are generated (\textbf{centre column}) by randomly retaining bonds with probability $p=0.2$. Additional disorder is then introduced (\textbf{right column}) by displacing each node from its lattice position by a random distance $d \in [0,\delta_{\max}/2]$ in a random direction, yielding a generic diluted network. Here, $\delta_{\max}=0.6$ is expressed in units of the lattice spacing (i.e. the initial bond length). Dangling ends and dangling clusters are removed in all diluted networks shown (as described in \cref{sec:PrunningNetwork}).
    }
    \label{fig:LatticeNetworks}
\end{figure}

While trivial to construct, such lattices lack the structural disorder typical of disordered networks, which exhibit no long-range order. This predictability is undesirable when modelling disordered networks, where the goal is to represent amorphous materials lacking long-range structure. Disorder is therefore introduced into these networks through two complementary approaches. First, a common method is bond dilution via the random removal of lattice bonds. In practice, each bond is retained with probability $p$ (and removed with probability $1-p$), where $p$ sets the desired bond occupation fraction \cite{broedersz_criticality_2011,head_distinct_2003}. This procedure is widely used when studying the role of geometric percolation in such models \cite{broedersz_criticality_2011,feng_percolation_1984,das_redundancy_2012}.

Although random bond removal may appear to introduce only limited structural disorder, it fundamentally alters the rigidity of the network (see \cref{sec:Rigidity}). In particular, the system exhibits a geometrical percolation transition at a critical bond occupation probability $p_c$, at which isolated clusters coalesce into a single system-spanning cluster \cite{jacobs_generic_1995,zhang_correlated_2019}. A related transition occurs at the isostatic point, where the network loses (or gains) first-order rigidity.

A further advantage of bond dilution in lattice-based networks is the disruption of long, system-spanning force chains \cite{ruiz-franco_force_2022} that arise from the perfect alignment of bonds in the underlying lattice geometry. Networks that do not undergo such pruning typically exhibit low-strain softening of the shear modulus $K$ \cite{licup_nonlinear_2016}, associated with the buckling of these extended fibres.

The effects of these system-spanning force chains can be mitigated by a second form of disorder commonly employed in lattice networks, in which the nodal positions are randomly perturbed and the corresponding bond rest lengths and angles are redefined \cite{jacobs_generic_1995,krushynska_multi-frequency_2025,sheinman_nonlinear_2012,head_predicting_2025}. Networks incorporating this positional disorder are commonly referred to as generic networks \cite{jacobs_generic_1995}. The magnitude of these perturbations is typically constrained to be less than half the minimum lattice spacing, $\delta_{\min}/2$, to prevent node overlap \cite{shivers_normal_2019}. We explore the physics of square-lattice-based networks in \cref{sec:SquareAppendix} in the context of composite network materials comprising two interconnected networks.

\subsubsection{Packing-Derived Networks} \label{sec:PackingBasedNetworks}

Packing-derived networks are computationally constructed from simulated packings of soft particles and are widely used as model disordered networks \cite{rens_rigidity_2019,wyart_elasticity_2008,lerner_unified_2012,shivers_normal_2019,shivers_scaling_2019,shivers_phase_2022}. In practice, these networks are generated numerically by minimising the total energy of a bidisperse ensemble of soft disks at densities above the jamming transition $\phi_J$ \cite{Hernjamming2003,BoyerJamming2011,ikeda_unified_2012}, after which the contact network of overlapping particles is extracted and used to define the network connectivity. This procedure provides an idealised disordered topology and does not correspond to a specific experimental realisation, but instead serves as a controlled model for studying generic features of fibre and elastic networks.

Each disk is initially placed at a random position within a square periodic domain of size $L_x=L_y=\sqrt{N}$, with radii drawn from a bidisperse distribution with values $R_0$ and $fR_0$ in equal proportion. The bidispersity parameter $f$ is chosen to suppress crystallisation and inhibit the formation of long-range positional order; a commonly adopted value, and the one used throughout this thesis, is $f=1.4$. Such bidisperse mixtures are well known to produce amorphous packings with minimal macroscopic ordering \cite{van_hecke_jamming_2010, chacko_slow_2019, desmond_random_2009, koezeMappingJammingTransition2016}. \DIFaddbegin \DIFadd{Bidispersity does not, however, make the packing completely structureless: the two discrete particle sizes introduce preferred microscopic length scales, and their size ratio and relative abundance influence short-range order, contact statistics, and the jamming point \mbox{%DIFAUXCMD
\cite{koezeMappingJammingTransition2016}}\hskip0pt%DIFAUXCMD
. A continuously polydisperse construction would replace these two length scales with a broader distribution of particle and bond lengths and could therefore alter the local coordination statistics and the resulting network topology, although the same contact-network construction procedure would apply.
}\DIFaddend 

The base radius $R_0$ is then rescaled so that the total particle area matches the prescribed packing fraction $\phi$ within the simulation box. While several forms of soft-disk interaction potentials have been proposed in the literature, we adopt the form used in Ref.~\cite{chacko_slow_2019}. In this model, overlapping disks interact via a purely repulsive pair potential given by
\begin{equation}
    U_{ij} \left(r_{ij}\right)
    = R_0^3\left(1-\frac{r_{ij}}{D_{ij}}\right)^2
    \Theta(D_{ij}-r_{ij}),
\end{equation}
where $r_{ij}$ is the centre-to-centre distance between disk $i$ and $j$, and $D_{ij}=R_i+R_j$ is the sum of their radii. The Heaviside function $\Theta(x)$ ensures that the interaction is non-zero only when particles overlap ($r_{ij}<D_{ij}$), and vanishes otherwise. 

Starting from this random initial condition, the system is relaxed via energy minimisation of the soft repulsive interactions until a mechanically stable packing is obtained. In Ref.~\cite{chacko_slow_2019}, this minimisation is performed by integrating an overdamped equation of motion equivalent to that presented in \cref{sec:NetworkDynamics}.

In the present work, however, our primary interest lies in the network structure obtained from the final minimised configuration rather than in the relaxation dynamics. We therefore generate packings using the Fast Inertial Relaxation Engine (F.I.R.E.) algorithm \DIFdelbegin \DIFdel{\mbox{%DIFAUXCMD
\cite{bitzek_structural_2006, guenole_assessment_2020}}\hskip0pt%DIFAUXCMD
}\DIFdelend \DIFaddbegin \DIFadd{\mbox{%DIFAUXCMD
\cite{bitzek_structural_2006,guenole_assessment_2020}}\hskip0pt%DIFAUXCMD
}\DIFaddend , which converges \DIFaddbegin \DIFadd{efficiently }\DIFaddend to local energy minima\DIFdelbegin \DIFdel{more efficiently. While this approach accelerates the energy minimisation, it does not capture physically accurate relaxation dynamics}\DIFdelend \DIFaddbegin \DIFadd{. The trajectory generated by the F.I.R.E. algorithm is a numerical optimisation path and should not be interpreted as physical time evolution. Nevertheless, the resulting configurations are mechanically stable, force-balanced minima of the chosen interaction energy, so their contact networks provide reasonable idealised models of disordered elastic networks}\DIFaddend .

\begin{figure}
\centering
\includegraphics[width=.32\textwidth]{chapters/2-Background/Figures/particles.pdf}\hfill
\includegraphics[width=.32\textwidth]{chapters/2-Background/Figures/Stacked.pdf}\hfill
\includegraphics[width=.32\textwidth]{chapters/2-Background/Figures/network.pdf}
\caption[Steps for generating a packing-derived network from 400 soft disks.]{
Steps for generating a packing-derived network from 400 soft disks at $\phi = 1$.\newline 
\textbf{Left:} A bidisperse disk packing with a radius ratio of $1:1.4$ in an energy-minimised state. \newline
\textbf{Middle:} Network bonds defined from the contact network of overlapping disks. \newline
\textbf{Right:} Disks removed, leaving only the bonds and nodes of the network.}
\label{fig:SingleNetworkGeneration}
\end{figure}

Network bonds are constructed from the contact network of an equilibrated packing \cite{lerner_simulations_2013,ellenbroek_non-affine_2009,mangal_topological_2023}. A network node is placed at the centre of each disk. For every overlapping pair of disks, the corresponding nodes are connected by a bond whose natural length is given by their centre-to-centre distance in the energy-minimised configuration. \DIFdelbegin %DIFDELCMD < 

%DIFDELCMD < %%%
\DIFdelend After all contacts have been processed, the disks themselves are removed, leaving a network composed solely of nodes and bonds. This construction procedure is illustrated in \cref{fig:SingleNetworkGeneration}. 

In this work, we take $\phi = 1 > \phi_J \approx 0.84$ unless stated otherwise \cite{chacko_slow_2019,ikeda_unified_2012}. For single networks, this corresponds to an average connectivity $z \approx 5.4$, which lies well above the isostatic threshold. In principle, networks with lower connectivity (while remaining above the isostatic threshold $z=4$) could be generated by preparing packings at $\phi$ closer to $\phi_J$. Instead, we obtain lower-connectivity networks by iteratively pruning bonds from an initially highly connected network constructed at $\phi=1$, following the procedure described below. This allows for tighter control of the network connectivity when studying systems with lower coordination.

These packing-derived networks provide a simple yet robust alternative to lattice-based networks. They exhibit emergent positional disorder, good numerical stability, and can be straightforwardly extended to three dimensions, as demonstrated in \cite{shivers_normal_2019, shivers_strain-controlled_2023, shivers_nonlinear_2020}. We therefore use them as the primary network structures considered in \cref{sec:Ultradelayed,sec:DoubleNetworks}.

\subsubsection{Pruning Networks} \label{sec:PrunningNetwork}
During network construction, the resulting graphs can contain disconnected clusters and dangling components (see \cref{fig:PruningNetworks}). These elements do not contribute to the static mechanical stress of the network and are therefore removed before simulation.

Dangling components, shown in panel (a) of \cref{fig:PruningNetworks}, consist of portions of the network attached to the bulk by a single bond. In graph theory, these correspond to biconnected components connected to the main network via a bridge edge. These are pruned by identifying bridge bonds and removing the associated dangling components. Disconnected clusters, illustrated in panel (b) of \cref{fig:PruningNetworks}, are eliminated by identifying all connected components of the network and retaining only the largest component, which corresponds to the bulk network. To perform this analysis efficiently, we represent the network as an undirected graph, with bonds as edges, and use the \texttt{NetworkX} Python package \cite{SciPyProceedings_11} to identify connected components, biconnected components, and bridges.

\begin{figure}
    \centering
    \begin{subfigure}{0.45\linewidth}
        \centering
        \includegraphics[width=\linewidth]{chapters/2-Background/Figures/network_dangling_motifs_400.pdf}
        \textbf{(a)}
    \end{subfigure}
    \begin{subfigure}{0.45\linewidth}
        \centering
        \includegraphics[width=\linewidth]{chapters/2-Background/Figures/network_disconnected_cluster_400.pdf}
        \textbf{(b)}
    \end{subfigure}
    \caption[Dangling and disconnected clusters.]{{\bf a)} A dangling cluster (blue) is defined as a cluster of bonds connected to the rest of the network by only one bond (a bridge, shown in red). These are detected and removed by removing bridges, i.e., bonds that disconnect the graph when cut. {\bf b)} A disconnected cluster (red) is a cluster of bonds that is disconnected from the bulk network. These are removed by identifying all connected components and retaining only the largest one.}
    \label{fig:PruningNetworks}
\end{figure}

Once a clean network has been generated, it can be pruned to reduce its average connectivity. This procedure is analogous to the bond dilution used in lattice-based networks to introduce disorder, although additional care is required due to the intrinsic positional disorder already present.

Bonds are removed selectively based on the local pair connectivity $z_{ij}=z_i+z_j$, where $z_i$ and $z_j$ denote the coordination numbers (i.e. the number of bonds) of the nodes connected by bond $(i,j)$. \DIFdelbegin \DIFdel{Bonds associated with the most highly connected node pairs are preferentially pruned to minimise }\DIFdelend \DIFaddbegin \DIFadd{At each pruning step, $z_{ij}$ is recalculated for every remaining bond and the maximum value $z_{\max}$ is identified. One bond is then selected uniformly at random from the set of bonds satisfying $z_{ij}=z_{\max}$ and removed. The selection is therefore deterministic with respect to pair connectivity but stochastic among equally ranked candidate bonds. This preferential removal of bonds between highly connected nodes limits }\DIFaddend spatial heterogeneity in the coordination field \DIFdelbegin \DIFdel{\mbox{%DIFAUXCMD
\cite{wyart_elasticity_2008, lerner_breakdown_2014}}\hskip0pt%DIFAUXCMD
. This process }\DIFdelend \DIFaddbegin \DIFadd{\mbox{%DIFAUXCMD
\cite{wyart_elasticity_2008,lerner_breakdown_2014}}\hskip0pt%DIFAUXCMD
. Following each removal, the coordination numbers are updated and any newly formed dangling or disconnected components are removed. The procedure }\DIFaddend is repeated until the target average connectivity is reached.
\DIFdelbegin \DIFdel{Any dangling or disconnected components that emerge during pruning are identified and removed as they are created.
}\DIFdelend 

\subsection{Calculating Stress} \label{sec:virialStress}

The concept of stress has been discussed extensively throughout this thesis and derived within the framework of continuum mechanics (see \cref{sec:StressAndStrain,sec:StressTensor}). In molecular dynamics simulations, however, stress is defined in terms of the underlying discrete particle interactions.

In principle, the macroscopic stress can be obtained via the principle of virtual work by numerically differentiating the system Hamiltonian with respect to an imposed strain. While formally exact, this approach is computationally expensive, because the derivative is typically evaluated using finite differences. Consequently, each stress evaluation requires relaxing the system following a small strain increment.

Instead, molecular dynamics simulations typically report the virial stress tensor, which can be computed efficiently from the instantaneous particle forces and positions. While there are many formulations of the virial stress tensor \cite{walton_pressure_1983,erpenbeck_einstein-kubo-helfand_1995,tsai_virial_1979}, we will present a brief derivation of the virial stress tensor derived by Irving and Kirkwood in Ref.~\cite{irving_statistical_1950} in two dimensions.


Following the arguments of Doi and Edwards in Ref.~\cite{doi_theory_1988}, we consider an arbitrary line perpendicular to the $j$-axis, which intersects the network at a position $h$ along that axis. The $ij$ component of the stress across this line, $\sigma_{ij}(h)$, is defined as the $i$ component of the force (per unit length) exerted by the portion of the network above the line on the portion below it, which we denote by $F_i(h)$. The macroscopic stress $\Sigma_{ij}$ is then obtained by averaging $\sigma_{ij}(h)$ along the $j$-direction,
\begin{gather}
    \Sigma_{ij} = \langle \sigma_{ij}(h) \rangle = \frac{\langle F_i(h) \rangle}{L_i}, \\
    \Sigma_{ij} = \frac{1}{L_i L_j}\int_0^{L_j} F_i(h)\, dh,
    \label{eq:virialDevInt}
\end{gather}
where $L_i$ is the system length in the $i$ direction and $L_j$ is the extent of the averaging along the $j$-axis.  Their product $L_iL_j \equiv A$ is the system area. For simplicity, we neglect solvent contributions and retain only the contribution to the stress from the inter-node forces.

The force $F_i(h)$ can be written as the sum of the $i$-components of the pairwise forces $f_{mn}$ exerted on node $m$ by node $n$,
\begin{equation}
    F_i(h) = \sum_{m,n} f_{mn,i}\,
    \Theta(h-r_{n,j})\Theta(r_{m,j}-h),
    \label{eq:virialDevForce}
\end{equation}
where $r_{n,j}$ is the $j$-coordinate of node $n$. The two Heaviside functions restrict node $n$ to lie below the line and node $m$ above it. Substituting \cref{eq:virialDevForce} into \cref{eq:virialDevInt} and integrating over $h$ yields
\begin{subequations}
\begin{align}
    \Sigma_{ij}
    &= \frac{1}{A}\sum_{m,n}
    f_{mn,i}(r_{n,j}-r_{m,j})
    \Theta(r_{n,j}-r_{m,j})
    \label{eq:virialDevTheta} \\
    \Sigma_{ij} &= \frac{1}{2A}\sum_{m,n}
    f_{mn,i}(r_{n,j}-r_{m,j})
    \label{eq:virialDevfinal} \\
    \Sigma_{ij} &= \frac{1}{2A}\sum_{m,n}
    f_{mn,i}\, r_{mn,j},
    \label{eq:virialStress}
\end{align}
\end{subequations}
where $\bm{r}_{mn}=\bm{r}_n-\bm{r}_m$. The Heaviside factor in \cref{eq:virialDevTheta} selects pairs with $r_{n,j}>r_{m,j}$. Using Newton's third law, $\bm{f}_{mn}=-\bm{f}_{nm}$, and symmetry under exchange of $m$ and $n$, this restriction can be replaced by the factor of $1/2$, which corrects for double counting in the summation, as seen in \cref{eq:virialDevfinal,eq:virialStress}.

The derivation above extends trivially to three dimensions, where instead force contributions above and below an arbitrary plane perpendicular to the $j$-axis are considered. Calculating the stress by this approach eliminates the computation of numerical derivatives and is more computationally efficient, because both $\bm{r}_{ij}$ and $\bm{f}_{ij}$ are already calculated at each timestep when evaluating the system's dynamics. A more thorough derivation of the virial stress tensor, including fluid effects, can be found in Ref.~\cite{yang_generalized_2012}.

\subsection{\DIFdelbegin \DIFdel{Generalised }\DIFdelend Lees--Edwards Boundary Conditions} \label{sec:Lees-Edwards}

When studying rheology, we aim to characterise a material's response to an imposed bulk stress or strain. In experimental settings, this is typically achieved by placing a sample of the material in a rheometer and imposing the desired loading or deformation through the motion of one or more boundaries. To ensure that the measured response reflects the bulk behaviour, samples are typically made sufficiently large that edge effects do not dominate the physics.

Achieving such system sizes in simulations, however, is often computationally prohibitive. A common strategy to mitigate this limitation is to employ periodic boundary conditions, whereby material leaving one side of the unit cell re-enters on the opposite side. This construction removes physical boundaries from the computational domain and thereby eliminates edge effects. Consequently, macroscopic strains must be imposed through alternative methods rather than via conventional boundary-driven forcing.

\begin{figure}[b]
    \centering
    \begin{subfigure}[t]{0.5\linewidth}
        \centering
        \includegraphics[width=.8\linewidth]{chapters/2-Background/Figures/TriClinic.tikz}
        \caption{Triclinic Boundary Conditions}
        \label{fig:Tri-Clinic}
    \end{subfigure}%
    \hfill
    \begin{subfigure}[t]{0.5\linewidth}
        \centering
        \includegraphics[width=.8\linewidth]{chapters/2-Background/Figures/LeesEdwards.tikz}
        \caption{Lees Edwards Boundary Conditions}
        \label{fig:Lees-Edwards}
    \end{subfigure}
    \caption{Illustration of two equivalent boundary conditions used in simulations to apply macroscopic stress to a system.}
    \label{fig:BoundaryConditions}
\end{figure}

In periodic simulations, macroscopic deformation is imposed by modifying the lattice vectors, which define the mapping between the unit cell and its periodic images. Simple loading protocols can therefore be implemented by appropriately deforming the simulation box and its lattice vectors, as illustrated in \cref{fig:Tri-Clinic}. This approach is commonly referred to as the use of triclinic boundary conditions, with the terminology derived from the triclinic crystal lattice.

While this method permits arbitrary homogeneous deformations, it can be computationally costly in practice due to the need to transform between bases when evaluating minimum-image distances. Moreover, at large strains, the unit cell can become highly skewed, which may lead to numerical inefficiencies. To address these issues, Lees and Edwards proposed a modified set of periodic boundary conditions that achieves the same macroscopic shear while remaining computationally efficient \cite{lees_computer_1972}.

To impose shear in the $x$ direction using Lees--Edwards boundary conditions, particle positions $\mathbf{r}_i=(x_i,y_i)$ are affinely transformed according to
\begin{equation}
    \DIFdelbegin %DIFDELCMD < \begin{pmatrix}x_i \\ y_i
%DIFDELCMD <     \end{pmatrix}%%%
\DIFdelend \DIFaddbegin \begin{pmatrix}
        x_i \\ y_i
    \end{pmatrix}\DIFaddend 
    \to
    \DIFdelbegin %DIFDELCMD < \begin{pmatrix}x_i + \gamma y_i \\ y_i
%DIFDELCMD <     \end{pmatrix}%%%
\DIFdelend \DIFaddbegin \begin{pmatrix}
        x_i+\gamma y_i \\ y_i
    \end{pmatrix}\DIFaddend \DIFaddbegin \DIFadd{,
}\DIFaddend \end{equation}
\DIFdelbegin \DIFdel{whilst periodic images are shifted by an amount }\DIFdelend \DIFaddbegin \DIFadd{while periodic images in adjacent rows are displaced by }\DIFaddend $\gamma L_y$\DIFdelbegin \DIFdel{between adjacent rows, as shown }\DIFdelend \DIFaddbegin \DIFadd{, as illustrated }\DIFaddend in \cref{fig:Lees-Edwards}. \DIFdelbegin \DIFdel{This construction is computationally advantageous because the }\DIFdelend \DIFaddbegin \DIFadd{We use this representation as it retains the orthogonal }\DIFaddend primary simulation cell \DIFdelbegin \DIFdel{remains orthogonal. Implementation }\DIFdelend \DIFaddbegin \DIFadd{and allows periodic bond vectors, including those crossing the cell boundaries, to be calculated using simple shear-dependent image shifts. This avoids the complex basis transformations and increasingly skewed simulation cell associated with an explicitly triclinic representation needed to calculate minimum-image distances. The imposed therefore only changes the geometry of each bond but not the network connectivity; bonds are removed only when the relevant failure criterion is satisfied. Further implementation }\DIFaddend details can be found in Ref.~\cite{allen_computer_2017}.

\section{Units and Parameters}
For the network simulations considered in this thesis, we work in units where the drag coefficient and filament stretching modulus are set to $\zeta = 1$ and $\mu = 1$, respectively, thereby fixing the characteristic viscous and elastic scales used throughout the network models. The reference length scale depends on the network construction: for lattice-based networks, the bond rest length is taken to be unity, whereas for packing-derived networks, the system size is taken to be $L_x = L_y = \sqrt{N}$, where $N$ is the number of nodes used to generate the initial packing. Unless stated otherwise, these packings are generated at packing fraction $\phi = 1$ with bidispersity parameter $f = 1.4$, which yields unpruned single networks with average connectivity $z \approx 5.4$. For the adaptive time-stepping scheme introduced in \cref{sec:EulerHeun}, we use an absolute error tolerance of $\tau = 10^{-4}$ and impose a maximum time step of $\delta t_{\max} = 0.1$.

\section{Conclusion}
% \red{
    In this chapter, we have introduced the theoretical framework used throughout the remainder of this thesis. We began with a continuum description of stress and strain, establishing the tensorial quantities used to characterise deformation and internal forces, and introduced the linear elastic constitutive relation that serves as a reference point for identifying yielding, plasticity, and fracture. We also outlined the strain-controlled deformation protocols used in later chapters, namely oscillatory strain, step strain, and strain startup, each of which probes a different aspect of material response.

    We then considered force balance in the Stokes limit, which provides the appropriate description for the overdamped systems studied here. Within this framework, local irreversible rearrangements may be represented through a plastic eigenstrain, and the resulting redistribution of stress is described by the Oseen tensor and the Eshelby stress propagator. These results provide the continuum basis for understanding how local events perturb the surrounding material and contribute to the bulk response.

    Finally, we introduced soft networks as model systems in which the coupling between structure, deformation, and failure can be resolved explicitly. We described the mechanics of semi-flexible filaments, the overdamped dynamics used to evolve the network, and the different notions of rigidity relevant to disordered systems. We then discussed how fracture emerges through heterogeneous stress concentration, non-affine rearrangements, and bond breakage, and introduced the network architectures and numerical methods used in later chapters, including lattice-based and packing-derived networks, pruning procedures, the virial stress tensor, and Lees--Edwards boundary conditions. Together, these concepts provide the basis for the analysis of damage accumulation under cyclic loading, delayed yielding following an imposed step strain, and stress redistribution and toughening in double-network materials.
% }

\newpage
\begin{subappendices}
    % !TeX root = ../../thesis.tex
\section{Fourier Derivation of the Oseen Tensor}
\label{app:OseenDerivation}

This appendix provides the intermediate derivation for the results quoted in \cref{sec:Oseen}. Starting from the Stokes equations for a plastic eigenstrain field,
\begin{subequations}
\label{eq:AppStokes}
\begin{align}
- \nabla p(\bm{r}) + \mu \nabla^2 \bm{u}(\bm{r}) &= 2\mu \nabla \cdot \bm{\epsilon}^{\mathrm{pl}}(\bm{r}), \\
\nabla \cdot \bm{u} &= 0,
\end{align}
\end{subequations}
which must be solved with the appropriate boundary conditions. The solution may formally be written as
\begin{equation}
    u_i(\bm{r}) = -2\mu \int \mathbb{O}_{ij}(\bm{r}-\bm{r}^\prime)\,\partial_l^\prime \epsilon_{jl}^{\mathrm{pl}}(\bm{r}^\prime)\, d\bm{r}^\prime.
\end{equation}
Here, $\mathbb{O}_{ij}(\bm{r}-\bm{r}^\prime)$ is the Oseen tensor, which is the Green's function for the Stokes equations in an infinite medium. To determine it, we consider the auxiliary problem of a point force $\bm{F}$ applied at $\bm{r}^\prime$. The Oseen tensor therefore satisfies
\begin{subequations}
\label{eq:AppStokesPointForce}
\begin{align}
- \nabla p(\bm{r}) + \mu \nabla^2 \bm{u}(\bm{r}) &= -\bm{F} \delta(\bm{r} - \bm{r}^\prime), \label{eq:AppStokesPointForce1} \\
\nabla \cdot \bm{u} &= 0. \label{eq:AppStokesPointForce2}
\end{align}
\end{subequations}

Following Ref.~\cite{kim_microhydrodynamics_1991}, we impose that the deformation field $\bm{u}(\bm{r})$ in \cref{eq:AppStokesPointForce} vanishes at infinity, i.e. $\lim_{r \to \infty}|\bm{u}(\bm{r})| = 0$. Because the Stokes equations are linear in $p(\bm{r})$ and $\bm{u}(\bm{r})$, we may write the pressure field as the scalar product of the point force $\bm{F}$ with a vector field $\bm{P}(\bm{r})$, and the deformation field as the tensor $\mathbb{O}_{ij}$ acting on $\bm{F}$:
\begin{equation}
    p(\bm{r}) = P_j F_j, \qquad
    u_i(\bm{r}) = \mathbb{O}_{ij} F_j.
    \label{eq:OseenProducts}
\end{equation}
Applying the Fourier transform to \cref{eq:AppStokesPointForce}, we obtain
\begin{equation*}
- i k_i \hat{P}_j F_j 
- \mu k^2 \hat{\mathbb{O}}_{ij} F_j 
= -F_i = -F_j \delta_{ij}.
\end{equation*}
Introducing the Kronecker delta $\delta_{ij}$ allows us to rewrite this as a tensor equation independent of $F_j$:
\begin{equation}
    - i k_i \hat{P}_j 
    - \mu k^2 \hat{\mathbb{O}}_{ij} 
    = -\delta_{ij}.
    \label{eq:StokesFourier2}
\end{equation}
The incompressibility condition $\nabla \cdot \bm{u} = 0$ reads in Fourier space
\[
    ik_i\hat{u}_i=ik_i\hat{\mathbb{O}}_{ij} F_j=0 \implies k_i\hat{\mathbb{O}}_{ij} = 0
\]
Multiplying \cref{eq:StokesFourier2} by $k_i$ and using the incompressibility relation eliminates the velocity-tensor contribution and gives
\begin{equation*}
    \hat{P}_j = -\frac{ik_j}{k^2}.
\end{equation*}
Substituting $\hat{P}_j$ back into \cref{eq:StokesFourier2}, we obtain the closed expression for the Oseen tensor in Fourier space,
\begin{equation}
    \hat{\mathbb{O}}_{ij} = \frac{1}{\mu k^2}\left( \delta_{ij} - \frac{k_ik_j}{k^2} \right).
    \label{eq:AppOseenPointForce}
\end{equation}
From \cref{eq:OseenProducts}, the displacement field is
\begin{equation}
    \hat{u}_i = \frac{1}{\mu k^2}\left( \delta_{ij} - \frac{k_ik_j}{k^2} \right)F_j.
    \label{eq:AppFourierDeformationPointForce}
\end{equation}
The corresponding deformation gradients ($\partial_l u_i$), expressed in Fourier space, are
\begin{equation}
    ik_l\hat{u}_i = \frac{ik_l}{\mu k^2}\left( \delta_{ij} - \frac{k_ik_j}{k^2} \right)F_j
    \label{eq:AppFourierDeformationGradientPointForce}
\end{equation}
which are the Fourier-space displacement and deformation-gradient fields used in \cref{sec:EshelbyPropogator}.

\DIFaddbegin \section{\DIFadd{Convergence of Numerical Integration}}
\label{app:stepStrainConvergence}

\DIFadd{We test the numerical convergence of the overdamped network dynamics described in }\cref{sec:EulerHeun} \DIFadd{following an imposed step strain. The comparison is restricted to $0\leq t\leq10$, over which all trajectories share the same checkpoints, separated by $\Delta t_{\mathrm{out}}=0.1$. Each trajectory is integrated directly to these checkpoints; no interpolation or nearest-time matching is used. The measured observable is the bulk virial shear stress $\Sigma_{xy}(t)$, consistent with the notation introduced in }\cref{sec:StressTensor,sec:virialStress}\DIFadd{.
}

\DIFadd{Fixed Euler trajectories with time steps $\delta t=10^{-6},10^{-5},10^{-4},10^{-3},10^{-2},$ and $10^{-1}$ are compared with a fixed reference trajectory using $\delta t_{\mathrm{ref}}=10^{-7}$. For each non-reference trajectory, we write the difference in bulk shear stress as
}\begin{equation}
    \DIFadd{\Delta\Sigma_{xy,n}=\Sigma_{xy,n}-\Sigma_{xy,n}^{\mathrm{ref}}.
}\end{equation}
\DIFadd{The normalised root-mean-square and maximum errors shown in }\cref{fig:stepStrainConvergence} \DIFadd{are then
}\begin{align}
 \DIFadd{\frac{\mathrm{RMS}(\Delta\Sigma_{xy})}{\Sigma_{\mathrm{ref}}^{\max}}
 }&\DIFadd{= \frac{1}{\Sigma_{\mathrm{ref}}^{\max}}
 \left[\frac{1}{N_t}\sum_{n=1}^{N_t}
 \left(\Delta\Sigma_{xy,n}\right)^2\right]^{1/2},}\\
 \DIFadd{\frac{\max_n|\Delta\Sigma_{xy,n}|}{\Sigma_{\mathrm{ref}}^{\max}}
 }&\DIFadd{= \frac{1}{\Sigma_{\mathrm{ref}}^{\max}}
 \max_n\left|\Delta\Sigma_{xy,n}\right|,
}\end{align}
\DIFadd{where
$\Sigma_{\mathrm{ref}}^{\max}=\max_n|\Sigma_{xy,n}^{\mathrm{ref}}|$. The fixed-step errors decrease linearly with $\delta t$, as expected for the first-order Euler method. Even the coarsest trajectory, $\delta t=10^{-1}$, gives a normalised RMS error of $9.20\times10^{-4}$ and a normalised maximum error of $4.12\times10^{-3}$ over this interval.
}

\DIFadd{We also test the adaptive Euler--Heun method using a rounded, approximately logarithmic grid of absolute tolerances,
}\begin{equation}
    \DIFadd{\tau=10^{-4},\;1.6\times10^{-4},\;2.5\times10^{-4},\;
    4\times10^{-4},\;6.3\times10^{-4},\;10^{-3}.
}\end{equation}
\DIFadd{The production parameters are $\tau_0=10^{-4}$ and $\delta t_{\max}=10^{-1}$. At the production tolerance, the normalised RMS and maximum errors are $4.80\times10^{-7}$ and $1.54\times10^{-6}$, respectively. The systematic reduction in both errors on reducing $\tau$ demonstrates convergence of the adaptive integration scheme.
}

\begin{figure}[t]
    \centering
    \includegraphics[width=0.95\linewidth]{chapters/2-Background/Figures/convergence_validation_mosaic.pdf}
    \caption[Time-step convergence of the step-strain dynamics]{\DIFaddFL{Numerical convergence of the overdamped network dynamics over $0\leq t\leq10$. }\textbf{\DIFaddFL{(a)}} \DIFaddFL{Bulk shear stress $\Sigma_{xy}(t)$ for all fixed Euler and adaptive Euler--Heun trajectories, including the fixed reference trajectory ($\delta t_{\mathrm{ref}}=10^{-7}$). }\textbf{\DIFaddFL{(b), inset}} \DIFaddFL{Corresponding pointwise normalised absolute errors $|\Delta\Sigma_{xy}(t)|/\Sigma_{\mathrm{ref}}^{\max}$ for every non-reference trajectory. }\textbf{\DIFaddFL{(c)}} \DIFaddFL{Normalised RMS and maximum errors for fixed Euler trajectories, demonstrating first-order convergence with $\delta t$. }\textbf{\DIFaddFL{(d)}} \DIFaddFL{The same normalised errors as functions of the absolute tolerance $\tau$ over the rounded logarithmic grid from $10^{-4}$ to $10^{-3}$.}}
    \label{fig:stepStrainConvergence}
\end{figure}

\DIFaddend \end{subappendices}
 \newpage 

 \newpage % !TeX root = ../../thesis.tex

\chapter{Absorption and Failure in Oscillatory Shear} \label{sec:Fatigue}

\def\ltilde{\tilde{l}}
\def\gammaMax{\gamma_0}
\def\lw{l_w}
\def\lc{\ell_c}
\def\Origin{\mathbb{G}_{xy}(\bm{0})}
\def\Norm{||\mathbb{G}_{xy}||_2}
\def\CStar{C^\ast}
\def\CycleNineFive{\CStar_{95\%}}
\def\Factive{f_{\mathrm{active}}}

\section{Introduction}

It is a common experience that materials can be repeatedly deformed over their lifetime, apparently without damage, until they ultimately fail. A paper clip, for example, which may eventually hold this document together, can be bent back and forth many times before suddenly snapping after a sufficient number of cycles. This phenomenon, known as fatigue, arises from the gradual accumulation of microscopic damage within the material, despite the absence of obvious macroscopic changes for much of its lifetime. Under repeated deformation, many solid materials progressively generate defects and microcracks that grow and interact, until a critical state is reached in which rapid crack propagation leads to catastrophic fatigue failure \cite{BasquinTheEL,Suresh1998,Anderson2017,DowlingMechanical2013}.

Consistent with the theme of this thesis, we will consider this phenomenon in the context of amorphous materials, which, unlike crystalline solids, lack any long-range order in the arrangements of their microscopic constituents (grains, bonds, droplets, etc.). Amorphous materials \cite{nicolas_deformation_2018,Berthier2011,Berthier2025} encompass a broad class of disordered solids, including soft materials such as gels and foams \cite{bonn_yield_2017}, as well as granular materials \cite{liuJammingNotJust1998} and metallic glasses \cite{SCHUH2007}. When subject to sufficiently small loads or deformations, such materials typically show an initial solid-like elastic response, with stress increasing approximately linearly with strain. At larger loads, however, this elastic response eventually breaks down, and the material yields plastically, undergoing irreversible rearrangements of its microscopic structure. Recent developments in experimental methods for tracking motion inside amorphous solids have made it possible to resolve the complex spatiotemporal dynamics that accompany yielding and flow \cite{porr_probing_2025,manneville_recent_2008,Besseling2009,Dutta2013,divoux_stress_2011,divouxTransientShearBanding2010,divoux_yielding_2012,grenard_timescales_2014,gibaud_heterogeneous_2010,kurokawa_avalanche-like_2015,Chan2013,Schott2025,Zhang2025,Shang2024,Aime2025}. Correspondingly, it has been shown that the yielding behaviour of amorphous materials can occur through different mechanisms. In some cases, yielding proceeds gradually and relatively homogeneously through the accumulation of many distributed plastic rearrangements \cite{Malandro1999,pollard_yielding_2022,falk_dynamics_1998,divoux_stress_2011,divoux_yielding_2012,divouxTransientShearBanding2010,grenard_timescales_2014,Gibaud2009,gibaud_heterogeneous_2010,kurokawa_avalanche-like_2015,Sentjabrskaja2015}. In other cases, yielding occurs more abruptly through the nucleation and propagation of localised shear bands that can span the system and precipitate failure \cite{divoux_stress_2011,pollard_yielding_2022,SCHUH2007,barlow_ductile_2020,Huang2016,bonnDelayedFractureInhomogeneous1998,leocmach_creep_2014,Tabuteau2009,fieldingTriggersSignaturesShear2016}. Despite the diversity of materials in which these phenomena occur, the similarities in their macroscopic behaviour have motivated the search for general theoretical descriptions of yielding in disordered solids \cite{cabriolu_precursors_2019}. 

Similar phenomenology is also observed in harder amorphous materials such as polymeric and metallic glasses \cite{hufnagel_deformation_2016,chen_mechanical_2008,boyce_large_1988,anand_large-deformation_2012,Argon1979}. These materials also typically exhibit elastic, solid-like behaviour at small strains before yielding at larger deformations, often accompanied by strain localisation and shear band formation \cite{greer_shear_2013,Maa2015}. \DIFdelbegin \DIFdel{However, unlike softer }\DIFdelend \DIFaddbegin \DIFadd{Softer }\DIFaddend amorphous solids, such as yield stress fluids, \DIFdelbegin \DIFdel{which }\DIFdelend \DIFaddbegin \DIFadd{can }\DIFaddend continue to flow \DIFdelbegin \DIFdel{and subsequently recover solidity when unloaded}\DIFdelend \DIFaddbegin \DIFadd{after yielding and recover solidity once unloaded. In metallic glasses, by contrast}\DIFaddend , shear banding \DIFdelbegin \DIFdel{in metallic glasses }\DIFdelend commonly leads to brittle failure and catastrophic fracture, limiting the material’s strength. The mechanisms underlying shear band formation in these hard materials remain an active topic of study \cite{greer_shear_2013,Cheng2011,Shi2006}. Homogeneous nucleation has been proposed as one route, in which correlated plastic rearrangements cooperatively evolve into a system-spanning shear band once a critical level of activity is reached \cite{Shi2006}. In practice, however, shear bands more often nucleate at pre-existing stress concentrations within the material, such as surface imperfections \cite{Vasoya2016,Shi2007}.

A key factor influencing the nature of yielding in all amorphous materials is their preparation history. The degree of annealing, which reflects how deeply the system resides within its energy landscape before deformation, affects the stability of a sample's microscopic constituents and the way plastic rearrangements develop under load \cite{Shi2005,Ozawa2018}. This annealing can be achieved either thermally or mechanically: thermal annealing involves equilibration through controlled temperature protocols such as slow cooling or aging \cite{yehGlassStabilityChanges2020,kumarMappingOutGlassy2022,mfielding_shear_2009,cochran_slow_2024}, while mechanical annealing arises from repeated deformation or cyclic loading, which can progressively reorganise the structure and drive the system toward stable steady states \cite{bhaumikRoleAnnealingDetermining2021,priezjevMolecularDynamicsSimulations2018,khandareElastoplasticModellingCyclic2026}. Well-annealed systems tend to exhibit more abrupt, localised yielding behaviour, while poorly annealed systems often yield more gradually through distributed plastic activity \cite{Leishangthem2017,Popovi2018,barlow_ductile_2020}. Accordingly, the mechanical response of amorphous materials under deformation depends heavily on their preparation history \cite{Shi2005,Ozawa2018,Vasisht2020,barlow_ductile_2020,Leishangthem2017,Popovi2018,nicolas_deformation_2018,SCHUH2007,Berthier2025}.

In this work, we focus on the behaviour of amorphous materials under cyclic deformation. Oscillatory shear has become a widely used protocol for probing the non-linear rheology of disordered solids and complex fluids. By repeatedly driving a material away from equilibrium and allowing it to relax during each cycle, cyclic deformation provides a controlled framework for investigating the interplay between elastic deformation and plastic relaxation events. 
% The repeated application of deformation over many cycles can lead to the gradual accumulation of microscopic damage within the material, as each loading cycle allows local plastic rearrangements and partial stress relaxation to occur. 
Although the early stages of this fatigue process often produce little observable macroscopic change, the progressive evolution of the microstructure can ultimately destabilise the material, leading to a higher-energy state characterised by chaotic, irreversible dynamics from cycle to cycle \cite{Regev_onset_2013,pineChaosThresholdIrreversibility2005,Das2020}.

As a result, oscillatory shear has been widely used as a powerful probe of the non-linear rheology of amorphous materials, revealing several characteristic signatures of their deformation and yielding dynamics. These include nonlinear stress responses \cite{yoshimuraResponseElasticBingham1987a,ewoldtLargeAmplitudeOscillatory2010,renouYieldingProcessesColloidal2010,rogersSequencePhysicalProcesses2012,blackwellSimpleThixotropicViscoelastic2014}, progressive structural changes \cite{Koumakis2013,Poulos2013,Poulos2015,vandoornStrandPlasticityGoverns2018}, and the localisation of deformation into shear bands \cite{radhakrishnan_shear_2016,radhakrishnan_shear_2018,rogersAgingYieldingShear2008,yehGlassStabilityChanges2020,Parmar2019}. \DIFaddbegin \DIFadd{Beyond characterising rheological behaviour, cyclic loading also forms the basis of accelerated fatigue-testing protocols, in which materials are subjected to controlled, often more severe, loading conditions so that damage develops over experimentally accessible timescales. The resulting dependence of the number of cycles to failure on loading amplitude can then be used to assess durability under less severe service conditions \mbox{%DIFAUXCMD
\cite{BasquinTheEL,Suresh1998,DowlingMechanical2013}}\hskip0pt%DIFAUXCMD
. This practical perspective motivates studying not only the eventual response under oscillatory shear, but also the number of cycles required to reach it and how the material evolves over those cycles.
}\DIFaddend 

Unlike in crystalline solids, where fatigue is often interpreted in terms of the gradual growth of microscopic cracks or defects, the mechanisms underlying failure in amorphous materials remain less well understood. Experiments on colloidal gels under oscillatory loading provide one such example \cite{Perge2014,Gibaud2016,SaintMichel2017}. These studies reveal a multistage yielding process: the material initially retains a solid-like response over many cycles, before wall slip develops at the boundaries. The deformation then becomes heterogeneous, with the emergence of coexisting solid–fluid shear bands in the bulk, before the system ultimately fluidises completely \cite{Gibaud2016}. Such observations highlight the complex and history-dependent processes by which amorphous materials evolve under cyclic deformation.

In addition to the initial state of the material, determined by its preparation history, the response of amorphous materials under oscillatory shear has been shown to depend strongly on the imposed strain amplitude $\gammaMax$ relative to a critical amplitude $\gamma_c$ \cite{bhaumikRoleAnnealingDetermining2021,Cort2008,Das2020,Fiocco2013,kawasakiMacroscopicYieldingJammed2016,KhirallahIntegerModel2021,knowltonMicroscopicViewYielding2014,kumarMappingOutGlassy2022,lavrentovichPeriodProliferationPeriodic2017,bhardwajSituClickChemistry2017,liuFateShearoscillatedAmorphous2022,munganMetastabilityMechanismYielding2021,himanagamanasaExperimentalSignaturesNonequilibrium2014,nessAbsorbingStateTransitionsGranular2020,pineChaosThresholdIrreversibility2005,priezjevMolecularDynamicsSimulations2018,Regev_onset_2013,sastryModelsYieldingBehavior2021,szulcForcedDeterministicDynamics2020,yehGlassStabilityChanges2020}. For sufficiently small strain amplitudes, $\gammaMax < \gamma_c$, the system can evolve towards a quiescent state in which plastic rearrangements eventually cease. In this regime, repeated deformation drives the system to self-organise into configurations that avoid further irreversible plastic events, such that after an initial period, the microscopic configuration repeats exactly after each deformation cycle \cite{Regev_onset_2013,kawasakiMacroscopicYieldingJammed2016,Leishangthem2017}. This quiescent regime is commonly described as an absorbing state, since once the system reaches this configuration, no further plastic activity can occur under the imposed oscillation. \DIFdelbegin %DIFDELCMD < 

%DIFDELCMD < %%%
\DIFdel{As the strain amplitude approaches the critical value, $\gammaMax\to\gamma_c^-$, the number of cycles required to reach this quiescent steady state increases rapidly and can diverge at the critical point $\gammaMax$ \mbox{%DIFAUXCMD
\cite{Regev_onset_2013,kawasakiMacroscopicYieldingJammed2016}}\hskip0pt%DIFAUXCMD
. This behaviour has been interpreted as critical slowing down in the dynamics associated with a nonequilibrium phase transition between reversible and irreversible dynamical states. Evidence for this phenomenon }\DIFdelend \DIFaddbegin \DIFadd{Similar reversible behaviour }\DIFaddend has been observed across a wide range of experimental and numerical systems. In periodically driven suspensions, for example, particles can self-organise under cyclic shear into configurations that avoid further collisions, producing a reversible steady state in which trajectories form closed, periodic loops \DIFdelbegin \DIFdel{, }\DIFdelend and diffusion vanishes \cite{pineChaosThresholdIrreversibility2005,Cort2008}. \DIFdelbegin \DIFdel{Similar behaviour has }\DIFdelend \DIFaddbegin \DIFadd{This type of reversible behaviour has also }\DIFaddend been reported in experiments on colloidal glasses and dense granular materials subjected to oscillatory shear, where the reversible regime is characterised by vanishing stroboscopic diffusion and the absence of sustained rearrangements \cite{himanagamanasaExperimentalSignaturesNonequilibrium2014,knowltonMicroscopicViewYielding2014}.

\DIFaddbegin \DIFadd{As the strain amplitude approaches the critical value, $\gammaMax\to\gamma_c^-$, the number of cycles required to reach this quiescent steady state increases rapidly and can diverge at the critical point \mbox{%DIFAUXCMD
\cite{Regev_onset_2013,kawasakiMacroscopicYieldingJammed2016}}\hskip0pt%DIFAUXCMD
. This increase reflects the progressively longer transient dynamics required for the system to organise into a reversible configuration as the transition is approached. This behaviour has been interpreted as critical slowing down in the dynamics associated with a nonequilibrium phase transition between reversible and irreversible dynamical states. Above the critical amplitude, the system instead remains dynamically active, with plastic rearrangements persisting indefinitely under repeated cycling.
}

\DIFaddend Particle-based simulations of amorphous solids reveal closely related behaviour. Athermal quasi-static simulations of model glasses subjected to cyclic shear show that, for strain amplitudes below yield, the system evolves towards periodic trajectories in configuration space, in which the same sequence of reversible rearrangements repeats every cycle \cite{Fiocco2013,Regev_onset_2013,Leishangthem2017,nessAbsorbingStateTransitionsGranular2020}. Near the transition, the transient time required for the system to reach this periodic state grows rapidly with strain amplitude, while above the transition, the system instead exhibits persistent plastic activity and diffusive particle motion \cite{Regev_onset_2013,kawasakiMacroscopicYieldingJammed2016}. In this irreversible regime, the continued accumulation of plastic rearrangements often leads to the localisation of deformation and the nucleation of shear bands, even in systems that remain homogeneous below yield \cite{Leishangthem2017,bhaumikRoleAnnealingDetermining2021}.

These observations have motivated interpreting cyclic yielding in amorphous materials within the framework of absorbing-state phase transitions, drawing on the analogy between reversible steady states and absorbing configurations \cite{nessAbsorbingStateTransitionsGranular2020,Cort2008,himanagamanasaExperimentalSignaturesNonequilibrium2014}. Within this picture, the reversible regime corresponds to an absorbing phase in which the system eventually reaches a configuration that no longer generates activity under the imposed oscillation. The dynamics then settle into a periodic limit cycle in which material configurations repeat exactly from cycle to cycle and the mean-squared displacement measured stroboscopically vanishes. In contrast, above a critical strain amplitude, the system fails to find such a repeating trajectory. The resulting persistent rearrangements lead to diffusive particle motion and irreversible structural evolution, marking the onset of yielding. 
% This perspective is consistent with recent work interpreting yielding itself, under imposed stress, as an absorbing phase transition \cite{JocteurYielding2024}.

To capture both yielding and absorbing-state behaviour under cyclic shear, many authors have turned to elastoplastic models (EPMs), which provide a mesoscopic, coarse-grained description of plastic deformation in amorphous solids. In these models, the material is partitioned into mesoscopic elements representing local coarse-grained volumes containing many microscopic constituents, rather than sharply defined physical subregions. Each element is assumed to be a mesoscopic sample of the material, each representing a coarse-grained region containing multiple microstructural constituents (for example, a few particles in a colloidal system \cite{Schall2007}, or several hundred metallic glass particles \cite{Pan2008}). This partition into mesoscopic elements should be understood as a conceptual coarse-graining procedure rather than as a division into sharply defined physical subregions; the elements do not correspond to fixed material boundaries, but instead denote representative volumes over which a local constitutive description is assumed to apply. Therefore, these element sizes are assumed to satisfy two conditions:
\begin{enumerate}
    \item Sufficiently small that the bulk response can be obtained by averaging over many such elements.
    \item Sufficiently large that, within each element, the deformation can be meaningfully characterised by local continuum elastic strain and stress fields.
\end{enumerate}

Within this class of model, each element deforms elastically under the applied load until its local stress (or equivalently strain) exceeds a local yield threshold. When this threshold is reached, the element undergoes a local plastic event, representing a rearrangement of its microscopic structure. This event leads to a drop in the local stress, and the element transitions to a new local configuration. The stress released during this rearrangement is then redistributed elastically throughout the surrounding medium, typically through a long-range interaction kernel reflecting the elastic response of the solid. As a result, plastic events are not independent but can trigger further yielding in nearby regions, giving rise to correlated cascades of rearrangements and complex spatio-temporal dynamics. A comprehensive review of these models can be found in Ref.~\cite{nicolas_deformation_2018}. 

Several different variations of mesoscopic elastoplastic models have been used to understand the deformation and yielding of amorphous materials under cyclic shear. These models aim to capture the collective dynamics that emerge from interacting local plastic rearrangements while remaining computationally tractable compared with fully microscopic particle simulations. In recent work, Cochran \etal applied the soft glassy rheology (SGR) framework to study fatigue and delayed yielding in amorphous materials subjected to oscillatory shear \cite{cochran_slow_2024}. Unlike the athermal threshold-based model considered here, SGR represents mesoscopic elements as traps in an energy landscape and yielding as activated hops between traps, with rates controlled by the barrier depth and an effective noise temperature $x$ \cite{sollich_rheological_1998,sollich_rheology_1997}. Within this framework, the material can accumulate strain heterogeneity over many cycles before yielding catastrophically through shear-band formation, with the number of cycles to failure depending strongly on the imposed strain amplitude, temperature, and initial degree of annealing \cite{cochran_slow_2024}.

Another important study of cyclic deformation is the tensorial elastoplastic model of Liu \etal~\cite{liuFateShearoscillatedAmorphous2022}, which extends the scalar picture adopted here by resolving the full local stress tensor and its evolution under oscillatory shear. Using an analytically constructed Eshelby kernel following Ref.~\cite{caoSoftModesStrain2018}, they showed that the large-cycle response depends on both annealing and strain amplitude: at small amplitudes, poorly annealed samples undergo progressive shear-annealing, whereas beyond a critical amplitude all samples transition to a mixed state consisting of a fluidised shear band embedded within a marginally stable solid.

Most recently, Khandare and Sastry introduced an athermal quasi-static elastoplastic model to study amorphous materials under athermal quasistatic cyclic shear \cite{khandareElastoplasticModellingCyclic2026}. In this framework, each mesoscopic element is represented as a mesostate, i.e. a local energy minimum, and elastic interactions are computed using a finite-element implementation of the Eshelby kernel \cite{Nicolas2015,Budrikis2017}. The model \DIFdelbegin \DIFdel{reproduces several established }\DIFdelend \DIFaddbegin \DIFadd{captures key }\DIFaddend features of amorphous-solid deformation under uniform shear, including the annealing-dependent \DIFdelbegin \DIFdel{brittle-to-ductile crossover }\DIFdelend \DIFaddbegin \DIFadd{crossover from abrupt, localised yielding to gradual, distributed plastic deformation, }\DIFaddend and the Bauschinger effect\DIFdelbegin \DIFdel{under load reversal, and provides a framework to characterise cyclic-shear behaviour}\DIFdelend \DIFaddbegin \DIFadd{, in which prior deformation produces an asymmetric response upon load reversal}\DIFaddend . In particular, poorly annealed samples undergo progressive relaxation towards lower-energy absorbing states at small strain amplitudes, whereas better-annealed samples remain close to their initial energy until a larger, annealing-dependent critical amplitude is reached. Beyond this amplitude, failure occurs via shear-band formation for all annealing levels, and the number of cycles required to reach the final state grows strongly as the transition is approached from either side. The model also exhibits a regime in which absorbing and yielding trajectories coexist, and the steady-state energies of these end states can be characterised in a manner comparable to Ref.~\cite{liuFateShearoscillatedAmorphous2022}.

Although the emergence of yielding and absorption has been extensively characterised in both experimental studies and particle-based simulations, most work on elastoplastic models has focused on characterising the resultant steady-state behaviour post-yielding \cite{liuFateShearoscillatedAmorphous2022} or the divergence in the number of cycles to failure \cite{cochran_slow_2024}. Recent work published \cite{khandareElastoplasticModellingCyclic2026} during the final stages of preparing this chapter has shown that a divergence in the number of cycles required to reach the final state can be captured in an athermal quasi-static elastoplastic model, and that a coexistence region exists in the vicinity of the yield point, where the probability of evolving to an absorbing state drops from $1$ to $0$ over a finite range. 

In this work, we performed a detailed numerical study to characterise this transition region in more detail and to examine the mechanisms by which systems approach an absorbing state under cyclic deformation. Rather than focusing solely on macroscopic observables, we consider the dynamics at the level of individual mesoscopic elements. In particular, we analyse how these systems exhibit less and less activity as the system evolves under repeated oscillation. Within this framework, the relaxation towards an absorbing state can be interpreted as a stochastic process in which each element evolves until it reaches a configuration from which no further plastic events occur.  This allows the problem to be naturally formulated as a first-passage problem. By adopting this perspective, we can relate the macroscopic relaxation of the system to the statistical properties of the first-passage dynamics of its constituent elements.

\section{Methodology} \label{sec:fatigueMethodology}

As a starting point\DIFdelbegin \DIFdel{for the model used in this work}\DIFdelend , we adopt the \DIFdelbegin \DIFdel{framework introduced }\DIFdelend \DIFaddbegin \DIFadd{scalar elastoplastic model introduced by Pollard and Fielding \mbox{%DIFAUXCMD
\cite{pollard_yielding_2022}}\hskip0pt%DIFAUXCMD
, formulated within the overdamped continuum framework described }\DIFaddend in \cref{sec:Hydrodynamics}. We briefly restate \DIFdelbegin \DIFdel{it }\DIFdelend \DIFaddbegin \DIFadd{the model }\DIFaddend here to fix notation \DIFaddbegin \DIFadd{and identify the modifications used in the present oscillatory-shear study}\DIFaddend . Physically, the governing equations describe the slow deformation of an incompressible viscoelastic material in the overdamped limit, where elastic and viscous stresses dominate over inertia. The total stress tensor is therefore decomposed as
\[
    \bm{\sigma} = \bm{\tau} - p\bm{I},
\]
where $\bm{\tau}$ denotes the deviatoric stress and $p$ is the hydrostatic pressure enforcing incompressibility.

In this low-Reynolds-number, overdamped limit, momentum conservation reduces to the Stokes force-balance equations
\begin{subequations}
\label{eq:fatigueForceBalance}
\begin{align}
\nabla \cdot (\bm{\tau} - p\bm{I}) &= 0, \\
\nabla \cdot \bm{v} &= 0.
\end{align}
\end{subequations}
Throughout this chapter, we further specialise to the quasi-static shear limit, so that elastic propagation and plastic relaxation are treated as instantaneous compared to the timescale of the applied deformation and the material regains mechanical equilibrium after each strain increment. This assumption motivates, as in \cref{sec:Hydrodynamics} the use of a displacement-based formulation in which the displacement field $\bm{u}(\bm{r},t)$ is defined through $\bm{v} = \partial_t \bm{u}$.

To solve the dynamics of \cref{eq:fatigueForceBalance}, we introduce a square lattice of $i=1,2,\ldots,N$ by $j=1,2,\ldots,N$ elements. While each element could be associated with a mesoscopic conformation tensor which fully defines the material's state, instead we consider a scalar model where each element is assigned a local elastic shear strain $l$, a corresponding local shear stress $\mu l$ and a local yield energy $E$. In this way we neglect normal stresses.
\DIFaddbegin 

\DIFaddend Each element in the lattice is assumed to follow the following set of rules. \DIFdelbegin \DIFdel{Between each plastic event, the strain of each element is assumed to follow affinely}\DIFdelend \DIFaddbegin \DIFadd{Following the scalar elastoplastic model of Pollard and Fielding \mbox{%DIFAUXCMD
\cite{pollard_yielding_2022}}\hskip0pt%DIFAUXCMD
, between plastic events we assume that the externally imposed shear loads all elements affinely, such that a change $\Delta\gamma$ in }\DIFaddend the globally applied \DIFdelbegin \DIFdel{elastic strain $\gamma$, corresponding to elastic loading in the material.
}\DIFdelend \DIFaddbegin \DIFadd{strain shifts the local strains uniformly as $l_{ij}\to l_{ij}+\Delta\gamma$. This treatment is consistent with the homogeneous linear-elastic formulation commonly adopted in elastoplastic models \mbox{%DIFAUXCMD
\cite{nicolas_deformation_2018}}\hskip0pt%DIFAUXCMD
. Within the small-deformation framework, the displacement gradient can be decomposed as $\nabla\bm{u}=\bm{\epsilon}+\bm{\omega}$, where $\bm{\epsilon}$ is the symmetric strain tensor and $\bm{\omega}$ is the antisymmetric rotational part. At linear order, the elastic deviatoric stress depends on $\bm{\epsilon}$, but not on $\bm{\omega}$. The heterogeneous residual stresses generated by previous plastic events are therefore retained through the heterogeneous local shear strain field $l_{ij}$, while couplings between these residual stresses and local rotations, which arise beyond the linear small-deformation description, are neglected. Although more complete tensorial elastoplastic models retain additional stress components neglected here, previous studies have found that, under uniform loading, their qualitative flow and avalanche behaviour is similar to that obtained from scalar models \mbox{%DIFAUXCMD
\cite{nicolas_deformation_2018,Budrikis2017,Nicolas2014}}\hskip0pt%DIFAUXCMD
.
}

\DIFaddend Each element is assumed to continue this elastic loading until its strain energy reaches the yield threshold
\[
    E=\frac{1}{2}\mu l_c^2.
\]
Beyond this threshold, the element is assumed to immediately yield. For this work we will work in units where $\mu=1$ and $E=0.5$, such that the yielding threshold for elements reduces to 
\[
    |l| > \lc := \sqrt{2E/\mu}=1.
\]
\DIFaddbegin \DIFadd{This choice of setting $l_c=1$ corresponds to measuring strains in units of the physical yield strain. Assuming that the physical yield strain is small, this remains consistent with the small-strain formulation.
}


%DIF >  Each element in the lattice is assumed to follow the following set of rules. Between each plastic event, the strain of each element is assumed to follow affinely the globally applied elastic strain $\gamma$, corresponding to elastic loading in the material. Each element is assumed to continue this elastic loading until its strain energy reaches the yield threshold
%DIF >  \[
%DIF >      E=\frac{1}{2}\mu l_c^2.
%DIF >  \]
%DIF >  Beyond this threshold, the element is assumed to immediately yield. For this work we will work in units where $\mu=1$ and $E=0.5$, such that the yielding threshold for elements reduces to 
%DIF >  \[
%DIF >      |l| > \lc := \sqrt{2E/\mu}=1.
%DIF >  \]
%DIF >  This choice of setting $l_c=1$ corresponds to measuring strains in units of the physical yield strain. Assuming that the physical yield strain is small, this remains consistent with the small-strain formulation.
\DIFaddend 

When an element yields, it adopts a new local strain drawn at random from a Gaussian distribution with mean $0$ and variance $\lw^2$, which we denote by $\mathcal{N}(0,\lw^2)$. This parameter does not alter whether the material ultimately yields or absorbs, and therefore does not modify the qualitative physics of the transition in the range studied here. It does, however, introduce a characteristic number of cycles, which will be examined and quantified in subsequent sections. 

Following a yielding event, the material is temporarily out of force balance, so the stress tensor no longer satisfies \cref{eq:fatigueForceBalance} and the local stress field has non-zero divergence. Force balance is then immediately reimposed via the Eshelby stress propagator in two dimensions \cite{picard_elastic_2004,eshelby_determination_1997} (see \cref{sec:EshelbyPropogator} for its derivation). As specified in \cref{sec:Hydrodynamics}, the displacement correction $\bm{u}^1$ generated by a stress perturbation $\bm{S}$ is, in Fourier space,
\begin{equation}
    \hat{\bm{u}}^1
    =
    \frac{1}{\mu k^2}
    \left(
        \bm{I}
        - \DIFdelbegin \DIFdel{\frac{\bm{k}\bm{k}}{k^2} }\DIFdelend \DIFaddbegin \DIFadd{\frac{\bm{k}\otimes\bm{k}}{k^2}
    }\DIFaddend \right)
    \DIFaddbegin \DIFadd{\cdot
    }\left(\DIFaddend i\DIFdelbegin %DIFDELCMD < \bm{k} %%%
\DIFdelend \hat{\bm{S}}\DIFaddbegin \DIFadd{\cdot}\bm{k}\right)
    \DIFaddend =
    \hat{\mathbb{G}}
    \DIFaddbegin \DIFadd{\cdot
    }\left(\DIFaddend i\DIFdelbegin %DIFDELCMD < \bm{k} %%%
\DIFdelend \hat{\bm{S}}\DIFaddbegin \DIFadd{\cdot}\bm{k}\right)\DIFaddend ,
    \label{eq:FatiguefourierDeformation}
\end{equation}
with wave vectors $\bm{k}$. Here\DIFaddbegin \DIFadd{, $\otimes$ denotes the dyadic product, while $\cdot$ denotes contraction over one pair of tensor indices. In index notation,
$(\hat{\bm{S}}\cdot\bm{k})_j=\hat{S}_{j\ell}k_\ell$, such that
$\hat{u}^1_i=\hat{G}_{ij}i\hat{S}_{j\ell}k_\ell$. Here }\DIFaddend and below, the superscript $1$ denotes the correction induced by any yielding events. Consistent with the infinitesimal strain tensor introduced in \cref{eq:linearStrain}, the corresponding strain correction is
\[
    \bm{\epsilon}^1 = \frac{1}{2}\left(\nabla \bm{u}^1 + (\nabla \bm{u}^1)^T\right).
\]
The associated deviatoric stress correction then follows from \cref{eq:elasticStress} as
\[
    \bm{\tau}^1 = 2\mu \bm{\epsilon}^1 = \mu\left(\nabla \bm{u}^1 + (\nabla \bm{u}^1)^T\right).
\]
In our model, we only consider the shear stress component of this correction $\tau^1_{xy}$, which is used to correct the distribution of local strains $l$. As noted in \cref{sec:Hydrodynamics}, equation \ref{eq:FatiguefourierDeformation} is only defined up to a constant at $\bm{k}=\bm{0}$. For strain-controlled protocols, this is fixed by imposing the conservation of the globally applied strain, which fixes the zero mode to $\hat{\mathbb{G}}_{xy}(\bm{0}) = -1$ \cite{nicolas_deformation_2018,martens_spontaneous_2012,jocteur_protocol_2025,lin_scaling_2014}.

\begin{figure}[t]
    \centering
    \includegraphics[width=0.99\textwidth]{chapters/Fatigue/Diagrams/Prop.pdf}
    \caption[Stress redistribution $\bm{\tau}^1_{xy}$ resulting from a point stress perturbation $\bm{S}_{xy}=1$ in $2D$]{{\bf a)} Stress redistribution $\bm{\tau}^1_{xy}$ resulting from a point stress perturbation $\bm{S}_{xy}=1$ in $2D$ {\bf b)} \DIFdelbeginFL \DIFdelFL{Convergence }\DIFdelendFL \DIFaddbeginFL \DIFaddFL{Finite-system-size convergence of the magnitude }\DIFaddendFL of the value at the origin\DIFdelbeginFL \DIFdelFL{$\Origin$ }\DIFdelendFL \DIFaddbeginFL \DIFaddFL{, $|\Origin|$, }\DIFaddendFL and \DIFaddbeginFL \DIFaddFL{the }\DIFaddendFL norm\DIFaddbeginFL \DIFaddFL{, }\DIFaddendFL $\Norm$\DIFaddbeginFL \DIFaddFL{, }\DIFaddendFL with increasing grid size $N$.}
    \label{fig:PropogatorConvergence}
\end{figure}

In \cref{fig:PropogatorConvergence}a, we \DIFdelbegin \DIFdel{plot the resulting shear stress correction $\tau^{1}_{xy}$. In }%DIFDELCMD < \cref{fig:PropogatorConvergence}%%%
\DIFdel{b we show the convergence, with grid size $N$, of two key quantities derived from $\mathbb{G}_{xy}$. We define the resulting stress drop from an element which yields from $l=1 \to l=0$ as
}\DIFdelend \DIFaddbegin \DIFadd{show the real-space response to a unit point-stress perturbation applied to the element at the origin, $\bm{r}=\bm{0}$. The magnitude at each element gives the size of its shear-stress correction after force balance has been restored, while the sign indicates whether its stress increases or decreases. The central value describes the stress change experienced by the perturbed element itself, while the four surrounding lobes show how the stress is redistributed through the rest of the system. For this unit perturbation, the central value is
}\DIFaddend \[
\Origin \approx -0.5708
\qquad \text{as } N \to \infty.
\]
\DIFdelbegin \DIFdel{This gives the stress correction at the yielded site per unit stress released}\DIFdelend \DIFaddbegin \DIFadd{The negative sign indicates a reduction in the stress of the perturbed element}\DIFaddend . Thus, if an element yields \DIFdelbegin \DIFdel{at }\DIFdelend \DIFaddbegin \DIFadd{from }\DIFaddend $l=1$ \DIFdelbegin \DIFdel{and is locally reset to $l=0$ before force balance is reimposed, the net stress drop of the element is not $1$ but $|\Origin| \approx 0.5708$; equivalently, the site retains a stress $1-|\Origin| \approx 0.4292$ after the force-balance correction}\DIFdelend \DIFaddbegin \DIFadd{with post-hop strain $\eta=0$, its net stress change after force balance is $\Origin \approx -0.5708$, leaving it with a stress $1+\Origin \approx 0.4292$}\DIFaddend . More generally, \DIFdelbegin \DIFdel{if the post-hop strain is $\eta$, }\DIFdelend the released stress is $1-\eta$, so the net stress \DIFdelbegin \DIFdel{drop is }\DIFdelend \DIFaddbegin \DIFadd{change at the yielding element is $\Origin(1-\eta)$, with magnitude }\DIFaddend $|\Origin|(1-\eta)$.
\DIFdelbegin \DIFdel{We also define }\DIFdelend \DIFaddbegin 

\DIFadd{Whereas $\Origin$ characterises the response of the element at the origin, }\DIFaddend the two-norm \DIFdelbegin \DIFdel{of the propagator as
}\DIFdelend \DIFaddbegin \DIFadd{measures the overall magnitude of the complete stress-redistribution pattern. It can be visualised by squaring the correction at every element in }\cref{fig:PropogatorConvergence}\DIFadd{a, summing these values, and taking the square root:
}\DIFaddend \[
\Norm
:=
\sqrt{
\sum_{i=1}^{N}
\sum_{j=1}^{N}
\left( \mathbb{G}_{xy}(i,j) \right)^2 
} \approx 0.6724
\qquad \text{as } N \to \infty.
\]
Accordingly, $\Norm$ quantifies the \DIFdelbegin \DIFdel{RMS (root-mean-square) }\DIFdelend \DIFaddbegin \DIFadd{root-sum-square }\DIFaddend magnitude of the redistributed shear-stress field generated by a unit local event. In particular, for a zero-mean random strain field with no spatial correlations, such as that used in our model, projection back into force balance rescales the standard deviation of the resulting field by a factor $\Norm$. \DIFaddbegin \Cref{fig:PropogatorConvergence}\DIFadd{b shows that both quantities converge towards their large-system limits as $N$ increases, consistent with a finite-size effect arising from the self-interaction of the long-range propagator through the
periodic boundaries. Numerical round-off from the double-precision FFTW3 calculations~\mbox{%DIFAUXCMD
\cite{Frigo2005} }\hskip0pt%DIFAUXCMD
is negligible on this scale.
}\DIFaddend 

To initialise the sample before any deformation, each element is assigned a local shear strain $l$, sampled at random from a distribution that has a width $l_0$ and zero mean. This distribution is naturally not in force balance; as such, we immediately impose it across the lattice using \cref{eq:FatiguefourierDeformation}. As a result of this imposed force balance, the resultant distribution of shear strain $l$ is reduced by a factor equal to $\Norm$. In the rare instances where an element remains above its yielding threshold after this projection, it is allowed to yield according to the rules specified \DIFdelbegin \DIFdel{later }\DIFdelend \DIFaddbegin \DIFadd{earlier }\DIFaddend in the chapter. We will only consider initial distributions sampled in this manner (from a Gaussian distribution) in this work. However, it should be noted that simulations in which the initial strain distribution was instead drawn from a beta distribution, as implemented in Ref.~\cite{pollard_yielding_2022}, showed little qualitative difference in the shear startup protocol (see \cref{sec:ShearStartup}). Alternative algorithms which enforce force balance have been used, such as the one discussed in Ref.~\cite{Popovi2018}.

Small values of the initial width $l_0$ represent well-annealed samples, whereas larger values correspond to poorly annealed ones. In experimental systems, strong annealing may result from equilibration at temperatures just above the glass transition, from slow cooling from high temperature to zero, from prolonged ageing at low temperature before quenching \cite{yehGlassStabilityChanges2020,kumarMappingOutGlassy2022,mfielding_shear_2009,cochran_slow_2024}, or from prior mechanical conditioning, such as repeated sub-yield deformation, that progressively relaxes internal stresses \cite{bhaumikRoleAnnealingDetermining2021,priezjevMolecularDynamicsSimulations2018,khandareElastoplasticModellingCyclic2026}. By contrast, weak annealing would be associated with equilibration at higher temperature before quench, more rapid cooling, shorter ageing times, or little prior mechanical conditioning. We do not attempt to model these preparation histories directly; instead, the single parameter $l_0$ serves as an effective measure of the degree of annealing.

For discussion later in this chapter, we introduce a decomposition,
\[
    l = \ltilde + \gamma
\]
where as stated previously, $\gamma$ is the globally applied elastic strain, and is therefore uniform across all elements. The quantity $\ltilde$ is therefore defined as the local deviation from this imposed contribution, $\ltilde := l - \gamma$ and quantifies the element-to-element variation about the externally applied deformation. Equivalently, $\ltilde$ represents the \DIFdelbegin \DIFdel{intrinsic (non-affine) }\DIFdelend \DIFaddbegin \DIFadd{heterogeneous residual }\DIFaddend strain field obtained after subtracting the uniform, externally imposed contribution $\gamma$. \DIFdelbegin \DIFdel{Since }\DIFdelend \DIFaddbegin \DIFadd{Consistent with the affine-loading assumption introduced above, }\DIFaddend $\ltilde$ \DIFdelbegin \DIFdel{is defined to be separate from the globally applied elastic strain, it }\DIFdelend \DIFaddbegin \DIFadd{remains fixed during elastic loading and }\DIFaddend evolves only through internal plastic events \DIFdelbegin \DIFdel{within the material }\DIFdelend and the resulting propagation of elastic stress. Although it only changes when an element yields, it should not be interpreted as the accumulated strain lost by that element over its yielding history. In particular, $\ltilde$ is not the sum of the local strain drops associated with past yielding events. Rather, it reflects the combined effect of the element’s initial strain state and the elastically redistributed strain generated by yielding activity throughout the material.

As discussed previously, we wish to consider the response of amorphous solids under oscillatory shear. The bulk elastic strain $\gamma$ is therefore taken to vary between $\gamma_{\max}$ and $-\gamma_{\max}$, with the system initially prepared at $\gamma = 0$. A single cycle is defined as the complete evolution of the imposed strain from $\gamma=0$ to $\gammaMax$, back to $0$, then to $-\gammaMax$, and finally returning to $0$. 

Since we impose the condition that the material remains in force balance at all times and that an element yields instantaneously once its local threshold is exceeded, the dynamics correspond to the athermal limit of zero temperature. Accordingly, the protocol studied here is that of \quotes{athermal quasi-static shear} (AQS), as considered widely in the literature on sheared amorphous materials \cite{liuCreepDynamicsAthermal2018,barlow_ductile_2020,pollard_yielding_2022,maloney_amorphous_2006,Popovi2018,ness_oscillatory_2017,Fiocco2013,Ruscher2021,khandareElastoplasticModellingCyclic2026}. As a result, the global deformation of the material is assumed to shear in small steps, stopping when any element yields. After each strain increment, the resulting unstable element, now just above threshold, is then yielded and force balance is reimposed. This may lead to further elements becoming unstable. Any such elements are then yielded, and force balance is reimposed. This procedure continues until no elements remain unstable and the avalanche has ceased. The system then continues to globally shear. Numerically, each elastic strain step is determined by the increment $\Delta\gamma$ required to bring the least stable element to its yielding threshold $\lc$. Since an elastic loading step shifts all local strains uniformly as $l_i \to l_i + \Delta\gamma$, this increment is given by
\begin{equation*}
\Delta \gamma =
\begin{cases}
\displaystyle \min_i \bigl( |\lc| - l_i \bigr),
& \text{for increasing $\gamma$}, \\[6pt]
\displaystyle \max_i \bigl( -|\lc| - l_i \bigr),
& \text{for decreasing $\gamma$}.
\end{cases}
\end{equation*}
The two cases correspond to shearing in the positive direction ($0 \to \gamma_{\max}$) and negative direction ($0 \to -\gamma_{\max}$), respectively; in each case, the selected increment is the value closest to zero, corresponding to the first element to reach its yielding threshold (the minimum of $|\lc|-l_i$ for increasing $\gamma$, and the maximum of $-|\lc|-l_i$ for decreasing $\gamma$).

The macroscopic stress of the system is defined as the average over the lattice of the local element stresses,
\begin{equation}
    \Sigma = \frac{\mu}{N^2}\sum_{i=1}^{N}\sum_{j=1}^{N} l_{ij}.
\end{equation}

As discussed previously, the large-cycle response of shear-oscillated amorphous solids is either to reach an absorbing state in which no further plastic events occur or to undergo a discontinuous transition to a yielded state. In the model described above, the absorbing state is realised when all elements have values of $\ltilde$ sufficiently close to zero that even the maximum imposed elastic strain, $\pm\gammaMax$, does not drive any element beyond its yielding threshold. Mathematically, this condition requires $|\ltilde|<1-\gammaMax$
for every element. We will refer to the region $-1+\gammaMax < \ltilde < 1-\gammaMax$ as the absorbed region, within which elements remain stable under the imposed oscillatory strain. Alternatively, the sample can undergo the formation of a shear band, which spreads quickly across the sample, and the sample fails suddenly. Within our models, these bands appear as localised regions of high plasticity when compared to the bulk of the material. Consequently, we detect the formation of these events by comparing the amount of plastic yielding across streamlines (rows or columns of our lattice) in our material. We compute the standard deviation of the number of yielded sites across streamlines, and if this exceeds a threshold value of $5 \times 10^{-2}$, we classify the system as having undergone macroscopic yielding. This provides a definition of macroscopic yielding that separates diffuse plastic activity from a clearly banded state; the results below should therefore be interpreted with this definition in mind. We will refer to this state as the yielded state, in which the system has undergone a transition to a higher-energy state characterised by persistent plastic activity and irreversible structural evolution. 

In \cref{fig:YieldCountSnapshots}, we plot maps of two representative systems: the top row shows a sample that evolves into an absorbing state, whereas the bottom row shows a sample that undergoes macroscopic yielding via the formation of a system-spanning shear band. In both cases, we observe only limited diffuse damage with no clear visual distinction between the samples until the final cycle, at which point a pronounced shear band emerges in the yielding case.

\begin{figure}
    \centering
    \includegraphics[width=0.99\linewidth]{chapters/Fatigue/Figures/YieldCountSnapshots.pdf}
    \caption[Maps of yielded sites in an absorbing and yielded system.]{Maps of sites that have yielded at least once (black) and sites that have never yielded (white), shown at the end of cycles 1, 4, 16, 64, and 256 (left to right). \newline {\bf Top row:} an absorbing state. {\bf Bottom row:} a yielding state. \newline Both systems reach their respective end states during cycle 256.  Each map has system size $N_x \times N_y=256\times256$, $l_0=0.075$, $\gamma_0=0.835$, $\lw=0.02$}
    \label{fig:YieldCountSnapshots}
\end{figure}

Recall that we work in units where $\mu=1$ and $E=0.5$, so that the critical yielding strain is $\lc=1$. The scale of the imposed amplitude $\gammaMax$ is therefore measured relative to this threshold, and elements yield whenever $|l|>\lc$, irrespective of sign.

The remaining parameters are the degree of initial annealing, set by $l_0$ (with smaller $l_0$ corresponding to more highly annealed samples), the system size $N$, and the width $\lw$ of the post-hop distribution. Unless stated otherwise, we take $N=1024$ and $\lw=0.02$.

\section{Shear Startup}
As a starting point for the results presented later in this chapter, we first examine the behaviour of the model under shear startup (see \cref{sec:ShearStartup}). This protocol has been analysed in detail by Pollard and Fielding in Ref.~\cite{pollard_yielding_2022}; our aim here is therefore to focus only on those aspects most relevant to the analysis developed below, which also serves as a check that our implementation of the model is working correctly. Although the primary focus of this chapter is oscillatory shear, the startup protocol provides the simplest setting in which to identify the elastic regime, the onset of yielding, and the approach to steady flow under a controlled deformation. These same ingredients govern the material response within each half-cycle of an oscillatory protocol, so analysing shear startup first isolates the basic mechanisms of stress build-up, yielding, and relaxation that will recur in the cyclic case.

\begin{figure}[b!]
    \centering
    \includegraphics[width=0.99\linewidth]{chapters/Fatigue/Figures/pollard.pdf}
    \caption[Stress-strain curves of the lattice elastoplastic model under shear startup.]{Stress-strain curves of the lattice elastoplastic model under shear startup. {\bf a)} Responses for a lattice size $N \times N = 1024 \times 1024$, with decreasing levels of sample annealing in curves downwards: $l_0=0.025, 0.05, 0.075, \ldots, 0.2$. {\bf b)} Finite-size dependence for two different annealing levels $l_0=0.05$ (dashed curves) and $l_0=0.125$ (solid curves) for system sizes $N=64, 128, 256, 512, 1024, 2048, 4096$.}
    \label{fig:ShearStarup}
\end{figure}

In \cref{fig:ShearStarup}a, we show stress-strain curves for a system of size $N \times N = 1024 \times 1024$ for various annealing levels. As can be seen, for highly annealed samples  (low $l_0$), the stress is initially linear in strain, corresponding to elastic loading of the elements. The stress then reaches a pronounced overshoot before dropping sharply towards a steady-state flow value. At low values of $l_0$, the distribution of element strains $l$ is narrow, and therefore a larger fraction of elements lie close to their local yield threshold at the point of the first element yielding. The stress redistributed by this initial plastic event is therefore sufficient to trigger further yielding events in rapid succession, promoting a collective cascade and an abrupt macroscopic stress drop. As $l_0$ increases, the height of this overshoot decreases with the system tending towards ductile behaviour. This correlates with the widening of the initial distribution, forcing the first yielding event to occur at a lower strain. Although for small enough $l_0$, the distribution is narrow enough that a small number of initial failures can trigger cascades of yielding events.

For less well-annealed samples (larger $l_0$), the overshoot becomes progressively smaller, and the subsequent stress drop more gradual. The failure of individual samples illustrated here occurs through a cumulative sequence of smaller discontinuous stress drops for the weakest annealed samples. However, Ref.~\cite{pollard_yielding_2022} shows that, when averaged over many runs with different initial random seeds, the decline is continuous. In these less well-annealed samples, plastic activity is initiated earlier in the deformation, because a broader initial distribution of local strains brings a finite fraction of elements close to threshold even at a small applied strain. As a result, stress relaxation is distributed over a wider strain interval, and the macroscopic response appears more ductile.

The role of the system size is illustrated in \cref{fig:ShearStarup}b. We consider two cases: a highly annealed sample ($l_0=0.05$) and a less annealed sample ($l_0=0.125$). In the highly annealed case, failure is already sharp, even at modest system sizes, but the peak strain shifts to lower values as $N$ increases. In the less annealed case ($l_0=0.125$), small systems display a more gradual relaxation composed of multiple smaller stress drops, whereas larger systems exhibit a progressively steeper decline \DIFdelbegin \DIFdel{that converges towards a single dominant drop. Increasing }\DIFdelend \DIFaddbegin \DIFadd{over the range of sizes studied. Across the finite system sizes considered here, increasing }\DIFaddend $N$ \DIFdelbegin \DIFdel{enhances the likelihood that a local plastic rearrangement triggers a system-spanning cascade. In the large-system limit (}\DIFdelend \DIFaddbegin \DIFadd{appears to sharpen the stress drop and shift the position of the stress maximum. Whether the responses under either oscillatory shear or shear startup converge to well-defined limiting curves as }\DIFaddend $N\to\infty$ \DIFdelbegin \DIFdel{), a finite stress overshoot signals the onset of a collective failure mode rather than a sequence of isolated plastic relaxations}\DIFdelend \DIFaddbegin \DIFadd{remains an open question. Because the present simulations are computationally limited to finite system sizes, these observations should be interpreted as finite-size trends rather than evidence for any particular thermodynamic-limit response}\DIFaddend .

While the work of Pollard and Fielding went on to characterise in detail the formation of shear bands and the subsequent approach to steady flow following failure \cite{pollard_yielding_2022}, the focus here is different. In the oscillatory shear protocol considered below, we do not resolve the post-banding dynamics once a shear band has formed. Instead, we restrict attention to the onset of instability itself, classifying trajectories according to whether they evolve towards failure or are absorbed into a stable absorbing state. This differs from the main emphasis of Khandare and Sastry in Ref.~\cite{khandareElastoplasticModellingCyclic2026}, who explicitly go on to analyse the spatio-temporal dynamics within the shear-banded regime.

\section{Oscillatory shear results}
We begin by outlining the basic physics that governs a sample’s response under the oscillatory shear protocol defined in \cref{sec:fatigueMethodology}. As the shear strain is cycled, plastic activity accumulates through successive local yielding events (up until an absorbing state is reached). 
Stress is then redistributed via the Eshelby propagator after each yielding event, therefore modifying the local strain field, altering the stability of neighbouring elements. 

Over repeated cycles, this accumulated plastic deformation can evolve in two qualitatively distinct ways. In one scenario, the redistributed stresses progressively organise into a correlated structure that ultimately forms a system-spanning shear band, leading to macroscopic failure. In the other, plastic events reorganise the system into a configuration in which all elements remain below threshold throughout subsequent cycles. In this latter case, the dynamics terminate in an absorbed state, with no further yielding possible under further cyclic shearing at the same strain amplitude.

Which of these outcomes occurs depends strongly on the imposed strain amplitude $\gammaMax$ relative to the width of the initial strain distribution, characterised by the annealing parameter $l_0$.
At sufficiently small amplitudes, it is possible for all elements $(i,j)$ in a finite system to satisfy
\[
    -1+\gammaMax < \ltilde_{ij} < 1-\gammaMax
\]
even in the initial configuration. 
In this limit, no element is driven beyond its local yield threshold during the first cycle or any subsequent cycle, so no plastic activity occurs and the system is immediately classified as absorbing.

However, after force balance is imposed, the initial distribution of local strain values is a Gaussian distribution of width $l_0\Norm$. In the thermodynamic limit $N \to \infty$, the extreme tails of the initial distribution are inevitably sampled. Although elements with very large initial strains $|l|>1$ would already have yielded during preparation and are therefore truncated, the probability of finding elements arbitrarily close to threshold $l_c = \pm 1$ remains finite. Consequently, in sufficiently large systems, near-critical elements are statistically unavoidable, so even at small imposed strain amplitudes the dynamics generally involve plastic activity, at least during the early cycles. \DIFaddbegin \DIFadd{The Gaussian form of the initial distribution is not essential here: its width determines the fraction of elements whose local strain $l$ lies sufficiently close to $l_c$ for the applied deformation to drive them beyond threshold during a cycle. Thus, the tails influence the initial level of plastic activity by setting the number of elements capable of yielding, rather than governing their subsequent dynamics.
}\DIFaddend 

At the opposite extreme, when the imposed amplitude $\gammaMax$ is large, the applied deformation immediately drives elements beyond their yield threshold within the first cycle. In this limit, the dynamics closely resemble those of shear start-up: plastic activity rapidly organises during monotonic deformation, damage accumulates cooperatively, and a shear band forms within the first cycle, leading to immediate failure. While in the large-$\gammaMax$ limit, this failure would be expected to occur during the first half cycle, $\gamma \to \gammaMax$, for intermediate values of $\gammaMax$ it may instead arise from statistical sampling of the initial distribution of $l$ during the second half-cycle, $\gamma \to -\gammaMax$.

Between these extremes lies an intermediate regime that is the primary focus of this work. In this regime, the system neither \DIFaddbegin \DIFadd{immediately }\DIFaddend attains an absorbing state nor fails immediately in the first cycle. Instead, a small amount of plastic activity occurs during each cycle, allowing the system to evolve gradually over many cycles. These trajectories can ultimately end either in a yielded state, via the formation of a shear band, or in an absorbing state in which all elements remain below threshold, with the outcome depending strongly on the imposed strain amplitude
$\gammaMax$.

\begin{figure}[p!]
    \centering
	\includegraphics[width=0.99\textwidth]{chapters/Fatigue/Figures/EndStates.pdf}
    \caption[Distribution of simulation end states as a function of imposed strain $\gammaMax$.]{Distribution of simulation end states as a function of imposed strain. For each strain amplitude on the $x$-axis, stacked bars show the fraction of realisations classified according to two binary criteria: whether failure occurs on the first cycle (immediate) or at later cycles (delayed), and whether the system ultimately becomes absorbed or yields. Results are shown for five values of the annealing parameter $l_0$ (arranged vertically), with smaller $l_0$ corresponding to better-annealed samples. Each bar represents the statistical outcome over all simulations at the corresponding $(\gamma, l_0)$ pair. \DIFdelbeginFL \DIFdelFL{For each value of $\gamma_0$ and $l_0$, the distributions are made up of at least $10{,}000$ samples}\DIFdelendFL \DIFaddbeginFL \DIFaddFL{Each parameter combination contains more than $2{,}500$ independent realisations}\DIFaddendFL , each \DIFdelbeginFL \DIFdelFL{sample }\DIFdelendFL \DIFaddbeginFL \DIFaddFL{generated }\DIFaddendFL with a different \DIFdelbeginFL \DIFdelFL{value of the random number }\DIFdelendFL \DIFaddbeginFL \DIFaddFL{random-number }\DIFaddendFL seed\DIFaddbeginFL \DIFaddFL{; the exact ensemble sizes are reported in }\cref{sec:FatigueSampleCounts}\DIFaddendFL .}
    \label{fig:EndStates}
\end{figure}

The distribution of these possible outcomes is summarised in \cref{fig:EndStates}, where we show the probability of each end state as a function of the imposed strain amplitude $\gamma_0$ for different levels of annealing. For both yielding and absorbing outcomes, we further distinguish between events that occur during the first cycle (“immediate”) and those that occur only after many cycles (“delayed”). Each bar, therefore, represents the statistical outcome of many independent realisations of the same $(\gamma_0,l_0)$ pair. \DIFaddbegin \DIFadd{Each parameter combination used in }\cref{fig:EndStates,fig:LogisticsL0,fig:FailureCyclel0} \DIFadd{was sampled with more than $2{,}500$ independent realisations. The precise ensemble size varies slightly across the sampled parameter values; the exact numbers used for each dataset are reported in }\cref{sec:FatigueSampleCounts}\DIFadd{.
}\DIFaddend 

Following the qualitative overview of possible outcomes in \cref{fig:EndStates}, we now quantify how the probability of reaching an absorbing state depends on the imposed strain amplitude and the level of annealing. In particular, we measure the fraction of simulations that terminate in an absorbing state, $P_{\mathrm{abs}}$, for a range of values of $\gamma_0$ and $l_0$. 

For each pair $(\gamma_0,l_0)$, we perform an ensemble of independent simulations, each with a different value of the random number seed, and classify the final state of each trajectory as either absorbing (brown and orange in \cref{fig:EndStates}) or yielding (blue and purple in \cref{fig:EndStates}). The resulting probability of absorption $P_{\mathrm{abs}}$ is then defined as the fraction of trajectories that reach an absorbing state before yielding. This provides a natural parameter for quantifying the transition between the two dynamical regimes. The resulting dependence of $P_{\mathrm{abs}}$ on $\gamma_0$ is shown in \cref{fig:LogisticsL0} for several values of the annealing parameter $l_0$. For each value of $l_0$, the probability of absorption decreases smoothly from unity at small strain amplitudes to zero at larger amplitudes. This crossover marks the coexistence region briefly discussed in Ref.~\cite{khandareElastoplasticModellingCyclic2026}, which lies between a regime in which trajectories almost always organise into absorbing states and one in which yielding becomes overwhelmingly likely. 

To characterise this transition more quantitatively, we fit the data for $P_{\mathrm{abs}}(\gamma_0)$ using a logistic function of the form
\[
    P_{\mathrm{abs}}(\gamma_0)=\frac{1}{1+\exp\!\left[k(\gamma_0-\gamma_{0,1/2})\right]}.
\]
In Ref.~\cite{khandareElastoplasticModellingCyclic2026}, this quantity is referred to as the yielding amplitude; however, yielding can still occur for imposed strain amplitudes below this value, while absorption can occur for amplitudes above it. The parameter $k$ then characterises the width of the crossover region.

\begin{figure}[t]
    \centering
	\includegraphics[width=0.99\linewidth]{chapters/Fatigue/Figures/logistics_l0.pdf}
    \caption[Logistic fits for absorption probability for various levels of annealing]{
    {\bf a} Absorption probability $P_{\mathrm{abs}}$ as a function of imposed shear amplitude $\gammaMax$ for increasing levels of annealing (curves from left to right: $l_0=0.025,0.0375,0.05,0.075,0.1,0.125$). For each $l_0$, the data are fitted with the logistic form $P_{\mathrm{abs}}(\gamma_0)=1/\left(1+\exp[k(\gamma_0-\gamma_{0,1/2})]\right)$, shown as dotted lines. \newline
    {\bf b} Fitted midpoint $\gamma_{0,1/2}$, defined by $P_{\mathrm{abs}}=0.5$, as a function of $l_0$, with a power-law guide $\gamma_{0,1/2}\propto l_0^{-0.10}$. \newline
    {\bf c} Fitted transition width $w$, defined as the strain interval between $P_{\mathrm{abs}}=0.01$ and $0.99$, giving $w=\ln(99)/|k|$, as a function of $l_0$. \DIFaddbeginFL \DIFaddFL{Each parameter combination contains more than $2{,}500$ independent realisations; exact ensemble sizes are reported in }\cref{sec:FatigueSampleCounts}\DIFaddFL{.}\DIFaddendFL }
    \label{fig:LogisticsL0}
\end{figure}

As shown in \cref{fig:LogisticsL0}, decreasing the annealing parameter $l_0$ shifts the transition to larger strain amplitudes, indicating that better-annealed systems are more resistant to yielding under cyclic deformation. The fitted midpoint $\gamma_{0,1/2}$ is plotted as a function of the annealing level $l_0$ (with stronger annealing corresponding to smaller $l_0$) in \cref{fig:LogisticsL0}b. Over the range of $l_0$ studied here, the data are consistent with a weak power-law dependence, $\gamma_{0,1/2} \propto l_0^{-0.10}$. Although this trend suggests that increasingly poorly annealed samples absorb only at smaller imposed strain amplitudes, the present data do not establish the asymptotic behaviour of $\gamma_{0,1/2}$ at large $l_0$. Interestingly, the transition width $w$ (defined as the strain interval over which $P_{\mathrm{abs}}$ decreases from $0.99$ to $0.01$) does not appear to depend strongly on the level of annealing $l_0$. However, as shown in \cref{fig:EndStates}, the character of the transition changes: it evolves from one between immediate absorption and immediate yielding to one between delayed absorption and delayed yielding.

\begin{figure}[t]
    \centering
    \includegraphics[width=1.0\textwidth]{chapters/Fatigue/Figures/Decay.pdf}
    \caption[Final cycle $\CStar$ as a function of imposed strain amplitude $\gammaMax$ for various annealing levels.]{Final cycle $\CStar$ as a function of the imposed strain amplitude $\gammaMax$ for various annealing levels $l_0$. Smaller values of $l_0$ correspond to better-annealed samples. We report the $95^\mathrm{th}$ percentile, $\CStar_{95\%}$, defined as the cycle at which the survival distributions in \cref{fig:SurvivalDistrbutions} intersect the red threshold line. Results \DIFaddbeginFL \DIFaddFL{conditioned on yielding and absorption }\DIFaddendFL are \DIFdelbeginFL \DIFdelFL{shown for systems that yield, absorb, or exhibit both behaviours, }\DIFdelendFL indicated by solid \DIFdelbeginFL \DIFdelFL{, dotted, }\DIFdelendFL and \DIFdelbeginFL \DIFdelFL{dashed }\DIFdelendFL \DIFaddbeginFL \DIFaddFL{dotted }\DIFaddendFL lines, respectively. \DIFdelbeginFL \DIFdelFL{For each pair of values for $\gamma_0$ and $l_0$}\DIFdelendFL \DIFaddbeginFL \DIFaddFL{Dashed lines labelled ``All'' pool all trajectories}\DIFaddendFL , \DIFdelbeginFL \DIFdelFL{the distributions are made up }\DIFdelendFL \DIFaddbeginFL \DIFaddFL{each }\DIFaddendFL of \DIFdelbeginFL \DIFdelFL{at least $10{,}000$ samples}\DIFdelendFL \DIFaddbeginFL \DIFaddFL{which ultimately either yields or reaches an absorbing state. Each parameter combination contains more than $2{,}500$ independent realisations}\DIFaddendFL , each \DIFdelbeginFL \DIFdelFL{run }\DIFdelendFL \DIFaddbeginFL \DIFaddFL{generated }\DIFaddendFL with a different \DIFdelbeginFL \DIFdelFL{value of the random number }\DIFdelendFL \DIFaddbeginFL \DIFaddFL{random-number }\DIFaddendFL seed\DIFaddbeginFL \DIFaddFL{; exact ensemble sizes are reported in }\cref{sec:FatigueSampleCounts}\DIFaddendFL . \DIFdelbeginFL \DIFdelFL{The curves }\DIFdelendFL \DIFaddbeginFL \DIFaddFL{Outcome-conditioned percentiles are }\DIFaddendFL shown \DIFdelbeginFL \DIFdelFL{in this plot include }\DIFdelendFL only \DIFdelbeginFL \DIFdelFL{those end states for which }\DIFdelendFL \DIFaddbeginFL \DIFaddFL{when }\DIFaddendFL more than \DIFdelbeginFL \DIFdelFL{$1000$ such samples }\DIFdelendFL \DIFaddbeginFL \DIFaddFL{$2{,}500$ realisations }\DIFaddendFL reach the corresponding \DIFdelbeginFL \DIFdelFL{outcome}\DIFdelendFL \DIFaddbeginFL \DIFaddFL{end state}\DIFaddendFL .}
    \label{fig:FailureCyclel0}
\end{figure}

Having identified the transition between absorbing and yielding behaviour through the probability $P_{\mathrm{abs}}$, we now turn to the dynamics of the system separately for the cases in which yielding or absorption arise after a delayed number of cycles. For each simulation, we define the termination cycle $\CStar$ as the number of cycles elapsed before the system either reaches an absorbing state or undergoes macroscopic yielding. \DIFaddbegin \DIFadd{In the figures below, ``All'' denotes the pooled ensemble of these two mutually exclusive outcomes: each trajectory ultimately either yields or reaches an absorbing state. }\DIFaddend The dependence of $\CStar$ on the level of annealing $l_0$ is shown in \cref{fig:FailureCyclel0}. Far from the transition, trajectories typically terminate after only a small number of cycles. For small strain amplitudes, the system rapidly relaxes to an absorbing state, while at sufficiently large amplitudes, yielding occurs almost immediately. In contrast, as the strain amplitude approaches the transition region, the number of cycles required to reach the final state increases dramatically.

This behaviour reflects the increasingly long-lived metastable dynamics that occur near the boundary between the two regimes. In this region, the system undergoes repeated plastic rearrangements over many cycles before eventually either organising into an absorbing state or nucleating a shear band that leads to failure. The strong growth in $\CStar$ is therefore consistent with critical-like slowing down across the transition regime. At the system sizes studied here, this regime remains broad, so we avoid identifying a unique critical point; however, the finite-size analysis below suggests that the transition may sharpen with increasing system size.

Such divergences in the number of cycles required to reach the final state have been widely observed in simulations of amorphous solids under oscillatory shear, as discussed previously. In particle simulations exploring this behaviour \cite{Fiocco2013,Regev_onset_2013,Leishangthem2017,kawasakiMacroscopicYieldingJammed2016}, a well-defined critical strain amplitude $\gamma_c$ is typically identified, with systems evolving to absorbing states for $\gammaMax < \gamma_c$ and yielding for $\gammaMax > \gamma_c$. In contrast, both this work and Ref.~\cite{khandareElastoplasticModellingCyclic2026} find a coexistence region in which both yielding and absorption occur. However, the presence of a finite region in which both outcomes are possible suggests that the transition may not be sharp, but instead occurs over a finite interval of strain amplitudes. Evidence for such mixed behaviour is not entirely unexpected: simulations and experiments have reported the coexistence of reversible and irreversible plastic events under oscillatory shear within a single sample \cite{LundbergReversiblePlasticEvents2008,KhirallahIntegerModel2021}.

Within this interpretation, the discrepancies reported in recent work between the power-law divergences obtained in particle simulations \cite{Maity2026,khandareElastoplasticModellingCyclic2026,kawasakiMacroscopicYieldingJammed2016}, elastoplastic models \cite{khandareElastoplasticModellingCyclic2026}, and those predicted by mean-field models \cite{parleyMeanFieldTheoryYielding2022,sarkar2025coarsegraineddescriptionsdynamics} may reflect the presence of this coexistence region. If the transition is a broad region rather than a sharp critical value, the apparent critical scaling will depend sensitively on how the transition point is defined and on which subset of trajectories, yielding or absorbing, is considered. Understanding this transition remains an open challenge for future work.

% Instead, we focus on the behaviour observed in the present model. 
In Ref.~\cite{khandareElastoplasticModellingCyclic2026}, the midpoint $\gamma_{0,1/2}$ was identified with the critical strain amplitude $\gamma_c$, and the subsequent analysis reported time-divergence data only outside the coexistence region, considering yielding cases for $\gammaMax > \gamma_{0,1/2}$ and absorbing cases for $\gammaMax < \gamma_{0,1/2}$. By contrast, our data in \cref{fig:FailureCyclel0} show that the number of cycles required for both absorption and yielding \DIFdelbegin \DIFdel{extends }\DIFdelend \DIFaddbegin \DIFadd{increases }\DIFaddend across the entire transition region. Consequently, this choice obscures part of the underlying physics of the transition.

Interestingly, \DIFdelbegin \DIFdel{the longest-lived trajectories tend to correspond to the least probableoutcomes at a given strain amplitude}\DIFdelend \DIFaddbegin \DIFadd{for both outcomes, the number of cycles to termination tends to increase as that outcome becomes less probable}\DIFaddend . For example, at strain amplitudes slightly above the midpoint of the transition $\gamma_{0,1/2}$, the probability of reaching an absorbing state becomes small. \DIFdelbegin \DIFdel{When, in }\DIFdelend \DIFaddbegin \DIFadd{In }\DIFaddend the rare cases \DIFdelbegin \DIFdel{, }\DIFdelend \DIFaddbegin \DIFadd{where }\DIFaddend absorption does occur in this regime, it \DIFdelbegin \DIFdel{therefore requires an unusually }\DIFdelend \DIFaddbegin \DIFadd{usually requires a }\DIFaddend long sequence of plastic rearrangements before the system eventually organises into a configuration that avoids further yielding. \DIFdelbegin \DIFdel{As a result, these rare absorbing trajectories are associated with the largest values of $\CStar$}\DIFdelend \DIFaddbegin \DIFadd{Consequently, the absorbing trajectories become increasingly long-lived as the strain amplitude is increased through the transition region}\DIFaddend . Conversely, \DIFdelbegin \DIFdel{for amplitudes slightly below }\DIFdelend \DIFaddbegin \DIFadd{below the midpoint}\DIFaddend , yielding becomes \DIFdelbegin \DIFdel{the less probableoutcome. In these cases, }\DIFdelend \DIFaddbegin \DIFadd{progressively less probable, and }\DIFaddend the trajectories that eventually fail \DIFdelbegin \DIFdel{are typically those that persist for many cycles while gradually }\DIFdelend \DIFaddbegin \DIFadd{tend to persist for an increasing number of cycles while }\DIFaddend accumulating damage before a shear band nucleates. \DIFaddbegin \DIFadd{Although the yielding trajectories remain shorter-lived than the absorbing trajectories overall, their characteristic lifetime also increases as yielding becomes rarer, in this case as $\gamma_0$ is decreased.
}

%DIF >  Interestingly, the longest-lived trajectories tend to correspond to the least probable outcomes at a given strain amplitude. For example, at strain amplitudes slightly above the midpoint of the transition $\gamma_{0,1/2}$, the probability of reaching an absorbing state becomes small. When, in the rare cases, absorption does occur in this regime, it therefore requires an unusually long sequence of plastic rearrangements before the system eventually organises into a configuration that avoids further yielding. As a result, these rare absorbing trajectories are associated with the largest values of $\CStar$. Conversely, for amplitudes slightly below, yielding becomes the less probable outcome. In these cases, the trajectories that eventually fail are typically those that persist for many cycles while gradually accumulating damage before a shear band nucleates. 
\DIFaddend 

As a consequence, sampling such rare events becomes increasingly challenging, particularly given the sensitivity of the dynamics to the initial strain configuration. To obtain reliable statistics, we therefore sample \DIFdelbegin \DIFdel{at least $10{,}000$ }\DIFdelend \DIFaddbegin \DIFadd{more than $2{,}500$ }\DIFaddend realisations for each pair of parameters ($l_0$, $\gammaMax$)\DIFaddbegin \DIFadd{, with exact ensemble sizes reported in }\cref{sec:FatigueSampleCounts}\DIFaddend . For strain amplitudes within the transition region identified in \cref{fig:LogisticsL0}, additional samples were generated to construct the survival functions shown in \cref{fig:SurvivalDistrbutions}. \DIFdelbegin \DIFdel{These survival distributions are computed }\DIFdelend \DIFaddbegin \DIFadd{We define the survival function as
}\[
\DIFadd{S(C)=1-F(C),
}\]
\DIFadd{where $F(C)=P(C^*\leq C)$ is the cumulative distribution of the termination cycle $C^*$. Accordingly, $S(C)=P(C^*>C)$ is the probability that a trajectory remains active beyond cycle $C$. We calculate these survival functions }\DIFaddend separately for trajectories that terminate in absorbing and yielding end states.

The value of $\CycleNineFive$, used in \cref{fig:FailureCyclel0} to quantify the number of cycles required for the sample to reach an end state, is defined as the value of $C^*$ at the $95^\textrm{th}$ percentile of the distribution of failure cycles across the ensemble. It is reported only when the corresponding end state is represented by more than $1{,}000$ realisations in the sampled ensemble. This threshold ensures that the reported statistics are not dominated by rare events with insufficient sampling. Accordingly, near the edges of the coexistence region, where yielding or absorption occurs only in a minority of samples relative to the competing end state, the reported curves are based on fewer contributing realisations than in regimes where all sampled seeds converge to the same outcome. Nevertheless, the underlying behaviours are robust.

\begin{figure}[p]
    \centering
    \includegraphics[width=1.0\textwidth]{chapters/Fatigue/Figures/Survival.pdf}
    \caption[Survival functions of the termination cycle $\CStar$]{Survival functions of the termination cycle $\CStar$, showing the fraction of simulations, for simulations performed at many values of random number seed, that remain active up to cycle $\CStar$ before reaching an absorbing state ({\bf Left}) or a yielding state ({\bf Right}).
    The survival probability $S(\CStar)=1-F(\CStar)$, resolved by end state and annealing parameter. From top to bottom, rows correspond to increasing annealing parameter $l_0 = 0.025, 0.05, 0.075, 0.1$. Within each panel, multiple curves represent different imposed strain amplitudes, specified in the right column for each $l_0$. Decay constants $\lambda$ obtained from exponential fits to the tails of the survival functions are shown in \cref{fig:SurvivalDistrbutionsFits}.}
    \label{fig:SurvivalDistrbutions}
\end{figure}

\DIFdelbegin \DIFdel{To examine the dynamics of these trajectories in more detail, we now consider the full distribution of the number of cycles to reach either a yielded or absorbed end state $\CStar$. In }\DIFdelend \DIFaddbegin \DIFadd{The resulting survival functions are shown in }\DIFaddend \cref{fig:SurvivalDistrbutions}\DIFdelbegin \DIFdel{we plot the survival function $S(\CStar)=1-F(\CStar)$, where $F(\CStar)$ is the cumulative distribution of $\CStar$. This survival function represents the probability that a trajectory remains active for at least $\CStar$ cycles before reaching its final state. We plot these distributions }\DIFdelend \DIFaddbegin \DIFadd{, }\DIFaddend separately for absorbing and yielding trajectories, allowing the \DIFaddbegin \DIFadd{termination }\DIFaddend statistics of the two end states to be compared directly. \DIFdelbegin %DIFDELCMD < 

%DIFDELCMD < %%%
\DIFdelend The survival distributions exhibit two distinct regimes. At small values of $\CStar$, many curves show a rapid initial drop, reflecting the significant fraction of trajectories that terminate within the first cycle. These correspond to cases in which immediate yielding or absorption \DIFdelbegin \DIFdel{occur in }\DIFdelend \DIFaddbegin \DIFadd{occurs during }\DIFaddend the first cycle, as identified in \cref{fig:EndStates}. The magnitude of this initial drop \DIFdelbegin \DIFdel{, therefore , }\DIFdelend \DIFaddbegin \DIFadd{therefore }\DIFaddend depends strongly on both the strain amplitude and the level of annealing, which \DIFdelbegin \DIFdel{control }\DIFdelend \DIFaddbegin \DIFadd{determine }\DIFaddend the probability of \DIFdelbegin \DIFdel{such }\DIFdelend \DIFaddbegin \DIFadd{these }\DIFaddend immediate end states.

Beyond this initial regime, the survival curves transition to a much more gradual decay. These samples correspond to delayed dynamics in which plastic activity persists over many cycles before the system ultimately reaches its final state. As indicated in \cref{fig:SurvivalDistrbutions}, the tails of the distribution, ignoring the immediate drop or the apparent initial plateau (arising from the absence of trajectories that end at small values of $\CStar$) in $S(\CStar)$, are well described by an exponential form,
\[
    S(\CStar) \sim \exp(-\lambda \CStar).
\]
The exponential form of the tail of $S(\CStar)$ suggests that, once the initial transient regime has passed, the process loses memory of how long the system has already survived. In this regime, the survival probability decreases by approximately the same factor with each additional cycle. Consequently, among the samples that have not yet failed, roughly the same fraction fails during each subsequent cycle.

\begin{figure}
    \centering
    \includegraphics[width=0.99\linewidth]{chapters/Fatigue/Figures/Survival_TailFits_DecayConstant.pdf}
    \caption[Decay constants $\lambda$ characterising the exponential tails of the survival functions shown in \cref{fig:SurvivalDistrbutions}.]{Decay constants $\lambda$ characterising the exponential tails of the survival functions shown in \cref{fig:SurvivalDistrbutions}. For each annealing parameter $l_0$ and strain amplitude, $\lambda$ is obtained by fitting the large-$\CStar$ regime to 
    $S(\CStar) \sim \exp(-\lambda \CStar)$. Results are shown separately for absorbing and yielding trajectories.}
    \label{fig:SurvivalDistrbutionsFits}
\end{figure}

The values of the decay constant $\lambda$ extracted from fits to the tails of the survival functions are shown in \cref{fig:SurvivalDistrbutionsFits}. For absorbing trajectories, the decay constants collapse onto a common curve when plotted as a function of the imposed strain amplitude. This indicates that the late-time dynamics are controlled primarily by the distance to the absorbing boundary set by $\gammaMax$, and are largely insensitive to the initial annealing once immediate absorption events are excluded, consistent with an effective first-passage picture, as discussed later.

In contrast, the yielding trajectories do not collapse when plotted as a function of $\gammaMax$, and exhibit a significant spread between different annealing levels. This suggests that the dynamics leading to yielding are not governed by a single characteristic scale set by the imposed strain amplitude, but instead depend on more complex, collective processes associated with shear-band formation.

Taken together, these results indicate that the transition between absorbing and yielding behaviour under cyclic deformation is governed by both the probability of sampling each end state and the number of cycles over which the corresponding trajectories take to evolve to that end state. The analysis done in \cref{fig:LogisticsL0,fig:EndStates} shows that the likelihood of reaching an absorbing state decreases smoothly with increasing strain amplitude, with the midpoint of this transition shifting systematically with the level of annealing. At the same time, the number of cycles to absorb $\CStar$ reveals that trajectories near this transition can persist for many cycles before terminating, leading to broad distributions of cycles $\CStar$. The longest delays occur for trajectories that realise the least probable outcome at a given strain amplitude, reflecting the need for an unusually long sequence of plastic rearrangements to occur before the system reaches an absorbing configuration or rearranges into a configuration which nucleates a shear band. These observations are consistent with previous studies of cyclic deformation in amorphous solids, which report a strong growth in the number of cycles for the sample to absorb or yield as the yielding transition is approached from either side.

\section{Robustness to \texorpdfstring{$\lw$ and $N$}{lw and N}} \label{sec:FatigueRobustness}

So far, we have considered a fixed system size $N$ and post-hop distribution $\lw$. However, it is well established that the system size plays a central role in the transition from ductile to brittle behaviour. In particular, Pollard and Fielding in Ref.~\cite{pollard_yielding_2022} demonstrated that increasing $N$ leads to increasingly abrupt, brittle-like stress drops, whereas smaller systems can exhibit behaviour that appears more ductile. As discussed previously, such system-size effects are therefore also expected to influence the fatigue behaviour investigated here.

\begin{figure}[b!]
    \centering
    \includegraphics[width=0.99\textwidth]{chapters/Fatigue/Figures/lw.pdf}
    \caption[Effect of changing the width $\lw$ of the post-hop distribution, with the annealing parameter fixed at $l_0 = 0.075$ and system size $N=1024$.]{Effect of changing the width $\lw$ of the post-hop distribution, with the annealing parameter fixed at $l_0 = 0.075$ and system size $N=1024$. \newline
    {\bf a)} The $95^\mathrm{th}$ percentile failure cycle $\CStar_{95\%}$ as a function of imposed strain $\gamma_0$. Results are shown for $\lw=0.01, 0.02, 0.03, 0.04$\DIFaddbeginFL \DIFaddFL{. Solid }\DIFaddendFL and \DIFdelbeginFL \DIFdelFL{for }\DIFdelendFL \DIFaddbeginFL \DIFaddFL{dotted lines are conditioned on }\DIFaddendFL trajectories that yield \DIFdelbeginFL \DIFdelFL{, }\DIFdelendFL \DIFaddbeginFL \DIFaddFL{and }\DIFaddendFL absorb, \DIFdelbeginFL \DIFdelFL{or exhibit either of these behaviours, indicated by solid, dotted}\DIFdelendFL \DIFaddbeginFL \DIFaddFL{respectively}\DIFaddendFL , \DIFdelbeginFL \DIFdelFL{and }\DIFdelendFL \DIFaddbeginFL \DIFaddFL{while }\DIFaddendFL dashed lines \DIFdelbeginFL \DIFdelFL{, respectively}\DIFdelendFL \DIFaddbeginFL \DIFaddFL{labelled ``All'' pool trajectories with either end state}\DIFaddendFL . The inset shows a rescaling of the data showing the failure cycle scales approximately as the square of the maximum strain. \newline
    {\bf b)} The corresponding distribution of end states: delayed yielding (blue), immediate yielding (purple), delayed absorption (orange), and immediate absorption (brown). \newline
    \DIFdelbeginFL \DIFdelFL{For each value of $\gamma_0$ and $\lw$, the distributions are constructed from at least $10{,}000$ }\DIFdelendFL \DIFaddbeginFL \DIFaddFL{Each parameter combination contains more than $2{,}500$ }\DIFaddendFL independent \DIFdelbeginFL \DIFdelFL{runs }\DIFdelendFL \DIFaddbeginFL \DIFaddFL{realisations generated }\DIFaddendFL with different \DIFdelbeginFL \DIFdelFL{random number }\DIFdelendFL \DIFaddbeginFL \DIFaddFL{random-number }\DIFaddendFL seeds\DIFaddbeginFL \DIFaddFL{; exact ensemble sizes are reported in }\cref{sec:FatigueSampleCounts}\DIFaddendFL . Quantiles are shown only for cases with more than $1{,}000$ samples.}
    \label{fig:FailureCyclelw}
\end{figure}

First, we consider the effect of the changing post-hop distribution, characterised by its width $\lw$. In the results presented so far, this width has been fixed at $\lw = 0.02$. Varying $\lw$ does not alter the eventual end state of the simulation, but it does control the number of cycles over which the system reaches that state. This follows from the definition of the post-hop distribution introduced in \cref{sec:fatigueMethodology}, where $\lw$ sets the characteristic magnitude of the strain increments following a yielding event.

This behaviour is illustrated in \cref{fig:FailureCyclelw}, where we show the number of cycles required to reach the final state for a fixed initial disorder $l_0 = 0.075$ and system size $N = 1024$, for several values of $\lw$. As shown in panel b, the fraction of samples reaching each end state remains approximately unchanged across the range $\lw=0.01-0.04$, indicating that, within this range, the post-hop width does not influence the selection of the final state. In contrast, the number of cycles required to reach these end states depends strongly on $\lw$: as $\lw$ decreases, trajectories require progressively more cycles before terminating.

The data further reveal a simple scaling relation between the characteristic cycle number and the post-hop width. When plotted in terms of $\sqrt{\CycleNineFive}\lw$, the curves approximately collapse up to a constant factor, implying the scaling $\CycleNineFive \propto \lw^{-2}$. This dependence is consistent with a diffusive first-passage process, in which $\lw$ sets the characteristic step size of the walk while $\gammaMax$ plays the role of a threshold displacement \cite{Redner_2001}. We return to discuss this interpretation later in the chapter in \cref{sec:DilluteModel}.

We now examine the effect of varying the system size $N$, while keeping the annealing level fixed at $l_0 = 0.075$, corresponding to an intermediate degree of annealing. The post-hop distribution width is fixed at $\lw=0.02$, and the system size is varied logarithmically, increasing in powers of two: $64, 128, 256, 512, 1024, 2048$.

\begin{figure}[t!]
    \centering
    \includegraphics[width=0.99\textwidth]{chapters/Fatigue/Figures/N.pdf}
    \caption[Effect of varying system size $N$ at fixed parameters $l_0 = 0.075$ and $\lw = 0.02$.]{Effect of system size $N$ at fixed parameters $l_0 = 0.075$ and $\lw = 0.02$. \newline
    {\bf a)} The $95^\mathrm{th}$ percentile failure cycle $\CStar_{95\%}$ as a function of imposed strain. Results are shown for system sizes $N = 64, 128, 256, 512, 1024, 2048$, indicated in blue, orange, green, and red, purple and brown, respectively (top to bottom). The inset shows the curves shifted to align the peak of $C^*_{95\%}$, evaluated \DIFdelbeginFL \DIFdelFL{across all samples and end states}\DIFdelendFL \DIFaddbeginFL \DIFaddFL{using the pooled set of trajectories that either yield or absorb}\DIFaddendFL , at $\gamma_0 - \gamma_{0,\mathrm{peak}} = 0$. \newline
    {\bf b)}  shows the corresponding value of $\CStar_{95\%}$ at $\gammaMax=0.799$, the closest sampled point to $0.8$ in our dataset, as a function of system size $N$. \newline
    {\bf c)} The curves in {\bf a)} shifted to align the peak of $C^*_{95\%}$, evaluated \DIFdelbeginFL \DIFdelFL{across all samples and end states}\DIFdelendFL \DIFaddbeginFL \DIFaddFL{using the pooled set of trajectories that either yield or absorb}\DIFaddendFL , at $\gamma_0 - \gamma_{0,\mathrm{peak}} = 0$.
    {\bf d)} The corresponding distribution of end states: delayed yielding (blue), immediate yielding (purple), delayed absorption (orange), and immediate absorption (brown). \newline
    \DIFdelbeginFL \DIFdelFL{For each value of $\gamma_0$ and $N \leq 1024$}\DIFdelendFL \DIFaddbeginFL \DIFaddFL{Unless stated otherwise}\DIFaddendFL , \DIFdelbeginFL \DIFdelFL{the distributions are constructed from at least $10{,}000$ }\DIFdelendFL \DIFaddbeginFL \DIFaddFL{each parameter combination contains more than $2{,}500$ }\DIFaddendFL independent \DIFdelbeginFL \DIFdelFL{runs }\DIFdelendFL \DIFaddbeginFL \DIFaddFL{realisations generated }\DIFaddendFL with different \DIFdelbeginFL \DIFdelFL{random number }\DIFdelendFL \DIFaddbeginFL \DIFaddFL{random-number }\DIFaddendFL seeds. Due to computational constraints, \DIFaddbeginFL \DIFaddFL{the }\DIFaddendFL $N=2048$ \DIFdelbeginFL \DIFdelFL{is constructed from $2500$ independent runs}\DIFdelendFL \DIFaddbeginFL \DIFaddFL{data contain $2{,}500$ realisations. Exact ensemble sizes are reported in }\cref{sec:FatigueSampleCounts}\DIFaddendFL . Quantiles are shown only for cases with more than $1{,}000$ samples.
    }
    \label{fig:EPSizeDependence}
\end{figure}

The resulting dependence of the characteristic number of cycles to reach an end state $\CycleNineFive$ on the imposed strain amplitude $\gammaMax$ is shown in \cref{fig:EPSizeDependence}a. As the system size increases, the range of strain amplitudes over which delayed dynamics are observed shifts to smaller values of $\gammaMax$. In particular, the strain amplitude at which trajectories exhibiting delayed absorption or yielding first appear decreases with increasing system size. This behaviour is consistent with improved sampling of the tails of the initial strain distribution in larger systems, which increases the likelihood of near-threshold elements.

For delayed absorption to occur, the system must contain at least one element with a sufficiently large initial value of $l$ such that the system does not immediately attain an absorbing state. With increasing $N$, the probability of sampling rare strain values in the tails of the initial distribution grows. Larger systems are therefore increasingly likely to contain elements whose initial strain lies close to the yielding threshold, allowing plastic activity to be triggered at smaller imposed strain amplitudes. Consequently, the transition from immediate absorption to delayed absorption shifts to smaller imposed strain amplitudes $\gammaMax$. In the thermodynamic limit $N\to\infty$, the system will almost surely contain elements with initial strains arbitrarily close to the threshold $l_c$, so that plastic activity can be triggered at arbitrarily small imposed strain amplitudes. 

To further quantify this effect, we examine the characteristic cycle number at a fixed strain amplitude of $\gammaMax=0.799$, chosen as a representative value in the absorbed region (the closest sampled point to $0.8$ in our dataset). \Cref{fig:EPSizeDependence}b shows that the maximum number of cycles required for absorption grows approximately logarithmically \DIFdelbegin \DIFdel{with system size $N$}\DIFdelend \DIFaddbegin \DIFadd{over the range of system sizes studied}\DIFaddend . A plausible interpretation of this trend is that the remaining active elements relax approximately independently, with an exponential number of cycles set by $\lw$, the origin of which we discuss later. The system's absorption cycle \DIFdelbegin \DIFdel{is then }\DIFdelend \DIFaddbegin \DIFadd{would then be }\DIFaddend controlled by the slowest surviving element \DIFdelbegin \DIFdel{, and is therefore expected to }\DIFdelend \DIFaddbegin \DIFadd{and would }\DIFaddend grow approximately as the maximum of many such waiting times\DIFdelbegin \DIFdel{, giving a logarithmic dependence on $N$. Consequently, $\CycleNineFive$ would be expected to diverge as }\DIFdelend \DIFaddbegin \DIFadd{. The accessible system-size range is, however, insufficient to determine whether this interpretation and the apparent logarithmic growth persist as }\DIFaddend $N\to\infty$\DIFdelbegin \DIFdel{. We return to the assumptions underlying this interpretation later in the chapter}\DIFdelend . When we shift the curves in \cref{fig:EPSizeDependence}c to align the peaks of $C^*_{95\%}$ across all system sizes, the yielding curves \DIFdelbegin \DIFdel{collapse onto a single scaling form, indicating universal behaviour near the transition}\DIFdelend \DIFaddbegin \DIFadd{approximately collapse over the range studied}\DIFaddend . In contrast, the absorbing branches do not exhibit this collapse, suggesting that the absorption process retains a \DIFdelbegin \DIFdel{system-size-dependent character even }\DIFdelend \DIFaddbegin \DIFadd{stronger system-size dependence }\DIFaddend after accounting for the shift in the transition location.

A similar statistical argument \DIFdelbegin \DIFdel{explains }\DIFdelend \DIFaddbegin \DIFadd{may explain }\DIFaddend the shift in the onset of yielding. Larger systems are increasingly likely to contain elements whose initial strain lies very close to the yield threshold, such that only a small additional strain is required to trigger instability. \DIFdelbegin \DIFdel{As a result}\DIFdelend \DIFaddbegin \DIFadd{Consistent with this argument}\DIFaddend , the strain amplitude at which yielding first appears \DIFdelbegin \DIFdel{also }\DIFdelend shifts to smaller values \DIFdelbegin \DIFdel{as the system size increases}\DIFdelend \DIFaddbegin \DIFadd{over the system sizes shown in }\cref{fig:EPSizeDependence}\DIFadd{d. The available data do not establish whether this shift converges, nor whether the coexistence region ultimately vanishes or remains finite as $N\to\infty$}\DIFaddend .

\begin{figure}[t]
    \centering
    \includegraphics[width=0.99\linewidth]{chapters/Fatigue/Figures/logistics_N.pdf}
    \DIFdelbeginFL %DIFDELCMD < \caption[Effect of varying system size $N$ at fixed parameters $l_0 = 0.075$ and $\lw = 0.02$.]{
%DIFDELCMD <     %%%
\DIFdelFL{Effect }\DIFdelendFL \DIFaddbeginFL \caption[Finite-size dependence of the transition between absorption and yielding.]{\DIFaddFL{Finite-size dependence }\DIFaddendFL of \DIFdelbeginFL \DIFdelFL{system size $N$ }\DIFdelendFL \DIFaddbeginFL \DIFaddFL{the transition between absorption and yielding }\DIFaddendFL at fixed parameters \DIFdelbeginFL \DIFdelFL{$l_0 = 0.075$ }\DIFdelendFL \DIFaddbeginFL \DIFaddFL{$l_0=0.075$ }\DIFaddendFL and \DIFdelbeginFL \DIFdelFL{$\lw = 0.02$}\DIFdelendFL \DIFaddbeginFL \DIFaddFL{$\lw=0.02$}\DIFaddendFL . \newline
    {\bf a)} \DIFdelbeginFL \DIFdelFL{The $95^\mathrm{th}$ percentile failure cycle $\CStar_{95\%}$ }\DIFdelendFL \DIFaddbeginFL \DIFaddFL{Absorption probability $P_{\mathrm{abs}}$ }\DIFaddendFL as a function of imposed strain \DIFaddbeginFL \DIFaddFL{amplitude }\DIFaddendFL $\gamma_0$ \DIFaddbeginFL \DIFaddFL{for the finite system sizes studied, with logistic fits shown as dotted curves}\DIFaddendFL . \DIFdelbeginFL \DIFdelFL{Results }\DIFdelendFL \DIFaddbeginFL \DIFaddFL{Curves }\DIFaddendFL are shown for \DIFdelbeginFL \DIFdelFL{system sizes $N = 64, 128, 256, 512, 1024, 2048$ (}\DIFdelendFL \DIFaddbeginFL \DIFaddFL{$N=64,128,256,512,1024$ and $2048$ in }\DIFaddendFL blue, orange, green, red, purple \DIFdelbeginFL \DIFdelFL{, }\DIFdelendFL and brown \DIFdelbeginFL \DIFdelFL{, }\DIFdelendFL respectively\DIFdelbeginFL \DIFdelFL{). The inset shows the curves shifted to align the peak of $C^*_{95\%}$, evaluated across all samples and end states, at $\gamma_0 - \gamma_{0,\mathrm{peak}} = 0$}\DIFdelendFL . \newline
    {\bf b)} \DIFdelbeginFL \DIFdelFL{The corresponding value of $\CStar_{95\%}$ at $\gammaMax = 0.799$, the closest sampled point to $0.8$, }\DIFdelendFL \DIFaddbeginFL \DIFaddFL{Fitted transition midpoint $\gamma_{0,1/2}$ }\DIFaddendFL as a function of system size $N$. \newline
    {\bf c)} \DIFdelbeginFL \DIFdelFL{The same data }\DIFdelendFL \DIFaddbeginFL \DIFaddFL{Fitted transition width $w$ }\DIFaddendFL as \DIFdelbeginFL \DIFdelFL{in }%DIFDELCMD < {\bf %%%
\DIFdelendFL a \DIFdelbeginFL \DIFdelFL{)}%DIFDELCMD < }%%%
\DIFdelFL{, shifted to align the peak of $C^*_{95\%}$ at $\gamma_0 - \gamma_{0,\mathrm{peak}} = 0$. }%DIFDELCMD < \newline
%DIFDELCMD <     {\bf %%%
\DIFdelFL{d)}%DIFDELCMD < } %%%
\DIFdelFL{The corresponding distribution of end states: delayed yielding (blue), immediate yielding (purple), delayed absorption (orange), and immediate absorption (brown). }%DIFDELCMD < \newline
%DIFDELCMD <     %%%
\DIFdelFL{For each value }\DIFdelendFL \DIFaddbeginFL \DIFaddFL{function }\DIFaddendFL of \DIFdelbeginFL \DIFdelFL{$\gamma_0$ and $N \leq 1024$, distributions are constructed from at least $10{,}000$ independent runs. Due to computational constraints, $N=2048$ uses $2500$ runs. Quantiles are shown only for cases with more than $1{,}000$ samples}\DIFdelendFL \DIFaddbeginFL \DIFaddFL{system size $N$}\DIFaddendFL .}
    \label{fig:LogisticsN}
\end{figure}

To \DIFdelbegin \DIFdel{characterise the resulting }\DIFdelend \DIFaddbegin \DIFadd{describe the finite-size }\DIFaddend shift in the \DIFdelbegin \DIFdel{distributions }\DIFdelend \DIFaddbegin \DIFadd{distribution }\DIFaddend of end states, we \DIFdelbegin \DIFdel{again analyse the probability of absorption $P_{\mathrm{abs}}$ using
logistic fits of the form
}\DIFdelend \DIFaddbegin \DIFadd{fit the absorption probability at each sampled system size using
}\DIFaddend \[
    P_{\mathrm{abs}}(\gamma_0)=\frac{1}{1+\exp\!\left[k(\gamma_0-\gamma_{0,1/2})\right]}.
\]
The \DIFdelbegin \DIFdel{resulting fits are shown in }%DIFDELCMD < \cref{fig:LogisticsN}%%%
\DIFdel{a, while the }\DIFdelend fitted midpoint $\gamma_{0,1/2}$ and transition width $w$ are shown in \cref{fig:LogisticsN}b and \cref{fig:LogisticsN}c, respectively. \DIFdelbegin %DIFDELCMD < 

%DIFDELCMD < %%%
\DIFdel{As the system size increases, the transition from predominantly absorbing trajectories to predominantly yielding trajectories shifts to lower strain amplitudes. This trend can also be observed directly in the underlying distributions of end states shown in }%DIFDELCMD < \cref{fig:EPSizeDependence}%%%
\DIFdel{d. Plotting the midpoint $\gamma_{0,1/2}$ as a function of the system size $N$ reveals a systematic drift that decreases with increasing $N$. Following the approach of Khandare and Sastry in Ref.~\mbox{%DIFAUXCMD
\cite{khandareElastoplasticModellingCyclic2026}}\hskip0pt%DIFAUXCMD
, we fit the form
}\[
    \DIFdel{\gamma_{0,1/2}(N)=aN^{-\nu}+\gamma_{0,1/2}(N\to\infty),
}\]%DIFAUXCMD
\DIFdel{obtaining $\gamma_{0,1/2}(N\to\infty)=0.78$, $\nu=0.16$, and $a=1.27\times10^{-1}$. The power-law decay of the finite-size correction indicates that the midpoint converges slowly towards a well-defined value in the thermodynamic limit. Physically, this reflects the gradual saturation of the tails of the initial distribution: as the system size increases, the probability of finding elements arbitrarily close to the instability threshold increases, but with diminishing returns as the system approaches the large-$N$ limit. The numerical values differ from those reported in Ref.~\mbox{%DIFAUXCMD
\cite{khandareElastoplasticModellingCyclic2026}}\hskip0pt%DIFAUXCMD
, which is expected given the different choice of units used here.
}%DIFDELCMD < 

%DIFDELCMD < %%%
\DIFdel{Of particular interest is the behaviour }\DIFdelend \DIFaddbegin \DIFadd{These quantities provide descriptive measures }\DIFaddend of the transition \DIFdelbegin \DIFdel{width $w$, shown in }%DIFDELCMD < \cref{fig:LogisticsN}%%%
\DIFdel{c. We find a power-law dependence $w\propto N^{-\nu}$ with $\nu=0.17$, indicating that the transition region between yielding and absorbing behaviour becomes increasingly sharp with increasing system size. This suggests that in the thermodynamic limit }\DIFdelend \DIFaddbegin \DIFadd{over the finite system sizes studied; we do not use them to infer an asymptotic scaling form or a limiting value as }\DIFaddend $N\to\infty$\DIFdelbegin \DIFdel{the transition may become effectively discontinuous. Such convergence was not reported in Ref.~\mbox{%DIFAUXCMD
\cite{khandareElastoplasticModellingCyclic2026}}\hskip0pt%DIFAUXCMD
. However, the present data do not allow us to determine whether, in the thermodynamic limit, the system exhibits a sharp transition directly between delayed absorption and delayed yielding, or whether the onset of yielding remains gradual, with a crossover from delayed to immediate yielding}\DIFdelend .

Taken together, these results show that the qualitative phenomenology of the model is robust to variations in both the post-hop distribution width $\lw$ and the system size $N$ \DIFaddbegin \DIFadd{over the ranges studied}\DIFaddend . Small changes to $\lw$ primarily rescale the characteristic number of cycles required for trajectories to reach their final state, consistent with a diffusive first-passage process. In contrast, varying the system size modifies the statistical likelihood of sampling rare initial strain configurations that can trigger yielding. \DIFdelbegin \DIFdel{As a result, the }\DIFdelend \DIFaddbegin \DIFadd{The }\DIFaddend transition between absorbing and yielding behaviour shifts systematically \DIFdelbegin \DIFdel{with $N$ and becomes increasingly sharp for larger systems, consistent with the emergence of a well-defined transition in the thermodynamic limit}\DIFdelend \DIFaddbegin \DIFadd{and appears to sharpen across the accessible system sizes, but its thermodynamic-limit behaviour remains unresolved}\DIFaddend .

Nevertheless, the key features of the dynamics, namely the coexistence of absorbing and yielding trajectories and the presence of long-lived delayed dynamics, are observed across all system sizes studied, but become increasingly confined to a narrowing region of strain amplitude as $N$ increases. These observations indicate that the fatigue behaviour described in the previous sections is not an artefact of the particular parameter choices used, but rather a robust feature of the underlying elastoplastic model.

\section{Path to a Single Active Element}

To understand in more detail the numerical results explored so far, we will now consider a simple model that considers the fraction of active elements $\Factive$ left in the system as a function of cycle number. We define here an active element to be one which has the possibility of yielding as a result of bulk elastic straining (it is not in the absorbed region). Mathematically, for our model, this 
corresponds to an element whose value of $\ltilde$ (the element's deviation in its strain from the bulk strain $\gamma$) is 
\[
    \ltilde \le -1+\gammaMax
    \quad \text{or} \quad
    \ltilde \ge 1-\gammaMax.
\]

The relaxation towards an absorbing state can be understood in terms of the progressive elimination of plastically active elements. In the early cycles, unless the strain amplitude is small, a large fraction of elements lies within the active region of strain space, leading to frequent yielding events and sustained plastic activity. As the system evolves, these events drive elements out of the active region, reducing the number of sites capable of further rearrangement.

\begin{figure}
    \centering
    \includegraphics[width=0.99\textwidth]{chapters/Fatigue/Figures/ActiveDecay_Path21Osc.pdf}
    \caption[Decay of number of active elements remaining in the system.]{Example decay of the number of active elements as a function of the number of cycles. Curves show the ensemble-averaged value of $\Factive$ over \DIFdelbeginFL \DIFdelFL{$1000$ }\DIFdelendFL \DIFaddbeginFL \DIFaddFL{$1{,}000$ }\DIFaddendFL trajectories. We distinguish between trajectories in which the system eventually reaches an absorbing state (solid lines) and those in which it eventually yields (dashed lines). \newline
    {\bf a)} $l_0 = 0.075$, $\lw = 0.02$, and $N = 1024$. \newline
    {\bf Inset)} binned mean of the fraction of active elements in the final cycle $\CStar$ taken across more than \DIFdelbeginFL \DIFdelFL{$10{,}000$ }\DIFdelendFL \DIFaddbeginFL \DIFaddFL{$1{,}000$ }\DIFaddendFL samples for each $\gammaMax$. \newline
    {\bf b)} Parameter sets for which the number of cycles to an absorbed state satisfies $\CycleNineFive \approx 200$. From bottom to top: $(l_0,\gamma_{\max})=(0.0375,0.882)$, $(0.05,0.848)$, $(0.075,0.802)$, $(0.1,0.784)$ and $(0.125,0.775)$. \newline
    }
    \label{fig:Pathto1}
\end{figure}

This process is quantified in \cref{fig:Pathto1}a and \cref{fig:Pathto1}b, which show the mean fraction of active elements $\Factive$ as a function of cycle number. At early cycles the decay is rapid, reflecting the large number of plastic events that occur as the system relaxes away from its initial distribution. Once an element enters the absorbing region it cannot leave through its own dynamics; it can only become active again if stress redistributed by the yielding of another active element pushes it back into the active region.

The inset of \cref{fig:Pathto1}a further illustrates that long-lived trajectories are characterised by a slow decrease in the fraction of active elements to
progressively smaller values. Samples that survive many cycles typically contain only a few remaining active elements immediately before reaching their final state. For absorbing trajectories, this typically corresponds to a single element remaining active immediately prior to absorption, consistent with the microscopic mechanism of absorption discussed earlier. In contrast, yielding trajectories can still contain many tens of active elements at the point of yielding, although this remains an order of magnitude fewer than in samples that yield after only a small number of cycles.

The same qualitative behaviour is observed across different levels of annealing, as shown in \cref{fig:Pathto1}b for parameter sets chosen such that $\CycleNineFive \approx 200$ (the closest value in the dataset). In all cases, the dynamics consist of a rapid initial depletion of active elements followed by a much slower decay at an approximately constant rate. The initial fraction of active elements is determined by the portion of the initial strain distribution that lies within the active region, which depends on both the oscillation amplitude $\gammaMax$ and the initial width of the distribution $l_0$. The ultimate decay rate is therefore controlled by three parameters: the imposed strain amplitude $\gammaMax$, which sets the width of the active region; the post-hop width $\lw$, which controls the typical step size; and the initial disorder $l_0$, which determines the width of the initial strain distribution.

This behaviour admits a simple physical interpretation. Once the number of active elements becomes sufficiently small, plastic events are well separated in space, so the elastic interactions between them are weak due to the $1/r^D$ decay of the Eshelby propagator. In this regime, the cumulative stress increments induced by distant events are small compared to the local stochastic updates from yielding, so the dynamics of different active elements become approximately independent. This justifies treating the late-time dynamics in a dilute limit, in which the decay of $\Factive$ reflects the stochastic evolution of effectively independent elements governed by the post-hop strain distribution.

\subsection*{Self-Interaction of an Active Element}

We therefore briefly examine the self-interaction experienced by an element immediately after it yields. This self-interaction plays an important role in determining the long-time dynamics when only a few active sites remain.

When an element yields, the stress released by the event is redistributed elastically throughout the system via the stress propagator defined in \cref{sec:EshelbyPropogator}. Because force balance is subsequently reimposed, a portion of this redistributed stress is recovered elastically. For an element that loses a stress of 1, the net change in stress it experiences corresponds to the value of the stress propagator at the origin, 
\[
\Delta l = \Origin \approx -0.5708 
\]
in the thermodynamic limit $N\to\infty$, although this quickly converges as shown in \cref{fig:PropogatorConvergence}.

In our model, the local yield thresholds correspond to $l_c=\pm 1$. For clarity, the discussion which follows considers the \DIFdelbegin \DIFdel{$l_c= 1$ }\DIFdelend \DIFaddbegin \DIFadd{$l_c=1$ }\DIFaddend threshold; the \DIFdelbegin \DIFdel{cases }\DIFdelend \DIFaddbegin \DIFadd{case }\DIFaddend for $l_c=-1$ \DIFdelbegin \DIFdel{follow }\DIFdelend \DIFaddbegin \DIFadd{follows }\DIFaddend by symmetry. When an element yields, its local strain $\eta$ is redrawn from the post-hop distribution, which we take to be Gaussian of width $\lw$ and mean $0$,
\[
 \eta \sim \mathcal{N}(0,\lw^2)
\DIFdelbegin \DIFdel{.
}\DIFdelend \]
\DIFaddbegin \DIFadd{as disscussed in }\cref{sec:fatigueMethodology}\DIFadd{. }\DIFaddend Immediately after the yielding event (before force balance), the element therefore has a strain $l=\eta$. The stress released by this event, which must be balanced, is therefore $1-\eta$, so the net stress loss by the yielding element after force balance is
\[
\Delta l = |\Origin|(1-\eta)
\]
and the element will then have a local strain
\DIFdelbegin \DIFdel{,
}\DIFdelend \[
l = 1 - |\Origin|(1-\eta).
\]

\DIFdelbegin \DIFdel{If this }\DIFdelend \DIFaddbegin \DIFadd{To understand when this element will yield again during the same oscillation cycle, it is useful first to consider the deterministic limit $\lw=0$. Consider an element with value $\ltilde_0$ at the start of a cycle, where $\gamma=0$, on the positive side of the strain distribution. During the forward stroke, the global strain increases to $\gammaMax$, and the element reaches the positive yield threshold if
}\[
    \DIFadd{\ltilde_0+\gammaMax \geq 1,
}\]
\DIFadd{or equivalently
}\[
    \DIFadd{\ltilde_0 \geq 1-\gammaMax.
}\]
\DIFadd{In the absence of post-hop noise, $\eta=0$, so after yielding the self-interaction shifts the element to
}\[
    \DIFadd{\ltilde_1=\ltilde_0-|\Origin|.
}\]
\DIFadd{As the applied strain is subsequently reversed, this element will reach the opposite threshold $l_c=-1$ provided
}\[
    \DIFadd{\ltilde_1-\gammaMax\leq -1.
}\]
\DIFadd{Combining these conditions, an element yields once on both the positive and negative strokes when
}\[
    \DIFadd{1-\gammaMax
    \leq \ltilde_0
    \leq -1+\gammaMax+|\Origin|.
}\]

\begin{figure}[t!]
    \centering
    \includegraphics[width=0.95\linewidth]{chapters/Fatigue/RandomWalk/YieldRegions.pdf}
    \caption[Regions governing the evolution of an active element under oscillatory strain.]{ \DIFaddFL{Distinct regions of $\ltilde$ governing the evolution of an element under oscillatory strain in the deterministic limit $\lw=0$. Elements in the central green region cannot be strained sufficiently to yield and are therefore absorbed. Elements in the blue regions yield on one half-cycle and are mapped into the opposite blue region, allowing them to yield again on the following half-cycle. Elements in the yellow regions yield but are mapped into the central absorbed region. Elements in the red regions have sufficiently large $|\ltilde|$ to yield more than once during a single stroke, with repeated yielding reducing $|\ltilde|$ until the element enters either a blue or yellow region. The Gaussian profile is shown schematically for clarity and is not intended to represent the instantaneous distribution of $\ltilde$, which will in general be non-Gaussian and vary throughout the oscillation cycle. }}
    \label{fig:YieldRegions}
\end{figure}

\DIFadd{The different possible trajectories are illustrated in }\cref{fig:YieldRegions}\DIFadd{. Elements in the central green region,
}\[
    \DIFadd{-1+\gammaMax < \ltilde < 1-\gammaMax,
}\]
\DIFadd{cannot reach either yield threshold under the imposed strain and are therefore absorbed. Elements in the blue regions can yield and are mapped by their self-interaction into the corresponding blue region of opposite sign. An element starting in the positive blue region therefore yields during the forward stroke and again during the negative stroke. In the deterministic limit, the two self-interactions have equal magnitude and opposite sign, returning the element to its initial value of $\ltilde$ after a complete cycle.
}

\DIFadd{By contrast, an element in one of the yellow regions yields but is mapped into the central absorbed region. For example, following a yielding event at the positive threshold, if
}\[
    \DIFadd{-1+\gammaMax < \ltilde_1 < 1-\gammaMax,
}\]
\DIFadd{the remaining negative stroke is insufficient to bring the element to $l_c=-1$. It therefore yields only on the first half-cycle and subsequently remains absorbed. Elements in the red regions have sufficiently large $|\ltilde|$ that a single }\DIFaddend yielding event does not \DIFdelbegin \DIFdel{cause the element to be absorbed , the element will yield again on the next half-cycle. This }\DIFdelend \DIFaddbegin \DIFadd{move them out of the active region on the same side. They can therefore yield multiple times during a single stroke, with each event reducing $|\ltilde|$ until they enter either a blue region, where they subsequently oscillate between the two thresholds, or a yellow region, where they become absorbed.
}

\DIFadd{For the long-lived trajectories of interest here, an element in the positive blue region therefore occupies the interval
}\[
    \DIFadd{1-\gammaMax
    \leq \ltilde
    \leq -1+\gammaMax+|\Origin|
}\]
\DIFadd{when viewed stroboscopically at the start of each cycle. In the deterministic limit, an element within this interval follows the same trajectory indefinitely. Crossing the lower boundary places it in the central region where it can no longer reach the positive threshold, while crossing the upper boundary causes the next yielding event to map it into the absorbed region. These two routes provide the two boundaries through which the single-element dynamics can terminate.
}

\DIFadd{For a finite post-hop width $\lw$, this simple mapping becomes stochastic. Following a yielding event at the positive threshold,
}\[
    \DIFadd{\ltilde_1
    =
    \ltilde_0-|\Origin|(1-\eta_1),
}\]
\DIFadd{so whether the element is mapped into the opposite active region or into the absorbed region depends on the particular post-hop sample $\eta_1$. An element close to the boundaries between the regions in }\cref{fig:YieldRegions} \DIFadd{can therefore either continue to yield or become absorbed. If the first event maps the element into the opposite active region, it reaches $l_c=-1$ during the subsequent negative stroke and yields for a second time. The second post-hop sample then introduces a further stochastic displacement and can likewise move the element towards, or into, the absorbed region.
}

\DIFadd{Conditional on the element yielding on both half-cycles, the }\DIFaddend second event occurs at the opposite threshold $l_c=-1$ \DIFdelbegin \DIFdel{, }\DIFdelend and therefore produces a stress change of opposite sign. Over one full oscillation cycle, the net stress increment is
\DIFdelbegin \DIFdel{therefore
}\DIFdelend \[
    \Delta l = |\Origin|(\eta_1\DIFdelbegin \DIFdel{-}\DIFdelend \DIFaddbegin \DIFadd{+}\DIFaddend \eta_2) \sim \mathcal{N}(0, 2|\Origin|^2\lw^2)\DIFaddbegin \DIFadd{,
}\DIFaddend \] 
where $\eta_1$ and $\eta_2$ are independent post-hop samples drawn on the positive and negative half-cycles, respectively. Since both are independent zero-mean Gaussian variables, their difference is also Gaussian with variance \DIFdelbegin \DIFdel{$2\lw^2$}\DIFdelend \DIFaddbegin \DIFadd{$2\lw^2|\Origin|^2$}\DIFaddend .

The key consequence is that the deterministic stress relaxation over the two half-cycles cancels, while the stochastic contributions from the post-hop distribution accumulate. The strain of \DIFdelbegin \DIFdel{the element }\DIFdelend \DIFaddbegin \DIFadd{an element that remains active }\DIFaddend therefore performs a random walk when viewed stroboscopically from cycle to cycle.

\begin{figure}[t!]
    \centering
    \includegraphics[width=0.99\linewidth]{chapters/Fatigue/RandomWalk/OscillatorRegions.pdf}
    \caption[Illustration of an ideal trajectory of a single independent element]{
    Illustration of an idealised trajectory of a single independent active element. The black curve shows the case without a post-hop distribution, where the stress changes from each half-cycle cancel exactly, leaving no net change over a full cycle. The orange curve shows the case with a Gaussian post-hop distribution (shown here with a large $l_w$ for illustration). \newline
    {\bf a)} Evolution of $\ltilde$ over several idealised cycles. {\bf b)} The value of $\ltilde$ viewed stroboscopically at the start of each cycle. Dotted lines are included to guide the eye and do not represent the actual evolution of $\ltilde$.
    }
    \label{fig:IdealElement}
\end{figure}

This behaviour is illustrated schematically in \cref{fig:IdealElement}. For a delta-function post-hop distribution ($l_w=0$), the stress changes over the two half-cycles cancel exactly, leaving the element at the same value of $\tilde l$ after each cycle (black trajectory in panel a and corresponding stroboscopic points in panel b). For a finite post-hop width $\lw$, however, the small stochastic shifts produced at each yielding event accumulate, giving rise to a discrete random walk in $\tilde l$, as illustrated by the orange trajectories. In both cases, the value of $\ltilde$ oscillates between positive and negative values within the (blue) active region while skipping over the absorbed (green) region. For the black trajectory, the element therefore always remains in the active region. By contrast, the orange trajectory performs a random walk in $\ltilde$ and will eventually enter the absorbed region.

\begin{figure}[p]
    \centering
    \includegraphics[width=0.99\textwidth]{chapters/Fatigue/Figures/RandomWalks.pdf}
    \caption[Stroboscopic evolution of element strain measured at the start of each cycle.]{Stroboscopic evolution of element strain $l$ measured at the start of each cycle. $\tilde{C}$ indicates the number of cycles after the random walk has begun.\newline
    {\bf a)} Example trajectory of an active element in a sample that ultimately absorbs. The red dashed line denotes the absorbing boundary; the blue trajectory crosses this threshold, leading to absorption. \newline
    {\bf b)} Example trajectory in a sample that ultimately yields. The blue curve shows the active element, while the orange curve shows a neighbouring element. In this case, the neighbouring element crosses the red boundary into the active region, triggering yielding. \newline
    {\bf c)} Ensemble of trajectories of the oscillating element, each terminated at the first crossing of the absorbing boundary (red dashed line). \newline
    {\bf d)} Standard deviation of an ensemble of $1{,}000$ active element displacements, $\sigma(\Delta l)$, as a function of cycle number, demonstrating the expected diffusive scaling. \newline
    Parameters: $l_0=0.075$, $\gammaMax=0.823$, $\lw=0.02$, $N=1024$
    }
    \label{fig:RandomWalkExamples}
\end{figure}

We observe the same behaviour in the full simulations, as shown in \cref{fig:RandomWalkExamples}. The trajectories shown correspond to cases in which only a single active element remains. In this situation, the dynamics must evolve as described above. However, similar random-walk behaviour can also occur in systems with more than one active element when the active sites are sufficiently spatially separated for their dynamics to be dominated by their own post-hop updates rather than the elastic interactions with other active sites.

For any active element in the dilute regime, in trajectories that ultimately reach an absorbing state, the dynamics reduce to a stochastic walk in $\tilde l$ (orange line in \cref{fig:RandomWalkExamples}a) until it first reaches the boundary separating the active and absorbed regions (red dashed line in \cref{fig:RandomWalkExamples}a). Once this boundary is crossed, the element can no longer yield because the bulk elastic strain $\gamma$ is insufficient to strain the element beyond either threshold $l_c=\pm1$. If this is the last remaining active element, the system therefore attains an absorbing state at this point.

An element’s path to absorption can therefore be viewed as a first-passage problem, in which the element performs a random walk in $\tilde l$ until it enters the absorbed region. The system’s path to absorption can then be regarded as an ensemble of such first-passage processes, with the number of cycles to absorption determined by the slowest element to reach the absorbed region.

Correspondingly, in trajectories that yield only after many cycles, the remaining active elements can become sufficiently dilute that their dynamics approximate independent random walks. Such trajectories are, however, less common than those that yield at early times. An example is shown in \cref{fig:RandomWalkExamples}b, where an active element undergoes a random walk (orange). We also show a neighbouring element on the same streamline in red, where a streamline refers to elements lying along the same vertical or horizontal axis. As a result of stress propagation, this absorbed element (which lies below the dotted line) undergoes a correlated random walk driven by the yielding of the active element, but with the opposite sign: the active element’s value of $\ltilde$ drifts downward toward the threshold, while the neighbouring absorbed element drifts upward toward it. For suitable initial configurations, the active element can therefore pull the absorbed element back into the active region, where it can subsequently yield. In the examples shown here, the surrounding elements are arranged such that this process nucleates the formation of a shear band, although this outcome does not occur in all realisations.

Finally, we confirm that the stochastic dynamics are consistent with a Gaussian random walk. In \cref{fig:RandomWalkExamples}c we show a sample of trajectories drawn from an ensemble that terminate either in absorption or yielding. For reference, we plot the absorption threshold; once a trajectory crosses this boundary, the walk terminates, as the element can no longer yield. For such a process, the standard deviation of the displacement grows diffusively as 
\[
    \sigma(\Delta l) \propto \sqrt{\tilde{C}}
\]
where $\tilde{C}$ is the number of cycles from the beginning of the random walk. The characteristic step size is set by the width of the cycle-to-cycle increment, $D=\sqrt{2}|\Origin|\lw$. In \cref{fig:RandomWalkExamples}d we plot $\sigma(\Delta l)$ averaged over $1000$ trajectories, showing the expected diffusive scaling while sufficient samples remain in the ensemble. Deviations at late times arise as trajectories terminate upon reaching either the absorbing boundary or the yielding state. Motivated by this random-walk description of the element dynamics, we now introduce a simplified model that isolates this behaviour.

\section{The Dilute Model} \label{sec:DilluteModel}

The random-walk behaviour identified in the previous section suggests that, once the number of active elements becomes sufficiently small, the dynamics of individual elements can be considered to be approximately independent. To explore this behaviour more directly, we therefore consider a simplified elastoplastic model in the dilute limit. Specifically, we use the elastoplastic model defined in \cref{sec:fatigueMethodology}, but neglect stress propagation between elements while retaining the self-interaction elastic stress recovered by force balance. When an element yields, its stress drops by an amount equal to $|\Origin|(1-\eta)$ where $\eta\sim\mathcal{N}(0,\lw^2)$; however, no other elements are affected by this yielding event.

The initial distribution of element strains in this dilute limit is taken to be the analytical form obtained after imposing force balance: a Gaussian of width $\Norm l_0$, where $\Norm$ arises from the rescaling associated with force balance. The system then evolves as usual, with elements yielding whenever the critical strain $l_c=\pm1$ is reached.

\begin{figure}[t!]
    \centering
    \includegraphics[width=0.99\linewidth]{chapters/Fatigue/Figures/ElementDecay.pdf}
    \caption[Comparison of the decay of the number of active elements remaining in the system against a dilute model, with an inset showing the decay gradient.]{Comparison of the decay of the number of active elements in the system with predictions from simplified models. Curves show the mean fraction of active elements $\Factive$ averaged over $1{,}000$ samples, and are plotted while at least $500$ samples remain active. Results from the full elastoplastic model (EP model) are compared with those from the dilute model. {\bf Inset:} The logarithmic decay gradient, $-d\ln \Factive/ dC$, for the same curves. Locations of curves are indicated on \cref{fig:Comparison}.}    
    \label{fig:ElementDecayComp}
\end{figure}

To verify that the dilute model reproduces behaviour qualitatively similar to the full elastoplastic model, we compare the predicted decay curves for the fraction of active elements. As shown in \cref{fig:ElementDecayComp}, the dilute model shows excellent agreement with the full elastoplastic model for high levels of annealing (low $l_0$). In this regime, the dilute model captures both the initial rapid drop and the subsequent slow decay. As the level of annealing decreases (increasing $l_0$), the numerical results from the dilute model begin to diverge from the full elastoplastic simulations. In this regime, the effective diffusion coefficients governing the rate at which elements are absorbed differ between the two models, leading to increasingly noticeable deviations.

In the dilute limit, the exponential decay observed in $\Factive$ depends only on the characteristic step size of the random walk, set here by $\lw$, and on the size of the active region determined by the imposed strain amplitude $\gammaMax$. This occurs because, once the dynamics become dilute, the remaining active elements evolve essentially independently, and their evolution reduces to a first-passage problem for a random walk between fixed boundaries. The decay rate is therefore controlled only by the diffusion scale of the walk and the distance to the absorbing boundary, and is insensitive to the initial distribution set by $l_0$. This behaviour can be seen in the inset of \cref{fig:ElementDecayComp}, which shows that the gradients for the two cases at $\gammaMax = 0.775$ converge independently of the level of annealing $l_0$. The initial fraction of active elements and the magnitude of the initial drop, however, remain dependent on $l_0$. 

In less dilute samples, shown here for the $l_0 = 0.125$ case, the two models no longer agree. In this regime, the higher density of active elements makes their dynamics correlated, violating the assumption of independent random walks that underlies the dilute model. As a result, the effective diffusion is reduced, and the decay of $\Factive$ is slower than that predicted by the dilute model.

As a final check of the dilute model, we compare the number of cycles required to reach an absorbing state, $\CStar$. The comparison between the two models is shown in \cref{fig:Comparison}. Excellent agreement between the predicted and simulated values is observed at intermediate levels of annealing $l_0$ with both models predicting similar values of $\CycleNineFive$. For less highly annealed samples, however, the predictions break down over the same range of imposed strain amplitudes 
$\gammaMax$ for which delayed yielding occurs. In this regime, yielding also occurs with a high probability in the full model. Because the dilute model does not capture yielding, which is inherently a many-body process involving spatially correlated plastic events, it cannot distinguish between initial configurations that ultimately yield and those that absorb. As discussed previously, the number of cycles to absorption increases with increasing $\gammaMax$; however, because the dilute model assumes that all trajectories eventually absorb, the predicted curves (e.g. the blue dashed curve for $l_0=0.025$) continue to increase over the entire range considered. By contrast, in the full model the corresponding curves are cut off once yielding becomes sufficiently probable.

\begin{figure}[b!]
    \centering
    \includegraphics[width=0.99\linewidth]{chapters/Fatigue/Figures/Comparison.pdf}
    \caption[Comparison of $\CStar$ against $\gammaMax$ for the full and dilute models at fixed annealing levels.]{Comparison of the final cycle $\CStar$ as a function of the imposed strain amplitude $\gammaMax$ for the full and dilute models at fixed annealing levels $l_0$. The full model corresponds to the absorbed cases shown in \cref{fig:FailureCyclel0}, and the dilute model is evaluated at the same set of parameters. Coloured markers indicate the position of the decay curves in \cref{fig:ElementDecayComp}. Parameters: $N=1024$ and $\lw=0.02$.}
    \label{fig:Comparison}
\end{figure}

Conversely, for poorly annealed samples (high $l_0$), the dilute model underpredicts the number of cycles to absorption. In this regime, the decay of active elements in the full model is slower than that predicted by the dilute model, as shown in the inset of \cref{fig:ElementDecayComp}. Elements, therefore, remain active for longer than expected, increasing the observed value of $C^*$ when compared to the dilute model. This discrepancy arises because a large number of active elements are initially present, so their dynamics are no longer independent. The resulting interactions introduce correlations between the random walks of different elements, which reduce the effective decay rate compared with the dilute prediction. We still however see reasonable qualitative agreement between the dilute and full model in this regime.

Overall, these results show that the dilute model captures the essential physics of the late-time dynamics when the number of active elements becomes small, and their behaviour is effectively independent. In this regime, relaxation towards absorption is well described by a stochastic random walk, with the remaining active elements behaving independently. Deviations arise when the assumptions of the dilute limit break down. For poorly annealed samples, the initially large number of active elements leads to correlated dynamics that slow the decay of activity, while for more highly annealed samples, the model cannot capture trajectories in which the system eventually yields and overestimates the number of cycles to absorption. Within its regime of validity, however, the dilute description provides a simple and physically transparent framework for understanding the approach to an absorbing state.

\section{Conclusion} \label{sec:FatigueConclusion}

In this chapter, we have investigated the behaviour of a simple mesoscopic elastoplastic model of amorphous materials subjected to large-amplitude oscillatory shear. The focus has been on understanding the delayed dynamics that arise near the transition between absorbing states, in which plastic activity eventually ceases, and yielding states, in which irreversible plastic activity persists, and the system ultimately fails. By systematically exploring the effects of the initial degree of annealing $l_0$, the characteristic width of the post-hop strain distribution $\lw$, and the system size $N$, we have characterised how these parameters control both the probability of reaching each final state and the number of deformation cycles required to do so.

Our results demonstrate that the dynamics close to the transition are governed by long-lived transient regimes where the system can remain active over many cycles before either relaxing into an absorbing configuration or undergoing catastrophic failure. The number of cycles required to reach these final states varies strongly with the imposed strain amplitude and with the level of annealing of the sample. In particular, more highly annealed systems exhibit behaviour more consistent with brittle, abrupt yielding, whereas more weakly annealed systems tend to show more gradual dynamics and, in parts of the transition region, longer-lived transients. These observations are consistent with the broader phenomenology observed in amorphous solids, in which the degree of disorder inherited from the preparation history plays a central role in determining whether deformation proceeds gradually or through abrupt localisation events such as shear banding \cite{pollard_yielding_2022,liuFateShearoscillatedAmorphous2022,khandareElastoplasticModellingCyclic2026}.

As was also observed in Ref.~\cite{khandareElastoplasticModellingCyclic2026}, we find a transition region in which the probability of reaching either an absorbing state or a yielding state varies continuously with the imposed strain amplitude. Within this region, the number of cycles required to reach absorption increases with increasing imposed strain amplitude $\gammaMax$, while the number of cycles required for yielding increases as the amplitude is reduced. At the same time, the probability of absorption decreases with increasing amplitude. Consequently, the strain amplitudes at which the number of cycles to absorption or yielding is largest correspond to regimes in which those outcomes are least probable.

A key observation emerging from our simulations is that, as the number of deformation cycles increases, the remaining plastic activity becomes progressively more dilute, with the behaviour of absorbing cases reducing to a first-passage problem. Early in the dynamics, many elements are active and interact strongly through the long-ranged Eshelby stress propagator, giving rise to correlated plastic events throughout the system. As the system evolves, however, the number of active elements decreases and for low levels of annealing the spatial separation between elements appears to grow. In this limit, the interactions between plastic events become weaker, and the dynamics of the remaining active elements increasingly resemble independent stochastic processes.

Motivated by this physical picture, we introduced a simplified model intended to capture the dilute limit of the dynamics. The evolution of each element is assumed to follow an ideal independent stochastic process governed only by the distribution of post-hop strains and the magnitude of the imposed strain amplitude. Despite its simplicity, this approach reproduces the key qualitative behaviour observed in the full elastoplastic simulations for low levels of annealing, including the overall decay of the fraction of active elements and the number of cycles to absorb. We propose that this simplified model, therefore, provides a useful framework for understanding the late-cycle relaxation of the system and for interpreting the emergence of delayed absorption observed in the full simulations.

More broadly, the results presented here illustrate how mesoscale elastoplastic models provide a useful bridge between microscopic particle simulations and macroscopic rheological behaviour. Such models capture the essential ingredients of plasticity in amorphous solids: local yielding events, long-range interactions, and structural disorder inherited from the sample's preparation history, while remaining sufficiently simple to allow insight into the dynamics.

% Several directions for further work follow naturally from this study. First, it would be valuable to place the simplified stochastic description developed here within a more formal analytical framework to make explicit the assumptions on which it relies. An initial formulation of such a model has already been developed by the author; however, in its present form, it remains informal and depends on assumptions that have not yet been systematically validated. Consequently, it does not currently provide additional physical insight beyond that already obtained from the numerical model.

Several directions for further study follow naturally from this work. First, it would be valuable to place the simplified stochastic description developed here within a more formal analytical framework to make explicit the assumptions on which it relies.
\DIFdelbegin \DIFdel{Once this framework has been formalised, a natural extension }\DIFdelend \DIFaddbegin \DIFadd{A natural extension of this simplified stochastic model }\DIFaddend would be to investigate whether it can \DIFdelbegin \DIFdel{be expanded to }\DIFdelend capture not only absorption but also yielding. Preliminary observations suggest that, in highly annealed samples, the interaction between an active element and neighbouring elements in the absorbed region may destabilise otherwise stable configurations. In particular, if a neighbouring element lies sufficiently close to the absorbing boundary, stress redistribution from a nearby plastic event may push it back into the active region. This process could, in turn, trigger a cascade of activity that ultimately leads to macroscopic failure. Incorporating such mechanisms within the simplified framework would therefore provide a route to describing both absorption and yielding within a unified stochastic model, at least in the limit of highly annealed systems.

More generally, although the present work has been able to quantify and characterise the transition region between absorbing and yielding behaviour, the physical origin of this transition remains unclear. In particular, the differences between the trajectories that lead to absorption and those that lead to yielding are not yet well understood, nor is the mechanism by which these distinct dynamical pathways can coexist within the same region of parameter space to give rise to delayed dynamics. Previous studies of amorphous materials under cyclic shear have suggested that yielding may be interpreted as a transition between an absorbing state, in which plastic activity eventually ceases, and an active state in which plastic rearrangements persist indefinitely \cite{JocteurYielding2024}. Developing a clearer microscopic understanding of how this transition emerges across all levels of annealing in the present model, therefore, remains an important direction for future work.

\DIFaddbegin \newpage
\begin{subappendices}

\section{\DIFadd{Simulation Ensemble Sizes}}
\label{sec:FatigueSampleCounts}

\DIFadd{For each parameter tuple $(N,l_0,l_w,\gamma_0)$, we define the ensemble size $n_{\mathrm{ens}}$ as the number of independently seeded simulation trajectories included in the analysis.The reported $n_{\mathrm{ens}}$ is the total number of trajectories for each parameter tuple, irrespective of whether they ultimately absorb or yield. The ensemble sizes used for the annealing, post-hop-width, and system-size sweeps are reported in }\cref{tab:FatigueSampleCountsL0,tab:FatigueSampleCountsLw,tab:FatigueSampleCountsN}\DIFadd{, respectively. These values give the number of samples available for each parameter tuple used in the corresponding figures. Here $N=N_x=N_y$ denotes the lattice size.
}

\begingroup
\footnotesize
\setlength{\tabcolsep}{4pt}
\renewcommand{\arraystretch}{0.9}
\setlength{\LTcapwidth}{\textwidth}
\begin{singlespace}
\begin{longtable}{@{}rrr@{\hspace{2.2em}}rrr@{\hspace{2.2em}}rrr@{}}
\caption[Exact ensemble sizes for the annealing sweep.]{\DIFadd{Exact ensemble sizes for the parameter combinations in }\texttt{\DIFadd{EndStates.toml}} \DIFadd{at $N=1024$ and $l_w=0.02$, used in }\cref{fig:EndStates,fig:LogisticsL0,fig:FailureCyclel0}\DIFadd{. The survival analyses in }\cref{fig:SurvivalDistrbutions,fig:SurvivalDistrbutionsFits} \DIFadd{use selected rows from the same ensembles. Entries are ordered by $l_0$ and then $\gamma_0$, read from left to right across each row.}}
\label{tab:FatigueSampleCountsL0}\\
\toprule
\DIFadd{$l_0$ }& \DIFadd{$\gamma_0$ }& \DIFadd{$n_{\mathrm{ens}}$ }& \DIFadd{$l_0$ }& \DIFadd{$\gamma_0$ }& \DIFadd{$n_{\mathrm{ens}}$ }& \DIFadd{$l_0$ }& \DIFadd{$\gamma_0$ }& \DIFadd{$n_{\mathrm{ens}}$ }\\
\midrule
\endfirsthead

\multicolumn{9}{c}{\tablename~\thetable\ continued} \\
\toprule
\DIFadd{$l_0$ }& \DIFadd{$\gamma_0$ }& \DIFadd{$n_{\mathrm{ens}}$ }& \DIFadd{$l_0$ }& \DIFadd{$\gamma_0$ }& \DIFadd{$n_{\mathrm{ens}}$ }& \DIFadd{$l_0$ }& \DIFadd{$\gamma_0$ }& \DIFadd{$n_{\mathrm{ens}}$ }\\
\midrule
\endhead

\midrule
\multicolumn{9}{r}{Continued on next page} \\
\endfoot

\bottomrule
\endlastfoot

\DIFadd{0.025 }& \DIFadd{0.9 }& \DIFadd{$10{,}115$ }& \DIFadd{0.025 }& \DIFadd{0.903 }& \DIFadd{$10{,}116$ }& \DIFadd{0.025 }& \DIFadd{0.906 }& \DIFadd{$23{,}156$ }\\
\DIFadd{0.025 }& \DIFadd{0.909 }& \DIFadd{$114{,}311$ }& \DIFadd{0.025 }& \DIFadd{0.912 }& \DIFadd{$42{,}120$ }& \DIFadd{0.025 }& \DIFadd{0.915 }& \DIFadd{$15{,}324$ }\\
\DIFadd{0.025 }& \DIFadd{0.918 }& \DIFadd{$12{,}516$ }& \DIFadd{0.025 }& \DIFadd{0.921 }& \DIFadd{$47{,}493$ }& \DIFadd{0.025 }& \DIFadd{0.924 }& \DIFadd{$10{,}116$ }\\
\DIFadd{0.025 }& \DIFadd{0.927 }& \DIFadd{$9{,}997$ }& \DIFadd{0.025 }& \DIFadd{0.93 }& \DIFadd{$9{,}997$ }& \DIFadd{0.0375 }& \DIFadd{0.864 }& \DIFadd{$12{,}517$ }\\
\DIFadd{0.0375 }& \DIFadd{0.867 }& \DIFadd{$13{,}472$ }& \DIFadd{0.0375 }& \DIFadd{0.87 }& \DIFadd{$23{,}559$ }& \DIFadd{0.0375 }& \DIFadd{0.873 }& \DIFadd{$26{,}503$ }\\
\DIFadd{0.0375 }& \DIFadd{0.876 }& \DIFadd{$17{,}955$ }& \DIFadd{0.0375 }& \DIFadd{0.879 }& \DIFadd{$49{,}492$ }& \DIFadd{0.0375 }& \DIFadd{0.882 }& \DIFadd{$24{,}997$ }\\
\DIFadd{0.0375 }& \DIFadd{0.885 }& \DIFadd{$74{,}099$ }& \DIFadd{0.0375 }& \DIFadd{0.888 }& \DIFadd{$24{,}997$ }& \DIFadd{0.0375 }& \DIFadd{0.891 }& \DIFadd{$21{,}565$ }\\
\DIFadd{0.0375 }& \DIFadd{0.894 }& \DIFadd{$10{,}118$ }& \DIFadd{0.05 }& \DIFadd{0.839 }& \DIFadd{$25{,}000$ }& \DIFadd{0.05 }& \DIFadd{0.842 }& \DIFadd{$25{,}000$ }\\
\DIFadd{0.05 }& \DIFadd{0.845 }& \DIFadd{$25{,}000$ }& \DIFadd{0.05 }& \DIFadd{0.848 }& \DIFadd{$25{,}000$ }& \DIFadd{0.05 }& \DIFadd{0.851 }& \DIFadd{$25{,}000$ }\\
\DIFadd{0.05 }& \DIFadd{0.854 }& \DIFadd{$25{,}000$ }& \DIFadd{0.05 }& \DIFadd{0.857 }& \DIFadd{$25{,}000$ }& \DIFadd{0.05 }& \DIFadd{0.86 }& \DIFadd{$25{,}000$ }\\
\DIFadd{0.05 }& \DIFadd{0.863 }& \DIFadd{$25{,}000$ }& \DIFadd{0.05 }& \DIFadd{0.866 }& \DIFadd{$25{,}000$ }& \DIFadd{0.05 }& \DIFadd{0.869 }& \DIFadd{$25{,}000$ }\\
\DIFadd{0.075 }& \DIFadd{0.805 }& \DIFadd{$83{,}959$ }& \DIFadd{0.075 }& \DIFadd{0.808 }& \DIFadd{$25{,}000$ }& \DIFadd{0.075 }& \DIFadd{0.811 }& \DIFadd{$25{,}000$ }\\
\DIFadd{0.075 }& \DIFadd{0.814 }& \DIFadd{$25{,}000$ }& \DIFadd{0.075 }& \DIFadd{0.817 }& \DIFadd{$25{,}000$ }& \DIFadd{0.075 }& \DIFadd{0.82 }& \DIFadd{$25{,}000$ }\\
\DIFadd{0.075 }& \DIFadd{0.823 }& \DIFadd{$25{,}000$ }& \DIFadd{0.075 }& \DIFadd{0.826 }& \DIFadd{$74{,}065$ }& \DIFadd{0.075 }& \DIFadd{0.829 }& \DIFadd{$84{,}148$ }\\
\DIFadd{0.075 }& \DIFadd{0.832 }& \DIFadd{$25{,}000$ }& \DIFadd{0.075 }& \DIFadd{0.835 }& \DIFadd{$25{,}000$ }& \DIFadd{0.1 }& \DIFadd{0.781 }& \DIFadd{$3{,}997$ }\\
\DIFadd{0.1 }& \DIFadd{0.784 }& \DIFadd{$3{,}997$ }& \DIFadd{0.1 }& \DIFadd{0.787 }& \DIFadd{$3{,}996$ }& \DIFadd{0.1 }& \DIFadd{0.79 }& \DIFadd{$3{,}997$ }\\
\DIFadd{0.1 }& \DIFadd{0.793 }& \DIFadd{$3{,}998$ }& \DIFadd{0.1 }& \DIFadd{0.796 }& \DIFadd{$8{,}179$ }& \DIFadd{0.1 }& \DIFadd{0.799 }& \DIFadd{$24{,}981$ }\\
\DIFadd{0.1 }& \DIFadd{0.802 }& \DIFadd{$7{,}241$ }& \DIFadd{0.1 }& \DIFadd{0.805 }& \DIFadd{$3{,}994$ }& \DIFadd{0.1 }& \DIFadd{0.808 }& \DIFadd{$3{,}997$ }\\
\DIFadd{0.1 }& \DIFadd{0.811 }& \DIFadd{$3{,}997$ }& \DIFadd{0.125 }& \DIFadd{0.775 }& \DIFadd{$3{,}438$ }& \DIFadd{0.125 }& \DIFadd{0.778 }& \DIFadd{$3{,}428$ }\\
\DIFadd{0.125 }& \DIFadd{0.781 }& \DIFadd{$3{,}955$ }& \DIFadd{0.125 }& \DIFadd{0.784 }& \DIFadd{$4{,}483$ }& \DIFadd{0.125 }& \DIFadd{0.787 }& \DIFadd{$3{,}412$ }\\
\DIFadd{0.125 }& \DIFadd{0.79 }& \DIFadd{$3{,}440$ }& \DIFadd{0.125 }& \DIFadd{0.793 }& \DIFadd{$3{,}463$ }& \DIFadd{0.125 }& \DIFadd{0.796 }& \DIFadd{$3{,}470$ }\\
\DIFadd{0.125 }& \DIFadd{0.799 }& \DIFadd{$3{,}474$ }& \DIFadd{0.125 }& \DIFadd{0.802 }& \DIFadd{$3{,}479$ }& \DIFadd{0.125 }& \DIFadd{0.805 }& \DIFadd{$3{,}477$ }\\
\end{longtable}
\end{singlespace}
\endgroup

\newpage
\begingroup
\footnotesize
\setlength{\tabcolsep}{14pt}
\renewcommand{\arraystretch}{0.9}
\setlength{\LTcapwidth}{\textwidth}
\begin{singlespace}
\begin{longtable}{@{}rrrrr@{}}
\caption[Exact ensemble sizes for the post-hop-width sweep.]{\DIFadd{Exact ensemble sizes for the post-hop-width sweep at $N=1024$ and $l_0=0.075$, used in }\cref{fig:FailureCyclelw}\DIFadd{. Entries are $n_{\mathrm{ens}}$; a dash denotes a parameter combination absent from the reported sweep.}}
\label{tab:FatigueSampleCountsLw}\\
\toprule
\DIFadd{$\gamma_0$ }& \multicolumn{4}{c}{$l_w$} \\
\cmidrule(l){2-5}
 & \DIFadd{$0.01$ }& \DIFadd{$0.02$ }& \DIFadd{$0.03$ }& \DIFadd{$0.04$ }\\
\midrule
\endfirsthead

\multicolumn{5}{c}{\tablename~\thetable\ continued} \\
\toprule
\DIFadd{$\gamma_0$ }& \multicolumn{4}{c}{$l_w$} \\
\cmidrule(l){2-5}
 & \DIFadd{$0.01$ }& \DIFadd{$0.02$ }& \DIFadd{$0.03$ }& \DIFadd{$0.04$ }\\
\midrule
\endhead

\midrule
\multicolumn{5}{r}{Continued on next page} \\
\endfoot

\bottomrule
\endlastfoot

\DIFadd{$0.751$ }& \DIFadd{$4{,}999$ }& \DIFadd{$4{,}999$ }& \DIFadd{$4{,}999$ }& \DIFadd{$4{,}999$ }\\
\DIFadd{$0.754$ }& \DIFadd{$12{,}520$ }& \DIFadd{$12{,}520$ }& \DIFadd{$12{,}520$ }& \DIFadd{$12{,}520$ }\\
\DIFadd{$0.757$ }& \DIFadd{$12{,}520$ }& \DIFadd{$12{,}520$ }& \DIFadd{$12{,}520$ }& \DIFadd{$12{,}520$ }\\
\DIFadd{$0.76$ }& \DIFadd{$12{,}520$ }& \DIFadd{$12{,}520$ }& \DIFadd{$12{,}520$ }& \DIFadd{$12{,}520$ }\\
\DIFadd{$0.763$ }& \DIFadd{$12{,}518$ }& \DIFadd{$12{,}518$ }& \DIFadd{$12{,}518$ }& \DIFadd{$12{,}518$ }\\
\DIFadd{$0.766$ }& \DIFadd{$12{,}520$ }& \DIFadd{$12{,}519$ }& \DIFadd{$12{,}520$ }& \DIFadd{$12{,}520$ }\\
\DIFadd{$0.769$ }& \DIFadd{$12{,}518$ }& \DIFadd{$12{,}518$ }& \DIFadd{$12{,}518$ }& \DIFadd{$12{,}518$ }\\
\DIFadd{$0.772$ }& \DIFadd{$12{,}520$ }& \DIFadd{$12{,}519$ }& \DIFadd{$12{,}520$ }& \DIFadd{$12{,}520$ }\\
\DIFadd{$0.775$ }& \DIFadd{$12{,}518$ }& \DIFadd{$12{,}518$ }& \DIFadd{$12{,}518$ }& \DIFadd{$12{,}518$ }\\
\DIFadd{$0.778$ }& \DIFadd{$12{,}521$ }& \DIFadd{$12{,}519$ }& \DIFadd{$12{,}520$ }& \DIFadd{$12{,}520$ }\\
\DIFadd{$0.781$ }& \DIFadd{$12{,}518$ }& \DIFadd{$12{,}518$ }& \DIFadd{$12{,}518$ }& \DIFadd{$12{,}518$ }\\
\DIFadd{$0.784$ }& \DIFadd{$12{,}521$ }& \DIFadd{$12{,}519$ }& \DIFadd{$12{,}520$ }& \DIFadd{$12{,}520$ }\\
\DIFadd{$0.787$ }& \DIFadd{$12{,}518$ }& \DIFadd{$12{,}519$ }& \DIFadd{$12{,}518$ }& \DIFadd{$12{,}518$ }\\
\DIFadd{$0.79$ }& \DIFadd{$12{,}520$ }& \DIFadd{$12{,}520$ }& \DIFadd{$12{,}521$ }& \DIFadd{$12{,}520$ }\\
\DIFadd{$0.793$ }& \DIFadd{$12{,}518$ }& \DIFadd{$12{,}519$ }& \DIFadd{$12{,}519$ }& \DIFadd{$12{,}519$ }\\
\DIFadd{$0.796$ }& \DIFadd{$12{,}521$ }& \DIFadd{$12{,}520$ }& \DIFadd{$12{,}521$ }& \DIFadd{$12{,}520$ }\\
\DIFadd{$0.799$ }& \DIFadd{--- }& \DIFadd{$84{,}087$ }& \DIFadd{$12{,}519$ }& \DIFadd{$12{,}519$ }\\
\DIFadd{$0.802$ }& \DIFadd{$12{,}521$ }& \DIFadd{$12{,}519$ }& \DIFadd{$12{,}521$ }& \DIFadd{$12{,}521$ }\\
\DIFadd{$0.805$ }& \DIFadd{$25{,}000$ }& \DIFadd{$83{,}959$ }& \DIFadd{$25{,}000$ }& \DIFadd{$25{,}000$ }\\
\DIFadd{$0.808$ }& \DIFadd{$25{,}000$ }& \DIFadd{$25{,}000$ }& \DIFadd{$25{,}000$ }& \DIFadd{$25{,}000$ }\\
\DIFadd{$0.811$ }& \DIFadd{$25{,}000$ }& \DIFadd{$25{,}000$ }& \DIFadd{$25{,}000$ }& \DIFadd{$25{,}000$ }\\
\DIFadd{$0.814$ }& \DIFadd{$24{,}997$ }& \DIFadd{$25{,}000$ }& \DIFadd{$25{,}000$ }& \DIFadd{$25{,}000$ }\\
\DIFadd{$0.817$ }& \DIFadd{$24{,}991$ }& \DIFadd{$25{,}000$ }& \DIFadd{$25{,}000$ }& \DIFadd{$25{,}000$ }\\
\DIFadd{$0.82$ }& \DIFadd{$24{,}984$ }& \DIFadd{$25{,}000$ }& \DIFadd{$25{,}000$ }& \DIFadd{$25{,}000$ }\\
\DIFadd{$0.823$ }& \DIFadd{$24{,}988$ }& \DIFadd{$25{,}000$ }& \DIFadd{$25{,}000$ }& \DIFadd{$25{,}000$ }\\
\DIFadd{$0.826$ }& \DIFadd{$73{,}988$ }& \DIFadd{$74{,}065$ }& \DIFadd{$74{,}080$ }& \DIFadd{$25{,}000$ }\\
\DIFadd{$0.829$ }& \DIFadd{$25{,}000$ }& \DIFadd{$84{,}148$ }& \DIFadd{$25{,}000$ }& \DIFadd{$25{,}000$ }\\
\DIFadd{$0.832$ }& \DIFadd{$25{,}000$ }& \DIFadd{$25{,}000$ }& \DIFadd{$25{,}000$ }& \DIFadd{$25{,}000$ }\\
\DIFadd{$0.835$ }& \DIFadd{$25{,}000$ }& \DIFadd{$25{,}000$ }& \DIFadd{$25{,}000$ }& \DIFadd{$25{,}000$ }\\
\DIFadd{$0.838$ }& \DIFadd{--- }& \DIFadd{$3{,}480$ }& \DIFadd{--- }& \DIFadd{--- }\\
\end{longtable}
\end{singlespace}
\endgroup

\begingroup
\footnotesize
\setlength{\tabcolsep}{10pt}
\renewcommand{\arraystretch}{0.9}
\setlength{\LTcapwidth}{\textwidth}
\begin{singlespace}
\begin{longtable}{@{}rrrrrrr@{}}
\caption[Exact ensemble sizes for the system-size sweep.]{\DIFadd{Exact ensemble sizes for the system-size sweep at $l_0=0.075$ and $l_w=0.02$, used in }\cref{fig:EPSizeDependence,fig:LogisticsN}\DIFadd{. Entries are $n_{\mathrm{ens}}$; a dash denotes a parameter combination absent from the reported sweep.}}
\label{tab:FatigueSampleCountsN}\\
\toprule
\DIFadd{$\gamma_0$ }& \multicolumn{6}{c}{$N$} \\
\cmidrule(l){2-7}
 & \DIFadd{$64$ }& \DIFadd{$128$ }& \DIFadd{$256$ }& \DIFadd{$512$ }& \DIFadd{$1024$ }& \DIFadd{$2048$ }\\
\midrule
\endfirsthead

\multicolumn{7}{c}{\tablename~\thetable\ continued} \\
\toprule
\DIFadd{$\gamma_0$ }& \multicolumn{6}{c}{$N$} \\
\cmidrule(l){2-7}
 & \DIFadd{$64$ }& \DIFadd{$128$ }& \DIFadd{$256$ }& \DIFadd{$512$ }& \DIFadd{$1024$ }& \DIFadd{$2048$ }\\
\midrule
\endhead

\midrule
\multicolumn{7}{r}{Continued on next page} \\
\endfoot

\bottomrule
\endlastfoot

\DIFadd{$0.751$ }& \DIFadd{$25{,}000$ }& \DIFadd{$25{,}000$ }& \DIFadd{$25{,}000$ }& \DIFadd{$25{,}000$ }& \DIFadd{$4{,}999$ }& \DIFadd{$5{,}000$ }\\
\DIFadd{$0.754$ }& \DIFadd{$25{,}000$ }& \DIFadd{$25{,}000$ }& \DIFadd{$25{,}000$ }& \DIFadd{$25{,}000$ }& \DIFadd{$12{,}520$ }& \DIFadd{$5{,}000$ }\\
\DIFadd{$0.757$ }& \DIFadd{$25{,}000$ }& \DIFadd{$25{,}000$ }& \DIFadd{$25{,}000$ }& \DIFadd{$25{,}000$ }& \DIFadd{$12{,}520$ }& \DIFadd{$5{,}000$ }\\
\DIFadd{$0.76$ }& \DIFadd{$25{,}000$ }& \DIFadd{$25{,}000$ }& \DIFadd{$25{,}000$ }& \DIFadd{$25{,}000$ }& \DIFadd{$12{,}520$ }& \DIFadd{$5{,}000$ }\\
\DIFadd{$0.763$ }& \DIFadd{$25{,}000$ }& \DIFadd{$25{,}000$ }& \DIFadd{$25{,}000$ }& \DIFadd{$25{,}000$ }& \DIFadd{$12{,}518$ }& \DIFadd{$5{,}000$ }\\
\DIFadd{$0.766$ }& \DIFadd{$25{,}000$ }& \DIFadd{$25{,}000$ }& \DIFadd{$25{,}000$ }& \DIFadd{$25{,}000$ }& \DIFadd{$12{,}519$ }& \DIFadd{$5{,}000$ }\\
\DIFadd{$0.769$ }& \DIFadd{$25{,}000$ }& \DIFadd{$25{,}000$ }& \DIFadd{$25{,}000$ }& \DIFadd{$25{,}000$ }& \DIFadd{$12{,}518$ }& \DIFadd{$5{,}000$ }\\
\DIFadd{$0.772$ }& \DIFadd{$25{,}000$ }& \DIFadd{$25{,}000$ }& \DIFadd{$25{,}000$ }& \DIFadd{$25{,}000$ }& \DIFadd{$12{,}519$ }& \DIFadd{$5{,}000$ }\\
\DIFadd{$0.775$ }& \DIFadd{$25{,}000$ }& \DIFadd{$25{,}000$ }& \DIFadd{$25{,}000$ }& \DIFadd{$25{,}000$ }& \DIFadd{$12{,}518$ }& \DIFadd{$5{,}000$ }\\
\DIFadd{$0.778$ }& \DIFadd{$25{,}000$ }& \DIFadd{$25{,}000$ }& \DIFadd{$25{,}000$ }& \DIFadd{$25{,}000$ }& \DIFadd{$12{,}519$ }& \DIFadd{$5{,}000$ }\\
\DIFadd{$0.781$ }& \DIFadd{$25{,}000$ }& \DIFadd{$25{,}000$ }& \DIFadd{$25{,}000$ }& \DIFadd{$25{,}000$ }& \DIFadd{$12{,}518$ }& \DIFadd{$5{,}000$ }\\
\DIFadd{$0.784$ }& \DIFadd{$25{,}000$ }& \DIFadd{$25{,}000$ }& \DIFadd{$25{,}000$ }& \DIFadd{$25{,}000$ }& \DIFadd{$12{,}519$ }& \DIFadd{$5{,}000$ }\\
\DIFadd{$0.787$ }& \DIFadd{$25{,}000$ }& \DIFadd{$25{,}000$ }& \DIFadd{$25{,}000$ }& \DIFadd{$25{,}000$ }& \DIFadd{$12{,}519$ }& \DIFadd{$5{,}000$ }\\
\DIFadd{$0.79$ }& \DIFadd{$25{,}000$ }& \DIFadd{$25{,}000$ }& \DIFadd{$25{,}000$ }& \DIFadd{$25{,}000$ }& \DIFadd{$12{,}520$ }& \DIFadd{$5{,}000$ }\\
\DIFadd{$0.793$ }& \DIFadd{$25{,}000$ }& \DIFadd{$25{,}000$ }& \DIFadd{$25{,}000$ }& \DIFadd{$25{,}000$ }& \DIFadd{$12{,}519$ }& \DIFadd{$5{,}000$ }\\
\DIFadd{$0.796$ }& \DIFadd{$25{,}000$ }& \DIFadd{$25{,}000$ }& \DIFadd{$25{,}000$ }& \DIFadd{$25{,}000$ }& \DIFadd{$12{,}520$ }& \DIFadd{$5{,}000$ }\\
\DIFadd{$0.799$ }& \DIFadd{$25{,}000$ }& \DIFadd{$25{,}000$ }& \DIFadd{$25{,}000$ }& \DIFadd{$25{,}000$ }& \DIFadd{$84{,}087$ }& \DIFadd{$5{,}000$ }\\
\DIFadd{$0.802$ }& \DIFadd{$25{,}000$ }& \DIFadd{$25{,}000$ }& \DIFadd{$25{,}000$ }& \DIFadd{$25{,}000$ }& \DIFadd{$12{,}519$ }& \DIFadd{$5{,}000$ }\\
\DIFadd{$0.805$ }& \DIFadd{$25{,}000$ }& \DIFadd{$25{,}000$ }& \DIFadd{$25{,}000$ }& \DIFadd{$25{,}000$ }& \DIFadd{$83{,}959$ }& \DIFadd{$5{,}000$ }\\
\DIFadd{$0.808$ }& \DIFadd{$25{,}000$ }& \DIFadd{$25{,}000$ }& \DIFadd{$25{,}000$ }& \DIFadd{$25{,}000$ }& \DIFadd{$25{,}000$ }& \DIFadd{$5{,}000$ }\\
\DIFadd{$0.811$ }& \DIFadd{$25{,}000$ }& \DIFadd{$25{,}000$ }& \DIFadd{$25{,}000$ }& \DIFadd{$25{,}000$ }& \DIFadd{$25{,}000$ }& \DIFadd{$5{,}000$ }\\
\DIFadd{$0.814$ }& \DIFadd{$25{,}000$ }& \DIFadd{$25{,}000$ }& \DIFadd{$25{,}000$ }& \DIFadd{$25{,}000$ }& \DIFadd{$25{,}000$ }& \DIFadd{$5{,}000$ }\\
\DIFadd{$0.817$ }& \DIFadd{$25{,}000$ }& \DIFadd{$25{,}000$ }& \DIFadd{$25{,}000$ }& \DIFadd{$25{,}000$ }& \DIFadd{$25{,}000$ }& \DIFadd{$5{,}000$ }\\
\DIFadd{$0.82$ }& \DIFadd{$25{,}000$ }& \DIFadd{$25{,}000$ }& \DIFadd{$25{,}000$ }& \DIFadd{$25{,}000$ }& \DIFadd{$25{,}000$ }& \DIFadd{$4{,}992$ }\\
\DIFadd{$0.823$ }& \DIFadd{$25{,}000$ }& \DIFadd{$25{,}000$ }& \DIFadd{$25{,}000$ }& \DIFadd{$25{,}000$ }& \DIFadd{$25{,}000$ }& \DIFadd{$4{,}999$ }\\
\DIFadd{$0.826$ }& \DIFadd{$25{,}000$ }& \DIFadd{$25{,}000$ }& \DIFadd{$25{,}000$ }& \DIFadd{$25{,}000$ }& \DIFadd{$74{,}065$ }& \DIFadd{$5{,}000$ }\\
\DIFadd{$0.829$ }& \DIFadd{$25{,}000$ }& \DIFadd{$25{,}000$ }& \DIFadd{$25{,}000$ }& \DIFadd{$25{,}000$ }& \DIFadd{$84{,}148$ }& \DIFadd{$5{,}000$ }\\
\DIFadd{$0.832$ }& \DIFadd{$25{,}000$ }& \DIFadd{$25{,}000$ }& \DIFadd{$25{,}000$ }& \DIFadd{$25{,}000$ }& \DIFadd{$25{,}000$ }& \DIFadd{$5{,}000$ }\\
\DIFadd{$0.835$ }& \DIFadd{$25{,}000$ }& \DIFadd{$25{,}000$ }& \DIFadd{$25{,}000$ }& \DIFadd{$25{,}000$ }& \DIFadd{$25{,}000$ }& \DIFadd{$5{,}000$ }\\
\DIFadd{$0.838$ }& \DIFadd{$25{,}000$ }& \DIFadd{$25{,}000$ }& \DIFadd{$25{,}000$ }& \DIFadd{$25{,}000$ }& \DIFadd{$3{,}480$ }& \DIFadd{$5{,}000$ }\\
\DIFadd{$0.841$ }& \DIFadd{$25{,}000$ }& \DIFadd{$25{,}000$ }& \DIFadd{$25{,}000$ }& \DIFadd{$25{,}000$ }& \DIFadd{$3{,}480$ }& \DIFadd{$5{,}000$ }\\
\DIFadd{$0.844$ }& \DIFadd{$25{,}000$ }& \DIFadd{$25{,}000$ }& \DIFadd{$25{,}000$ }& \DIFadd{--- }& \DIFadd{$3{,}480$ }& \DIFadd{--- }\\
\DIFadd{$0.847$ }& \DIFadd{$25{,}000$ }& \DIFadd{$25{,}000$ }& \DIFadd{$25{,}000$ }& \DIFadd{--- }& \DIFadd{$3{,}480$ }& \DIFadd{--- }\\
\DIFadd{$0.85$ }& \DIFadd{$25{,}000$ }& \DIFadd{$25{,}000$ }& \DIFadd{--- }& \DIFadd{--- }& \DIFadd{$3{,}480$ }& \DIFadd{--- }\\
\DIFadd{$0.853$ }& \DIFadd{$25{,}000$ }& \DIFadd{$25{,}000$ }& \DIFadd{--- }& \DIFadd{--- }& \DIFadd{--- }& \DIFadd{--- }\\
\DIFadd{$0.856$ }& \DIFadd{$25{,}000$ }& \DIFadd{$25{,}000$ }& \DIFadd{--- }& \DIFadd{--- }& \DIFadd{--- }& \DIFadd{--- }\\
\DIFadd{$0.859$ }& \DIFadd{$25{,}000$ }& \DIFadd{$25{,}000$ }& \DIFadd{--- }& \DIFadd{--- }& \DIFadd{--- }& \DIFadd{--- }\\
\DIFadd{$0.862$ }& \DIFadd{$25{,}000$ }& \DIFadd{$25{,}000$ }& \DIFadd{--- }& \DIFadd{--- }& \DIFadd{--- }& \DIFadd{--- }\\
\DIFadd{$0.865$ }& \DIFadd{$25{,}000$ }& \DIFadd{$25{,}000$ }& \DIFadd{--- }& \DIFadd{--- }& \DIFadd{--- }& \DIFadd{--- }\\
\DIFadd{$0.868$ }& \DIFadd{$25{,}000$ }& \DIFadd{$25{,}000$ }& \DIFadd{--- }& \DIFadd{--- }& \DIFadd{--- }& \DIFadd{--- }\\
\DIFadd{$0.871$ }& \DIFadd{$25{,}000$ }& \DIFadd{$25{,}000$ }& \DIFadd{--- }& \DIFadd{--- }& \DIFadd{--- }& \DIFadd{--- }\\
\DIFadd{$0.874$ }& \DIFadd{$25{,}000$ }& \DIFadd{$25{,}000$ }& \DIFadd{--- }& \DIFadd{--- }& \DIFadd{--- }& \DIFadd{--- }\\
\DIFadd{$0.877$ }& \DIFadd{$25{,}000$ }& \DIFadd{$25{,}000$ }& \DIFadd{--- }& \DIFadd{--- }& \DIFadd{--- }& \DIFadd{--- }\\
\DIFadd{$0.88$ }& \DIFadd{$25{,}000$ }& \DIFadd{$25{,}000$ }& \DIFadd{--- }& \DIFadd{--- }& \DIFadd{--- }& \DIFadd{--- }\\
\end{longtable}
\end{singlespace}
\endgroup


\end{subappendices}

\DIFaddend % \begin{itemize}
    % \item Consider the low Reynolds limit, where inertia is negligible 
    % \item We define a lattice of $i=1,2,\ldots,N_x$ by $j=1,2,\ldots,N_y$, and take $N_x=N_y=N$ for simplicity.
    % \item Each lattice site/element is considered to be a mesoscopic region of the material, typically approximately the size of a rearrangement event (bond or cluster of particles)
    % \item Each element has a local strain $l$, a corresponding local shear stress $\mu l$, and a local yield energy $E$.
    % \item The material is assumed to undergo a globally applied bulk strain $\gamma$, which is assumed to be applied affinely to the material.
    % \item Each element is assumed to follow the same set of rules
    % \begin{itemize}
        % \item Between plastic events, the strain of each element affinely follows the globally applied strain $\gamma$, corresponding to elastic loading.
        % \item Each element has a local yield criterion ($\frac{1}{2} \mu l^2=E$) which determines the onset of a plastic event.
        % \item Work in units where the elastic constant $\mu=1$ and $E=0.5$ such that the critical strain $l_c = \pm 1$
        % \item Upon yielding a element adopts a new local strain chosen at random from a Gaussian distribution of width $\lw$ and mean $0$.
        % \item The plastic event forces a rearrangement of particles; the system is not in force balance, which needs to be reimposed. This is done by yielding the element and applying force balance using \cref{eq:eshelbyReal}.
        % \item The elastic stress lost from a plastic event is redistributed to neighbouring blocks through the Eshelby stress propagator (defined in \cref{sec:EshelbyPropogator})
        % \item Define the propagator here, make it clear the constant at $\hat{G}[0,0] = -1$, such that the stress that would be lost at a single element is instead spread across the system.
    % \end{itemize}
    % \item define both $\Origin$ and $\Norm$ here for use later.
    % \begin{figure}[hbt]
    %     \centering
    %     \includegraphics[width=0.9\textwidth]{chapters/Fatigue/Diagrams/Prop.pdf}
    %     \caption{{\bf Left)} Average stress redistribution around a shear transformation in $2D$ {\bf Right)} Convergence of two important parameters from $\mathbb{G}_{xy}$ with grid size $N$. The value at the origin $\Origin$ and norm $\Norm$.}
    %     \label{fig:PropogatorConvergence}
    % \end{figure}
    % \item We assign the initial distribution of local strains $l$ from a Gaussian distribution of width $l_0$, this sets the level of annealing. We then immediately impose force balance on the system, which has the effect of scaling the distribution of local strains by a factor $\Norm$.
    % \item Introduce the concept of $\ltilde$, the deviation in $l$ from the bulk strain $\gamma$.
    % \item We impose cyclic shear deformation with amplitude $\gammaMax$.
    % \item The deformation is applied quasi-statically: the strain is incremented in small steps, and after each increment, the system is fully relaxed to mechanical equilibrium, allowing any plastic events to complete before further loading.
    % \item Describe the implementation of jumping to the next plastic event and then evolving any avalanches. 

    % \item define macrosopic shear stress

    % \item Need to explicitly define end states of the system. 
    % \begin{itemize}
    %     \item Yielded states are where the displayed in a shear band.
    %     \item Absorbed states where the system has no element which can yield.
    % \end{itemize}

    
    % \[\Sigma = \mu\sum_{i=1}^{N_x}\sum_{j=1}^{N_y} l_{ij}.\]
    % \item State default parameters $N=1024$, $\lw=0.02$
    % \item Reiterate the $l_0$ is an important parameter
% \end{itemize}

% \section{Shear Startup}
% \begin{itemize}
    % \item I think we need to discuss the shear startup here, as it gives some background to each half cycle.
    % \todo[inline]{Make Figure from Pollard paper - already have data}
    % \item As $l_0$ increases at fixed system size $N$, the system becomes more ductile, and the failure strain decreases.
    % \item As $N$ increases, the system also becomes more brittle
    % \item Systems ultimately fail through shear banding, which is more akin to failures in metallic glasses, rather than the coexistence of macroscopic bands of different strain rates.
    % \todo[inline]{This needs discussing - Do I need to include it?}
% \end{itemize}

% \section{Elasto-Plastic Model Results}

% \begin{itemize}
    % \item We ultimately do not resolve the system any further once shear banding occurs, considering only the path to failure or absorption.
    % \item We begin with the basic physics. As the system oscillates, it will end up in one of two states. Either it will shear band as a result of collective plastic events. Or will absorb with no elements able to yield.
    % \todo[inline]{Discuss max $l_0$ with Suzanne}

    % \item If we consider the extreme cases:
    % \begin{itemize}
    %     \item If $\gamma_0$ is not large enough, such that $|\ltilde_{ij} \pm \gamma_0| < l_c \forall \{ i = 1, \ldots, N_x;\; j = 1, \ldots, N_y \}$ then no yielding happens.
    %     \item If $\gamma_0$ is too large, then the system acts like the shear startup cases and shear bands during the first cycle. 
    % \end{itemize}
    % \item it is in the intermediate regime where delayed failure and absorption can occur over many cycles, we illustrate this in \cref{fig:EndStates}. This is the main focus of this work
    % \item For each level of annealing, we transition from always absorbing in the first cycle, to delaying this absorption as elements $\ltilde$ decay towards $0$ as we increase $\gamma_0$. We then see yielding occurring at a late number of cycles, before the system always forms a shear band in the first cycle. 
    % \todo[inline]{Need to work out how to describe this better - some notation? + quantify it?}
    % \item As the level of annealing increases, the strain point at which the system immediately forms a shear band decreases as predicted in the shear startup case, and the range of maximum strains where delayed failures and absorption occurs increase.
    % \begin{figure}[h!]
    %     \centering
    %     \includegraphics[width=0.9\textwidth]{chapters/Fatigue/Figures/Decay.pdf}
    %DIF <      \caption{Final cycle $\CStar$ as a function of the annealing parameter $l_0$. Smaller values of $l_0$ correspond to better annealed samples. We report the $95^\mathrm{th}$ percentile, $\CStar_{95\%}$, defined as the cycle at which the survival distributions in \cref{fig:SurvivalDistrbutions} intersect the red threshold line. Results are shown for systems that yield, absorb, or exhibit both behaviours, indicated by solid, dotted, and dashed lines, respectively. For each value of $\gamma_0$ and $l_0$, the distributions are made up of at least $10{,}000$ runs, each run with a different value of the random number seed.}
    %DIF >      \caption{Final cycle $\CStar$ as a function of the annealing parameter $l_0$. Smaller values of $l_0$ correspond to better annealed samples. We report the $95^\mathrm{th}$ percentile, $\CStar_{95\%}$, defined as the cycle at which the survival distributions in \cref{fig:SurvivalDistrbutions} intersect the red threshold line. Results are shown for systems that yield, absorb, or exhibit both behaviours, indicated by solid, dotted, and dashed lines, respectively. For each value of $\gamma_0$ and $l_0$, the distributions are made up of at least $1{,}000$ runs, each run with a different value of the random number seed.}
    %     \label{fig:FailureCyclel0}
    % \end{figure}
    % \item As we increase strain away from the immediate absorption cases, we see the delay in absorption correspondingly rise, peaking at a maximum strain, around where we begin to observe yielding occurring. Here we see yielding delayed to a long number of cycles, although occurring sooner than their absorbing counterparts. We show this in \cref{fig:FailureCyclel0}
    % \item The peak of these distributions shifts according to where the transition to immediate banding occurs.
    % \begin{figure}[hp!]
    %     \centering
    %     \includegraphics[width=1.0\textwidth]{chapters/Fatigue/Figures/Survival.pdf}
    %     \caption{Survival functions of the termination cycle $\CStar$, showing the fraction of initial seed values that remain active up to cycle $\CStar$ before reaching an absorbing state ({\bf Left}) or a yielding state ({\bf Right}).
    %     Panels (a–h) show the survival probability $S(\CStar)=1-F(\CStar)$, conditioned on end state and annealing parameter. From top to bottom, rows correspond to increasing annealing parameter $l_0 = 0.025, 0.05, 0.075, 0.1$. Within each panel, multiple curves represent different imposed strain amplitudes, specified in the right column for each $l_0$.}
    %     \label{fig:SurvivalDistrbutions}
    % \end{figure}
    % \item For completeness, we show a set of the corresponding survival functions for several $l_0$ and $\gamma_0$. In a) and b) for $l_0=0.025$, we note the vertical form of the distribution indicating the large number of early failures, where $\approx 95\%$ of samples either immediately absorb or fail. However, of those $5\%$ who do survive the first cycle, see note exponential scaling of $\CStar$. 
    % \item As $l_0$ increase we see a transition in the absorbed curves from immediate to delayed behaviour. As both $l_0$ and correspondingly $\gamma_0$ increases we see the minimum time to absorption increases; however, the tails of the distribution remain exponential. 
    % \item The corresponding patterns occur in failure where there exist regimes at high enough $l_0$ where yielding only occurs in a delayed fashion.
    % \item The exponential tails here indicate that failure can occur at any number of cycles; however, with an ever-decreasing likelihood 
    % \todo[inline]{Extract Decay constant}
% \end{itemize}

% \section{Robustness to $\lw$ and $N$} \label{sec:FatigueRobustness}

% \begin{itemize}
    % \item So far, we have only considered the case where we have a fixed system size $N$ and post hop distribution $\lw$. We know from previous work that the system size $N$ affects the transition from brittle to ductile behaviour, which plays an important role in the physics we have seen. 
    % \begin{figure}[h]
    %     \centering
    %     \includegraphics[width=0.9\textwidth]{chapters/Fatigue/Figures/lw.pdf}
    %DIF <      \caption{Final cycle $\CStar$ as a function of the parameter $\gamma_0$, with the annealing parameter fixed at $l_0 = 0.075$. Multiple curves correspond to different values of $\lw$, increasing from top to bottom. We report the $95^\mathrm{th}$ percentile, $\CStar_{95\%}$, defined as the cycle at which the survival distributions in \cref{fig:SurvivalDistrbutions} intersect the red threshold line. Results are shown for systems that yield, absorb, or exhibit both behaviours, indicated by solid, dotted, and dashed lines, respectively. The values of $\lw$ considered are $0.01, 0.02, 0.03, 0.04$. The insert shows a rescaling of the data showing the failure cycle scales approximately as the square of the maximum strain. For each value of $\gamma_0$ and $\lw$, the distributions are constructed from at least $10{,}000$ independent runs, each with a different random number seed.}
    %DIF >      \caption{Final cycle $\CStar$ as a function of the parameter $\gamma_0$, with the annealing parameter fixed at $l_0 = 0.075$. Multiple curves correspond to different values of $\lw$, increasing from top to bottom. We report the $95^\mathrm{th}$ percentile, $\CStar_{95\%}$, defined as the cycle at which the survival distributions in \cref{fig:SurvivalDistrbutions} intersect the red threshold line. Results are shown for systems that yield, absorb, or exhibit both behaviours, indicated by solid, dotted, and dashed lines, respectively. The values of $\lw$ considered are $0.01, 0.02, 0.03, 0.04$. The insert shows a rescaling of the data showing the failure cycle scales approximately as the square of the maximum strain. For each value of $\gamma_0$ and $\lw$, the distributions are constructed from at least $1{,}000$ independent runs, each with a different random number seed.}
    %     \label{fig:FailureCyclelw}
    % \end{figure}
    % \item In \cref{fig:FailureCyclelw} we plot the number of cycles to failure at $l_0=0.075$ and $N=1024$ for various widths of post-hop distributions. Here we note an effect with smaller values of $\lw$ corresponding to more delayed failures/absorption. 
    % \item The overall shapes of the distribution appear consistent, indicating the dynamics of the system are not dependent on this parameter. However, the characteristic timescale of the system does depend on this. 
    % \item In the insert, we plot the same data using a rescaled parametrisation showing the failure cycle scales as the square of the maximum strain up to a constant.
    % \begin{figure}[h]
    %     \centering
    %     \includegraphics[width=0.99\textwidth]{chapters/Fatigue/Figures/N.pdf}
    %     \caption[Effect of system size $N$ at fixed parameters $l_0 = 0.075$ and $\lw = 0.02$.]{Effect of system size $N$ at fixed parameters $l_0 = 0.075$ and $\lw = 0.02$. \newline
    %     {\bf a)} The $95^\mathrm{th}$ percentile failure cycle $\CStar_{95\%}$ as a function of imposed strain. Results are shown for system sizes $N = 128, 256, 512, 1024$, indicated in blue, orange, green, and red, respectively. \newline
    %     {\bf b)} The corresponding distribution of end states. 
    %     }
    %     \label{fig:EPSizeDependence}
    % \end{figure}
    % \item in \cref{fig:EPSizeDependence} we plot both the failure cycle $\CStar$ and end states. 
    % \item In a) we show that the time to absorb both decreases with decreasing system size, but also the strain, which increases (as seen in Pollard's paper). This is effectively how well the initial distribution is sampled.
    % \item The cycles to failure results appear to shift in strain, however reach a maximum at about the same value. Indicating different physics occurring for absorbed vs yielding cases.
    % \item b) is just there for completeness, I find it easier to understand with it. 
% \end{itemize}



% \section{Path to a single oscillation}

% \begin{itemize}
    % \item If we consider the final number of active elements just before the point of failure. 
    % \begin{figure}[h]
    %     \centering
    %     \includegraphics[width=0.9\textwidth]{chapters/Fatigue/Figures/1Osc.pdf}
    %     \caption[A sample of the final number of active elements remaining in the oscillating set immediately before the system reaches its end state.]{A sample of the final number of active elements remaining in the oscillating set immediately before the system reaches its end state. We distinguish between runs whose end states end in yielding and absorption. \newline 
    %     Parameter: $l_0 = 0.075$, $\lw = 0.02$, and $N = 1024$.}
    %     \label{fig:Pathto1}
    % \end{figure}
    % \item We see an overall trend for runs which end in yielding, where systems which take a long time to yield end up with only a few active elements remaining
    % \item This is definitely true for runs which end in absorption, which have very few elements remaining. 
    % \item If we consider the extreme case where there is only one element remaining, then the system is very simple and becomes solely driven by a random walk determined by the post hop distribution (we will discuss in \cref{sec:EPSimpleModel})
    % \begin{figure}[hp]
    %     \centering
    %     \includegraphics[width=0.9\textwidth]{chapters/Fatigue/Figures/RandomWalks.pdf}
    %     \caption[Stroboscopic evolution of element lengths measured at the start of each cycle.]{Stroboscopic evolution of element lengths measured at the start of each cycle. $\tilde{C}$ indicates the number of cycles after the random walk is assumed to have begun.\newline
    %     {\bf a)} Example trajectory of the oscillating element in a realisation that ultimately absorbs. The red dashed line denotes the absorbing boundary; the blue trajectory crosses this threshold, leading to absorption. \newline
    %     {\bf b)} Example trajectory in a realisation that yields. The blue curve shows the oscillating element, while the orange curve shows a neighbouring element. In this case, the neighbouring element crosses the red boundary first, triggering yielding. \newline
    %     {\bf c)} Ensemble of trajectories of the oscillating element, each terminated at the first crossing of the absorbing boundary (red dashed line). \newline
    %     {\bf d)} Standard deviation of the oscillating element displacement, $\sigma(\Delta l_{\mathrm{osc}})$, as a function of cycle number, demonstrating the expected diffusive scaling.}
    %     \label{fig:RandomWalkExamples}
    % \end{figure}
    % \item We illustrate an example of a random walk observed in \cref{fig:RandomWalkExamples} for {\bf a)} an absorbing case and  {\bf b)} a yielding case. 
    % \item These random walks occur when observed stroboscopically at the start of the cycle. Here, yielding has occurred twice, and $\ltilde$ has shifted by a small amount. This random walk continues to evolve until the element enters the absorbed region, where it is unable to yield (subject to no other elements yielding). Once all elements enter this region.
    % \item Alternatively, as a result of this walk and stress propagation, a neighbouring element which follows a coupled random walk can be displaced from the absorbed region and, as a result, begin to yield. This commonly initiates a shear banding event. 
    % \item We illustrate an ensemble of such random walks in {\bf c)}, resulting from an oscillating element, and plot the standard deviation of the ensemble's displacement in {\bf d)}, which grows as expected, proportional to $\sqrt{\tilde{C}}$.
% \end{itemize}

% \section{Origin of Random Walk} \label{sec:EPSimpleModel}
% \begin{itemize}
    % \item Concerning a simple model where elements are diffusely scattered across the system, then the subsequent iterations resulting from stress propagation are small. Therefore, elements can, in the simplest model, be considered independent.
    % \begin{figure}[hbtp]
    %     \centering
    %     \includegraphics[width=0.8\linewidth]{chapters/Fatigue/RandomWalk/SimpleFirstCycle.pdf}
    %     \caption{Illustration of how the distribution of $l$ changes in the first cycle. We display the distribution for several key points in the cycle a)-f).}
    %     \label{fig:SimpleFirstCycle}
    % \end{figure}
    % \item To construct a simple model, we first consider the evolution of an element in a single cycle. Use \cref{fig:SimpleFirstCycle} to explain the evolution of elements and how they flip signs in each half cycle.
    % \begin{figure}[hbtp]
    %     \centering
    %     \includegraphics[width=0.8\linewidth]{chapters/Fatigue/RandomWalk/OscillatorRegions.pdf}
    %     \caption{Illustration of how the distribution of $l$ changes in the first cycle. We display the distribution for several key points in the cycle a)-f).}
    %     \label{fig:RegionsLw}
    % \end{figure}
    % \item Use \cref{fig:RegionsLw} to explain where elements yield based on their value of $\tilde{l}$. The green region is the absorbed region, the blue the active region, and elements in the yellow region yield to the absorbed state. Follow this up with a discussion on how $\lw$ slightly changes this picture as the dynamics are not entirely deterministic.
    % \item Finally, describe the origin of the apparent random walk using \cref{fig:RegionsLw}b.
    % \item For systems which go onto yield, we therefore can consider a first passage time problem. 
    % \item To the authors’ knowledge, no analytical solution exists for $\rho(\tilde{c})$ for a process with continuous spatial states but discrete time steps. In these cases, it is common to instead consider a continuous-time approximation, which can be modelled as a diffusion process of the form
    % \[
    %     \frac{\partial c(x, t)}{\partial t} = D \frac{\partial^2 c(x, t)}{\partial x^2}
    % \]
    % where $D$ is a diffusion coefficient and $c(x,t)$ the concentration.
    % \item We construct D based on the expected values $\mathbb{E}\left[(X_n - X_0)^2\right] = n \sigma^2 = \mathbb{E}\left[(x(t) - x_0)^2\right] = 2 D t$.
    % \item we consider the case where there is a single absorbing BC on the green-blue boundary, two absorbing BC on the green-blue and blue-yellow boundary. Additionally, adding a robin boundary on the green-blue boundary to account for not perfectly absorbing BC on that side. The comparisons are plotted in \cref{fig:CtildeDistributions}.
    % \begin{figure}[h]
    %     \centering
    %     \includegraphics[width=0.9\linewidth]{chapters/Fatigue/RandomWalk/Ctilde_Models.pdf}
    %     \caption{Cumulative distribution functions (CDFs) of $\tilde{C}$ comparing the empirical simulation distribution with model predictions (dashed curves: one-boundary, two-boundary, and two-boundary with Robin condition). \newline
    %     Parameters $N=1024$, $\lw=0.02$. {\bf Left:} $l_0=0.0375$,$\gamma_0=0.885$. {\bf Right:} $l_0=0.075$, $\gamma_0=0.826$.}
    %     \label{fig:CtildeDistributions}
    % \end{figure}
    % \item The one BC case clearly does not fit the data, with the two BC solutions providing a better fit. The use of a Robin boundary condition models the data better than the purely absorbing case, but does not provide a qualitative match to the data. 
% \end{itemize}

% \section{A simple Model}
% \begin{itemize}
    % \item While it might be possible to construct a set of boundary conditions which match the physics observed, we instead decide to use a Monte-Carlo method. We run the same simulation dynamics as in the full Elastoplastic, but instead do not resolve the propagator. Instead, when an element yields, it reduces its local stress $l$ by a factor $\Origin$. This models the dynamics of a diffuse set of elements but allows for the same initial starting configurations to be considered. 
    % \begin{figure}
    %     \centering
    %     \includegraphics[width=0.9\linewidth]{chapters/Fatigue/RandomWalk/Combined_Comparison_Plot.pdf}
    %     \caption{Comparison of simulations using the full elastoplastic mode (solid) and diffuse model without propagation (dashed) for the case of absorption. \newline
    %     {\bf a)} Comparison of absorption time $\CStar$ a function of annealing parameter $l_0=0.025,0.0375,0.05,0.075,0.1,0.125$ from right to left. \newline
    %     {\bf b)} Mean activity (fraction of elements yielding) versus cycle for highlighted parameter sets in {\bf a)}. \newline
    %     Parameters: $N_x=1024$, $\lw=0.02$}
    %     \label{fig:SimpleModelComparison}
    % \end{figure}
    % \item This simple model shows remarkable similarities to the full simulation, particularly at low levels of annealing. Here, elements act independently of each other, with the process being driven by a set of independent random walks. 
    % \item This indicated that with increasing system size, the time to absorb increases, as the systems can support more of these random walks. As the system only absorbs once all the walks have ended 
% \end{itemize}

% \section{Conclusion}
% \begin{itemize}
%     \item We have modelled the delayed failure and absorption of these systems under large amplitude oscillatory shear.
%     \item Captured the transition from absorption to yielding across a range of annealing 
%     \item Quantified the delayed failure/absorption observed in these systems as a function of annealing, post-hop and System size. 
%     \item Captured the trend to diffuse behaviour as the system progresses to a higher number of cycles. 
%     \item Began to quantify this diffuse behaviour both analytically and using a very simple method that captures the overall behaviour of the system accurately. 
% \end{itemize}

% \subsection{Further Work}
% \begin{itemize}
%     \item Extend the simple model to try and explore if yielding can be modelled in this context. Early work shown in \cref{fig:RandomWalkExamples}b shows that a neighbour in the absorbed region, close enough that it can be displaced from the absorbed region, is enough in brittle systems to force the system to yield. 
%     \item Is it possible to model the BC as a diffusion process to get an analytical form of the failure time.
% \end{itemize}

% \input{chapters/Fatigue/RandomWalk}
 \newpage 
 \newpage % !TeX root = ../../thesis.tex

\chapter{Ultradelayed Fracture After Step Strain in Networks} \label{sec:Ultradelayed}

\section{Introduction} \label{sec:delayedInto}
In this chapter, we investigate the behaviour of soft network materials subjected to the step strain protocol outlined in \cref{sec:StepStrain}. In particular, we focus on delayed failure occurring long after the initial deformation. This focus is motivated by previous theoretical studies by Lockwood \etal\cite{lockwood_ultradelayed_2025}, which demonstrated ultra-delayed instabilities after step strain in continuum and mesoscopic constitutive models. While those results compellingly showed that dramatic late-time failure is possible in principle, they did not show how related delayed-fracture physics may emerge in an explicit disordered network model. By introducing thermal bond-breaking into the simple network model described in \cref{sec:Networks}, we show in this chapter how delayed fracture can emerge within a mesoscopic description of soft networks.

The formation of strain localisation and material failure has been widely explored both experimentally and theoretically, as discussed in \cref{sec:ShearBanding}, and these phenomena have been observed under a range of deformation protocols. Typically, protocols such as shear startup, outlined in \cref{sec:ShearStartup}, where at some time $t=0$ the system is subjected to the switch on of a shear rate $\dot{\gamma}$, which is then held constant for all time $t>0$, have been used to study these phenomena \cite{moorcroft_shear_2014,moorcroft_criteria_2013,pollard_yielding_2022,moorcroft_age-dependent_2011,divouxTransientShearBanding2010,divoux_stress_2011,kurokawa_avalanche-like_2015}. In this protocol, shear stress $\Sigma_{xy}$ accumulates approximately linearly with the applied strain $\gamma = \dot{\gamma}t$, until it reaches a maximum value where yielding initiates and the stress decays until the sample obtains a flowing steady state, or the material catastrophically fails. 

Alternatively, in a stress-controlled protocol, the material is subjected to a fixed stress of amplitude $\Sigma$ for $t \geq 0$. Here, typically the material slowly creeps initially with a strain $\gamma \sim t^{1-\alpha}$ that increases sublinearly in time and a rate that decays as $\dot{\gamma} \sim t^{-\alpha}$, with $\alpha \in \left( 0,1 \right)$. On longer timescales the material may ultimately fail, with the measured shear rate suddenly increasing as the material yields \cite{miguel_dislocation_2002,gibaud_heterogeneous_2010,leocmach_creep_2014,ballesta_creep_2016,liuCreepDynamicsAthermal2018,cabriolu_precursors_2019,siebenburger_creep_2012}.

In both cases, damage in the material accumulates over time, either through externally applied deformation that triggers yielding events or through stress propagation from a previous yielding event. During the initial stages of deformation, the shear response is often statistically uniform with any plastic events spread diffusely across the sample. With accumulating damage, this homogeneous state typically breaks down and deformation becomes spatially heterogeneous, with localised clusters of plastic events emerging and spatial gradients developing in the shear response across the sample. In some materials and protocols this localisation sharpens into shear bands, while in others it develops into more strongly fracture-dominated failure. These localisation mechanisms have been captured theoretically within a minimal set of continuum rheological constitutive assumptions \cite{fielding_shear_2014,moorcroft_age-dependent_2011,moorcroft_shear_2014,moorcroft_criteria_2013,adams_transient_2011,fieldingTriggersSignaturesShear2016}.

\begin{figure}[t]
    \centering
    \includegraphics[width=.8\linewidth]{chapters/UltraDelayedFracture/Figures/TrivialRelax.tikz}
    \caption[Illustration of response curves under a step strain deformation]{Schematic Illustration of typical response curves under a step strain deformation. The material relaxes to a finite residual stress as the material elastically or plastically rearranges (blue). Alternatively, when the imposed deformation is sufficiently large, it undergoes a catastrophic stress drop associated with macroscopic failure (red).}
    \label{fig:trivialRelaxation}
\end{figure}

Here, instead, we consider a step strain protocol in which, at time $t=0$, the material is rapidly deformed to a finite fixed shear strain $\gamma_0$, which is then held constant thereafter. Unlike the shear startup and creep protocols, in which the strain can increase indefinitely such that the system attains a fully fluidised steady state, the strain remains fixed in this step strain protocol. Consequently, the subsequent relaxation processes of these viscoelastic materials have generally been assumed to be straightforward. The stress is assumed either to relax gradually towards a finite residual value or to drop abruptly upon catastrophic failure, depending on the magnitude of the imposed strain $\gamma_0$. A schematic representation of these two limiting responses is shown in \cref{fig:trivialRelaxation}, with the slow relaxation indicated in blue and the catastrophic failure in red.

In the athermal limit, the emergence of either rapid failure or mechanical arrest follows directly from purely mechanical stability. Because any sources of thermal fluctuations are absent, there is no intrinsic timescale for delayed failure: deformation proceeds through discrete, avalanche-like rearrangements that redistribute stress \cite{nicolas_deformation_2018,agoritsas_relevance_2015,ness_oscillatory_2017,lin_hydrodynamic_2015}, so the system either fails promptly following the imposed step strain or relaxes to a mechanically stable arrested state. A single yielding event can therefore trigger a cascade of further events, potentially leading to system-spanning failure \cite{barlow_ductile_2020}, while in the absence of such cascades no further rearrangements occur.

% In athermal networks, the slow sequence of microscopic rearrangements that could build toward an eventual avalanche and produce delayed banding is therefore less likely than either rapid catastrophic failure or complete arrest. Relaxation times have been shown to diverge at the strain-stiffening point \cite{shivers_strain-controlled_2023}, suggesting that the material may be able relax extremely slowly toward its eventual failure point even in the absence of thermal activation. We expect this behaviour to be highly network-dependent and, therefore, in this work we will only consider the effect of thermally driven failure in highly connected networks.

In previous work by Lockwood \etal~\cite{lockwood_ultradelayed_2025}, constitutive models in the thermal regime were shown to develop ultra-delayed instabilities long after the imposed deformation \cite{lockwood_ultradelayed_2025}. In these systems, catastrophic failure can occur only after extremely long delay times at low temperatures and for imposed strains below that required to trigger an immediate instability, with the delay time increasing strongly as the strain amplitude is reduced, the temperature is lowered, and the degree of annealing is increased. During this interval, the material appears macroscopically stable, before a sudden shear banding instability leads to rapid localisation and a sharp drop in the stress. In the absence of knowledge of the system’s prior evolution, the onset of failure might therefore appear abrupt to an external observer.

Because the step strain protocol is less commonly studied compared to other protocols, relatively little work exists on such late-time behaviour before the study by Lockwood \etal \cite{lockwood_ultradelayed_2025}. Some alternative theoretical studies focused on short-timescale localisation or immediate failure \cite{moorcroft_shear_2014,moorcroft_criteria_2013,adams_nonmonotonic_2009}. Experimentally, similar short-timescale failure processes have been observed in polymer melts \cite{boukany_step_2009,barroso_evaluation_2002,einaga_stress_1973}. However, studies examining much longer-time dynamics after the system has reached an apparently quiescent state remain limited. A small number of works have reported delayed yielding or delayed fracture, in which materials subjected to subcritical loading remain macroscopically unchanged for extended periods before undergoing sudden failure \cite{boukany_step_2009,jacob_rheological_2019,van_der_kooij_laser_2018}.

\DIFaddbegin \DIFadd{Throughout this chapter, we use packing-derived networks constructed by energy-minimising bidisperse soft-disk packings at $\phi=1$ and replacing the disk contacts with Hookean bonds, as described in }\cref{sec:PackingBasedNetworks}\DIFadd{. Unless otherwise stated, these networks have an average coordination $z\approx5.4$; networks with lower coordination are obtained by pruning bonds to the specified value of $z$.
}

\DIFaddend \begin{figure}
    \centering
    \includegraphics[width=1.0\linewidth]{chapters/UltraDelayedFracture/Figures/Relax.pdf}
    \caption[Stress relaxation of networks after a step strain of amplitude $\gamma_0$]{Stress relaxation of \DIFaddbeginFL \DIFaddFL{packing-derived }\DIFaddendFL networks \DIFaddbeginFL \DIFaddFL{pruned to the specified average coordination $z$ }\DIFaddendFL after a step strain of amplitude $\gamma_0$. \newline
    {\bf Left:} Stress relaxation response under an imposed shear strain of $\gamma_0$ at $z=3.5$. \newline
    {\bf Right:} Phase space of normalised steady state stress $\Sigma_\infty$ under changing $z$ and $\gamma_0$. We observe, in particular, a critical point $z_c=4.0$, where for $z \geq z_c$, all networks can hold a state of self-stress; below this, we observe the introduction of a strain stiffening transition.}
    \label{fig:networkRelaxation}
\end{figure}

The step strain protocol has previously been studied in the context of fibre networks consisting of unbreakable fibres by Shivers \etal in Ref.~\cite{shivers_strain-controlled_2023}. To check our simulation method, we reproduce the relevant features of this behaviour in \cref{fig:networkRelaxation}. Although their work includes a detailed analysis of relaxation timescales and dynamics, these aspects are not central to the focus of this chapter. Instead, the key feature for our purposes is the presence of a strain-stiffening transition in sub-isostatic networks with connectivity $z<4$. At strains less than the strain stiffening point $\gamma_c$, the network is naturally floppy and possesses many zero-energy modes. Beyond a critical strain $\gamma_c(z)$, these modes are eliminated, and the network can sustain stress in a manner comparable to networks above the isostatic point ($z>4$). We examine the origins of this transition in more detail in \cref{sec:Rigidity}.

% \red{
This work aims to follow naturally from the work of Ref.~\cite{lockwood_ultradelayed_2025} and explore how introducing thermally induced bond breakage allows for the eventual delayed failure of these network materials. In attractive colloidal gels subjected to sustained loading, delayed catastrophic yielding under creep is well documented, with failure times that decrease rapidly with applied stress and often show an exponential sensitivity to the applied load and other control parameters \cite{sprakel_stress_2011,blindstrom_structures_2012,Landrum2016,Brenner2013}. Although these studies are most commonly performed under constant stress, the underlying mechanism, thermally activated bond rupture within a heterogeneous network, is intuitively aligned with the ideas we consider in this chapter. The same literature emphasises that microscopic dynamics and structural evolution can change substantially, a long time before the eventual macroscopic rupture, even when the bulk rheological signal is comparatively smooth \cite{blindstrom_structures_2012}. This suggests that step strain experiments in colloidal gels or soft glasses provide a realistic experimental setting in which the delayed-fracture scenario studied here could be observed. In particular, rheological X-ray photon correlation spectroscopy (rheo-XPCS) \cite{Leheny2015} has been implemented to directly link nanoscale dynamics to yielding and relaxation under step deformations, including in step-strain protocols that induce yielding and flow \cite{Chen2020,Bellour2003}. Notably, data reported by Harden \etal \cite{Chen2020} exhibit features consistent with strongly delayed failure, although this behaviour was not explicitly identified or analysed in that study.

A second class of materials that shows promise for delayed fracture under imposed strain comprises soft network materials. Experiments on notched elastomers and soft polymer networks under step strain demonstrate delayed fracture following a substantial induction period, with characteristic times that decrease sharply as the imposed strain is increased \cite{van_der_kooij_laser_2018,wangDelayedFractureGels2012,Ju2023}. More generally, slow crack growth and delayed rupture under a fixed deformation have long been observed in highly elastic materials, where failure is governed by stress-induced bond breaking \cite{Fan2024,Sanoja2021,Creton2016} and thermally activated bond rupture \cite{Wang2023,Persson2024,Baumberger2009,Hui2004,Creton2016}. Crucially, full-field strain imaging via laser speckle strain imaging and related optical methods reveals the gradual development of an internal pre-fracture deformation zone well before catastrophic failure is evident in the global stress signal \cite{van_der_kooij_laser_2018,Wang2023,ju_role_nodate,Ju2023}. In these materials, a non-linear process zone forms at the crack tip and evolves gradually prior to rupture, even though the system appears macroscopically quiescent \cite{van_der_kooij_laser_2018,Wang2023}.

The relative lack of explicit reports of ultra-delayed failure under step strain is therefore unlikely to reflect its absence in real materials. Rather, the phenomenon is easily obscured in practice and therefore difficult to characterise systematically, owing to its strong sensitivity to both internal material parameters and external experimental conditions. First, the required experimental protocol of maintaining a fixed strain for very long durations while observing the response at multiple scales is much less common than creep or oscillatory measurements. Most rheological tests are abandoned once the material enters an apparently stable state. Second, boundary and edge effects can dominate the macroscopic response. While simulations can impose periodic boundary conditions and effectively represent an infinite system, experiments are necessarily performed on finite samples. As a result, measured bulk behaviour is always influenced to some degree by surfaces, interfaces and sample geometry. Wall slip and thin lubrication layers are ubiquitous in yield-stress fluids and soft particulate materials, such as colloidal gels, near smooth boundaries. They can dominate the apparent deformation near and below yielding, decoupling the imposed motion from the bulk response \cite{Divoux2024,Walls2003,Ballesta2013,Jung2021}. Separately, edge fracture is a well-known free surface instability that can invade the sample at the perimeter \cite{ElastomersEdge2024,Hemingway2017}. Third, delayed failure or yielding is strongly dependent on the history of the sample: the microstructure (and therefore the delay time) depends on the preparation, deformation history, and ageing of the sample, which can be further complicated by stochastic nucleation mechanisms that can further complicate the response even under nominally identical conditions \cite{blindstrom_structures_2012,Divoux2024,bouzidNetworkTopologySoft2018,Cipelletti2000}. Without careful control of the system's initial configuration and repeated measurements, such behaviour can seem inconsistent from one test to another or be mistaken for ordinary experimental variability rather than reflecting an intrinsic delayed instability.

Taken together, the existing literature indicates that the predicted ultra-delayed failure should be experimentally accessible, but only in a carefully designed step strain experiment that precisely controls the system's initial configurations and the experimental conditions before the material fails catastrophically.
%DIF <  }

\DIFdelbegin \DIFdel{For computational efficiency}\DIFdelend \DIFaddbegin \DIFadd{To model thermally activated bond rupture}\DIFaddend , we adopt a \DIFdelbegin \DIFdel{thermal breakage condition analogous to the }\DIFdelend \DIFaddbegin \DIFadd{breakage condition motivated by established descriptions of thermally activated fracture in soft networks and elastomers \mbox{%DIFAUXCMD
\cite{Hui2004,Baumberger2009,Creton2016,Wang2023}}\hskip0pt%DIFAUXCMD
, together with the activated }\DIFaddend yielding rules used in soft glassy rheology \cite{sollich_rheological_1998,sollich_thermodynamic_2012} and thermally activated elasto-plastic models \cite{nicolas_deformation_2018,bonn_yield_2017,Rodney2011}. Specifically, we assign to each bond a failure rate governed by an Arrhenius (Boltzmann) factor\DIFdelbegin \DIFdel{that depends on its local stress. Within this framework, we explore how failure can be shifted to }\DIFdelend \DIFaddbegin \DIFadd{, in which the elastic energy stored in a strained bond lowers the effective energetic barrier to failure and thereby increases its rupture rate. Rather than modelling thermal fluctuations explicitly through Brownian or Langevin dynamics (as in \mbox{%DIFAUXCMD
\cite{tauber_role_2020}}\hskip0pt%DIFAUXCMD
), we coarse-grain their effect directly into this stochastic bond-breakage rate. This provides a simple representation of the competition between elastic loading and thermal activation while making the }\DIFaddend extremely long waiting times \DIFaddbegin \DIFadd{associated with delayed fracture computationally accessible}\DIFaddend . In the remainder of this chapter, we introduce this thermally activated bond-failure rule and use it to study delayed fracture at both high and low strain. In particular, we examine how the characteristic failure time depends on the imposed strain $\gamma_0$ and temperature $T$.

% To explore the ultra-delayed failure of networks, we extend the network model outlined in \cref{sec:Networks}, specifically focusing on the Jamming-based networks defined in \cref{sec:PackingBasedNetworks}. We will introduce temperature through two methods, a Brownian simulation where node positions are randomly \quotes{kicked}, and a Boltzmann-like breakage rate for the network bonds inspired by the soft glassy rheology model\cite{sollich_rheology_1997,sollich_soft_2006}. 

% The bulk of the results in this work will use the latter method, which we introduced for its computational simplicity and ability to reach long time scales. In both models, shear is applied through applying Lees-Edwards boundary conditions (outlined in \cref{sec:Lees-Edwards}). This imposes effective periodic boundary conditions in all directions and enforces a simple shear in the $xy$ plane. The jamming-based networks used in this chapter are generated as outlined in \cref{sec:PackingBasedNetworks}, with $\phi = 1.0$ and the box size $L_x=L_y = \sqrt{N}$ determining the length scale of the system. Unless otherwise states this gives networks with connectivities of $z \approx 5.4$. Further, we also take units where $\mu_m=1.0$. 

% \section{Brownian Dynamics} \label{sec:UltraDelayedBrownian}
% The simplest model, which introduces a temperature $T$ is Brownian dynamics. We amend the dynamics defined in \cref{sec:NetworkDynamics} to be of the following form
% \begin{equation}
%     \dot{\bm{r}_i} = \frac{1}{\zeta}\bm{F}_{n,i} + \bm{v}_f(\bm{r}_i) + \sqrt{\frac{2k_BT}{\zeta}}R(t). \label{eq:Brownian}
% \end{equation}
% We work in unit where drag coefficient $\zeta = 1$, Boltzmann constant $k_B=1$ and $R(t)$ is a Gaussian white noise term with $\langle R(t)\rangle = 0$ and $\langle R(t) R(t^\prime) \rangle = \delta(t-t^\prime)$. This is equivalent to the limiting cases of the dynamics explored in Ref.\cite{tauber_role_2020}. We integrate \cref{eq:Brownian} using a simple Euler algorithm with time-step $\delta t = 10^{-3}$. During the simulation, we will take bonds to break irreversibly when their strain $\epsilon_b = (l-l_0)/l_0$ exceeds a threshold $\lambda$, which is evaluated after each dynamics step.

% \begin{figure}
%     \centering
%     %\includegraphics[width=0.5\linewidth]{}
%     \missingfigure{Left time series / Right $t^*$ extracted}
%     \caption{Caption}
%     \label{fig:placeholder}
% \end{figure}

\section{Rate Dependent Fracture} \label{sec:BoltzmanFracture}
To explore the ultra-delayed failure of networks, we extend the network model outlined in \cref{sec:Networks}, specifically focusing on the packing-derived network model defined in \cref{sec:PackingBasedNetworks}. It is common when exploring fractures in mesoscopic network models to introduce a finite breakage strain $\lambda$, where bonds with axial strain $\epsilon_b=(l-l_0)/l_0 \geq \lambda$ break \cite{tauber_stretchy_2022}. While this model could be used to understand athermal fracture in a step strain regime, it is unable to resolve the slow thermal activation processes that underlie delayed fracture, as discussed in Ref.~\cite{lockwood_ultradelayed_2025} and Ref.~\cite{lockwood_yielding_nodate}.

% Taking inspiration from the Soft Glassy Rheology (SGR) model \cite{sollich_rheological_1998,sollich_thermodynamic_2012}, we will consider each bond to have a Threshold energy for breaking $E$. Each bond then has a local strain energy $E_b = \frac{\mu}{2}\epsilon_b^2$. Through thermal fluctuations, the bond is therefore allowed to explore its own energy well, defined by its strain energy $E_b$ and failure energy $E$. These thermal fluctuations can be thought of as arising from both the surrounding material and the bond's internal structure. Small fluctuations from the equilibrium point are therefore more common, while large excursions away from the equilibrium point are increasingly less common due to the higher energy costs.

% As the local strain energy $E_b$ increases towards $E$, the extra energy required from thermal fluctuations decreases and therefore the probability of escaping the well increases. This idea leads to an Arrhenius distribution for the rate of a bond failing
% \begin{equation*}
%     \Gamma(\epsilon_b) \sim \exp \left[ -\dfrac{E - \frac{\mu}{2}\epsilon_b^2}{x} \right],
% \end{equation*}
% where $T$ is an effective temperature, controlling the amplitude of the allowed fluctuations. In systems where the thermal energy is of the same order as the depth of the energy well, the temperature parameter can be interpreted as the true thermodynamic temperature, with $T = k_B T_{\mathrm{therm}}$ \cite{sollich_thermodynamic_2012}. However, as in the SGR model, the energy depths in our network are much larger than the thermal contributions, the effective noise temperature is more closely related to how stress would be distributed in the microscopic structure of the bond.

% We will also adopt a slight variation to this form of $\Gamma(\epsilon_b)$, where an upper limit is placed on the rate of failing \cite{sollich_thermodynamic_2012},
% \begin{equation*}
%     \Gamma(\epsilon_b) \sim \begin{cases}
%         \exp \left[ -\dfrac{E - \frac{\mu}{2}\epsilon_b^2}{T} \right] & \text{if } E_b < E \\
%         1 & \text{if } E_b \geq E.
%     \end{cases}
% \end{equation*}
% This decouples the rate of failing from the bond's strain after the bond has escaped the well, with the probability remaining constant rather than rising indefinitely. It also has the additional benefit of increasing the numerical stability of the system, as an infinitely small time step is no longer required to accurately resolve the rate $\Gamma$. 

% To implement this explicitly, we introduce a microscopic rate $\Gamma_0$ that sets the characteristic attempt frequency for bond failure. The failure rate of an individual bond is then given by
% \begin{equation} \label{eq:SGRRate}
%     \Gamma(\epsilon_b) = \Gamma_0 \begin{cases}
%         \exp \left[ -\dfrac{E - \frac{\mu}{2}\epsilon_b^2}{T} \right] & \text{if } E_b < E \\
%         1 & \text{if } E_b \geq E.
%     \end{cases}
% \end{equation}
% The factor $\Gamma_0$ corresponds to the microscopic hopping rate of the bonds, for which we will work in units where $\Gamma_0=1$. We illustrate an example of $\Gamma(\epsilon_b)$ in \cref{fig:FailureRate} for a bond with $E=1.0$ and $T=0.1$.

\DIFdelbegin \DIFdel{Taking inspiration from the Soft Glassy Rheology (SGR) model }\DIFdelend \DIFaddbegin \DIFadd{Following treatments of thermally activated bond rupture in soft materials and elastomers \mbox{%DIFAUXCMD
\cite{Hui2004,Baumberger2009,Creton2016,Wang2023}}\hskip0pt%DIFAUXCMD
, and adopting the strain-lowered Arrhenius form used in soft glassy rheology }\DIFaddend \cite{sollich_rheological_1998,sollich_thermodynamic_2012}, we \DIFdelbegin \DIFdel{will consider each bond to have }\DIFdelend \DIFaddbegin \DIFadd{assign each bond }\DIFaddend a threshold energy for breaking $E$ and \DIFaddbegin \DIFadd{a }\DIFaddend local strain energy $E_b = \frac{\mu}{2}\epsilon_b^2$, where $\mu$ is the bond stiffness. \DIFdelbegin \DIFdel{Through thermal fluctuations, each bond then has a rate of breaking }\DIFdelend \DIFaddbegin \DIFadd{The rate of bond breaking is then }\DIFaddend given by an Arrhenius rate,
\begin{equation*}
    \Gamma(\epsilon_b) \sim \exp \left[ -\dfrac{E - \frac{\mu}{2}\epsilon_b^2}{T} \right],
\end{equation*}
 where $T$ is an effective temperature, controlling the magnitude of the thermally activated fluctuations.

 We will also adopt a slight variation to this form of $\Gamma(\epsilon_b)$, where an upper limit is placed on the rate of failing \cite{sollich_thermodynamic_2012},
\begin{equation} \label{eq:SGRRate}
    \Gamma(\epsilon_b) = \Gamma_0 \DIFdelbegin %DIFDELCMD < \begin{cases}
%DIFDELCMD <         \exp \left[ -\dfrac{E - \frac{\mu}{2}\epsilon_b^2}{T} \right] & \text{if } E_b < E \\
%DIFDELCMD <         1 & \text{if } E_b \geq E.
%DIFDELCMD <     \end{cases}%%%
\DIFdelend \DIFaddbegin \begin{cases}
        \exp \left[ -\dfrac{E - \frac{\mu}{2}\epsilon_b^2}{T} \right], & \text{if } E_b < E, \\
        1, & \text{if } E_b \geq E.
    \end{cases}\DIFaddend 
\end{equation}
We introduce a microscopic attempt rate for bond breakage $\Gamma_0$ that sets the characteristic attempt frequency for bond failure. By limiting the breakage rate of a bond, we decouple the rate of breaking from the bond's strain after the bond has exceeded its energy threshold, with the probability remaining constant rather than rising indefinitely with increasing $\epsilon_b$. \DIFdelbegin \DIFdel{It also has the additional benefit of increasing the numerical stability of the simulation, as an arbitrarily small time step is no longer required to accurately resolve the rate }\DIFdelend \DIFaddbegin \DIFadd{For the fixed-time-step algorithm used here, bounding $\Gamma$ by $\Gamma_0$ also relaxes the time-step constraint needed to resolve bond failure accurately, since it is sufficient to choose $\Delta t$ such that $\Gamma_0\Delta t\ll1$ rather than accommodate an unbounded }\DIFaddend $\Gamma$. \DIFaddbegin \DIFadd{This is a practical feature of the present implementation; an adaptive or event-driven algorithm could avoid the same restriction. }\DIFaddend We illustrate an example of $\Gamma(\epsilon_b)$ in \cref{fig:FailureRate} for a bond with $E=1.0$ and $T=0.1$.

 
\begin{figure}[b!]
    \centering
    \includegraphics[width=.9\linewidth]{chapters/UltraDelayedFracture/Figures/Rate.tikz}
    %\missingfigure{An illustration of the rate failure rate $\Gamma$ of a bond as a function of bond strain $\epsilon_b$}
    \caption[Illustration of the failure rate of bonds $\Gamma(\epsilon_b)$]{
    Bond failure rate $\Gamma(\epsilon_b)$ as a function of bond strain $\epsilon_b$, shown for $E=1.0$ and $T=0.1$. The failure rate increases rapidly as bonds approach their failure threshold.}
    \label{fig:FailureRate}
\end{figure}

Bond failure is assumed to occur as a Poisson process with rate $\Gamma(\epsilon_b)$. At the end of each discrete time step $\delta t$, each bond undergoes a Bernoulli trial with failure probability
\begin{equation*}
    p(E_b) = 1 - \exp\left[ -\Gamma(\epsilon_b)\delta t \right].
\end{equation*}
This construction preserves the correct continuous-time Poisson limit while remaining numerically stable for finite time steps. \DIFaddbegin \DIFadd{Alternatively a Gillespie-type event-driven implementation could improve computational efficiency in this long-waiting-time regime by advancing directly between bond-breaking events, with the failure rates recalculated after each event and the ensuing mechanical relaxation.
}\DIFaddend 

We choose units in which the microscopic attempt rate $\Gamma_0=1$, the drag coefficient $\zeta=1$ (see \cref{sec:NetworkDynamics}), the bond stiffness $\mu=1$ and the system size $L_x=L_y=\sqrt{N}$ such that the characteristic bond length $\sim O(1)$ (detailed in \cref{sec:PackingBasedNetworks}). Additionally, the form of the fracture rate given in \cref{eq:SGRRate} gives two free variables $E$ and $T$. Analogous to the failure threshold $\lambda$ used in Refs.~\cite{tauber_role_2020,tauber_sharing_2021,tauber_microscopic_2020}, the parameter $E=\frac{\mu}{2}\epsilon_b^2$, sets the units of strain $ \sim \sqrt{2E / \mu}$ at which a bond is most likely to break. Since the rigidity transition point $\gamma_c(z)$ is fixed by structure alone, rather than any parameters specific to the dynamics ($\zeta$, $\mu$, $\sqrt{N}$) or failure rate ($\Gamma_0$, $E$, $T$), the relevant control parameter governing the point of failure is the relative magnitude of $E$ with respect to $\gamma_c$. The qualitative failure behaviour therefore depends on whether the strain scale $\sqrt{2E / \mu}$ is reached in the floppy or rigid regime of the network and by the nature of the load-bearing structures that have emerged. Furthermore, the temperature $T$ then in combination with the failure threshold $E$ sets the characteristic failure time of the network's bonds.

In this work, we examine two representative strain amplitudes: a small strain ($\gamma_0=0.01$) and a large strain ($\gamma_0=1.0$), using networks containing $N=64{,}000$ nodes. To obtain initial estimates for the free parameters in the breakage rate, we first approximate the network immediately after the imposed step strain by an affine deformation. In this approximation, a bond that initially makes an angle $\theta$ to the horizontal has strain $\epsilon_b(\theta|\gamma_0)=\sqrt{(\cos\theta+\gamma_0\sin\theta)^2+\sin^2\theta}-1$. Assuming the bond orientations are approximately uniformly distributed, we estimate the largest bond energy from the maximum value of $\mu\epsilon_b^2/2$ and then evaluate it at the midpoint of the strain range of interest, multiplying by a factor of two to allow for the non-affine relaxations neglected by this estimate. We then estimate $T$ from the mean affine bond energy by substituting it into the rate \cref{eq:SGRRate} and choosing the corresponding typical bond lifetime to be very long, before refining the parameters through direct exploration of the response. 
\DIFdelbegin \DIFdel{This gives, for }\DIFdelend \DIFaddbegin 

\DIFadd{This parameter refinement is necessary because non-trivial delayed failure occurs only within a narrow region of parameter space. If $E$ is too small relative to the strain energy, or $T$ is too large, widespread bond breaking produces rapid failure immediately after the step strain. Conversely, if $E$ is too large or $T$ is too small, no macroscopic failure is observed within the computationally accessible time window. We therefore tune $E$ and $T$ to access the intermediate regime in which the network first reaches an apparently stable state before failing after a long but observable delay. The resulting non dimensional parameter values should be regarded as representative points within this regime rather than as unique material-specific values. For the system considered, }\DIFaddend the small-strain case\DIFaddbegin \DIFadd{, }\DIFaddend $\gamma_0=0.01$, \DIFdelbegin \DIFdel{the values }\DIFdelend \DIFaddbegin \DIFadd{gives }\DIFaddend $E=5.025\times10^{-5}$ and $T=10^{-6}$, \DIFdelbegin \DIFdel{and for }\DIFdelend \DIFaddbegin \DIFadd{while }\DIFaddend the large-strain case\DIFaddbegin \DIFadd{, }\DIFaddend $\gamma_0=1.0$, \DIFdelbegin \DIFdel{the values }\DIFdelend \DIFaddbegin \DIFadd{gives }\DIFaddend $E=0.35$ and $T=0.0062$.

\section{Stress Response} \label{sec:lowStrainResponse}

The step-strain protocol \DIFdelbegin \DIFdel{, in which a }\DIFdelend \DIFaddbegin \DIFadd{was simulated using packing-derived networks containing $N=64{,}000$ nodes with average coordination $z\approx5.4$, prepared at $\phi=1$ as described in }\cref{sec:PackingBasedNetworks}\DIFadd{. A }\DIFaddend shear strain of amplitude $\gamma_0$ \DIFdelbegin \DIFdel{is }\DIFdelend \DIFaddbegin \DIFadd{was }\DIFaddend instantaneously applied and then held constant for time \DIFdelbegin \DIFdel{$t\geq 0$, was simulated within our network model }\DIFdelend \DIFaddbegin \DIFadd{$t\geq0$, }\DIFaddend using the dynamics described in \cref{sec:NetworkDynamics} \DIFdelbegin \DIFdel{, }\DIFdelend together with the bond failure rate defined in \cref{sec:BoltzmanFracture}. Results for a fixed strain amplitude $\gamma_0 = 0.01$ and a set of temperatures spaced logarithmically in the range $T \in [5\times10^{-7}, 5\times10^{-6}]$ are shown in \cref{fig:LowGammaTempResponse}.

We report the decay of the shear stress $\Sigma_{xy}(t)$ as a function of time $t$ since the imposition of the strain, normalised by its initial value $\Sigma_0 = \Sigma_{xy}(t=0^+)$. At $t=0^+$, the stress is given by the linear elastic response, $\Sigma_0 = G_{xy}\gamma_0$, where $G_{xy}$ is the shear modulus of the network. This modulus depends on the exact network topology and the bond stiffness $\mu=1$. For brevity, we write $\Sigma \equiv \Sigma_{xy}$ below.

\begin{figure}[b!]
    \centering
    \includegraphics[width=0.9\linewidth]{chapters/UltraDelayedFracture/LowGammaTemp/Stress.pdf}
    \caption[Decay of normalised stress $\Sigma/\Sigma_0$ following a step strain $\gamma_0=0.01$ at different temperatures.]
    {Decay of normalised stress $\Sigma/\Sigma_0$ as a function of time $t$ following the imposition of a step strain of amplitude $\gamma_0=0.01$, with breakage parameter $E=5.025\times10^{-5}$. The colour scale indicates the temperature $T$ associated with each curve. The red dashed line indicates the bond failure threshold. A network is considered to have failed once $\Sigma/\Sigma_0 \leq 0.75$, with the failure time $t^*$ defined to be the time at which the curves cross this threshold. The symbols show the locations of snapshots in \cref{fig:LowGammaTempSnapShots}. \textbf{Inset:} The same data plotted on a linear horizontal scale.}
    \label{fig:LowGammaTempResponse}
\end{figure}

At short times, the stress relaxation is nearly identical across all temperatures. This reflects the fact that thermally activated bond breakages are negligible in this early-time regime, with relaxation instead governed by non-affine network deformations. At later times, the relaxation curves for different temperatures become distinct from each other as networks undergo failure driven by thermally activated bond breakages.

For low temperatures, the relaxation closely resembles that shown in \cref{fig:networkRelaxation}, with the stress reaching a plateau that persists for long times. Although these networks ultimately fail in a brittle manner at times $t \gg 0$, we do not observe an obvious precursor in the global stress response before failure. While the failure appears instantaneous in the main figure due to the scale of the plot, the inset shows the same data on a linear horizontal time axis, revealing qualitatively similar failure behaviour across the failure regime, albeit still characterised by an abrupt onset.

\begin{figure}[b!]
    \centering
    \includegraphics[width=0.9\linewidth]{chapters/UltraDelayedFracture/LowGammaTemp/BondCount.pdf}
    \caption[Fraction of broken bonds $B$ following a step strain $\gamma_0=0.01$ at different temperatures.]{Fraction of broken bonds $B$ as a function of time $t$ following the imposition of a step strain of amplitude $\gamma_0=0.01$, with bond breakage energy $E=5.025\times10^{-5}$. Here, $B:= N_b/N$ denotes the fraction of broken bonds relative to the total initial number of bonds \DIFdelbeginFL \DIFdelFL{$N_{\mathrm{tot}}$}\DIFdelendFL \DIFaddbeginFL \DIFaddFL{$N_{\mathrm{tot}}=64{,}000$}\DIFaddendFL . The colour scale indicates the temperature $T$ associated with each curve. The symbols show the location of snapshots in \cref{fig:LowGammaTempSnapShots}.}
    \label{fig:LowGammaTempBondResponse}
\end{figure}

In previous work on delayed failure in elastoplastic models \cite{lockwood_ultradelayed_2025}, displacement profiles were used to diagnose the developing instability. Here, instead, we consider the fraction of broken bonds in the network, $B:= N_b/N$, as a complementary measure. This directly measures the number of bond-breakage events that have occurred in the system, although it does not retain spatial information.

We plot the corresponding bond-breakage response curves in \cref{fig:LowGammaTempBondResponse}, showing the fraction of bonds that have broken as a function of $t$ since the imposition of the strain. These curves display behaviour consistent with the stress decays shown in \cref{fig:LowGammaTempResponse}: the networks that exhibit rapid stress relaxation early in the response are those that show larger values of $B$. Conversely, curves that appear to reach a quasi-steady state exhibit little bond failure prior to an abrupt increase at the point of macroscopic failure.

\begin{figure}[t!]
    \centering
    \includegraphics[width=0.99\linewidth]{chapters/UltraDelayedFracture/LowGammaTemp/LowGammaTempSnapshots.png}
    \caption[Snapshot of bond failures at low imposed strain $\gamma_0$ at a high and low temperature $T$]{Snapshot of bond failures after an imposed step strain of $\gamma_0=0.01$, shown in its undeformed configuration. The colour scale indicates the order in which bonds fail from dark to light. Locations of {\bf a)}, {\bf b)}, {\bf c)} and {\bf d)} are indicated by a square, triangle, star and diamond, respectively, on the stress-time curves in \cref{fig:LowGammaTempResponse} and the bond failure curves in \cref{fig:LowGammaTempBondResponse}. \newline
    Temperature {\bf a)} $T=5 \times 10^{-6}$, {\bf b)} $T=5.35 \times 10^{-6}$, {\bf c)} $T=5.65 \times 10^{-7}$ and, {\bf d)} $T=5 \times 10^{-7}$. \newline
    Parameters: $E=5.025\times10^{-5}$, $\gamma_0=0.01$ and $N=64{,}000$.
    }
    \label{fig:LowGammaTempSnapShots}
\end{figure}

To gain further insight into the failure mechanisms of each network, we present representative snapshots of the microstructure at both low and high temperature in \cref{fig:LowGammaTempSnapShots}. Here, bond failures are shown in the undeformed configuration of the network, corresponding to the initial configuration before strain imposition. We begin by considering the high-temperature cases, shown in \cref{fig:LowGammaTempSnapShots}a and b. In this regime, failure is characterised by extensive diffuse damage: bond breakage occurs throughout the system in a largely spatially uncorrelated manner. As a result, the network rapidly loses rigidity, which is reflected in the corresponding drop in stress. This has been seen in a related model studied by Tauber \etal in Ref.~\cite{tauber_role_2020}. In this work, the authors use an equivalent central force model; however, they consider a topology derived from a two-dimensional triangular lattice (see \cref{sec:LatticeBasedNetworks}) at varying levels of connectivity. The temperature-dependent dynamics are implemented explicitly via Langevin equations of motion. As in our approach, bond failure is governed by a stochastic, thermally activated process, so that rupture events arise from noise-driven fluctuations rather than from a purely deterministic threshold. 

By contrast, at lower temperatures, our networks predominantly fail through a highly localised rupture event, shown in \cref{fig:LowGammaTempSnapShots}c and d. Here, the initial failure of a single bond redistributes stress to its neighbours, sufficiently loading them to bring them close to their own failure thresholds. This initiates a cascade of correlated bond breakage events that propagates through the material, ultimately culminating in the macroscopic failure of the network.

In this small-strain regime, macroscopic failure in the snapshots considered here propagates either vertically or horizontally when viewed in the undeformed configuration of the network, with no clear directional preference. We attribute this to the fact that, at sufficiently small strains, the network remains largely homogeneous and isotropic, with its structure only weakly perturbed by the imposed strain prior to failure. Consequently, the stress redistribution induced by a single bond failure is qualitatively consistent with the Eshelby stress propagator, as discussed in \cref{sec:EshelbyPropogator}. This predicts a quadrupolar stress redistribution, in which bonds aligned with the principal directions of the perturbation (which coincide with the vertical and horizontal axes in this frame) experience an increase in stress, whereas bonds oriented approximately diagonally are comparatively unloaded.

\section[Dependence of the Failure Time \texorpdfstring{$t^*$}{t*} on  Temperature \texorpdfstring{$T$}{T}]{Dependence of the Failure Time \texorpdfstring{$t^*$}{t*} on \\ Temperature \texorpdfstring{$T$}{T}} \label{sec:LowGammaTemp}

To quantify the delay in fracture following the imposition of the step strain, we define a characteristic failure time $t^*$ for the material. We take $t^*$ to be the time at which the stress $\Sigma(t)$ first falls below $\Sigma^*=0.75\Sigma_0$, as indicated by the red line in \cref{fig:LowGammaTempResponse}. This threshold lies below the plateau stress attained via non-affine relaxation at early times, and is therefore only crossed once significant bond failures have occurred. As such, it provides a reliable measure of the onset of macroscopic failure. In practice, the trends reported below are robust to moderate variations in the chosen threshold, provided it lies below the initial plateau reached after early-time relaxation. The resulting failure times are reported in \cref{fig:LowGammaTempFailure}a as a function of $1/T$, consistent with the form of the failure rate $\Gamma$, for which \DIFdelbegin \DIFdel{$t^* \sim \exp(1/T)$}\DIFdelend \DIFaddbegin \DIFadd{$t^* \sim \exp(A/T)$ where $A$ is an apparent activation parameter}\DIFaddend . Consequently, small values of $T$ (large $1/T$) lie to the right of the plot, while increasing temperature corresponds to moving left along the horizontal axis.

\begin{figure}[b!]
    \centering
    \includegraphics[width=0.99\linewidth]{chapters/UltraDelayedFracture/LowGammaTemp/FailureJoin.pdf}
    \caption[Failure time $t^*$ and the corresponding fraction of broken bonds $B^*$ as a function of temperature $T$ at low imposed strain $\gamma_0=0.01$.]{Failure time $t^*$ {\bf a)} and the corresponding fraction of broken bonds $B^*$ {\bf b)} as a function of temperature $T$ at low imposed strain $\gamma_0=0.01$. \newline
    Parameters: $E=5.025\times10^{-5}$ and $N=64{,}000$.}
    \label{fig:LowGammaTempFailure}
\end{figure}

This distinction can be interpreted more quantitatively in terms of how $B^*$ scales with system size. If failure occurs through the formation of a single approximately linear crack spanning the system, the number of broken bonds at failure should scale as the crack length, $N_b^* \sim L \sim \sqrt{N}$ in two dimensions. Since $B^*$ is normalised by the total number of bonds, which is extensive in system size, this implies $B^* \sim 1/\sqrt{N} \approx 3.95 \times 10^{-3}$ for brittle single-crack failure. By contrast, if damage remains diffuse throughout the sample, the number of broken bonds at failure is expected to scale extensively, $N_b^* \sim N$, so that $B^*$ remains of order unity in the thermodynamic limit. Large values of $B^*$ are therefore consistent with diffuse failure, whereas small values that decrease with increasing system size would indicate failure dominated by a single crack.

In \cref{fig:LowGammaTempFailure}a we observe a high-temperature regime in which the failure time increases rapidly with decreasing temperature. In this regime, the stress decay following strain imposition is driven by diffuse bond breakage across the system, rather than by non-affine relaxation of the network. This is reflected in \cref{fig:LowGammaTempFailure}b, where the fraction of broken bonds at failure $B^*$ is large compared with that at lower temperatures. With each bond that breaks, the system loses stress, which may be partially recovered locally as the network deforms to restore force balance. However, at these high temperatures the bond-failure rate is sufficiently large that the network cannot compensate for the structural damage, leading to catastrophic failure at small $t^*$. 

We next consider a temperature low enough that the corresponding rates of bond breakage $\Gamma(\epsilon_b)$, determined by \cref{eq:SGRRate}, are small compared with the characteristic time for viscoelastic relaxation of the network. In this regime, the temperature and strain are sufficiently low that no damage occurs during the initial relaxation, thereby preventing diffuse damage from developing early in the response. Because all systems considered here are subjected to the same imposed shear strain $\gamma_0=0.01$, each network relaxes to a statistically similar configuration determined primarily by the network structure rather than by any macroscopic control parameter. Consequently, each bond has a well-defined and approximately constant failure rate $\Gamma$ determined by its strain in the relaxed state. In this mechanically equilibrated regime, the response depends primarily on the next bond to break, which is statistically most likely to be one of the most highly strained bonds. The eventual failure time for a given network structure therefore depends on the temperature $T$ and the stochastic nature of the underlying Poisson breakage process.

The mechanism that determines the outcome of an individual bond breakage is not yet fully characterised. In some instances, the breaking of a single bond results only in a small, localised deformation that does not precipitate further damage. In other cases, however, a single bond breaking can act as a nucleus for brittle fracture, initiating a propagating crack that ultimately leads to an almost complete loss of bulk stress. Signatures of this behaviour are evident in the step-like increases in $B(t)$ observed in \cref{fig:LowGammaTempBondResponse} for all low-temperature cases, which reflect discrete bond-breaking events separated by long waiting times. Before catastrophic fracture, isolated bond breakages or small cluster failures may occur and persist for extended periods before the next yielding event is triggered.

In this low temperature mechanically equilibrated regime, the lower bound on the failure time of these systems is ultimately set by the rate defined in \cref{eq:SGRRate} for the given parameters $T$ and $E$ that prescribes the average time for the first bond, and subsequent uncorrelated bonds, to break. Network topology can further delay catastrophic failure, as the material must wait for a bond with a sufficiently susceptible local environment to fail and nucleate brittle fracture. As discussed in \cref{sec:Rigidity}, predicting local deformations without explicitly simulating the dynamics is computationally demanding \cite{damavandi_energetic_2022,damavandi_energetic_2022-1}. Consequently, the bonds most likely to nucleate brittle fracture cannot be determined a priori from the initial configuration. The behaviour is therefore well characterised in an ensemble sense, but less predictable for a given realisation.

\section[Dependence of the Failure Time on the Imposed Strain \texorpdfstring{$\gamma_0$}{gamma 0}]{Dependence of the Failure Time on the \\ Imposed Strain \texorpdfstring{$\gamma_0$}{gamma 0}} \label{sec:DelayedGamma}
\subsection{Low Strain Regime}

When varying the temperature $T$ in \cref{sec:LowGammaTemp}, we effectively varied how steeply the bond-breaking rate increased as a bond approached its energy threshold for breakage. Here, instead, by varying the imposed strain, we effectively control how close each bond’s strain energy is to the failure threshold $E$. In both cases, we tune the distribution of breakage rates $\Gamma$ within the network. Under different imposed step strains $\gamma_0$, the network relaxes into slightly different configurations, such that distinct subsets of bonds lie close to the failure threshold $E$.

\begin{figure}[b!]
    \centering
    \includegraphics[width=0.99\linewidth]{chapters/UltraDelayedFracture/LowGammaStrain/LowGammaStrainJoint.pdf}
    \caption[Network response for different values of $\gamma_0$ in the small-strain regime]{
    Network response for different values of imposed strain $\gamma_0$ in the small-strain regime ($\gamma_0 \ll 1$). The symbols in {\bf a)} and {\bf b)} show the location of snapshots in \cref{fig:LowGammaStrainSnapshots}. \newline
    {\bf a)} Normalised stress decay $\Sigma/\Sigma_0$ as a function of time $t$ since the shear strain was imposed. \newline
    {\bf b)} Failure time $t^*$. \newline
    {\bf c)} Fraction of broken bonds $B$ as a function of time $t$ since the shear strain was imposed. \newline
    {\bf d)} Fraction of broken bonds at failure $B^*$. \newline
    Parameters: $E=5.025\times10^{-5}$, $T=10^{-6}$, $N=64{,}000$.
    }
    \label{fig:LowGammaStrainJoint}
\end{figure}

\begin{figure}[t!]
    \includegraphics[width=0.99\linewidth]{chapters/UltraDelayedFracture/LowGammaStrain/LowGammaStrainSnapshots.png}
    \caption[Snapshot of bond failures at low strains.]{Snapshot of bond failures after an
    imposed step strain of {\bf a)} $\gamma_0=0.009$ and {\bf b)} $\gamma_0=0.011$. The colour scale indicates the order in which bonds fail from dark to light. Locations of {\bf a)} and {\bf b)} are indicated by a square and triangle, respectively, on the response curves in \cref{fig:LowGammaStrainJoint}a and \cref{fig:LowGammaStrainJoint}c. \newline
    Parameters: $E=5.025\times10^{-5}$, $T=10^{-6}$, $N=64{,}000$.
    }
    \label{fig:LowGammaStrainSnapshots}
\end{figure}

Despite these microscopic differences, the macroscopic response of the system and the associated failure times $t^*$, when examined as a function of $\gamma_0$, increase with decreasing $\gamma_0$, resembling the manner in which $t^*$ increases with decreasing temperature $T$ (see \cref{fig:LowGammaTempFailure}a). We report the stress-time responses $\Sigma(t)$ for various imposed step strains $\gamma_0$ in \cref{fig:LowGammaStrainJoint}a and the corresponding failure times $t^*$ in \cref{fig:LowGammaStrainJoint}b. The bond-breakage response $B(t)$ and the corresponding fraction of broken bonds at failure $B^*$ are shown in \cref{fig:LowGammaStrainJoint}c and \cref{fig:LowGammaStrainJoint}d, respectively. In this regime, following the arguments made when considering the effect of temperature $T$ discussed above, we fix the bond breakage energy $E=5.025\times10^{-5}$ (see \cref{sec:approxEx} for details), temperature $T=10^{-6}$ and vary the step strain amplitude linearly in the range $\gamma_0\in[0.009, 0.011]$.

As in the temperature-controlled case for this low-strain regime, we observe a transition from more brittle to more diffuse failure with increasing strain. Here, brittle failure refers to crack-like failure for which $N_b^*$ scales as $\sim O(\sqrt{N})$, whereas diffuse failure refers to spatially distributed damage $\sim O(N)$. As indicated by both the spatial distribution of bond failures found at the highest strain considered in this regime ($\gamma_0=0.011$), shown in \cref{fig:LowGammaStrainSnapshots}a, and the fraction of broken bonds at failure $B^*$ reported in \cref{fig:LowGammaStrainJoint}d, the parameter range considered does not exhibit the widespread diffuse damage observed at high temperatures (see \cref{fig:LowGammaTempSnapShots}a). Instead, we see a spatial distribution of bond breakage similar to that shown in \cref{fig:LowGammaTempSnapShots}b, where the system forms several smaller macroscopic cracks rather than a single system-spanning rupture. The corresponding number of bonds that fail in this regime, $N_b^* \approx 340$, remains comparable to $N_b^*$ in the brittle regime, $\approx 250 \approx \sqrt{N}$, and much smaller than in the highly diffuse high-temperature regime, where $N_b^* \approx 1500$.

In both cases, the initial step strain brings multiple bonds close to their failure threshold, leading to spatially distributed breakage events. These bond failures occur on a similar time-scale, allowing for multiple failures to propagate simultaneously in an avalanche-like manner until multiple macroscopic cracks have formed, terminating when the failure of a subsequent bond does not release sufficient stress to trigger continued bond breaking.

Notably, there is no obvious plateau in the fraction of broken bonds at failure $B^*$ reported in \cref{fig:LowGammaStrainJoint}d at low strains, in contrast to the clear plateau in \cref{fig:LowGammaTempFailure}b, despite reaching longer failure times $t^*\approx 5\times 10^{5}$ at the smallest considered $\gamma_0=0.009$ and observing similar failure modes (see \cref{fig:LowGammaTempSnapShots}d and \cref{fig:LowGammaStrainSnapshots}a). Despite this, within the limits of either low temperature or low strain, the underlying physics is the same: the system remains intact in a period of apparent stability until the late failure of a single critical bond. 

Although not currently characterised in this work, the fraction of broken bonds at failure $B^*$ is expected to grow strongly, as in the high-temperature regime, since a critical strain must exist above which all bonds exceed the threshold energy $E$ and therefore fail rapidly. Here, the network fails diffusely across the system with no correlation between breakage events. Conversely, there must exist a lower strain bound below which the imposed deformation induces a finite bond breakage rate, yet is sufficiently small that failure occurs only after the material has relaxed to a quasi-steady state. The response in this regime would therefore again be controlled by the next bond failure event, and the total number of broken bonds would plateau to the value required for brittle failure to reduce the stress below $\Sigma^*$. Although this value depends on the network topology, it is expected to be broadly consistent with failures shown in the low temperature regime discussed in \cref{sec:LowGammaTemp}.

\subsection{High Strain Regime}

So far, we have focused on the regime in which the imposed step strain $\gamma_0$ is small, such that the initially imposed deformation remains small throughout most of the network prior to failure. We now consider the response of the system under large imposed strains, $\gamma_0 \sim 1$, sufficient to induce substantial non-affine deformations within the network \cite{shivers_scaling_2019}. We therefore expect qualitatively different relaxation and failure behaviour in this regime. In what follows, we present results obtained using $E = 0.35$, temperature $T = 0.0062$, and step strains in the range $\gamma_0 \in [0.9,1.0]$. The value of the breakage threshold $E$ forces bonds to fail once they exceed a critical strain of $\pm\sqrt{0.9}\approx\pm0.9487$ and follows from the prediction of a purely affinely deformed network immediately after a strain $\gamma_0=1.0$ is imposed, as detailed in \cref{sec:approxEx}.

\begin{figure}[b!]
    \centering
    \includegraphics[width=0.95\linewidth]{chapters/UltraDelayedFracture/HighGammaStrain/HighGammaStrainJoint.pdf}
    \caption[Response under varying $\gamma_0$ in the high-strain regime]{
    Network response to different values of imposed strain $\gamma_0$ in the high-strain regime ($\gamma_0 \sim 1$). The symbols show in {\bf a)} and {\bf c)} the location of snapshots in \cref{fig:HighGammaStrainSnapShots}. \newline
    {\bf a)} Normalised stress decay $\Sigma/\Sigma_0$ as a function of time $t$ since the shear strain was imposed.\newline
    {\bf b)} Failure time $t^*$. \newline
    {\bf c)} Fraction of broken bonds $B$ as a function of time $t$ since the shear strain was imposed. \newline
    {\bf d)} Fraction of broken bonds at failure $B^*$. \newline
    Parameters: $E=0.35$, $T=0.0062$ and $N=64{,}000$.
    }
    \label{fig:HighGammaStrain}
\end{figure}

\begin{figure}[b!]
    \centering
    \includegraphics[width=0.99\linewidth]{chapters/UltraDelayedFracture/HighGammaStrain/HighGammaStrainSnapshots.png}
    \caption[Snapshot of bond failures at high imposed strains $\gamma_0$]{Snapshot of bond failures after an imposed step strain of {\bf a)} $\gamma_0=0.9$, {\bf b)} $\gamma_0=0.9125$, {\bf c)} $\gamma_0=0.975$ and, {\bf d)} $\gamma_0=1.0$. The colour scale indicates the order in which bonds fail from dark to light. Locations of {\bf a)}, {\bf b)}, {\bf c)} and {\bf d)} are indicated by a square, triangle, star and diamond, respectively, on the response curves in \cref{fig:HighGammaStrain}a and \cref{fig:HighGammaStrain}c. \newline
    Parameters: $E=0.35$, $T=0.0062$, $N=64{,}000$.
    }
    \label{fig:HighGammaStrainSnapShots}
\end{figure}

The stress-time responses in \cref{fig:HighGammaStrain}a-b resemble those found in the low-strain regime (see \cref{fig:LowGammaStrainJoint}a-b), with the failure time $t^*$ increasing systematically as $\gamma_0$ decreases over the range considered, as observed previously. A characteristic feature of failure at the high strains considered here is the rapid increase in bond breakage at early times during relaxation, as shown in \cref{fig:HighGammaStrain}c. The curves closely resemble those observed in the high-temperature cases in \cref{fig:LowGammaTempBondResponse}. This behaviour indicates that many bonds are driven close to their failure threshold and therefore break diffusely and independently. Such diffuse damage is evident in all cases shown in \cref{fig:HighGammaStrainSnapShots}. In each case, early bond failures are distributed throughout the network, largely independent of the imposed shear strain, and do not immediately coalesce into a macroscopic crack, regardless of the eventual mode of failure. Instead, they appear as a spatially dispersed background of isolated events at early stages of the response, before global failure sets in.

Despite this level of diffuse damage, which is not seen in the small-strain regime, networks subjected to the lower step strain $\gamma_0=0.9$ can survive for several more decades in time than networks subjected to the highest considered strain, $\gamma_0=1$.

As indicated in \cref{fig:HighGammaStrain}c, a substantial number of bonds have already broken in all cases before the earliest observed failure time, $t^*\approx10^{2}$. Moreover, the plateau value of $B^*$ reached upon decreasing $\gamma_0$ remains comparable to, or even exceeds, that of the most diffuse failure case (largest $\gamma_0$) shown in \cref{fig:LowGammaStrainJoint}d. Because the imposed strain is still small enough for the network to relax into a quasi-steady state, rather than failing immediately, the macroscopic stress response in \cref{fig:HighGammaStrain}a remains approximately constant in time prior to failure, despite the progressive accumulation of bond breakages evident in \cref{fig:HighGammaStrain}c. This behaviour may reflect stabilising effects associated with strain stiffening in the network \cite{sharma_strain-controlled_2016}, which enable the system to maintain its macroscopic load-bearing capacity even as local damage accumulates. Consequently, the network can accommodate a higher density of bond breakage events than in the low-strain regime while exhibiting only minimal changes in its sustained stress.

In these high-strain cases, the eventual failure, when it occurs long after the imposed strain, is characterised by a single macroscopic fracture that propagates perpendicular to the principal strain axis. This preferred orientation is discussed further in \cref{sec:DirectionOfFracture}. Failure is triggered when a critical bond breaks in a configuration that permits large-scale crack propagation. Whereas at sufficiently small strain or temperature this event can occur among the earliest bond failures, here it emerges only after substantial diffuse damage has accumulated throughout the network. This may indicate that only a small subset of bonds within the system are capable of initiating system-spanning failure, although preliminary investigations (not shown in this thesis) have not conclusively confirmed this. Alternatively, another property not quantified in this work, such as structural changes or the distribution of strains relative to $E$, may control when and where these brittle failures nucleate out of diffuse damage.

The plateau in $B^*$ with decreasing $\gamma_0$ suggests convergence towards a brittle-failure regime in which the number of bonds broken at failure approaches $N_b^*\approx 700 \approx 3 \sqrt{N}$, albeit with a larger variance than in networks that fail early in the response. As indicated in \cref{fig:HighGammaStrainSnapShots}a and \cref{fig:HighGammaStrainSnapShots}b, which show failures within this plateau, the failure modes are similar across network structures and strain in this regime. 

Consistent with the low-strain regime, when the imposed strain is large relative to the breakage energy $E$, many bonds are brought close to their breaking threshold. In these cases, bond failure is more rapid and diffuse than at smaller strains, and the network either fails through diffuse bond loss or through the formation of several smaller macroscopic fractures, as illustrated in \cref{fig:HighGammaStrainSnapShots}b.

\section{Direction of Fracture Propagation} \label{sec:DirectionOfFracture}

As discussed in \cref{sec:lowStrainResponse}, at small strain, the direction of macroscopic crack propagation in the examples studied here is either horizontal or vertical, with no clear directional preference. In contrast, at high strains $\gamma_0\sim 1$, as shown in \cref{fig:HighGammaStrainSnapShots}, macroscopic failure consistently propagates perpendicular to the principal strain axis, and approximately along a direction running from the top left to the bottom right of the sample.

For simplicity, we consider a network immediately after it has undergone an affine deformation with shear strain $\gamma_0$. At this instant, the imposed step strain has displaced every node according only to the externally applied shear field, before the internal forces in the disordered network have had time to drive any additional non-affine relaxation. The network configuration at this stage can therefore be derived analytically.

Consider a bond that makes an angle $\theta$ with the horizontal axis in the network’s initial configuration. Under a shear deformation of magnitude $\gamma_0$ in the $xy$ plane characterised by the matrix
\begin{equation}
    \left(\begin{matrix}
        1 & \gamma_0 \\
        0 & 1 \\
    \end{matrix}\right),
    \label{eq:ultraShearMatrix}
\end{equation}
take, without loss of generality, a bond of length $l_0$ with one node fixed at the origin and the other located at $(l_0 \cos\theta,\; l_0 \sin\theta)$. Under the shear transformation defined in \cref{eq:ultraShearMatrix}, the node at the origin remains unchanged, while the second node is mapped affinely to
\begin{equation*}
    \left(\begin{matrix}
        1 & \gamma_0 \\
        0 & 1 \\
    \end{matrix}\right)
    \left(\begin{matrix}
        l_0\cos\theta \\
        l_0\sin\theta \\
    \end{matrix}\right)
    =
    \left(\begin{matrix}
        l_0(\cos\theta + \gamma_0\sin\theta) \\
        l_0\sin\theta \\
    \end{matrix}\right).
\end{equation*}

The post-shear bond therefore has length
\begin{equation}
l(\theta \mid \gamma_0)
= l_0 \sqrt{\bigl(\cos\theta + \gamma_0 \sin\theta\bigr)^2 + \sin^2\theta},
\end{equation}
and makes an angle $\theta^\prime$ with the horizontal axis, where
\begin{equation}
\tan\left( \theta^\prime(\theta \mid \gamma_0) \right)
= \frac{l_0\sin\theta}{l_0(\cos\theta + \gamma_0\sin\theta)}
= \frac{\sin\theta}{\cos\theta + \gamma_0\sin\theta}.
\end{equation}
The corresponding bond strain is
\begin{align}
\epsilon_b(\theta \mid \gamma_0)
&= \frac{l(\theta \mid \gamma_0) - l_0}{l_0} \\
&= \sqrt{\bigl(\cos\theta + \gamma_0 \sin\theta\bigr)^2 + \sin^2\theta} - 1.
\end{align}

\begin{figure}[t]
    \centering
    \includegraphics[width=0.9\linewidth]{chapters/UltraDelayedFracture/CrackDirections/BondStrain.pdf}
    \caption[Illustrations of bond state under affine deformation]{
    Illustrations of {\bf a)} the strain $\epsilon_b$ and {\bf b)} the angle $\theta^\prime$ resulting from an affine deformation of a bond initially oriented at an angle $\theta$. We plot the two curves corresponding to two strains considered in this work $\gamma_0=0.01$ and $\gamma_0=1$. The principal angle $\theta\in[0,\pi)$ is measured relative to the horizontal $x$-axis. The zero of each curve, indicated with a dashed vertical line, shows the angle $\theta\in(\pi/2,\pi)$ at which the applied strain transitions from extension to compression.}
    \label{fig:AffineBondStrain}
\end{figure}

This expression for $\epsilon_b(\theta \mid \gamma_0)$ neglects any non-affine relaxation of the surrounding network and should therefore be regarded as an approximation of the bond strain in the relaxed network. We plot in \cref{fig:AffineBondStrain}a, $\epsilon_b(\theta| \gamma_0)$ for the two strain amplitudes considered in this work, $\gamma_0=0.01$ and $\gamma_0=1$. 

For small $\gamma_0$, the resultant affine strains remain small and close to zero with a roughly equal split of bonds in both compression and extension (red solid line in \cref{fig:AffineBondStrain}a). In the limit of vanishing affine strain, this distribution would be exactly symmetric about $\epsilon_b = 0$. The transition is indicated by the vertical dashed line associated with the red ($\gamma_0=0.01$) curve. The magnitudes of the resultant strains also remain of similar size, although with a slight bias to elongated bonds, as is exemplified by the $\gamma_0=1$ case. The resultant distribution both immediately following strain imposition (dotted red line) and in the relaxed configuration (solid red line) is plotted in \cref{fig:PDFBondDistributions}a. When force balance is reimposed via the network relaxing to an equilibrium configuration, the affine distribution progressively broadens compared to the distribution at $t=0^+$, though importantly remains roughly symmetric around $\epsilon_b=0$. In this low-strain regime, because the breakage condition considers the bond strain energy $E=\frac{\mu}{2}\epsilon_b^2$ rather than the bond strain $\epsilon_b$, the network's fracture is agnostic to the strain being compressive or extensional. As a result, fracture propagation at low strains does not show a preferred orientation at which bonds fail, with bonds failing under both extension and compression.

\begin{figure}[t]
    \centering
    \includegraphics[width=0.9\linewidth]{chapters/UltraDelayedFracture/CrackDirections/StrainDistributions.pdf}
    \caption[Normalised strain distributions $\rho(\epsilon_b/\gamma_0)$ for a relaxed network and affine (dashed) networks]{Normalised strain distributions $\rho(\epsilon_b/\gamma_0)$ for a relaxed network (solid line), affinely deformed (dashed line) networks and analytical predictions (dotted line), for {\bf a)} $\gamma_0=0.01$ and {\bf b)} $\gamma_0=1.0$. Affine network distributions are taken immediately after a step strain of magnitude $\gamma_0$ is imposed affinely. The analytical form of the affine prediction is derived in \cref{sec:affineStrainDist}.}
    \label{fig:PDFBondDistributions}
\end{figure}

In contrast, the distribution of bond strain in the affinely deformed network at $\gamma_0=1$ is asymmetric about $\epsilon_b=0$, with a bias towards positive (extensional) strains. When the network subsequently relaxes to restore force balance, this distribution broadens further, producing a larger population of bonds with high strain magnitude. Consequently, the bonds most likely to fail are those whose initial orientation in the undeformed frame is $\theta\approx58^{\circ}$ (dotted blue line in \cref{fig:AffineBondStrain}a), corresponding to the angle that produces the largest affine deformation. This \DIFdelbegin \DIFdel{crack failure mechanism is again }\DIFdelend \DIFaddbegin \DIFadd{predicted directional bias in bond failure is }\DIFaddend illustrated in \cref{fig:CrackDirections}b. We illustrate in black the directions parallel and perpendicular to the maximally strained bonds ($\theta\approx58^{\circ}$), which show qualitative agreement with the observed failure pattern in the network, for the eventual failure directions observed in both strain regimes ($\gamma \ll 1$ in \cref{fig:CrackDirections}a and $\gamma \sim 1$ in \cref{fig:CrackDirections}b). The exact fracture direction will depend strongly on the particular network topology and the local level of non-affine arrangement.

\begin{figure}[t]
    \centering
    \includegraphics[width=0.99\linewidth]{chapters/UltraDelayedFracture/CrackDirections/CrackSnapshots.png}
    \caption[Snapshot of crack features plotted in its undeformed frame]{
    Snapshots of systems in the rest frame at failure. Arrows indicate directions parallel and perpendicular to lines at $45^\circ$ (solid) and $58^\circ$ (dashed) relative to the horizontal. The colour scale indicates the order in which bonds fail from dark to light. \newline
    Parameters {\bf a)} $\gamma_0=0.009$, $T=10^{-6}$, $E=5.025\times10^{-5}$ and {\bf b)} $\gamma_0=1.0$, $T=0.0062$, $E=0.35$, $N=64{,}000$.\newline
    }
    \label{fig:CrackDirections}
\end{figure}

Finally, of particular interest are the side cracks that branch off from the initial failure and form in the small-$\gamma_0$ regime, shown in \cref{fig:LowGammaTempSnapShots}c, \cref{fig:LowGammaTempSnapShots}d, \cref{fig:LowGammaStrainSnapshots}a and, for comparison to the high-$\gamma_0$ regime, in \cref{fig:CrackDirections}a. While cracks at these small strains initially propagate in both the horizontal and vertical directions, a defining feature of these single brittle events is the splitting at the crack tip. The primary fracture bifurcates into secondary branches that propagate along diagonal directions through the network. In the undeformed reference frame, these branches appear to be oriented at approximately $45^\circ$ to the principal axes.

% Although we do not directly resolve the local stress or rigidity during branching, two plausible mechanisms may contribute to the observed crack-tip splitting. First, progressive bond loss near the crack tip likely reduces the local stiffness, permitting larger non-affine rotations and shear concentrations that deflect the advancing crack tip away from the principal axes. Second, once a crack forms, stress propagation across the failed region is suppressed, so the stress redistribution is expected to differ from the Eshelby-like redistribution characteristic of an intact elastic medium and this altered redistribution may favour the loading of bonds along oblique directions and hence branching. Testing these hypotheses would require direct measurements of the evolving local stiffness (or rigid cluster structure) and of the spatial pattern of stress/strain energy release around the tip, as is done experimentally under normal deformation in Refs. \cite{ju_role_nodate,qiMappingDeformationDissipation2024,zheng_how_2021,ducrot_structure_2015,wangCharacterizationDeformationFracture2025}.

Taken together, these observations suggest that the direction of fracture in these networks is controlled by a competition between the initially imposed affine strain field and the subsequent redistribution of stress during relaxation and failure. At small imposed strains, the affine strain distribution and resulting stress field render the network effectively isotropic with respect to fracture. This allows cracks to nucleate and propagate along either horizontal or vertical directions with comparable likelihood, with subsequent branching arising from local deformations at the crack tip. At high strain, however, the anisotropic strains resulting from the affine deformation strongly bias the population of highly strained bonds toward a well-defined orientation. After relaxation and force balance, these bonds select a preferred fracture direction perpendicular to these maximally loaded bonds. While non-affine deformations and network structure introduce variability into the precise failure path, the overall fracture direction is nevertheless governed by the underlying strain-induced anisotropy of the bond loading.

\section{Conclusion}
In this chapter, we have investigated the delayed failure of soft network materials under a simple set of assumptions. Our central aim has been to show that the broad class of ultra-delayed post-strain instabilities highlighted by Lockwood \etal~\cite{lockwood_ultradelayed_2025} has a clear fracture analogue in direct simulations of network materials. Within a minimal thermal model with thermally activated bond breakage events, we have shown that the failure time of the network can be pushed to extremely long timescales following the imposed deformation. In this regime, the networks can appear macroscopically stable over these long timescales, exhibiting little to no obvious precursor in the bulk material properties, before the eventual rupture of a critical bond triggers a catastrophic, brittle failure of the system.

At low imposed step strains, we found that networks that ultimately undergo ultra-delayed failure initially relax in a manner that is almost indistinguishable from an equivalent athermal network in which bonds are not permitted to break. Following this initial non-affine relaxation, the system enters a long-lived stable state in which further mechanical deformation is minimal. In this state, bond failure is rare and controlled by thermal fluctuations in the bonds rather than by deterministic mechanical processes. The rupture of bonds is therefore well described by a stochastic Poisson process, modelling the process in which each bond explores its local energy well and occasionally escapes over the energy barrier that impedes its breakage. 

The resulting macroscopic fractures in these networks typically nucleate from the failure of a single critical bond whose local configuration can trigger a cascade of correlated failures. Once initiated, this cascade of breakages leads to a brittle, system-spanning fracture. In the low-strain regime, these cracks predominantly propagate along horizontal or vertical directions in the network's rest frame, consistent with the near-isotropic distribution of bond strains. During crack formation, we observe the development of secondary branching, causing the fracture to change direction and split, with the direction appearing to depend on the local network structure and the imposed strain. 
% We argued that this behaviour arises from bond failure, resulting in both a local loss of rigidity at the crack tip and modifications to the stress propagator as the network becomes damaged.

In contrast, at large imposed strains, the network is driven far beyond the regime of linear elasticity, and substantial non-affine deformations occur during relaxation. Here, a large number of aligned bonds are pushed close to their failure thresholds immediately after the step strain, leading to widespread, spatially diffuse bond breakage at early times. Despite this early damage, we found that networks can still exhibit ultra-delayed failure, persisting essentially intact for many decades in time before ultimately fracturing. Interestingly, in this regime, the material can sustain a much larger number of broken bonds before macroscopic failure, compared to the low-strain case, suggesting that the strain stiffening exhibited by highly deformed networks allows the network to accommodate progressive damage without immediate collapse.

\DIFdelbegin \DIFdel{Taken together, these results suggest that, despite the simplicity of }\DIFdelend \DIFaddbegin \DIFadd{The correspondence with Lockwood }\etal\DIFadd{~\mbox{%DIFAUXCMD
\cite{lockwood_ultradelayed_2025} }\hskip0pt%DIFAUXCMD
extends beyond the generic occurrence of delayed failure. Their constitutive models and the network model considered here employ the same basic strain-assisted activated dynamics, in which stored local elastic energy lowers the barrier to a yielding or bond-breaking event. In }\DIFaddend the \DIFdelbegin \DIFdel{model, we have captured the essential physical ingredients underlying ultra-delayed fracture driven by thermal effects. In these disordered networks, the interplay between non-affine deformations, thermally activated bond failures, and stress redistribution controls their mechanical response. The fact that qualitatively similar behaviour emerges across different parameter regimes, as well as in the broader class of models studied by Lockwood }%DIFDELCMD < \etal%%%
\DIFdel{, indicates that delayed post-strain failure may be a robust phenomenon}\DIFdelend \DIFaddbegin \DIFadd{small-strain affine limit, the most highly loaded bonds in our network have $|\epsilon_b|\simeq\gamma_0/2$, so }\cref{eq:SGRRate} \DIFadd{gives the estimate $\log t^*\sim(E-\mu\gamma_0^2/8)/T$. This has the same functional dependence, $\log t^*\propto-\gamma_0^2/T$, as the scaling reported for the SGR model in Ref.~\mbox{%DIFAUXCMD
\cite{lockwood_ultradelayed_2025}}\hskip0pt%DIFAUXCMD
. The difference in the numerical coefficient is likely to reflect both the projection of the applied shear strain onto the axial strain of an individual bond and differences in how stress is redistributed following a local plastic event. The critical-bond cascade described above is also the network counterpart of their positive-feedback instability, in which local stress relaxation redistributes load and promotes further nearby events. In their constitutive models this feedback produces a shear band, whereas irreversible bond removal here converts it into a macroscopic crack}\DIFaddend .

At the same time, the complex and often unpredictable nature of failure, from failure time to crack structure, raises new questions about the microscopic mechanisms that ultimately govern when and where catastrophic failure nucleates. We therefore suggest that future work in this area could include:
\begin{itemize}
    \item {\bf Alternative network models} -- In this work, thermal effects are incorporated through a rate of bond failure that depends on a noise temperature $T$. An alternative and widely used approach in molecular dynamics is to introduce temperature via a Brownian or Langevin thermostat \cite{allen_computer_2017}, in which thermal fluctuations act directly on the node dynamics through a stochastic equation of motion. Work by Tauber et al. in Ref.~\cite{tauber_role_2020} showed that increasing temperature in such models leads to more homogeneous stress redistribution and enhances thermally activated bond failure, in qualitative agreement with the behaviour captured by the rate-based formulation used here (see \cref{eq:SGRRate}).

    We therefore expect the two approaches to yield broadly similar phenomenology, although they differ in how thermal effects enter the model. In the stochastic formulation, the temperature $T$ sets the scale of Brownian fluctuations relative to the elastic energy of the bonds, rather than appearing explicitly as an activation parameter in the breakage rate. This description may provide a more direct representation of materials such as elastomers and hydrogels, where thermal fluctuations and entropic elasticity play a central role \cite{arruda_three-dimensional_1993,james_theory_1943}. Such models are, however, significantly more computationally demanding, typically requiring much smaller timesteps than the dynamics used elsewhere \cite{leimkuhler_robust_2013}.

    [Early simulations by the author into this model proved successful, although they were unable to reach the system sizes and timescales needed to fully characterise the dynamics.]

    \item {\bf Athermal fracture} -- As discussed in \cref{sec:delayedInto}, we have focused on delayed failure arising from thermally activated bond breaking, treating purely athermal failure as unlikely in the regime studied here. In the constitutive models explored in Ref.~\cite{lockwood_ultradelayed_2025}, delayed failure is absent in the athermal limit because there is no intrinsic timescale associated with local plastic rearrangements.

    However, in the networks considered here, the relaxation time increases as the connectivity is reduced towards the isostatic point $z_c$, diverging as a power law at $z_c$ \cite{shivers_strain-controlled_2023}. This could, in principle, provide an alternative athermal route to ultra-delayed failure if the material relaxes sufficiently slowly. However, in this regime, the sample-to-sample variance of the fracture strain also grows significantly \cite{zhang_fiber_2017}, so failure becomes strongly dependent on the specific network realisation. Moreover, it is not clear that a marginally stable network is able to store sufficient elastic energy to fail a long time after the imposed strain.  As a result, although an athermal mechanism cannot be excluded, it would require fine tuning of connectivity and imposed strain (if it is possible), and is therefore unlikely to provide a robust route to ultra-delayed failure.

    \item {\bf Stress redistribution and Critical bonds} -- A key question raised by the results presented in this work concerns how stress is redistributed during fracture. Much of the existing literature on this problem considers the evolution of a crack tip under predominantly elongational loading perpendicular to the fracture plane \cite{sanner_less_2025, bolander_fracture_1998, shekhawat_damage_2013, slepyan_models_2002, ju_role_nodate}, a regime in which the local deformations near the crack tip are always large. In contrast, in the low-strain regime studied here, we observe an abrupt transition in fracture behaviour at some point during the failure. This suggests that the local stress redistribution evolves qualitatively during propagation. 

    Further analytical progress may be possible by extending classical treatments of the Eshelby stress propagator \cite{picard_elastic_2004, eshelby_determination_1997} to account for the presence of non-homogeneous fracture surfaces or more complex wall geometries resembling those created by pre-existing damage. Such approaches could provide deeper insight into how stress is redistributed in the vicinity of evolving cracks in disordered networks.

    Finally, we introduced the concept of a critical bond, whose failure resulted in the immediate catastrophic failure of the network. The existence of such critical bonds or regions of heightened susceptibility have been widely observed in amorphous materials, both experimentally and computationally \cite{mattsson_time-dependent_2017,cipelletti_microscopic_2020,yu_shortest_2024,membrillo_solis_tracking_2022,ozawaRareEventsDisorder2022,aimeMicroscopicDynamicsFailure2018,nabizadeh_network_2024}. Continued advances in experimental techniques \cite{porr_probing_2025} and numerical or data-driven methods \cite{mangal_predicting_2024,mangal_topological_2023,cubuk_identifying_2015} may, in the future, enable the identification of where such failures nucleate and ultimately provide routes to predicting or controlling their dynamics.

    \item {\bf Relation to gel collapse in colloidal networks} -- A potentially related phenomenon occurs in colloidal gels and arrested particle networks, for which experiments have shown that the materials can support their own weight for long times before ultimately collapsing catastrophically \cite{harich_gravitational_2016,bartlett_sudden_2012,manley_gravitational_2005,starrs_collapse_2002}. It has been shown in these systems that the network never truly reaches a static state: individual bonds continuously break and reform due to thermal fluctuations, and the microstructure slowly evolves. Therefore, this form of delayed yielding has been interpreted as a form of thermally activated creep rupture, where rare microscopic rearrangements progressively weaken load-bearing strands until a critical failure event nucleates \cite{blindstrom_structures_2012,teece_gels_2014,sprakel_stress_2011}.

    These behaviours suggest a natural analogy with the ultra-delayed failures observed here under step strain. In both cases, the systems appear macroscopically quasi-stationary over long times while microscopic damage accumulates, with rare bond- or strand-level events ultimately nucleating an instability. A useful direction for future work is to consider the model explored in this chapter, to compare the resulting statistics of delay times and precursor activity against experiments.

    Ultimately, this may be able to separate the purely dynamics-based delay routes from the solvent-mediated mechanisms, alternatively proposed as a source of delayed failure \cite{harich_gravitational_2016,varga_modelling_2018,padmanabhan_gravitational_2018}. Here, hydrodynamical erosion of the gel forms fluidised channels (\quotes{streamers}), which have been proposed to destabilise the gel's structure. Establishing which ingredients are essential across these systems would clarify whether ultra-delayed instabilities are a generic consequence of thermally activated rupture in disordered networks, or whether additional hydrodynamical couplings are required in the colloidal case.

    \item {\bf Experimental studies with the protocol of step shear strain} -- At present there remains limited direct experimental evidence for ultra-delayed fracture or failure under a true step strain protocol. Most work has been limited to short timescale failure in polymer melts \cite{boukany_step_2009,barroso_evaluation_2002,einaga_stress_1973} or simple constitutive models \cite{moorcroft_shear_2014,moorcroft_criteria_2013,adams_nonmonotonic_2009}. With consequences for material manufacturing, it is important to test this experimentally, in both network materials and, more generally, amorphous materials. The major problem with observing this experimentally is the parametrisation of the system such that the failure time ultimately occurs within experimental timescales, but also late enough to be defined as delayed rather than immediate failure. Identifying such a parameter regime is already non-trivial at the computational level, and is therefore expected to pose an even greater challenge in experimental realisations. While a true step strain protocol is inaccessible experimentally, a rapid strain ramp, at a rate quick enough to avoid any failure before the maximum strain $\gamma_0$ is reached, is achievable \cite{barroso_evaluation_2002}. 
\end{itemize}

\newpage
\begin{subappendices}

\section{Calculating Approximate free parameters in Breakage Rate} \label{sec:approxEx}
Predicting the relaxed post-strain configurations of disordered networks is inherently challenging. This problem has been extensively studied in the context of rigidity (see \cref{sec:Rigidity}). Recent work \DIFdelbegin \DIFdel{exploring }\DIFdelend \DIFaddbegin \DIFadd{on }\DIFaddend the high-dimensional energy landscapes of such networks has \DIFdelbegin \DIFdel{demonstrated that predicting their evolution is computationally intractable without direct numerical simulation }\DIFdelend \DIFaddbegin \DIFadd{shown that their mechanically relaxed state cannot generally be inferred using analytical or reduced-order approximations, but must instead be determined by numerically minimising the full network energy }\DIFaddend \cite{damavandi_energetic_2022,damavandi_energetic_2022-1}. These difficulties arise primarily from the large non-affine deformations characteristic of semi-flexible networks. Because the final parameter ranges will ultimately be refined through systematic parameter space searches, we require only a reasonable initial estimate rather than an exact prediction.

Based on the discussion in \cref{sec:DirectionOfFracture} for a network immediately after it has undergone an affine deformation with shear strain $\gamma_0$, a bond that, in the network’s initial (rest) configuration, makes an angle $\theta$ with the horizontal axis, has a strain defined by
\begin{equation*}
     \epsilon_b(\theta| \gamma_0) = \sqrt{(\cos(\theta) + \gamma_0\sin(\theta))^2 + \sin^2(\theta)} - 1.
\end{equation*}

For the large, disordered networks considered here, we assume that bond orientations are approximately uniformly distributed. We therefore consider the maximum value of $\epsilon_b(\theta|\gamma_0)$ over $\theta \in [0,\pi)$ and define
\begin{equation}
\epsilon_{\max}(\gamma_0) = \max_{\theta \in [0,\pi)} \left\{ \epsilon_b(\theta|\gamma_0) \right\}.
\end{equation}
The corresponding maximum expected bond energy at strain $\gamma_0$ is then $\mu \epsilon_{\max}^2/2$. To estimate a suitable central value for the strain parameters, we evaluate this energy at the midpoint of the parameter range of interest, multiplying by a factor of two to allow for the non-affine relaxations neglected by this estimate. This provides a practical starting point for the parameter exploration.

We adopt a similar strategy to obtain an approximate starting value for the temperature-like parameter $T$. Here we consider the mean bond energy, defined as
\begin{equation*}
    \langle \epsilon_b \rangle = \int_0^{2\pi} \frac{\mu}{2} \epsilon_b(\theta | \gamma_0)^2 d\theta.
\end{equation*}
Substituting this average strain into the failure rate given in \cref{eq:SGRRate} yields an average bond failure rate $\langle \Gamma \rangle$ for a typical bond in the system. Assuming that bond failures follow a Poisson process, such that the mean lifetime equals $1/\langle \Gamma \rangle$, we choose $T$ such that the average bond survives for a time $\langle \Gamma \rangle \approx 10^{20}$ in simulation units.

These estimates \DIFdelbegin \DIFdel{provide a reasonable initial parameter set. However, additional tuning is required to identify parameter combinations that best capture the relevant physics across the full range of behaviours studied in the main text}\DIFdelend \DIFaddbegin \DIFadd{identify a starting region in parameter space. As discussed in }\cref{sec:BoltzmanFracture}\DIFadd{, we then explore nearby values of $E$ and $T$ to locate the narrow intermediate regime in which failure is delayed but remains observable within the accessible simulation time}\DIFaddend .

\section{Strain Distribution for an Affine Deformation} \label{sec:affineStrainDist}
As discussed in \cref{sec:approxEx}, the expected strain following an affine deformation of magnitude $\gamma_0$, for a bond with principal angle $\theta\in[0,\pi)$ is
\begin{equation}
         \epsilon_b(\theta| \gamma_0) = \sqrt{(\cos(\theta) + \gamma_0\sin(\theta))^2 + \sin^2(\theta)} - 1. \label{eq:affine1}
\end{equation}
We define $u\coloneq\epsilon_b+1$ and square \cref{eq:affine1} to obtain
\begin{align*}
        u^2 
        &=(\cos(\theta) + \gamma_0\sin(\theta))^2 + \sin^2(\theta) \\
        &= 1+\gamma_0^2\sin^2(\theta) + \gamma_0\sin(2\theta) \\
        &= 1 + \gamma_0^2\left[\frac{1}{2}(1-\cos(2\theta))\right] + \gamma_0\sin(2\theta)\\
        &= 1+ \frac{\gamma_0^2}{2} - \frac{\gamma_0^2}{2}\cos(2\theta)+\gamma_0\sin(2\theta)
\end{align*}
By defining $C\coloneq1+ \frac{\gamma_0^2}{2}$, $q\coloneq- \frac{\gamma_0^2}{2}$ and using the identity $\gamma_0\sin(2\theta)+q\cos(2\theta)=R\sin(2\theta + \phi)$, where 
\begin{equation*}
    R=\sqrt{\gamma_0^2+q^2}=\sqrt{\gamma_0^2+\frac{\gamma_0^4}{4}},
\end{equation*}
we can express that
\begin{equation*}
    u^2=C+R\sin(2\theta+\phi)
\end{equation*}
where $\phi$ solves $\tan(\phi) = -\gamma_0/2$.

If we assume that the initial bond angle takes the form of a uniform distribution defined by $\theta \sim \text{Unif}[0,\pi)$, then $\alpha\coloneq2\theta+\phi$ also takes the form of a uniform distribution on $[0,2\pi)$. The distribution $U$ of $\sin(\alpha)$ has a known form 
\[
    f(u)=\frac{1}{\pi\sqrt{1-u^2}}, \qquad u\in[-1,1].
\]
We can consider the affine transformation of $U$ to $u^2$, $u^2 = C + R U$.
Inverting this relation gives
\[
    U = \frac{u^2 - C}{R} = \frac{(\epsilon+1)^2 - C}{R}.
\]
A change of variables then yields the probability density function of $\epsilon$ as
\begin{align*}
    f_\epsilon(\epsilon)
    &= f\left(\frac{(\epsilon+1)^2-C}{R}\right)
      \left|\frac{\mathrm{d}}{\mathrm{d}\epsilon}
      \left(\frac{(\epsilon+1)^2-C}{R}\right)\right| \\
    &= \frac{2(\epsilon+1)}{\pi R\sqrt{1-\left(\frac{(\epsilon+1)^2-C}{R}\right)^2}},
\end{align*}
with bounds determined by $U\in[-1,1]$. This implies
\[
    \epsilon \in \left[\sqrt{C-R}-1,\; \sqrt{C+R}-1\right].
\]

\section{Verification of Numerics} \label{sec:DelayedConvergenceChecks}

\begin{figure}[t]
    \centering
    \includegraphics[width=0.9\textwidth]{chapters/UltraDelayedFracture/Convergence/Combined_Plot.pdf}
    \caption[RMS velocity after a strain step and convergence of failure time with network size.]{{\bf Left:} Root mean squared velocity against time $t$ for a system strained to $\gamma_0=1$ at $t=0$. We plot the response for four values of absolute error tolerance defined in \cref{sec:NetworkDynamics}. {\bf Right:} Convergence of failure time $t^*$ with increasing network size at $\gamma_0=1$, $E=1.0$, $T=0.0062$.}
    \label{fig:DelayedConvergence}
\end{figure}

When varying the absolute error tolerance $\tau$ in the adaptive time-stepping scheme defined in \cref{sec:NetworkDynamics}, we observe little quantitative change in the relaxation dynamics. As shown in \cref{fig:DelayedConvergence}a, the response remains essentially unchanged as $\tau$ is varied across four orders of magnitude, from $10^{-2}$ to $10^{-5}$. We therefore adopt the smallest tolerance that still permits reasonable computation times, which we take to be $\tau = 10^{-4}$. The Bernoulli trials governing bond breaking, defined using the rate in \cref{eq:SGRRate}, remain numerically accurate across all time step sizes. Nevertheless, we impose an upper bound on the adaptive time step of $\delta t = 0.1$ to ensure that individual bond failure times are resolved reliably.

In \cref{fig:DelayedConvergence}b, we further demonstrate that the characteristic failure time $t^*$ converges with increasing system size. All results presented in this chapter were therefore obtained using networks containing $N = 64{,}000$ nodes.


\end{subappendices}
 \newpage 
 \newpage % !TeX root = ../../thesis.tex

\setcounter{equation}{0}
\chapter{\DIFdelbegin \DIFdel{Toughness }\DIFdelend \DIFaddbegin \DIFadd{Failure and Toughening }\DIFaddend of Double Network Hydrogels} \label{sec:DoubleNetworks}

\section{Introduction}

In the previous chapter, we discussed the delayed thermal failure of soft network materials that occurs a long time after an applied deformation. These materials, such as rubbers or gels, exhibit a predominantly elastic mechanical response, deforming reversibly up to the point of failure. This has motivated their use in applications, including biomedical tissue engineering \cite{leeHydrogelsTissueEngineering2001}, stretchable electronics \cite{rogers_materials_2010}, and soft robotics \cite{martinezSoftActuatorsRobots2014}. However, many of these materials typically exhibit abrupt and brittle failures at low strains, which lead to inferior mechanical performance \cite{Creton50th2017}. One of the most promising strategies for overcoming this limitation is to introduce a second polymer network to the material, to create a so-called double network. These composite materials have proved to provide remarkable material properties in double network hydrogels \cite{gong_why_2010, nakajima_synthesis_2013, sun_highly_2012, xin_effect_2013}, elastomers \cite{ducrot_toughening_2014, millereauMechanicsElastomericMolecular2018} and macroscopic composites \cite{kingMacroscaleDoubleNetworks2019,takahashiCreatingStiffTough2018}.

The work in this chapter can be considered to focus on double network hydrogels, motivated by the breadth of experimental data on these materials. However, because the model contains only generic ingredients, we expect it to provide qualitative insight into other composite materials constructed from two network types. The pioneering work which has inspired the work in this chapter was performed by Gong \etal \cite{gong_double-network_2003}, who combined a stiff and weak \quotes{sacrificial} network with a soft and extensible \quotes{matrix} network. The resultant double network hydrogel retained the stiffness of the sacrificial component while acquiring the large failure strain of the matrix component.

Experimentally, double network hydrogels have been studied extensively, and there is, therefore, already a wealth of insights into their failure mechanics. Small-angle neutron scattering probing the characteristic length scales of internal non-affine rearrangements has suggested that stress becomes delocalised in double networks \cite{ducrot_structure_2015, tominaga_molecular_2007}. Birefringence imaging further reveals that non-optimised double network structures exhibit large stress concentration when compared to more optimised structures, which distribute the stress more uniformly \cite{gong_double-network_2003, zheng_how_2021}. Dynamic light scattering experiments have demonstrated a strong coupling between fluctuations in the two networks consistent with mutual entanglements and the ability for the two network components to share load \cite{huang_importance_2007,Orr2026}. Additionally, increasing the molar fraction of the matrix network relative to the sacrificial network favours the formation of diffuse microcracks indicative of ductile behaviour, rather than the propagation of a single macroscopic crack associated with brittle failure \cite{fukao_effect_2020, ju_role_nodate}.

\DIFaddbegin \DIFadd{An important distinction in discussing the toughness of double network materials is between their response under homogeneous deformation and their resistance to the propagation of an existing crack. In fracture experiments, toughness is commonly characterised through a fracture energy, corresponding to the energy required to advance a crack per unit newly created area. This differs from a measure based on the integrated area under a stress-strain curve up to failure. Experimentally, optimised double network hydrogels perform strongly according to both measures: they can sustain tensile strains of order $1000$--$2000\%$ \mbox{%DIFAUXCMD
\cite{gong_why_2010,ahmedBrittleDuctileTransition2014} }\hskip0pt%DIFAUXCMD
and exhibit fracture energies substantially larger than those of single networks formed from either constituent \mbox{%DIFAUXCMD
\cite{tominaga_molecular_2007,xin_effect_2013,Creton50th2017}}\hskip0pt%DIFAUXCMD
.
}

\DIFadd{This enhancement in fracture resistance is commonly associated with the development of an extended damage zone around the crack tip \mbox{%DIFAUXCMD
\cite{brown_model_2007,tanaka_local_2007,zheng_how_2021,huang_importance_2007,ju_role_nodate,sanner_less_2025,Persson2024}}\hskip0pt%DIFAUXCMD
. Within this region, sacrificial bonds rupture over an extended volume while the matrix preserves connectivity, allowing energy to be dissipated away from the immediate crack tip and reducing the tendency for a single crack to propagate catastrophically \mbox{%DIFAUXCMD
\cite{gong_why_2010,fukao_effect_2020,millereauMechanicsElastomericMolecular2018}}\hskip0pt%DIFAUXCMD
. The resulting process therefore involves physics across several length and time scales, from individual bond failures to a macroscopic crack and its surrounding damage zone.
}

\DIFadd{The simulations considered in this chapter address a more restricted part of this problem. We consider initially unnotched periodic systems undergoing approximately homogeneous deformation and therefore do not directly calculate a fracture energy or resolve the propagation of a pre-existing macroscopic crack. Instead, we focus on a network-scale mechanism that can contribute to the experimentally observed toughening: how coupling to the matrix modifies the stress redistribution following sacrificial-bond failure and allows damage in the sacrificial component to remain diffuse rather than localising immediately into a catastrophic fracture.
}

\DIFaddend Early work in continuum models undertaken by Tanaka~\cite{tanaka_local_2007} and Brown~\cite{brown_model_2007} looked at the breakage of the sacrificial bonds, which were stabilised by a background matrix network. These models assumed that the stress field in the material can be described by continuum mechanics, with both networks acting like a homogeneous elastic solid. These models were later extended to consider the evolution of damage in the material and have been able to describe the hysteresis in the stress-strain response, \DIFdelbegin \DIFdel{the progressive softening}\DIFdelend \DIFaddbegin \DIFadd{progressive softening, }\DIFaddend and the Mullins effect\DIFaddbegin \DIFadd{, a history-dependent stress softening in which subsequent loading follows a lower-stress path until the previous maximum strain is exceeded, }\DIFaddend observed in double network hydrogels \cite{vernereyStatisticalDamageMechanics2018, wang_pseudo-elasticity_2011, bacca_model_2017, lavoie_continuum_2019,okumura_toughness_2004,zhao_theory_2012}. While these models can capture the macroscopic behaviours of the materials, they cannot resolve the detailed underlying microphysics.

This microphysics can instead be investigated using molecular dynamics simulations \cite{jang_mechanical_2007,wang_simulational_2017,higuchi_fracture_2018,mugnaiInterspeciesInteractionsDual2025,tauber_sharing_2021}. However, these models often fail to reach the necessary system sizes to access fracture or explore enough of the required parameter space due to the high computational costs associated with these techniques. Therefore, mesoscale models, such as the one presented later in this chapter, have proven successful in capturing a simplified representation of microphysics while still allowing access to large enough systems or parameter ranges. Previous mesoscopic models, which we use as inspiration for our work, have examined fracturing in single-component materials \cite{dussi_athermal_2020, lei_mesoscopic_2021}. Alternative double network frameworks have already been proposed; however, they either fail to resolve the difference in characteristic length scales between the two networks \cite{tauber_microscopic_2020}, consider a sparse unbreakable network \cite{noguchi_fracture_2024}, or focus on the edge/necking effects of non-periodic systems with non-standard network constructions \cite{li_fracture_2024}. 

This work aims to expand on these pre-existing models and explore how a simple mesoscopic model can be used to analyse the effect of differing characteristic length scales between the two networks. In our tough double networks, these characteristic lengths occur due to the distinct structural scales associated with the two interpenetrating networks. The stiff but weak sacrificial network is sparse with a large mesh size when compared to the soft, extensible but dense matrix. In particular, we will focus on how these two length scales control both where stress is initially stored elastically and where it is redistributed following a local bond failure. By explicitly considering the double network structure in a controlled manner, we aim to capture the relevant macro- and microphysics while including only minimal ingredients. In the remainder of this chapter, we will introduce our minimal model, describe its numerical implementation, and use it to investigate the fracture behaviour of double network hydrogels across a range of parameters.

\section{Methodology}
To capture the salient physics of double network materials, we must extend the network structures introduced in \cref{sec:NetworkStructure} and used in the previous chapters to model simple networks. In particular, in this chapter, we seek to model both a sacrificial component and a matrix component of a double network and how they interact with each other. As a starting point for our model, we consider a worm-like chain model with the Hamiltonian defined in \cref{eq:NetworkHamiltonian}, with only central forces acting on the nodes. We will take each node to evolve under the dynamics described in \cref{sec:NetworkDynamics}, both within each component, but also in interactions between the two components. To connect the two components, we will consider a subset of nodes in each component to be shared, allowing for the two networks to interact via these inter-network connections. These inter-network connections may correspond to either chemically bonded interactions or physical entanglement points between the two networks, and in either case, they enable load sharing, which has been observed experimentally \cite{ahmedBrittleDuctileTransition2014}. Each bond in our system will then be given a modulus $\mu_s$ and $\mu_m$ and a breakage threshold $\lambda_s$ and $\lambda_m$ for the sacrificial and matrix components, respectively. We will consider a bond to break immediately when its own strain $l/l_0 - 1$ exceeds its corresponding threshold: $\lambda_s$ for sacrificial bonds or $\lambda_m$ for matrix bonds.

We will focus mainly on packing-derived networks as introduced in \cref{sec:NetworkStructure}. Square lattice-derived networks will also be discussed in \cref{sec:SquareAppendix}. These two constructions provide distinct approaches for generating double network geometries by defining how each constituent network is constructed relative to the other and how the subset of nodes shared between the two networks is defined. Importantly, both methods define a relative length scale difference $M$ between the two components (sacrificial and matrix), which will allow us to capture both the relevant microphysics as well as the macroscopic physics of bulk failure. 

\subsection{Packing-Derived Double Networks}
For packing-derived double networks, we will generate the sacrificial component using the method described in \cref{sec:PackingBasedNetworks} by minimising the energy of an ensemble of $N_s$ soft disks with bidispersed radii $R_s$ and $1.4R_s$, each in equal numbers at a packing fraction $\phi = 1.0$. We set the typical length scale of a sacrificial bond by defining the initial size of our periodic unit cell to be $L_{x0} = L_{y0} = \sqrt{N_s}$. Overlapping pairs of disks are then connected by a bond with its natural length $l_0$ being set to the distance between the centres of the two disks. 

The matrix network is constructed similarly, but now starting with an ensemble of $N_m$ disks with radii $R_m$ and $1.4R_m$, where $R_m \leq R_s$, in equal numbers at the same packing fraction as the sacrificial component $\phi = 1.0$. The ratio of $R_s$ to $R_m$ defines $M = R_s / R_m \in \mathbb{R}_{>0}$, which can take any positive real value, because by construction the networks do not need to be commensurate to define the internetwork connections. A different choice of packing fraction could be used to generate different network structures for the two components. However, we suggest that if this property is desired, it is instead achieved through pruning each component individually. Unless otherwise stated, the main-text results are obtained for systems with matrix network size $N_m = 320^2$.

\begin{figure}[bt]
\centering
\resizebox{0.99\textwidth}{!}{\begin{tikzpicture}[
    node distance=1.2cm,
    % Style for the boxes
    imgbox/.style={
        draw, 
        line width=1mm, 
        inner sep=0.5\pgflinewidth
    },
    % Style for the text labels
    labeltext/.style={
        font=\sffamily\Large\bfseries, 
        align=center,
        anchor=center % Ensures rotation happens around the true center
    }
]

% === Column 1: Raw Networks ===
\node[imgbox] (A1) {\includegraphics[width=4cm]{chapters/DoubleNetworkFracture/NetworkGeneration/Sacrificial.pdf}};
\node[imgbox, below=1.2cm of A1] (A2) {\includegraphics[width=4cm]{chapters/DoubleNetworkFracture/NetworkGeneration/Matrix.pdf}};

% === Column 2: Colored Networks ===
\node[imgbox, right=1.5cm of A1] (B1) {\includegraphics[width=4cm]{chapters/DoubleNetworkFracture/NetworkGeneration/Sacrificial_Red.pdf}};
\node[imgbox, right=1.5cm of A2] (B2) {\includegraphics[width=4cm]{chapters/DoubleNetworkFracture/NetworkGeneration/Matrix_Blue.pdf}};

% === Column 3: The Combined Network ===
% This places the west (left) edge of C1 exactly 1.5cm to the right of the Column 2 boundary
\coordinate (Col2Center) at ($(B1.east)!0.5!(B2.east)$);
\node[imgbox, anchor=west] (C1) at ($(Col2Center) + (1.5cm, 0)$) {
    \includegraphics[width=6cm]{chapters/DoubleNetworkFracture/NetworkGeneration/Matrix_Blue.pdf}%
    \llap{\includegraphics[width=6cm]{chapters/DoubleNetworkFracture/NetworkGeneration/Sacrificial_Red.pdf}}%
};

% === LABELS ===
% Using 'at (Node.center)' with a negative xshift ensures perfect Y-axis centering
\node[labeltext, rotate=90] at ($(A1.center) + (-3.2cm, 0)$) {Sacrificial};
\node[labeltext, rotate=90] at ($(A2.center) + (-3.2cm, 0)$) {Matrix};

% Final Label (Centered above the large box)
\node[labeltext, above=0.3cm of C1] {Double};

% === Arrows ===
\draw[->, line width=1mm] (A1.east) -- (B1.west);
\draw[->, line width=1mm] (A2.east) -- (B2.west);

% Vertical-aware horizontal arrows
\draw[->, line width=1mm] (B1.east) -- (C1.west |- B1.east);
\draw[->, line width=1mm] (B2.east) -- (C1.west |- B2.east);

\end{tikzpicture}}
\caption[Figure detailing how a packing-derived double network is constructed]{Figure detailing how a packing-derived double network is constructed. 
{\bf Top row:} Construction of a packing-based network, as described in \cref{sec:PackingBasedNetworks}, by minimising the energy of an ensemble of soft disks of two different radii and connecting the centres of all overlapping disks with a bond. This produces a sacrificial network, with bonds shown in red. \newline
{\bf Bottom row:} Construction of the second matrix network, by performing a second energy minimisation of a new bidispersed ensemble of soft disks. The positions of existing rescaled disks from the sacrificial network are held fixed during this process. This produces a second matrix network with bonds illustrated in blue. \newline
{\bf Right:} The two networks are combined to form the double network on the right, with inter-network connections at all shared sacrificial nodes.}
\label{fig:DoubleJammedNetwork}
\end{figure}

To define the inter-network connections between the sacrificial and matrix networks, we take the existing nodes from the sacrificial component. These nodes are scaled by a factor of $1/M$ and constrained to coincide with their corresponding sacrificial nodes, while the remaining $N_m - N_s$ disks are randomly initialised. The matrix ensemble is then equilibrated, \DIFdelbegin \DIFdel{taking care to ensure a homogeneous distribution of disks despite the inclusions introduced by the constrained sacrificial disks, }\DIFdelend with matrix bonds formed in the same manner as in the sacrificial network. The double network is formed by taking the sacrificial and matrix networks defined in the same domain ($L_x$ and $L_y$ are the same for both network components) and combining them into a single network. Nodes that were fixed during the construction of the matrix network coincide with nodes in the sacrificial network and are therefore shared by both networks. These shared nodes therefore provide the inter-network connections, allowing stress to be shared between the two components. This process is illustrated in \cref{fig:DoubleJammedNetwork}. For the packing-derived networks considered here, the corresponding coordinations for both components are $z_s = z_m \approx 5.4$.

For $M < 2$, the presence of constrained sacrificial disks renders energy minimisation computationally challenging. In such cases, alternative strategies such as swap Monte Carlo may need to be considered \cite{NiklasSwapMonte2023}. All results presented in this thesis correspond to systems with $M \geq 2$, for which energy minimisation remains computationally tractable using the F.I.R.E. algorithm \cite{guenole_assessment_2020}.

\section{Stretching Protocol} \label{sec:NetworkQuasistatic}
Unlike the previous chapter, where we considered systems deformed under simple shear, here we consider the system to be subjected to elongational stretching, with the principal extension axis aligned with the $y$-direction (vertical axis). We deform the square unit cell with initial dimensions $L_{x,0} = L_{y,0}=L_0$ along the y-axis with periodic boundary conditions on both axes, preserving the area $L_0^2$ of the unit cell through a corresponding compression in the $x$-direction. The deformation of the system is measured by the global engineering strain along the $y$-axis, defined as $\epsilon = L_{y}/L_{y,0} - 1$. The strain rate $\dot{\epsilon}$ is taken to be quasi-static, i.e. $\dot{\epsilon} \to 0$, such that the system remains in mechanical equilibrium throughout the deformation.

Numerically, this is implemented by applying incremental deformation steps of size $\Delta \epsilon$, which by default are taken as $\Delta \epsilon = 10^{-3}$, followed by relaxation of the system until the root-mean-square (r.m.s.) force on the nodes falls below $10^{-6}$. When no bond breakages occur and $\Delta \epsilon$ is sufficiently small, this procedure approximates the quasi-static deformation path.

However, taking a step of size $\Delta\epsilon=10^{-3}$ may be large enough to overstep the precise strain $\epsilon$ at which a single bond fails, obscuring the true ordering of failure events and thereby modifying the resulting fracture of the system from the true quasi-static response. To avoid this, we treat each deformation step as a trial step consisting of an affine deformation of size $\Delta\epsilon$, followed by relaxing the system while suppressing any bond breakages. If the relaxed configuration contains no bonds with strains greater than their threshold, the trial step is accepted, and the process continues.

If, instead, the relaxed configuration contains more than one bond that exceeds its failure threshold, the trial identifies that the step was too large and a smaller step needs to be considered. To do this, we define a function
\begin{equation*}
    f(\Delta\epsilon) = \max_{i\in B} \left[ \frac{l_i(\Delta\epsilon)-l_{i,0}}{l_{i,0}} - \lambda_i \right],
\end{equation*}
where $B$ denotes the set of all remaining bonds, $l_i(\Delta\epsilon)$ is the relaxed length of bond $i$ after a trial step of size $\Delta\epsilon$, $l_{i,0}$ its natural length, and $\lambda_i$ its failure threshold. $f(\Delta\epsilon)$ therefore tracks the bond which is closest to failure, which may not always be the same, but by construction will be continuous, but not necessarily differentiable. The sign of $f(\Delta\epsilon)$ therefore determines whether a step is acceptable:

\begin{align*}
    f(\Delta\epsilon) < 0 &\iff \text{no bond exceeds its threshold (trial step is too small)}, \\
    f(\Delta\epsilon) = 0 &\iff \text{at least one bond is exactly at threshold (critical trial step size)}, \\
    f(\Delta\epsilon) > 0 &\iff \text{at least one bond exceeds its threshold (trial step is too large)},
\end{align*}
Thus, identifying the correct step size reduces to finding the smallest $\Delta\epsilon >0$ such that $f(\Delta\epsilon) = 0$, which can be achieved by bisection on the range $(0,\Delta\epsilon_{Max}]$ where $\Delta\epsilon_{Max}=10^{-3}$ is the default trial step size. For all simulations, we converge the window for $\Delta\epsilon$ to within a tolerance of $\pm 10^{-8}$ and take the right-hand boundary to ensure a bond breaks.  

Once a suitable step size is found, the dynamics of the fracture are resolved under overdamped dynamics described in \cref{sec:NetworkDynamics} where the system is assumed not to strain over the course of the relaxation ($\dot{\epsilon}=0$). The resultant equation of motion of a node $i$ is therefore 
\begin{equation}
    \dot{\bm{r}}_i = \sum_{j}\mu_{ij} \frac{l_{ij}-l_{ij,0}}{l_{ij,0}}\frac{\bm{r}_{ij}}{\left\lVert \bm{r}_{ij} \right\rVert}\DIFaddbegin \DIFadd{,
    }\DIFaddend \label{eq:DoubleMotion}
\end{equation}
in which the sum is taken over all neighbouring nodes $j$. The spring modulus $\mu_{ij}$ takes either $\mu_s$ or $\mu_m$, and time units are set by the drag coefficient $\zeta=1$. We integrate \cref{eq:DoubleMotion} with an adaptive \DIFdelbegin \DIFdel{Euler-Heun }\DIFdelend \DIFaddbegin \DIFadd{Euler--Heun }\DIFaddend algorithm (described in \cref{sec:EulerHeun}) with an absolute error tolerance of $10^{-3}$. Bond breakages are evaluated at the start of each time step, and if any bond breaks, the time step $\delta t$ is reset to $10^{-4}$ to account for the discontinuity resulting from the removed bond. The system's dynamics evolve in this manner until the r.m.s. force at the nodes again falls below $10^{-6}$, where another trial step of size $\Delta\epsilon=10^{-3}$ is attempted. In this manner, we have implemented a quasi-static elongational stretching protocol. We show in \cref{sec:DoubleRate} that stretching at a finite rate $\dot{\epsilon} > 0$ recovers this protocol as $\dot{\epsilon} \to 0$.

In all cases, we choose units of stress such that $\mu_s = 1.0$ and the characteristic length of our sacrificial bonds is set by our network structure. The length scale of the matrix bonds is set by $M$, the typical number of matrix-bond lengths per sacrificial-bond length, with larger $M$ corresponding to a denser matrix network relative to the sacrificial network. For notational convenience, we will define $M=0$ to denote a system with only a sacrificial component.

In line with the experimental literature discussed previously, and in particular, the guidelines for making tough double networks found in Ref.~\cite{gong_why_2010}, we will consider cases where the sacrificial network is sparse when compared to the matrix ($M>1.0$), stiffer than the matrix ($\mu_m < 1.0$), and paired with a more extensible matrix ($\lambda_m > 0.5$). The data we present are obtained by averaging over $5$ samples, and for each quantity, we report the mean and standard deviation over the $5$ samples. Each sample corresponds to a separate realisation of a double network with a different initial seed.

\section{Stress-Strain Response}

\begin{figure}[b!]
    \centering
    \includegraphics[width=0.99\linewidth]{chapters/DoubleNetworkFracture/Result Figures/StressStrainCurves.pdf}
    \caption[Stress-strain curves of a packing-derived network.]{Stress-strain curves of a packing-derived network. $\Sigma$ is the $yy$ component of the virial stress tensor defined in \cref{sec:virialStress} and $\epsilon$ is the engineering strain along the $y$ direction defined as $\epsilon=L_y/L_{y,0}-1$. \newline 
    {\bf a)} Single sacrificial network (red solid line) and single matrix network (blue solid line). {\bf b)} Composite double network (black solid line), with the contributions of the component sacrificial network (red dashed line) and matrix network (blue dotted line) shown separately. {\bf Inset:} modulus $G = d\Sigma / d\epsilon$ at small strain.\newline 
    Parameters: $M=3.0, \mu_m=0.2, \lambda_m=3.0$.}
    \label{fig:JammedStressStrain}
\end{figure}

This section explores how differing network length scales affect the fracture and failure behaviour of double networks. In particular, we focus on how stress is redistributed across the networks following a plastic event. The first key result of this work is shown in \cref{fig:JammedStressStrain}, where we compare the stress-strain response of a double network material with its sacrificial and matrix components simulated individually. We illustrate that the simple model described above, under a quasi-static elongational stretching protocol, can capture \DIFdelbegin \DIFdel{the }\DIFdelend qualitative features of the stress-strain response \DIFdelbegin \DIFdel{found in experimental systems }\DIFdelend \DIFaddbegin \DIFadd{observed experimentally, }\DIFaddend such as those shown in Ref.~\cite{gong_why_2010}. \DIFdelbegin \DIFdel{By comparing the stress-strain response of the double network shown in }%DIFDELCMD < \cref{fig:JammedStressStrain}%%%
\DIFdel{b to the single networks }\DIFdelend \DIFaddbegin \DIFadd{For the representative parameter set }\DIFaddend shown in \cref{fig:JammedStressStrain}a, \DIFdelbegin \DIFdel{we illustrate that it is possible within our model to create a double network that retains }\DIFdelend \DIFaddbegin \DIFadd{the double network retains approximately }\DIFaddend the failure strain of the matrix network while maintaining a stiffness comparable to that of the sacrificial network, particularly during the early elastic regime of deformation. \DIFaddbegin \DIFadd{Experimental double networks can exhibit an even stronger synergistic response, with both a greater toughness, corresponding to a larger work to fracture, and a failure strain exceeding that of either constituent network alone. The present model therefore captures the combined stiff and ductile response of experimental double networks, but not the full enhancement in failure strain observed relative to both isolated constituents.
}


\DIFadd{Experimental double networks can exhibit a synergistic response, with both a greater toughness (work to fracture) and a larger failure strain than either constituent network alone \mbox{%DIFAUXCMD
\cite{gong_double-network_2003,gong_why_2010,ducrot_toughening_2014,millereauMechanicsElastomericMolecular2018}}\hskip0pt%DIFAUXCMD
. For the representative parameter set shown in }\cref{fig:JammedStressStrain}\DIFadd{, however, the model gives $\epsilon_s^*<\epsilon_d^*<\epsilon_m^*$, while the initial stiffness of the double network remains comparable to that of the sacrificial network. This example therefore captures part of the combined stiff and ductile response of experimental double networks, but not the full enhancement of failure strain beyond that of both isolated constituents.
}\DIFaddend 

Firstly, we consider the stress-strain response of a single matrix network (blue curve in \cref{fig:JammedStressStrain}a) and separately a single sacrificial network (red curve in \cref{fig:JammedStressStrain}a), where we observe a non-linear elastic stretching response with strain stiffening, often found in the network materials \cite{feng_nonlinear_2016,licup_nonlinear_2016,storm_nonlinear_2005}, followed by abrupt brittle or quasi-brittle failure \cite{dussi_athermal_2020, tauber_stretchy_2022} at its failure strain $\epsilon^*_{s/m}$. This response is to be expected and follows the same principle discussed in \cref{sec:Ultradelayed}. Next, considering the response of the double network (solid black curve in \cref{fig:JammedStressStrain}b), we see a similar feature with an initial non-linear elastic response, again commonly found in elastomers~\cite{millereauMechanicsElastomericMolecular2018, ducrot_toughening_2014}, hydrogels~\cite{gong_why_2010} and macroscopic composites~\cite{kingMacroscaleDoubleNetworks2019, takahashiCreatingStiffTough2018} up until the first sacrificial bond breaks, after which the network enters an extended softening regime before eventual failure at $\epsilon^*_d$. 

In the intermediate regime, after the first sacrificial bond breaks but before macroscopic failure, the double network undergoes an irreversible softening response starting at the first sacrificial failure at a strain $\epsilon \approx \epsilon^*_{s}$. Throughout this interval, the sacrificial component contributions to the overall response plateau, as the sacrificial bonds fail (red dashed line in \cref{fig:JammedStressStrain}b) with these failure events arising diffusely across the system. This is seen experimentally in Ref.~\cite{millereauMechanicsElastomericMolecular2018}, where mechanoluminescence reveals that bond scission in the sacrificial network is initially distributed homogeneously throughout the sample before becoming localised near the onset of necking. The double network \DIFdelbegin \DIFdel{then ultimately fails when the response becomes dominated by }\DIFdelend \DIFaddbegin \DIFadd{ultimately fails once }\DIFaddend the matrix component \DIFdelbegin \DIFdel{through the same macroscopic failureobserved in the single networks (blue dotted line }\DIFdelend \DIFaddbegin \DIFadd{can no longer stabilise the accumulated damage and a localised rupture develops. Although the matrix carries most of the load immediately before failure, this process should not be identified with the failure of the isolated matrix network: coupling between the two components changes the distribution of stress and damage \mbox{%DIFAUXCMD
\cite{ducrot_structure_2015,tominaga_molecular_2007,huang_importance_2007}}\hskip0pt%DIFAUXCMD
. For the representative parameters }\DIFaddend in \cref{fig:JammedStressStrain}\DIFdelbegin \DIFdel{b)}\DIFdelend \DIFaddbegin \DIFadd{, this coupling delays failure relative to the isolated sacrificial network, but does not increase the failure strain beyond that of the isolated matrix network}\DIFaddend .

In the results that follow, we will explore the underlying physics that drives this difference in response between single and double networks. In particular, we consider how stress delocalisation suppresses macroscopic fracture and how $M$ and $\mu$ can be tuned to create an optimal double network.

\begin{figure}[p]
    \centering
    \includegraphics[width=0.99\linewidth]{chapters/DoubleNetworkFracture/Result Figures/StressStrainExtract.pdf}
\caption[Failure strains $\epsilon^*$, stresses $\Sigma^*$, and sacrificial network stiffness contributions $G_s/G_d$ for packing-derived double networks.]{
Failure strains $\epsilon^*$, stresses $\Sigma^*$, and sacrificial network stiffness contributions $G_s/G_d$ for packing-derived double networks.\newline
{\bf Top panels:} The failure strain $\epsilon^*$ of the double network (solid curves) plotted as a function of a) the relative length scale $M$ and b) the matrix bond stiffness $\mu_m$. In each case, results are shown for matrix bond failure thresholds $\lambda_m = 1.0,\; 2.0,\; 3.0$, displayed from lowest to highest curve (circles, triangles, and squares, respectively). For comparison, the failure strains of the corresponding single sacrificial and single matrix networks are indicated by dashed and dotted curves, respectively.\newline
{\bf Middle panels:} Failure stress $\Sigma^*$ of the double network, presented analogously to panels a) and b).\newline
{\bf Bottom panels:} Sacrificial contribution to the double-network modulus, quantified by $G_s/G_d$.\newline
{\bf Panels a), c), and e):} Fixed matrix bond stiffness $\mu_m = 0.2$.\newline
{\bf Panels b), d), and f):} Fixed relative length scale $M = 3.0$.
}
    \label{fig:JammedStrainExtract}
\end{figure}
\pagebreak

To quantify how effectively a double network characterised by the parameters $M$, $\mu_m$, and $\lambda_m$ resists failure, we characterise its mechanical response by extracting the failure strain $\epsilon^*_d$, the failure stress $\Sigma^*_d$ (defined in \cref{fig:JammedStressStrain}), and the initial moduli $G_s$, $G_m$, and $G_d$ of the sacrificial network, matrix network, and combined double network, respectively. These moduli are defined via the initial stress-strain slopes in \cref{fig:JammedStressStrain}. Correspondingly, for each component of every double network, we realise the associated single network and quantify its mechanical response in the same way, obtaining either $\epsilon_s$ and $\Sigma_s$ or $\epsilon_m$ and $\Sigma_m$ depending on the origin of the single network.

\begin{figure}[t!]
    \centering
    \begin{subfigure}{0.485\linewidth}
        \centering
        \includegraphics[width=\linewidth]{chapters/DoubleNetworkFracture/Result Figures/Fail.pdf}
    \end{subfigure}
    ~
    \begin{subfigure}{0.485\linewidth}
        \centering
        \includegraphics[width=\linewidth]{chapters/DoubleNetworkFracture/Result Figures/StiffnessRatio.pdf}
    \end{subfigure}
    \caption[Phase diagrams of double-network failure strain $\epsilon^*_d$ and relative sacrificial contribution $G_s/G_d$ for packing-derived networks]{Phase diagrams of the double network showing a) the failure strain $\epsilon^*_d$ and b) the relative sacrificial contribution $G_s/G_d$, as functions of $M$ and $\mu_m$.}
    \label{fig:JammedPhaseDiagrams}
\end{figure}

\DIFdelbegin \DIFdel{Based on the experimental guidelines outlined in Ref.~\mbox{%DIFAUXCMD
\cite{gong_why_2010}}\hskip0pt%DIFAUXCMD
, a well-designed tough double network has two key requirements }\DIFdelend \DIFaddbegin \DIFadd{For the purposes of the present model, we characterise a successful double network using two requirements motivated by the design principles discussed in Ref.~\mbox{%DIFAUXCMD
\cite{gong_why_2010}}\hskip0pt%DIFAUXCMD
}\DIFaddend :
\begin{enumerate}[label=\Roman*)]
\item It inherits the large failure strain of \DIFaddbegin \DIFadd{its }\DIFaddend soft and ductile matrix \DIFdelbegin \DIFdel{components}\DIFdelend \DIFaddbegin \DIFadd{component}\DIFaddend , such that the overall failure strain satisfies \DIFdelbegin \DIFdel{$\epsilon^*_D \approx \epsilon^*_m \gg \epsilon^*_s$}\DIFdelend \DIFaddbegin \DIFadd{$\epsilon^*_d \gtrsim \epsilon^*_m \gg \epsilon^*_s$}\DIFaddend .
\item It \DIFdelbegin \DIFdel{inherits }\DIFdelend \DIFaddbegin \DIFadd{retains }\DIFaddend the stiffness of its stiff but brittle sacrificial component, with modulus $G_d \approx G_s \gg G_m$.
\end{enumerate}
Together, these \DIFdelbegin \DIFdel{requirements }\DIFdelend \DIFaddbegin \DIFadd{provide }\DIFaddend serve as the criteria \DIFdelbegin \DIFdel{by which we assess the success of the double networks across parameter space.
}\DIFdelend \DIFaddbegin \DIFadd{for assessing the balance between stiffness and extensibility within the model considered here. They should not be interpreted as a complete definition of an optimal experimental double network. In particular, experimental double networks can exhibit failure strains and fracture energies significantly greater than those of either constituent single network \mbox{%DIFAUXCMD
\cite{gong_why_2010,ahmedBrittleDuctileTransition2014}}\hskip0pt%DIFAUXCMD
. Fully capturing this enhancement beyond the response of the isolated matrix would require resolving fracture across a broad range of length and time scales, from local bond breaking to the development of an extended damage zone and the propagation of a macroscopic crack. The present simulations are not able to access all of these scales simultaneously and should therefore be viewed as addressing only part of this multiscale fracture process.
}

\DIFaddend We will address \DIFdelbegin \DIFdel{these two requirements }\DIFdelend \DIFaddbegin \DIFadd{the two model criteria }\DIFaddend in turn. We first focus on property I), examining how \DIFaddbegin \DIFadd{effectively the matrix prevents premature failure of }\DIFaddend the \DIFdelbegin \DIFdel{failure strain of the double network depends on the choice of model parameters }\DIFdelend \DIFaddbegin \DIFadd{sacrificial component and allows the double-network failure strain to approach that of the isolated matrix as }\DIFaddend $M$, $\mu_m$ and $\lambda_m$ \DIFaddbegin \DIFadd{are varied}\DIFaddend .

We illustrate property I) in \cref{fig:JammedStrainExtract}a-b and \cref{fig:JammedPhaseDiagrams}a where we show how the failure strain $\epsilon^*$ changes as a function of our control parameters $M$, $\mu_m$ and $\lambda_m$. In \cref{fig:JammedPhaseDiagrams}a, we plot a colour map of $\epsilon^*$ for the plane of $M$ and $\mu_m$ for fixed $\lambda_m=3.0$. This shows an increasing failure strain as either $M$ or $\mu_m$ increases. This trend is shown more clearly in the upper solid line of \cref{fig:JammedStrainExtract}a, which represents a horizontal slice of \cref{fig:JammedPhaseDiagrams}a at $\mu_m=0.2$. For all values of $\lambda_m$ in \cref{fig:JammedStrainExtract}a we observe a consistent trend: as $M \to \infty$, the network's failure strain transitions from that of the single sacrificial network, with $\epsilon^*_d \approx \epsilon^*_s$, towards the failure strain of the matrix, with $\epsilon_d^* \to \epsilon_m^*$. We see a less pronounced upwards trend for $\lambda_m=1.0$, where the matrix and sacrificial bonds have similar breakage thresholds. In this case, early failure events in the sacrificial networks propagate enough stress to result in premature failure events in the matrix. In contrast, the curves corresponding to $\lambda_m=2.0$ or $3.0$ show that allowing enough separation in the failure of the two bond types ensures that failure requires a significant weakening in the sacrificial network. We will return to this behaviour in more detail when we discuss the role of stress propagation.

Property II) is considered in \cref{fig:JammedPhaseDiagrams}b, where we show the evolution of the relative contribution of the sacrificial component to the double network's stiffness, quantified by the ratio $G_s/G_d$. As either $M$ or $\mu_m$ increases, we see a decrease in $G_s/G_d$, indicating an increasing contribution of the matrix to the overall stiffness $G_d$. This behaviour is expected and indicates the trivial limit at $M \to \infty$ or $\mu_m \to \infty$ in which the double network effectively just becomes a single matrix network thereby inheriting its softness. While this trivial limit does therefore formally satisfy property I), it does so at the expense of losing any contribution from the sacrificial component and is therefore unable to satisfy property II).

This trend is also shown in \cref{fig:JammedStrainExtract}e-f where the decrease in $G_s/G_d$ is clearly apparent. We do not plot the values relative to $\lambda_m$ since the reported moduli are taken as the initial slope of the stress-strain response, where no bonds have failed. As a result, this is independent of $\lambda_m$, $\lambda_s$, and we are unable to report the single network moduli in this context. 

Considering the trends shown in \cref{fig:JammedPhaseDiagrams} and \cref{fig:JammedStrainExtract}, it becomes clear that achieving the desired \DIFdelbegin \DIFdel{optimal }\DIFdelend double network properties I) and II) requires a trade-off. Optimising property I), for a given $\lambda_m$, requires increasing $M$ and $\mu_m$ until the double network material fails \DIFaddbegin \DIFadd{increasingly }\DIFaddend like the single matrix network\DIFdelbegin \DIFdel{(}\DIFdelend \DIFaddbegin \DIFadd{, }\DIFaddend whose failure point \DIFdelbegin \DIFdel{$\epsilon^*_m$ only depends }\DIFdelend \DIFaddbegin \DIFadd{$\epsilon_m^*$ depends only }\DIFaddend on $\lambda_m$\DIFdelbegin \DIFdel{). At this point}\DIFdelend \DIFaddbegin \DIFadd{. In this limit, however}\DIFaddend , the sacrificial component contributes very little \DIFaddbegin \DIFadd{to the mechanical response}\DIFaddend , and the double network offers \DIFdelbegin \DIFdel{no advantage over a }\DIFdelend \DIFaddbegin \DIFadd{little advantage over the }\DIFaddend single matrix network \DIFdelbegin \DIFdel{, whose failure strain $\epsilon_m^*$ is set solely by $\lambda_m$}\DIFdelend \DIFaddbegin \DIFadd{in terms of its initial stiffness}\DIFaddend .

Conversely, decreasing $M$ or $\mu_m$ enhances the contribution of the sacrificial component to the overall stiffness but drives the failure strain towards that of the sacrificial network, as indicated by the red dashed lines in panels (a) and (b). \DIFdelbegin \DIFdel{An optimal }\DIFdelend \DIFaddbegin \DIFadd{A suitable }\DIFaddend choice of parameters \DIFdelbegin \DIFdel{, therefore , }\DIFdelend \DIFaddbegin \DIFadd{therefore }\DIFaddend requires a compromise: for a given $\lambda_m$ sufficiently large to avoid \DIFdelbegin \DIFdel{an }\DIFdelend immediate failure of the matrix, $M$ and $\mu_m$ should be chosen such that $\epsilon_d^* \approx \epsilon_m^*$ (Property I)\DIFdelbegin \DIFdel{. At the same time, the stiffness must remain }\DIFdelend \DIFaddbegin \DIFadd{, while the stiffness remains }\DIFaddend dominated by the sacrificial network (Property II). In practice, we find that the region of parameter space exemplified by \DIFdelbegin \DIFdel{$M = 3.0$, $\mu_m = 0.2$, and $\lambda_m = 3.0$}\DIFdelend \DIFaddbegin \DIFadd{$M=3.0$, $\mu_m=0.2$, and $\lambda_m=3.0$}\DIFaddend , shown in \cref{fig:JammedStressStrain} \DIFdelbegin \DIFdel{, }\DIFdelend provides a sensible compromise \DIFaddbegin \DIFadd{between these competing requirements.
}

\DIFadd{This should not, however, be interpreted as identifying the full optimum observed experimentally. For the representative parameters considered here, $\epsilon_s^<\epsilon_d^<\epsilon_m^*$, whereas experimental double networks can exhibit failure strains and fracture resistance exceeding those of the corresponding matrix-only material. It instead illustrates the internal compromises required within one part of the overall fracture mechanism to achieve the remarkable combination of stiffness and extensibility exhibited by these materials}\DIFaddend .

\section{Stress Propagation} \label{sec:doubleStressProp}

To understand these trends, and in particular the enhanced toughness of double networks, we first consider the effect of a single bond breaking in a single network material. When such a network is deformed along the $y$ axis until a single sacrificial bond fails, surrounding nodes undergo non-affine rearrangements as the system relaxes to a new minimum-energy state.
This relaxation redistributes some of the stress previously carried by the broken bond, increasing the stress on bonds aligned along the $x$ axis and decreasing it on those aligned along the $y$ axis, as shown in the left column of \cref{fig:JammedPropogators}.

\begin{figure}[p]
    \centering
    %\missingfigure{Propogator M =0, 2, 5 + contour}
    \includegraphics[width=0.99\linewidth]{chapters/DoubleNetworkFracture/Result Figures/Propogators.pdf}
    \caption[Heat map of Stress propagation in Double Network]{Stress propagation to neighbouring sacrificial bonds from the breakage of a sacrificial bond at the origin. We display each quantity for the case of a single sacrificial network in the left column and a double network with $M=5$ in the right column. In the double network cases, we take $\mu_m = 0.2$. Each quantity is averaged over all sacrificial bonds at each grid square. \newline
    Colour scales show the change due to stress propagation of the: \newline
    {\bf Top panels:} $yy$ component of the virial stress $\Sigma_{yy}$. \newline
    %{\bf Middle panels:} $xx$ component of the virial stress $\Sigma_{xx}$. \newline
    {\bf Bottom panels:} sacrificial bond stress $\delta\sigma_{sb} = \mu_s\left(l/l_0 -1 \right)$.
    }
    \label{fig:JammedPropogators}
\end{figure}

Unless otherwise stopped, this stress redistribution makes neighbouring bonds more susceptible to failure, either immediately or under continued deformation. In either case, these subsequent breakages may result in further rearrangements and increased stress on neighbouring bonds. If unchecked, this process develops into an avalanche of plastic breakages that propagates along the $x$ axis, ultimately leading to catastrophic material failure, as illustrated in \cref{fig:JammedSacrifical}. When a sacrificial bond breaks, stress is redistributed into the matrix; for sufficiently large $\lambda_m$, the matrix remains intact, lowering the stress in the sacrificial network and suppressing further bond failure.

\begin{figure}[b!]
    \centering
    \begin{subfigure}[t]{0.485\textwidth}
         \centering
         \includegraphics[width=\textwidth]{chapters/DoubleNetworkFracture/Result Figures/Sacrifical.pdf}
         \caption{Initiation of a macroscopic crack in a stretched single sacrificial network, which will result in a catastrophic material failure.}
         \label{fig:JammedSacrifical}
     \end{subfigure}
     ~
     \begin{subfigure}[t]{0.485\textwidth}
        \centering
        \includegraphics[width=\textwidth]{chapters/DoubleNetworkFracture/Result Figures/Double.pdf}
        \caption{Stretched double network at large strain, showing diffuse microcracking in the sacrificial network (red) that creates small, spatially localised voids arising from the failure of a limited number of bonds spread diffusely through the material.}
        \label{fig:JammedDouble}
    \end{subfigure}
    \caption[Portions of stretched sacrificial and double network whose stress-strain response is shown in \cref{fig:JammedStressStrain}.]{Portions of stretched sacrificial and double networks whose stress-strain response is shown in \cref{fig:JammedStressStrain}. The point on the stress-strain curve for the sacrificial and double networks is illustrated by a triangle and a circle, respectively, in \cref{fig:JammedStressStrain}.}
    \label{fig:JammedImages}
\end{figure}

Optimally, to satisfy properties I) and II) discussed earlier, the matrix must be sufficiently stiff and tough to absorb enough stress to arrest the chain of sacrificial bond failures, while remaining compliant enough that it does not dominate the initial modulus $G_d$. This reduction in local stresses on neighbouring sacrificial bonds allows for other sacrificial bonds elsewhere in the material to fail instead, promoting more diffuse damage in the sacrificial component. 

To quantify the reduction in the stress propagated to the neighbouring bonds caused by the matrix, we plot the average change in stress resulting from the failure of the first sacrificial bond during deformation in \cref{fig:JammedPropogators}. To construct a detailed enough mapping, we plot binned averages taken over an ensemble of $100$ network realisations. Each system is strained as described in \cref{sec:NetworkQuasistatic} until the failure of the first sacrificial bond. To capture the resultant propagation of stress after this failure event, we now deviate from the method outlined in \cref{sec:NetworkQuasistatic}. The dynamics are resolved using the same overdamped equation; however, following the failure of the first sacrificial bond, we disallow any further bond failures, and the system is then allowed to relax to a local energy minimum. Each network is then rotated such that the two nodes of the broken bonds lie vertically aligned along the $y$ axis. Bonds are then assigned a bin based on the position of their centre relative to the centre of the broken bond, which is defined to be at $(0,0)$, and the reported quantities are taken as the average value over all bonds in each bin. 

The change in the $yy$ component of the virial stress tensor (defined in \cref{sec:virialStress}) is plotted in \cref{fig:JammedPropogators}a-b for the case of single sacrificial networks and double networks with $M=5$. The plot is consistent with the quadrupolar structure characteristic of the Eshelby stress propagator discussed in \cref{sec:EshelbyPropogator}.
% We note that $\delta \Sigma_{yy}$ also exhibits a strongly negative stress relaxation region along the vertical axis, not typically present in the Eshelby propagator. We hypothesise that this feature arises from the non-isotropic nature of our networks at a strain $\epsilon \gg 0$. 

\begin{figure}[b!]
    \centering
    \includegraphics[width=0.99\linewidth]{chapters/DoubleNetworkFracture/Result Figures/Lines.pdf}
    \caption[Stress propagation in Double networks]{
    Stress propagation to neighbouring sacrificial bonds from the breakage of a sacrificial bond at the origin. We show the change in bond stress as a function of distance along $x$ at $y=0$ for (left panel) several $M$ at fixed $\mu_m = 0.2$, and (right panel) several $\mu_m$ at fixed $M=3.0$. Lengths are expressed in units such that $\Delta x=1$ corresponds to the typical sacrificial bond length. The propagator of a single sacrificial network is recovered for $M=0$ or $\mu_m = 0$.
    }
    \label{fig:JammedPropogatorsLines}
\end{figure}

% Correspondingly, panels c)-d) plot the $\delta \Sigma_{xx}$ and show the loading of neighbouring bonds along the $x$ axis, as discussed previously as a driving cause of macroscopic failure. 
% This is more directly emphasised by panels e)-f) which directly plots the change in scalar stress $\delta \sigma_{sb}$ in each sacrificial bond along its own axis. 
Correspondingly, panels c)-d) plot the change in scalar stress $\delta \sigma_{sb} = \mu_s\left(l/l_0 -1 \right)$ in each sacrificial bond along its own axis. Due to our choice in units, with $\mu_s=1.0$, $\delta \sigma_{sb}$ can be interpreted directly as the change in strain along each bond. Although this quantity alone does not determine if a bond fails, because the failure criterion also depends on the initial strain of each bond, it is nevertheless a reliable indicator of the underlying physics being driven by this reduction in stress propagation. 

By plotting slices of $\delta \sigma_{sb}$ along the $x$ axis in \cref{fig:JammedPropogatorsLines} as a function of a) $M$ at fixed $\mu_m$ and b) $\mu_m$ at fixed $M$, we observe a decrease in the amplitude of stress propagation. In each curve, the propagator decays as $x^{-2}$, consistent with the long-range propagator in \cref{eq:eshelbyReal}, and consistent with a linear elastic response. With increasing $M$ or $\mu_m$, the magnitude of $\delta \sigma_{sb}$ decreases as the stress released by a single bond failure is redistributed into an increasingly stiff matrix component, either through a larger matrix stiffness $\mu_m$ or a greater number of matrix bonds sharing the load. As a result, each bond experiences a smaller stress perturbation, and the amplitude of the stress propagator is reduced.

In each inset, we extract the value $\delta \sigma_{sb}$ at $x = 1$ and plot its decay as a function of $M$ and $\mu_m$ respectively, suggesting an approximate power-law dependence for $M \gg 0$ or $\mu_m \gg 0$. We note that a different dependence must emerge as $M/\mu_m \to 0$, where the response must recover the finite value associated with the single sacrificial network.
% This is consistent with the plateaus shown earlier, which show convergence to matrix-dominated behaviour at large $M$ or $\mu_m$.

\section{Damage in the Sacrificial component}

If our understanding is correct, the underlying physics, driven by a reduction in the Eshelby stress propagator, allows for diffuse rather than concentrated damage in the sacrificial component of the double network. This behaviour should be observable by examining the fraction of bonds broken in each component. Importantly, catastrophic failure can occur at any value of $B$, the fraction of bonds broken in the network. We begin again by considering the specific set of network parameters for which the stress-strain curves were shown in \cref{fig:JammedStressStrain}. In \cref{fig:JammedBrokenStrain}, we now plot the fraction of bonds broken $B$ for the single matrix and single sacrificial networks in panel a), and for the double network and its individual components in panel b). 

\begin{figure}[t!]
    \centering
    \includegraphics[width=0.99\linewidth]{chapters/DoubleNetworkFracture/Result Figures/BrokenStrainCurves.pdf}
    \caption[Breakage-strain curves of a double packing-derived network.]{The fraction of broken bonds $B$ is shown as a function of increasing strain for the networks whose stress–strain responses are presented in \cref{fig:JammedStressStrain}. {\bf a)} Single sacrificial network (red solid line) and single matrix network (blue solid line). {\bf b)} Double network (black solid line), sacrificial component (red dashed line), and matrix component (blue dotted line).}
    \label{fig:JammedBrokenStrain}
\end{figure}

First, considering the single matrix network (solid blue line), \cref{fig:JammedBrokenStrain}a shows a simple, highly brittle response with an abrupt increase in bond failures at the failure strain $\epsilon_m^*$. A similar but slightly more ductile response is observed in the single sacrificial network (solid red line), where a small number of bonds fail before the network's ultimate brittle failure, illustrated in \cref{fig:JammedSacrifical}. This behaviour is consistent with previous observations reported in Ref.~\cite{dussi_athermal_2020}.

In contrast, the double network curve (solid black line) in \cref{fig:JammedBrokenStrain}b demonstrates the markedly more ductile response of the double network. This behaviour has been discussed previously and is illustrated in \cref{fig:JammedDouble}, which shows the state of the strained network at a deformation well above $\epsilon^*_s$. At small strains $\epsilon \approx \epsilon_s^*$, the number of breakages in the sacrificial component (red dashed line) tracks closely with the single network curves, indicating that the presence of the matrix does not significantly alter the onset of failure in the sacrificial component. 

\begin{figure}[b!]
    \centering
    \includegraphics[width=0.9\linewidth]{chapters/DoubleNetworkFracture/Result Figures/BrokenStrainExtract.pdf}
    \caption[Fraction of bonds broken $B^*$ at failure.]{
    {\bf Top panels:} The fraction of broken bonds $B^*$ of the double network (solid curves) is plotted as a function of a) the matrix content $M$ and b) the matrix bond stiffness $\mu_m$. In each case, we show results for matrix bond failure thresholds $\lambda_m = 1.0,\; 2.0,\; 3.0$, displayed from the lowest to highest curve (circles, triangles and squares, respectively). For comparison, the fraction of broken bonds of the corresponding single sacrificial network and single matrix network is indicated by dashed and dotted curves, respectively.\newline
    {\bf Bottom panels:} For the solid lines shown in the corresponding plot above, we show the fraction of sacrificial and matrix bonds broken, indicated by dashed and dotted curves, respectively.
    }
    \label{fig:JammedBrokenExtract}
\end{figure}

However, as shown earlier, the sacrificial component ultimately survives to higher strains in a double network due to the mechanical support provided by the matrix. Correspondingly, the curve extends and begins to plateau at higher strains, reaching $B \approx 0.41$. At the large but still subcritical strain illustrated in \cref{fig:JammedDouble}, most of the sacrificial component remains intact and continues to percolate vertically despite diffuse bond breakage, and can therefore still contribute to the stiffness of the double network. At larger strains, the double network eventually fails in the same manner as the single matrix network, through an abrupt and brittle failure of the matrix component (blue dotted line).

% At the large but subcritical strain illustrated in \cref{fig:JammedDouble}, most of the sacrificial component remains intact and continues to percolate vertically, despite diffuse bond breakage. At larger strains, the double network eventually fails via abrupt failure of the matrix component.

To see how this effect depends on $M$ and $\mu_m$, we plot the fraction of broken bonds at the point of failure $B^*$ in \cref{fig:JammedBrokenExtract}. Importantly, the value of $B^*$ does not indicate whether the network failed catastrophically; however, a large $B^*$ indicates that the material may have accrued damage diffusely before failing.

The top panels a) and b) are equivalent to \cref{fig:JammedStrainExtract}a-d and show $B^*$ for the double network (solid curves) as a function of a) $M$ and b) $\mu_m$. For comparison, the corresponding single sacrificial and single matrix networks are shown by red dashed and blue dotted curves, respectively. The bottom panels c) and d) show, for the same parameters, the fraction of broken sacrificial and matrix bonds in the double network, indicated by dashed and dotted curves, respectively. In each panel, we again plot the response for different values of the matrix bond failure threshold, $\lambda_m=1.0$, $2.0$, and $3.0$, from bottom to top. 

The dependence on $M$ shown in \cref{fig:JammedBrokenExtract}a for the double network is a consequence of the construction of our networks: as $M$ increases, the fraction of sacrificial bonds relative to the total number of bonds decreases. This dilution causes the total fraction of bonds $ B^* \to 0.0$ as $M\to\infty$. The non-monotonic behaviour at $M=2$, however, is a consequence of the sacrificial component itself. As shown in \cref{fig:JammedBrokenExtract}c, the fraction of broken sacrificial bonds at failure increases, indicating increasingly diffuse damage within the sacrificial network.

Panels \cref{fig:JammedBrokenExtract}b and d show the corresponding dependence on $\mu_m$. In panel b, the total fraction of broken bonds in the double network increases monotonically, for $\lambda_m = 2.0, \; 3.0$. This again suggests a trend towards stress delocalisation, reflected in the increased extent of sacrificial bond failure that can occur before rupture of the matrix network. For $\lambda_m =1.0$ we observe a small decrease in the $B^*$ prior to the expected monotonic increase, together with an increase in the standard deviation in $B^*$ at small $\mu_m$. At low $\mu_m$ the matrix bonds are too weak to resist the deformation caused by the breaking of the sacrificial bonds, which leads to failure in the matrix network. This increases the total number of broken bonds but does not change the physics of the sacrificial network because little stress is propagated when the weak matrix bonds fail. Such behaviour is observed only when $\lambda_m \sim \lambda_s$, where the failure thresholds of the two networks are comparable.
%For $\lambda_m =1.0$ we see a small decrease in the $B^*$ before the expected monotonic increase. This does correspond with an increase in the standard deviation of our results over samples and may correspond to slight under-convergence in our athermal quasi-static algorithm (convergence was checked for major reported quantities $\epsilon^*$, $\Sigma^*$, and moduli).
% \todo{Need to put this in an appendix}

In each plot, we show results for $\lambda_m=1.0$, $2.0$, and $3.0$, which converge to different values of $B^*$ under both increasing $M \to \infty$ and $\mu_m \to 1$. This is consistent with the work done by Tauber \etal \cite{tauber_microscopic_2020} in a simple multi-spring model with $M=1$. While a theoretical extension of this model to include $M > 1$ would provide a useful analytical framework for interpreting these trends, constructing such a model lies beyond the scope of this thesis. 

\section{Conclusion}
In this chapter, we studied the increased toughness of double network materials in a simple mesoscopic model. Our central aim was to understand how stress delocalisation in the sacrificial component of a double network allows for the double network to inherit both the stiffness of its stiff and brittle sacrificial component and the ductility of its soft and ductile matrix component. We \DIFdelbegin \DIFdel{show }\DIFdelend \DIFaddbegin \DIFadd{have shown }\DIFaddend that this simple model reproduces several key qualitative features of double network materials observed experimentally.

\DIFaddbegin \DIFadd{The mechanism identified here also clarifies the sense in which the double network outperforms both single networks within the present model. The matrix-only network is highly extensible, but its low stiffness means that it carries comparatively little stress throughout deformation. By contrast, the sacrificial-only network is initially much stiffer, but fails at a substantially lower strain. The double network combines these two advantages, retaining much of the stiffness provided by the sacrificial network while the matrix maintains connectivity as sacrificial bonds progressively fail. Consequently, the double network can sustain comparatively large stresses over strains approaching those accessible to the matrix-only network, increasing the area under the stress--strain curve and hence the mechanical energy absorbed before failure.
}

\DIFadd{There is, however, an important limitation to this comparison. In the present model, the failure strain of the double network remains approximately bounded by that of the single matrix network, and we therefore do not reproduce the additional enhancement in failure strain that can occur in experimental double network materials. The increase over the matrix-only network obtained here instead arises primarily from combining the large extensibility of the matrix with the greater stiffness and distributed damage supplied by the sacrificial network. Thus, within the scope of this model, stress delocalisation explains how sacrificial-bond damage can occur without premature brittle failure, allowing the material to retain an elevated stress over a matrix-like strain range and thereby increase the work required to reach failure.
}

\DIFaddend This framework has allowed us to systematically explore how stiffness and network length scales control the mechanical response of double network materials. Our results align with previous experimental findings: the enhanced ductility of double networks can be understood in terms of a reduction in the Eshelby stress propagator between neighbouring sacrificial bonds. This reduction suppresses the cascade of bond failures that would otherwise drive brittle fracture in a single sacrificial network, instead promoting diffuse damage and delaying macroscopic failure.

The model presented in this work should be viewed as an idealised representation of real double network hydrogel systems. Our analysis has been restricted to two-dimensional networks governed by purely central-force interactions, with no energetic cost associated with bond bending. \DIFdelbegin \DIFdel{Although we do not anticipate that the introduction of bending interactions would qualitatively alter the key behaviours identified here, prior }\DIFdelend \DIFaddbegin \DIFadd{In this work, we have considered networks above the isostatic point, with an average network coordination $z \geq 4=2d$. By contrast, real gels inhabit three dimensions and can have connectivities below the corresponding central-force isostatic threshold $z_c=2d=6$. Extending the present model to lower connectivities or to three dimensions would therefore require some care, since a central-force network with $z<z_c$ is mechanically floppy. In this regime, additional stabilising mechanisms, including bond-bending interactions or pre-existing stress, become important.
}

\DIFadd{Previous }\DIFaddend studies have demonstrated that \DIFdelbegin \DIFdel{a }\DIFdelend finite bending rigidity can act as a stabilising field, particularly in networks with low connectivity $z$ \cite{shivers_strain-controlled_2023}. \DIFdelbegin \DIFdel{Such }\DIFdelend \DIFaddbegin \DIFadd{Although we do not anticipate that the introduction of bending interactions would qualitatively remove the load-sharing mechanism identified here, such }\DIFaddend stabilisation may enhance the ability of the sacrificial network to sustain diffuse damage \DIFdelbegin \DIFdel{, thereby potentially further increasing the toughness of double network materials. }\DIFdelend \DIFaddbegin \DIFadd{before macroscopic failure.
}\DIFaddend 

\DIFdelbegin \DIFdel{In this work, we have considered networks above the isostatic point with an average network coordination $z \geq 4=2d$. By contrast, real gels inhabit three dimensions, and typically have connectivity below the critical point $z_c=2d=6$. In experimental }\DIFdelend \DIFaddbegin \DIFadd{We nevertheless expect the basic load-sharing mechanism identified in the present two-dimensional model to remain relevant in three dimensions. Following a sacrificial-bond failure, coupling to the matrix would still provide additional pathways through which the released load can be redistributed, thereby reducing the perturbation experienced by neighbouring sacrificial bonds. Quantitative differences should, however, be expected. The spatial structure of the elastic propagator changes with dimensionality, a crack becomes a two-dimensional surface rather than a line, and the mechanical response may additionally depend on bending rigidity and prestress of the fibres. These differences may alter the extent and morphology of diffuse damage and shift the values of $M$ and $\mu_m$ at which the different failure regimes occur.
}

\DIFadd{Experimental }\DIFaddend tough double network hydrogels \DIFdelbegin \DIFdel{, the stiff sacrificial component}\DIFdelend \DIFaddbegin \DIFadd{also contain structural differences between the two constituent networks that are represented only indirectly in the present model. The stiff sacrificial network}\DIFaddend , which is synthesised first, \DIFaddbegin \DIFadd{typically }\DIFaddend has a higher \DIFdelbegin \DIFdel{density of cross-linkers when compared to }\DIFdelend \DIFaddbegin \DIFadd{cross-link density than }\DIFaddend the soft, weakly cross-linked matrix network. \DIFdelbegin \DIFdel{Moreover, during the }\DIFdelend \DIFaddbegin \DIFadd{During }\DIFaddend synthesis, the sacrificial network is \DIFdelbegin \DIFdel{often swollen to accommodate the }\DIFdelend \DIFaddbegin \DIFadd{also often swollen prior to }\DIFaddend formation of the second network, \DIFdelbegin \DIFdel{resulting in a naturally prestretched sacrificial network}\DIFdelend \DIFaddbegin \DIFadd{leaving it in a prestressed state}\DIFaddend . In our \DIFdelbegin \DIFdel{current model, we have used }\DIFdelend \DIFaddbegin \DIFadd{model, each network is instead characterised by }\DIFaddend a single stiffness parameter $\mu$\DIFdelbegin \DIFdel{for each network, and taken $\mu_s > \mu_m$ as a proxy for both increased }\DIFdelend \DIFaddbegin \DIFadd{, with $\mu_s>\mu_m$ used as a simplified representation of these differences in }\DIFaddend cross-link density and \DIFdelbegin \DIFdel{prestretch in the sacrificial network. This simplified approach has enabled us }\DIFdelend \DIFaddbegin \DIFadd{pre existing stresses. This approach allows the model }\DIFaddend to capture the \DIFdelbegin \DIFdel{relevant physics at play while retaining only }\DIFdelend \DIFaddbegin \DIFadd{resulting contrast between the stiff sacrificial and soft matrix networks while retaining }\DIFaddend a minimal number of ingredients\DIFdelbegin \DIFdel{in our model }\DIFdelend \DIFaddbegin \DIFadd{.
}

\DIFadd{A further simplification is that the elasticity of the present model is represented through deterministic central-force interactions, whereas the elasticity of real polymer gels is predominantly entropic in origin. In a polymer network, the force carried by a strand depends on its molecular conformation and becomes strongly non-linear as the strand approaches failure strain. At moderate deformation, this behaviour may be represented approximately through an effective elastic modulus, but this approximation becomes less complete at the very large strains accessible to double network hydrogels.
}

\DIFadd{Including entropic elasticity would therefore be expected to modify the quantitative redistribution of load between the two components. In particular, strongly stretched matrix strands would stiffen as deformation increases and could consequently carry a greater fraction of the load released by sacrificial-bond failure. This may further stabilise damaged regions, although it could also concentrate the stress onto a smaller population of highly extended matrix bonds and thereby influence the eventual failure of the matrix component}\DIFaddend .

\begin{figure}[b!]
    \centering
    \includegraphics[width=0.9\linewidth]{chapters/DoubleNetworkFracture/Result Figures/StressStrainCurves_lowZ.pdf}
    \caption[Stress-strain curves of a packing-derived network at low $z_m$.]{Stress-strain curves of a packing-derived network at a low matrix connectivity $z_m=3.5$. $\Sigma$ is the $yy$ component of the virial stress tensor defined in \cref{sec:virialStress} and $\epsilon$ is the engineering strain along the $y$ direction defined as $\epsilon=L_y/L_{y,0}-1$. \newline 
    {\bf a)} Single sacrificial network (red solid line) and single matrix network (blue solid line). {\bf b)} Composite double network (black solid line), with the contributions of the component sacrificial network (red dashed line) and matrix network (blue dotted line) shown separately. {\bf Inset:} modulus $G = d\Sigma / d\epsilon$ at small strain.\newline 
    Parameters: $M=3.0, \mu_m=0.2, \lambda_m=3.0$, $z_s=5.42$, $z_m=3.5$.}
    \label{fig:LowZStressStrain}
\end{figure}

In \cref{sec:SquareAppendix} we briefly explore the role of prestretching bonds as a stabilising field in square lattice-based networks. We also explore the role of subcritical networks in \cref{fig:LowZStressStrain}. Here, the matrix network is pruned such that initial connectivity is reduced to $z_m=3.5 < z_c$. The sacrificial component retains its high connectivity $z_s=5.4$, and keeps all other parameters the same as in \cref{fig:JammedStressStrain}. The double network remains significantly tougher than the single sacrificial network; however, its stiffness $G_d$ is now almost entirely dominated by the sacrificial component, consistent with experimental observations.

In experimental systems, the sacrificial network is often synthesised first, resulting in it being naturally prestressed in the double network's natural condition. In our current construction, we do not consider this, apart from the constraint that $\mu_s > \mu_m$. Incorporating controlled pre-stress into the sacrificial network would be straightforward within our framework, following the approach outlined for square double networks in \cref{sec:SquareAppendix}. Further adding disorder in this pre-stress may allow for more complex behaviour as the components transition from hyperstatic to sub-isostatic during fracture. These two modifications to our network geometries could be further extended by alternative potentials such as a rope-like potential \cite{arzashStressstabilizedSubisostaticFiber2019a} or generating more physically accurate network geometries using more complex molecular dynamics simulations \cite{mugnaiInterspeciesInteractionsDual2025,rijns_synthetic_2024}.  

In our framework, bond failure is treated as irreversible, with no formation of new bonds during deformation. However, a growing body of experimental work on double network hydrogels has demonstrated that incorporating self-healing, through the reformation of bonds, provides an additional route to toughness and durability \cite{wangRapidSelfstrengtheningDoublenetwork2025,yasuiRateIndependentSelfHealingDouble2022,yinSelfHealingHydrogelsSynthesis2023}. This has been particularly studied in the context of cyclic shear loading, where experiments \cite{zhang_fatigue_2018, chen_improvement_2016, chen_novel_2015} have shown that double-network hydrogels can endure many loading cycles with little loss of stiffness. 

Although the monotonic strain-startup simulations studied here do not address cyclic loading directly, the framework may provide a useful starting point for future studies of cyclic loading. It would be interesting to test whether repeated sacrificial-bond breakage and load transfer to the matrix network can be captured within this model over many cycles. This has not yet been theoretically explored due to the long timescales and system sizes. Our model, however, remains sufficiently simple to access the large number of cycles and parameter values required, while still explicitly resolving the essential physics. Additionally, recent experimental work done by You \etal \cite{you_damage_2025} has been able to establish a simple connection between monotonic and cyclic loading behaviours through damage modelling, which could be directly explored within this framework.

In summary, this simplified approach has enabled us to isolate \DIFdelbegin \DIFdel{the key physical mechanisms }\DIFdelend \DIFaddbegin \DIFadd{a physical mechanism }\DIFaddend governing fracture in double networks using minimal ingredients. We hope that this framework will serve as a foundation for future studies addressing more complex and realistic systems \DIFaddbegin \DIFadd{modelling the full fracture mechanism of double network materials}\DIFaddend .

\newpage
\begin{subappendices}
\section{Square Lattice-Based Network} \label{sec:SquareAppendix}
The simplest form of a single network structure is a lattice-based network. We extend this to a double network by constructing commensurate lattices with different length scales, connected where lattice sites overlap. By using this alternative method of network construction, we provide a test of the overall robustness of the observed physics.

\subsection{Lattice-Based Double Networks}
 For our lattice-based double network, the sacrificial network is constructed using the same method described in \cref{sec:LatticeBasedNetworks}, by defining a set of lattice vectors and rules for connecting them. The matrix network can also be generated in the same way, but with a scaled set of lattice vectors. 

 \begin{figure}[t!]
    \centering
    \includegraphics[width=0.99\linewidth]{chapters/DoubleNetworkFracture/NetworkGeneration/SquareTypes-4.pdf}
    \caption[Figure detailing the types of network constructions for a square lattice-based double network with $M=4$.]{Schematic illustrating the different network construction types for a square lattice–based double network with $M=4$. The constructions differ in how nodes are permitted to interact with both the sacrificial and matrix networks. In the overlapping network, matrix and sacrificial bonds are allowed to co-exist in the same spatial region. In contrast, networks with subdivisions split sacrificial bonds into smaller segments, enabling connections to matrix nodes that they pass over.\newline
    We illustrate the three forms of square lattice-based networks considered: \newline
    {\bf Top panels:} single matrix networks. \newline
    {\bf Middle panels:} single sacrificial networks. \newline
    {\bf Bottom panels:} combined double networks.
    }
    \label{fig:DoubleSquareTypes}
\end{figure}

For simplicity, we choose $M\in\mathbb{Z}$, such that the sacrificial and matrix lattices are commensurate: the matrix lattice vectors are integer multiples of the sacrificial ones. This ensures that the nodes of the sacrificial lattice lie exactly on a subset of the nodes of the fine matrix lattice. The natural length of each bond $l_0$ is then set in this configuration. 

During the network generation, we optionally decide whether the double network's matrix and sacrificial bonds are permitted to lie directly on top of one another in the network’s natural configuration. Additionally, sacrificial bonds can be subdivided such that they connect at every lattice point of the matrix bonds they traverse.

These are two binary choices, allowing overlap and allowing subdivision, which would, in principle, yield four distinct lattice-based network geometries. However, the case in which overlaps are disallowed and sacrificial bonds are not subdivided results in a complete decoupling of the matrix and sacrificial networks, and therefore does not constitute a double network. The three remaining cases are illustrated in \cref{fig:DoubleSquareTypes} for a square lattice with $M=4$.

Square lattice-based networks sit on the isostatic point $z_c$ with connectivity $z=z_c=4.0$ and are therefore marginally rigid, remaining rigid only in their initial configuration \cite{Head_marginally_rigid_2022}; any bond failures which occur during deformation will consequently locally drop the network's connectivity below $z_c$, causing the network to become floppy. As discussed in \cref{sec:Rigidity}, the network may strain-stiffen during its deformation and behave as if it has a $z > z_c$ even after bond failures. To ensure that the network remains in the strain-stiffened regime throughout its entire deformation, we pre-stress the network by reducing each bond's natural length $l_0$ such that its initial local bond strain $\epsilon_{b}$ is sampled from a normal distribution $\mathcal{N}(\epsilon_p, \sigma_p^2)$ with mean $\epsilon_p$ and variance $\sigma_p^2$. This pre-stress acts as a stabilising field allowing the network to hold states of self-stress at lower strains \cite{tauber_sharing_2021}. By varying each bond's local strain, we address the problems described in \cref{sec:LatticeBasedNetworks}, which is usually solved by pruning and perturbing the natural positions of each network's nodes. 

\subsection{Stress-Strain Response}

We present the stress-strain response of three networks in \cref{fig:SquareStressStrain}, one for each type of construction described above. As in \cref{fig:JammedStressStrain}, we plot in a) the corresponding single sacrificial and matrix responses and in b) the response of the combined double network. We will use a consistent set of colours to denote the relevant constructions with \quotes{Overlap With Subdivisions}, \quotes{Overlap Without Subdivisions} and \quotes{No Overlap With Subdivisions} denoted in blue, green and red, respectively. In \cref{fig:SquareStressStrain}a, the sacrificial component of the overlap with subdivisions case is identical to the sacrificial component of networks with no overlaps but with subdivisions. Likewise, its matrix component is equivalent to the matrix component of cases with overlaps but without Subdivisions.

\begin{figure}[b!]
    \centering
    %\missingfigure{Stress-strain curves of a prestressed Square Lattice Network}
    \includegraphics[width=0.95\linewidth]{chapters/DoubleNetworkFracture/SquareFigures/StressStrain.pdf}
    \caption[Stress-strain curves of a prestressed square lattice-based network.]{Stress-strain curves of a prestressed square lattice-based network. 
    We plot the three network constructions discussed in \cref{fig:DoubleSquareTypes} with \quotes{Overlap With Subdivisions}, \quotes{Overlap Without Subdivisions} and \quotes{No Overlap With Subdivisions} in blue, green and red, respectively.\newline 
    {\bf a)} Single sacrificial networks (dashed lines) and single matrix networks (dotted lines). {\bf b)} Composite double network (solid lines), with the contributions of the component sacrificial network (dashed line) and matrix network (dotted line) shown separately. \newline
    Parameters: $M=4, \mu_m=0.4$.}
    \label{fig:SquareStressStrain}
\end{figure}

Compared to the packing-derived networks discussed in the main text, we will consider a slightly different set of units now with $\lambda_s=1.0$. For simplicity, we will only consider the case of $\lambda_m=3.0$ and set the two disorder parameters discussed above to be $\epsilon_p=0.75$ and $\sigma_p=0.1$ unless otherwise stated. In \cref{fig:SquareStressStrain}, the sacrificial (red) and matrix (blue) curves display the same behaviour as the single packing-derived networks in \cref{fig:JammedStressStrain}, with a non-linear stress-strain response followed by brittle failure. All curves presented show a non-zero strain at $\epsilon=0.0$ due to the initial pre-stress in the system $\epsilon_p$, with a small deviation due to the allowed initial relaxation. 


\begin{figure}[b!]
    \centering
    %\missingfigure{Square Strain Extract}
    \includegraphics[width=0.99\linewidth]{chapters/DoubleNetworkFracture/SquareFigures/SquareStrainExtract.pdf}
    \caption[Failure strains of square lattice-based double networks]{
    Failure strains $\epsilon^*$ of square lattice-based double networks. The three network constructions discussed in \cref{fig:DoubleSquareTypes} with \quotes{Overlap With Subdivisions}, \quotes{Overlap Without Subdivisions} and \quotes{No Overlap With Subdivisions} in blue, green and red, respectively. The failure strain $\epsilon^*$ of the double network (solid curves) is plotted as a function of a) the matrix content $M$ and b) the matrix bond stiffness $\mu_m$. For comparison, the failure strains of the corresponding single sacrificial network and single matrix network are indicated by dashed and dotted curves, respectively.\newline
    Panel a) uses a fixed matrix bond stiffness $\mu_m = 0.4$, whereas panel b) fixes the relative length scales at $M = 4$.
    }
    \label{fig:SquareStrainExtract}
\end{figure}

We plot curves analogous to those presented in \cref{fig:JammedStrainExtract}a-b in \cref{fig:SquareStrainExtract}, extracting the failure strain $\epsilon^*$ as a function of a) $M$ and b) $\mu_m$. The behaviour of the square double network cases follows the overall trends explored in the packing-derived systems, diverging from that of the sacrificial cases at low $M$ or $\mu_m$ to failure strains closer to that of the matrix component with increasing $M$ or $\mu_m$. For the chosen parameters, increasing $M$ at $\mu_m=0.4$ does not lead to failure strains that reach that of the matrix. Instead, $\epsilon^*$ plateaus at a lower value that depends on the network construction.

Overall, we see that the blue network performs the worst; we expect, based on our exploration in the packing-derived case, that this is due to the lack of variation in length scales between the two components, forcing sacrificial bonds to rearrange alongside the matrix bonds, resulting in early failures. The responses of the cases with overlap and without subdivisions and cases with no overlaps but with subdivisions are nearly indistinguishable, both plateauing at $M \approx 4$. Both alternative constructions allow their sacrificial components to deform almost independently of their matrix components. This is achieved either by constraining only each bond’s endpoints, as in the packing-derived networks, or by weakly constraining each node, where all nodes remain at a connectivity of $z=4$. Importantly, they both appear to allow for the necessary separation between the dynamics of the two components. 

\begin{figure}[bt]
    \centering
    %\missingfigure{Square Strain Extract}
    \includegraphics[width=0.99\linewidth]{chapters/DoubleNetworkFracture/SquareFigures/HeatMaps.pdf}
    \caption[Phase diagrams of double-network failure strain $\epsilon^*_d$ and relative sacrificial contribution $G_s/G_d$ for square lattice-based networks]{Phase diagrams representing the square lattice-based networks. For each quantity, we plot the three network constructions discussed in \cref{fig:DoubleSquareTypes} separately. Colour scales show the change of the: \newline
    {\bf Top panels:} failure strain $\epsilon^*_d$ in the $(M,\mu_m)$ plane.\newline
    {\bf Bottom panels:} relative sacrificial contribution $G_s/G_d$ in the $(M,\mu_m)$ plane. \newline
    Parameters: $\lambda_s=1.0,\; \lambda_m=3.0,\; \mu_s=1.0, \; \epsilon_p=0.75$ 
    }
    \label{fig:SquareHeatMaps}
\end{figure}

Panel b) further illustrates the transition from sacrificial- to matrix-dominated behaviour. The double-network failure strain $\epsilon^*$ for all constructions begins to deviate from the corresponding single sacrificial network at approximately $\mu_m \gtrsim 0.15$. At this point, the matrix network becomes sufficiently stiff to suppress sacrificial bond failure and thereby enhance the overall toughness of the system. In most cases, the curves reach a plateau before converging to the pure matrix response. We hypothesise that this premature plateau arises from a progressive loss of local rigidity as bonds break and the initial pre-stress relaxes, permitting increased deformation and earlier failure. Although this mechanism has not yet been directly tested, it may account for the higher failure strains $\epsilon^*$ observed in networks with overlaps and without subdivisions (Fig.~\cref{fig:SquareHeatMaps}a–c). These networks more closely resemble the packing-derived structures and therefore appear better able to preserve their local structure at larger strains.

The role of stiffness is briefly explored in \cref{fig:SquareHeatMaps}d–f, which exhibit behaviour very similar to that observed for the packing-derived networks shown in \cref{fig:JammedPhaseDiagrams}b. We observe that the rate of reduction in sacrificial contribution to $G_d$ with increasing $M$ and $\mu_m$ is consistent with that of the packing-derived networks. However, as explored earlier, $\epsilon^*$ plateaus at larger strains, implying a tighter window of parameters over which the double network satisfies both properties I) and II): namely, retaining the large failure strain of the matrix ($\epsilon^*_D \approx \epsilon^*_m \gg \epsilon^*_s$) while simultaneously preserving the high modulus of the sacrificial network ($G_d \approx G_s \gg G_m$).

\subsection{Effective Propagator}

To explore the same physics that was discussed in \cref{sec:doubleStressProp}, but at a significantly reduced computational cost, we introduce a simplified but related approach. For each set of parameters, we evolve a single network using the method described in \cref{sec:doubleStressProp} and record the number of sacrificial bonds whose strain $\epsilon_{sb}$ increases by more than a threshold of $\delta\epsilon_{sb} \geq 10^{-2}$. As discussed in \cref{sec:doubleStressProp}, this is not a direct correlation with the number of bonds which will break in a disordered system. However, it does provide a useful insight into the magnitude of a single plastic event.

\begin{figure}[bt]
    \centering
    \includegraphics[width=0.9\linewidth]{chapters/DoubleNetworkFracture/SquareFigures/AboveThreshold.pdf}
    \caption[Phase space diagrams showing the change in bond strains under changing $M$ and $\mu_m$]{Phase space diagrams showing the number of bonds whose strain has increased by at least $10^{-2}$ as a result of single bond failure. For networks with Subdivisions in the sacrificial bonds, we normalise by $M$. Each data point is taken from a single network realisation with:
    \newline
    {\bf Top panels:} network disorder parameter $\sigma_p=0.1$. \newline
    {\bf Bottom panels:} network disorder parameter $\sigma_p=0.0$. \newline
    Parameters: $\lambda_s=1.0,\; \lambda_m=3.0,\; \mu_s=1.0, \; \epsilon_p=0.75$ 
    }
    \label{fig:SquarePropogators}
\end{figure}

Along the top row of \cref{fig:SquarePropogators}, we show results for a single network realisation with disorder parameter $\sigma_p = 0.1$. These exhibit a markedly less structured phase space when compared to the bottom row with $\sigma_p = 0.0$, where all vertical bonds would fail simultaneously, but only one is selected to fail. Comparing the two rows, we observe an order of magnitude reduction in the number of bonds exceeding the threshold, indicating a decrease in the stress propagated due to disorder. In both cases, the construction without subdivisions shows the least stress propagation, with only a minimal number of bonds crossing the threshold, even at low $M$ and $\mu_m$. 

We will not directly compare network construction here due to subdivisions dividing each sacrificial bond into multiple segments, and we have not yet fully considered or quantified the effect of losing a segment of a sacrificial bond compared to the whole bond; therefore we will not comment on results in \cref{fig:SquarePropogators} and how they map to the physics in \cref{fig:SquareHeatMaps} when comparing constructions. 

The global trend in this data remains consistent; increasing $M$ or $\mu_m$ systematically reduces the magnitude of the stress propagator generated by a single bond failure, thereby reducing the redistribution of stress through the network. However, the precise manner in which this reduction occurs, as well as its quantitative magnitude, depends sensitively on the underlying network topology. While these trends are clearly resolvable in ordered lattices, the introduction of disorder significantly complicates the interpretation of individual failure events. In disordered, packing-derived networks, the response to a single bond break becomes highly heterogeneous, making it difficult to extract general physical insights from isolated realisations. As a result, understanding stress propagation in these systems requires analysis based on ensemble-averaged behaviour, as performed for the packing-derived networks, rather than on the dynamics of a single event.

In summary, the square lattice-based double network reproduces the key qualitative features which we observe in the packing-derived networks. For the parameters presented here, they show consistent behaviours for their stress-strain response, failure-strain trends under changing both $M$ and $\mu_m$ and how stress propagates. They continue to illustrate the important role of controlling stress propagation between the two networks in the increased toughness of double networks. Network constructions that constrain sacrificial bonds too strongly to the matrix component suppress this separation and fail prematurely. This causes them to diverge from the constructions that allow more independent deformation. A deeper exploration into the role of $\epsilon_p$, along with other lattice constructions (such as a triangular lattice), may help deepen the understanding of the double network more generally. The agreement already demonstrated between lattice-based and packing-derived networks suggests that the underlying physics is robust to the details of network generation. This reinforces the generality of the conclusions drawn in the main text.

\section{Rate Dependency} \label{sec:DoubleRate}

In the main text, we carried out simulations in the quasi-static limit $\dot{\epsilon} \to 0$ using the algorithm described in \cref{sec:NetworkQuasistatic}. Here, we instead consider deformation at a finite stretching rate $\dot{\epsilon} > 0$. The results shown in \cref{fig:DoubleRate} demonstrate that as $\dot{\epsilon}$ decreases, the stress--strain curves progressively converge towards the quasi-static limit, indicating that introducing a finite stretching rate does not qualitatively alter our central result: a substantial enhancement of toughness in materials with a double-network structure.

\begin{figure}[bt]
    \centering
    \includegraphics[width=0.99\linewidth]{chapters/DoubleNetworkFracture/Result Figures/DoubleRateDependency.pdf}
    \caption[Stress-strain curves for a double network stretched at a finite rate.]{Stress-strain curves for a double network stretched at a finite rate $\dot{\epsilon} > 0$, showing the approach to the quasi-static limit $\dot{\epsilon} \to 0$ explored in the main text. Parameters used are the same as in \cref{fig:JammedStressStrain}}
    \label{fig:DoubleRate}
\end{figure}

\section{Dependence on matrix failure threshold \texorpdfstring{$\lambda_m$}{lambda m}}  \label{sec:JammedlambdaDependecy}

In \cref{fig:JammedStrainExtract,fig:JammedBrokenExtract} we briefly explore the effect of the matrix failure threshold $\lambda_m$ by plotting characteristic quantities of the failure behaviour at distinct values $\lambda_m = 1.0,\; 2.0,\; 3.0$. We extend this analysis by plotting intermediate values of $\lambda_m$ for a double network with parameters $M = 3$, $\mu_m = 0.2$, and $\lambda_s = 0.5$, for the characteristic quantities $\epsilon^*$ and $B^*$.

\begin{figure}[bt]
    \centering
    \includegraphics[width=0.99\linewidth]{chapters/DoubleNetworkFracture/Result Figures/lambda.pdf}
    \caption[Failure strain $\epsilon^*$ dependence on $\lambda_m$]{Failure behaviour of a packing-derived double network with $M=3$, $\mu_m=0.2$, and $\lambda_s=0.5$ as a function of $\lambda_m$. \newline
    {\bf a)} Failure strain for the double network (black), single matrix network (blue), and single sacrificial network (red). \newline
    {\bf b)} Fraction of bonds broken at failure for the double network (black), single matrix network (blue), and single sacrificial network (red). \newline
    {\bf c)} Fraction of bonds broken at failure in the matrix (blue) and sacrificial (red) components of the double network.}
    \label{fig:DoubleLambda}
\end{figure}

In \cref{fig:DoubleLambda}a, we plot the failure strain $\epsilon^*$. As expected, $\epsilon^*$ increases with $\lambda_m$ for both single matrix and double networks. For the single matrix networks, $\epsilon^*$ grows linearly with $\lambda_m$. However, the variance across samples also increases. We attribute this to the growing magnitude of non-affine rearrangements, which causes matrix bonds to fail earlier or later than a linear fit would predict. The double networks show a similar, roughly linear dependence on $\lambda_m$ for $\lambda_m \gtrsim 1.5$, where we do not observe any early matrix failures. 

In \cref{fig:DoubleLambda}b-c, we show the analogue of \cref{fig:JammedBrokenExtract}. Specifically, \cref{fig:DoubleLambda}b indicates the number of failed bonds at $\epsilon^*$ in the single sacrificial (red) and single matrix (blue) networks. In both cases, we observe little or no dependence on $\lambda_m$. This is trivial for the sacrificial network. For the matrix network, it reflects the abrupt brittle failures discussed in the main text.

The double network material exhibits a non-monotonic response in the number of failures $B^*$ for $\lambda_m < 1$, which can be understood as follows. For $\lambda_m \approx \lambda_s=0.5$, bond breakages occur uniformly across the components. Each network shows a slightly ductile response, similar to the sacrificial case in \cref{fig:JammedBrokenStrain}a. Consequently, $\epsilon_d^* \approx \epsilon_s^*$. However, the stress redistributed by broken sacrificial bonds accumulates locally in the matrix. Once sufficient local damage builds up, this stress triggers premature matrix failure. The behaviour for larger $\lambda_m$ then follows the trends discussed in the main text.

\end{subappendices}
 \newpage 
 \newpage % !TeX root = ../../thesis.tex

\setcounter{equation}{0}
\chapter{Conclusion}
In this thesis, we have used numerical simulations to study how soft disordered materials deform, yield, and fracture when subjected to different externally applied deformation protocols. Although the systems, methods, and deformation protocols considered differ between our three results chapters, the question has been the same: how does the interplay between structural disorder and the redistribution of stress following local failure events control the ultimate response of the material? In \cref{sec:Fatigue} this took the form of the collective organisation of mesoscopic elastoplastic elements, which evolved towards either a shear band or an absorbing state. In \cref{sec:Ultradelayed} it was reflected in the collective organisation of the network structure into an apparently quasi-stable state that ultimately failed much later through the breakage of a critical bond, and the knock-on consequences of that failure. \Cref{sec:DoubleNetworks} showed how two networks can interact to share load, resulting in a composite material that is both stiff and can be stretched to large strains without failure. This final chapter aims to draw together the main ideas developed across the thesis and to summarise the primary conclusions of each chapter in that broader context.

\section*{Absorption and Failure in Oscillatory Shear}
% In \cref{sec:Fatigue}, we examined a mesoscopic elastoplastic model under large-amplitude oscillatory shear and characterised the evolution of the system towards either an absorbing state, in which plastic activity eventually ceases, or a yielded state, in which deformation localises into a system-spanning shear band. We quantified the number of cycles, $C^*$, required to reach either outcome. In doing so, we identified a crossover region in strain amplitude, separating the absorbing and yielding regimes, in which both outcomes remain accessible, and the number of cycles required to reach them becomes large. Within this region, $C^*$ increases with increasing amplitude for absorption, but with decreasing amplitude for yielding. 

% The transition region itself was only very recently characterised by Khandare and Sastry \cite{khandareElastoplasticModellingCyclic2026}, and our results are in close agreement with this picture. At finite system sizes, as the imposed amplitude increases, the probability of reaching an absorbing state decreases continuously across a finite interval rather than at a single sharply defined threshold. This behaviour is broadly consistent with earlier particle-based studies, which often described cyclic yielding in terms of a critical amplitude separating reversible and irreversible dynamics \cite{Fiocco2013,Regev_onset_2013,kawasakiMacroscopicYieldingJammed2016,Leishangthem2017}.

% In this work, we further characterise the survival curves across both the absorbing and yielding regimes, allowing the transition between them to be described in full. An important observation that emerges from this analysis is that the longest-lived trajectories occur near the boundaries of the transition region. In these regimes, the less probable outcome, whether absorption or yielding, requires the largest number of cycles to occur, leading to markedly longer lifetimes than trajectories deep within either regime.

% Perhaps the most important mechanistic result of this chapter is the characterisation of the dynamics in the limit of increasingly dilute activity. For highly annealed samples, the dynamics towards an absorbing state can be interpreted as a weakly interacting first-passage process, with elements acting independently. This picture is supported by two observations. First, the characteristic cycle number scales approximately with the inverse square of the characteristic post-hop width, $C^* \sim l_w^{-2}$, which is consistent with a diffusive process. Second, the large-cycle decay rates for absorbing trajectories are the same for highly annealed samples in both the full and simplified models, suggesting that the simplified stochastic framework captures the essential physics of the late-cycle relaxation towards absorption in this limit.

% In \cref{sec:Fatigue}, we examined a mesoscopic elastoplastic model under large-amplitude oscillatory shear and characterised the evolution of the system towards either an absorbing state, in which plastic activity eventually ceases, or a yielded state, in which deformation localises into a system-spanning shear band. By measuring both the probability of reaching each outcome and the associated number of cycles $C^*$ required to do so, we showed that the response is nota single sharply defined threshold, but by a finite transition region separating absorption from yielding. Within this interval, both outcomes remain possible, and the probability of absorption decreases continuously with increasing strain amplitude, consistent with the recent picture proposed by Khandare and Sastry \cite{khandareElastoplasticModellingCyclic2026}. In addition to identifying a finite transition region, we showed that its dynamics are strongly asymmetric: near either boundary, the less probable outcome, whether absorption or yielding, is also the slowest to occur.

In \cref{sec:Fatigue} we examined a mesoscopic elastoplastic model under large-amplitude oscillatory shear and characterised the evolution of the system towards either an absorbing state, in which plastic activity eventually ceases, or a yielded state, in which deformation localises into a system-spanning shear band. By measuring both the probability of reaching each outcome and the associated number of cycles $C^*$ required to do so, we demonstrated that the response is not controlled by a single sharply defined threshold, but by a transition region separating absorption from yielding. In this regime, both outcomes remain possible, and the probability of absorption decreases continuously with increasing strain amplitude, consistent with the recent picture proposed by Khandare and Sastry \cite{khandareElastoplasticModellingCyclic2026}. In addition to identifying a finite transition region, we showed that the evolution of the system is strongly asymmetric: near either boundary of the transition region, the less probable outcome, whether absorption or yielding, is also the slowest to occur.

%The main mechanistic insight of the chapter is that, for highly annealed samples, the late-time progression towards absorption can be understood in the limit of increasingly dilute activity. In this regime, the remaining active elements behave approximately as weakly interacting stochastic processes, so that delayed absorption can be interpreted as a first-passage problem. This provides a simple physical picture for delayed absorption in terms of the stochastic relaxation of a small number of weakly interacting active elements

The main mechanistic insight of the chapter is that, for highly annealed samples, the late-time progression towards absorption can be understood in the limit of increasingly dilute activity. In this regime, the remaining active elements behave approximately as a non-interacting stochastic processes, so that delayed absorption can be interpreted as a first-passage problem. This provides a simple physical picture for delayed absorption in terms of the stochastic relaxation of a small number of weakly interacting active elements. More broadly, it shows that the long-lived dynamics of systems that eventually reach an absorbing state do not necessarily remain many-body all the way to the final state. Instead, after an initially collective regime dominated by long-ranged Eshelby interactions between active elements, the system crosses over to a much simpler stochastic regime. This is important for the field more generally because it identifies a possible route by which complex mesoscopic yielding models can be reduced, at late times, to effective stochastic descriptions with a clear physical picture.

%As discussed in \cref{sec:FatigueConclusion}, an important direction for future work would be to investigate whether the ideas behind the dilute-limit picture can be extended to describe not only absorption, but also yielding. While yielding is inherently a many-body process, it is possible that the transition to yielding is controlled by the dynamics of a small number of spatially localised elements. By extending the framework to include the possibility that a plastic event can destabilise a neighbouring element lying close to the absorbing boundary, it may be possible to capture the reactivation of the system and the cascade of events that culminates in shear-band formation. This would provide a route towards a unified stochastic framework for both delayed absorption and delayed yielding, and would help clarify the mesoscopic origin of the transition region between these two outcomes.

At the same time, the results also make clear what a study such as this does not yet capture. The model is deliberately minimal: it is athermal, quasi-static, scalar, and does not include explicit finite-rate effects, tensorial stress couplings between different deformation components, wall slip, or free surfaces. These simplifications were necessary to isolate the competition between absorption and yielding, but they also limit the extent to which the results can be mapped quantitatively onto experiments \cite{pineChaosThresholdIrreversibility2005,Cort2008,tapadiaBandingEntangledPolymer2006,Koumakis2011,Petekidis2004,knowltonMicroscopicViewYielding2014}. In particular, the dilute-limit framework developed here only captures delayed absorption. Yielding remains inherently more collective because it depends on the interaction of many elements and the eventual emergence of a localised shear band. As discussed in \cref{sec:FatigueConclusion}, an important direction for future work would therefore be to investigate whether the ideas behind the dilute-limit picture can be extended to describe not only absorption, but also yielding. By extending the framework to include the possibility that a plastic event can destabilise a neighbouring element lying close to the absorbing boundary, it may be possible to capture the reactivation of the system and the cascade of events that culminates in shear-band formation. This would provide a route towards a unified stochastic framework for both delayed absorption and delayed yielding, and would help clarify the mesoscopic origin of the transition region between these two outcomes.

These findings also have several implications for experimental work and for real amorphous materials. Experimentally, they suggest that the response under cyclic deformation need not always be interpreted purely in terms of a single yielding amplitude, but rather in terms of a finite regime in which absorption and yielding are both possible. In particular, the results suggest that bulk rheology alone may not be sufficient to distinguish trajectories that are slowly approaching absorption from those that are slowly evolving towards failure; spatially resolved measurements of local rearrangements or strain localisation would be especially valuable in testing this picture \cite{divoux_stress_2011,divouxTransientShearBanding2010,divoux_yielding_2012,grenard_timescales_2014,gibaud_heterogeneous_2010}. More broadly, for materials used in applications, the results highlight that fatigue resistance under repeated loading is controlled not only by the amount of plastic activity generated in each cycle, but by whether that activity tends to attenuate with increasing cycle number or instead to organise into a localised failure event. This is relevant to a wide range of soft and amorphous materials, from colloidal gels and emulsions to polymeric and metallic glasses, in which repeated deformation, vibration, or cyclic loading can gradually reorganise the material long before macroscopic failure becomes apparent \cite{bonn_yield_2017,SCHUH2007,hufnagel_deformation_2016,greer_shear_2013}.

\section*{Ultradelayed Fracture After Step Strain in Networks}
% In \cref{sec:Ultradelayed}, we shifted focus from generalised constitutive descriptions of viscoelastic materials to the fracture of soft disordered networks. Motivated by the work of Lockwood \etal~\cite{lockwood_ultradelayed_2025}, we imposed a step strain on a disordered network model with thermally activated bond failure and characterised the resulting delayed fracture dynamics. The central result of this chapter is that the ultra-delayed post-strain instabilities identified in Ref.~\cite{lockwood_ultradelayed_2025} have a clear fracture analogue in direct simulations of disordered networks. The material can remain apparently stable for a very long time after the imposed deformation before ultimately failing catastrophically, consistent with delayed-fracture and delayed-yielding phenomenology reported in soft gels, elastomers, and polymer networks \cite{sprakel_stress_2011,blindstrom_structures_2012,van_der_kooij_laser_2018,wangDelayedFractureGels2012,Ju2023}.

% At low imposed strains, networks that fail only after a long delay initially relax almost indistinguishably from their athermal counterparts before entering a long-lived quasi-stationary state in which further deformation is minimal, and bond breaking is rare. In this regime, failure is controlled by the stochastic rupture of a critical bond whose local environment allows it to trigger a cascade of correlated breakages, producing a system-spanning crack \cite{mattsson_time-dependent_2017,cipelletti_microscopic_2020,aime_microscopic_2018}.

% By contrast, at higher strains, the behaviour changes qualitatively. The network accumulates substantial diffuse damage during, and shortly after, relaxation, yet can still survive for many decades before macroscopic fracture occurs. This shows that large amounts of damage are not, by themselves, sufficient to cause immediate collapse. Instead, the system can retain its load-bearing capacity while broken bonds accumulate in the background, as both strain stiffening and elastic stress redistribution continue to stabilise the remaining structure \cite{tauber_role_2020,sharma_strain-controlled_2016}.

% A further interesting result is that the direction of brittle fracture depends strongly on the imposed strain regime. Low-strain failures are approximately isotropic in the network frame, whereas high-strain failures become biased by the anisotropic bond loading generated by the imposed deformation. This suggests that the route to catastrophic failure is controlled not only by when damage accumulates, but also by how the imposed strain and subsequent stress redistribution determine the geometry of the final instability.

In \cref{sec:Ultradelayed} we shifted focus from generalised constitutive descriptions of viscoelastic materials to the fracture of soft disordered networks. Motivated by the work of Lockwood \etal~\cite{lockwood_ultradelayed_2025}, we studied the case of step strain imposed on a disordered network model with thermally activated bond failure. We characterised in particular the resulting delayed fracture dynamics. The central result of this chapter is that the ultra-delayed post-strain instabilities identified in Ref.~\cite{lockwood_ultradelayed_2025} have a clear analogue in direct simulations of disordered networks manifesting here as delayed fracture. The material can remain apparently stable for a very long time after the imposed deformation before ultimately failing catastrophically, consistent with delayed-fracture and delayed-yielding phenomenology reported in soft gels, elastomers, and polymer networks \cite{sprakel_stress_2011,blindstrom_structures_2012,van_der_kooij_laser_2018,wangDelayedFractureGels2012,Ju2023}. At the same time, the ultra-delayed regime occupies only a narrow region of parameter space, in which imposed strain, thermally activated bond rupture, and the residual load-bearing stability of the network are finely balanced. Outside this window, the system either remains effectively stable after relaxation or fails on much shorter timescales.

The mechanism depends on the imposed strain regime. At low strains, delayed-failure trajectories relax almost indistinguishably from their athermal counterparts before entering a long-lived quasi-stable state in which deformation and bond breaking are both minimal. Catastrophic failure later occurs when the stochastic rupture of a critical bond triggers a cascade of correlated breakages \cite{mattsson_time-dependent_2017,cipelletti_microscopic_2020,aime_microscopic_2018}. At higher strains, by contrast, the network accumulates substantial diffuse damage during and shortly after relaxation. However, it can still survive for many decades before fracture, showing that widespread bond breaking is not, by itself, sufficient to cause immediate collapse, because the remaining structure continues to carry load. Consistent with this change in mechanism, low-strain fractures are approximately isotropic in the network frame, whereas high-strain fractures become biased by the anisotropic bond loading induced by the imposed deformation.

These results also have implications for how soft network materials are tested, processed, and used in practice. Most directly, they show that the absence of immediate failure after loading should not be taken as evidence of long-term stability. A network may relax to an apparently quiescent state and yet remain mechanically metastable for a very long time before ultimately failing catastrophically. This suggests that experimental protocols based only on the short-time response after loading may systematically underestimate the risk of failure, because samples that appear stable on accessible observation times might potentially still fracture at a much later time. The results, therefore, point to the importance of continued measurements, particularly in the narrow metastable window just below the failure condition where ultra-delayed fracture is possible.

This may also be relevant for applications involving hydrogels and elastomers used in tissue engineering, stretchable electronics, and soft robotics \cite{leeHydrogelsTissueEngineering2001,rogers_materials_2010,martinezSoftActuatorsRobots2014}. In such settings, deformation applied during fabrication, handling, assembly, or storage may place a material into a metastable damaged state whose consequences only emerge much later. Two samples that appear very similar immediately after loading may therefore have very different long-time mechanical integrity. Understanding how such delayed instabilities emerge, and how their timescales depend on the imposed strain and mesoscopic failure processes, may therefore help identify regimes in which late-time failure is most likely and ultimately inform strategies for preventing such failures.

Several directions for future work follow from this work. It would first be valuable to test whether the same ultra-delayed fracture phenomenology persists in more realistic thermal network models, such as those in which temperature enters through Brownian or Langevin fluctuations of the nodes rather than through an activated bond-breaking rule \cite{tauber_role_2020}. More generally, a more detailed understanding of how the local structure of the network selects critical bonds and failure paths would help clarify when and where catastrophic failure nucleates. Related questions of local susceptibility have been explored in other amorphous materials \cite{mattsson_time-dependent_2017,cipelletti_microscopic_2020,aime_microscopic_2018}. This would provide deeper insight into the mechanisms that control the response both in this chapter and in \cref{sec:DoubleNetworks}.

\section*{Toughness of Double Network Hydrogels}
In \cref{sec:DoubleNetworks}, we considered how fracture can be suppressed in composite double-network materials. The primary result of this chapter is that a minimal model of a stiff, brittle sacrificial network embedded in a softer and more extensible matrix reproduces the central qualitative hallmark of tough double networks \cite{gong_double-network_2003,gong_why_2010,ducrot_toughening_2014,millereauMechanicsElastomericMolecular2018,kingMacroscaleDoubleNetworks2019,takahashiCreatingStiffTough2018}: the combined material can retain the high initial stiffness of the sacrificial component while failing at strains set by the much more ductile matrix.

The mechanism underlying this increased toughness is identified in this chapter as stress delocalisation. By increasing either the density of the matrix relative to the sacrificial network, or the stiffness of the matrix bonds, the load released by the failure of a sacrificial bond is distributed over a larger and more compliant background. This reduces the magnitude of the stress perturbation experienced by nearby sacrificial bonds and therefore suppresses the cascade of local failures that would otherwise produce brittle fracture. This is consistent with experimental suggestions from neutron- and light-scattering measurements, together with birefringence imaging \cite{ducrot_structure_2015,tominaga_molecular_2007,huang_importance_2007,zheng_how_2021,Orr2026}. Our results show that this behaviour can emerge from a simple and physically interpretable mechanism: the matrix does not merely carry load after the sacrificial network has been damaged, but suppresses brittle failure by delocalising the stress released by local bond rupture. That principle is relevant not only for double-network hydrogels used in load-bearing and biomedical contexts, but also for tougher elastomers, soft adhesives, and other multi-network soft materials in which resistance to catastrophic fracture is required \cite{gong_why_2010,millereauMechanicsElastomericMolecular2018,kingMacroscaleDoubleNetworks2019,takahashiCreatingStiffTough2018}.

This change in stress propagation leads directly to the diffuse damage that is characteristic of tough double networks \cite{millereauMechanicsElastomericMolecular2018,fukao_effect_2020,ju_role_nodate}. Rather than failing immediately after the first sacrificial bond breaks, the composite can tolerate extensive distributed damage within the sacrificial network while preserving much of its stiffness. Macroscopic failure occurs only later, when the matrix itself can no longer sustain the imposed deformation. The chapter thus provides a simple mesoscopic explanation for how double-network architectures combine stiffness and toughness: not by preventing damage, but by spreading its mechanical consequences broadly enough that catastrophic localisation is delayed.

At the same time, the model necessarily omits a number of features that will matter in real materials. It does not include pre-stress in the sacrificial network, heterogeneous network topology, viscoelastic or poroelastic effects, or bond reformation after damage. Some of these omissions are secondary to the central mechanism identified here, but others are likely to be important for quantitative comparison with experiments and for understanding durability. In particular, reversible bond breaking and healing may substantially alter the way damage accumulates under repeated loading, while network heterogeneity may influence how easily failure localises in the first place. Extending the present framework to include such ingredients would therefore help determine how robust the stress-delocalisation mechanism remains once more realistic material physics is introduced \cite{gong_why_2010,shivers_strain-controlled_2023,wangRapidSelfstrengtheningDoublenetwork2025,yasuiRateIndependentSelfHealingDouble2022,yinSelfHealingHydrogelsSynthesis2023}.

A natural direction for future work would be to extend this minimal framework towards more realistic double-network hydrogels by incorporating additional ingredients such as pre-stress in the sacrificial network or controlling the connectivity of each network \cite{gong_why_2010,shivers_strain-controlled_2023}, as well as bond reformation \cite{wangRapidSelfstrengtheningDoublenetwork2025,yasuiRateIndependentSelfHealingDouble2022,yinSelfHealingHydrogelsSynthesis2023}. It would be particularly interesting to test whether the same stress-delocalisation mechanism continues to govern toughness under cyclic loading, where repeated sacrificial-bond failure, stress redistribution, and partial healing may together control long-term durability \cite{zhang_fatigue_2018,chen_improvement_2016,chen_novel_2015,you_damage_2025}.

\section{Overall Perspective}
Taken together, the chapters of this thesis support a single broader idea: the mechanical outcome of soft disordered materials is determined by how local irreversible events interact through a heterogeneous and evolving stress field. When stress redistribution promotes stress localisation, the material eventually fails through abrupt yielding, brittle fracture, or the nucleation of a system-spanning instability. When redistribution instead becomes dilute or delocalised, the material can redistribute the stress from yielding events, sustain long-lived metastable states, or accommodate extensive diffuse damage before failing. The key distinction is therefore not simply whether microscopic failures occur, but how the surrounding material responds to them.

In oscillatory shear, delayed absorption emerges when the remaining active regions become effectively isolated. Under an applied step-strain deformation, ultra-delayed fracture emerges when rare thermally activated failures must first occur at a mechanically critical state from which they can trigger collapse. In double networks, toughness is enhanced because the matrix suppresses the local stress perturbations generated by sacrificial bond failure, allowing damage to accumulate without immediately causing catastrophic localisation. Across all of these cases, the inherent disorder of the material sets the landscape of possible responses, but it is the redistribution of stress that ultimately determines the outcome: whether activity dies out, damage accumulates diffusely, or the system undergoes catastrophic failure.

Viewed from this perspective, a central open question for the field is whether the dynamics can be predicted from the local structure with sufficient accuracy to anticipate the eventual mode of failure. In other words, can one identify, from a disordered configuration and its loading history, which bonds, sites, or regions are genuinely mechanically critical, how their failure will redistribute stress, and whether the resulting plastic activity will attenuate, remain diffuse, or develop into a correlated avalanche that produces yielding or fracture? This question joins together several issues that are often treated separately: why one bond is critical, whereas many apparently similar bonds are not; how local susceptibility is encoded by structure and pre-existing damage; and how initially isolated events become correlated strongly enough in space and time to generate a shear band, a crack, or some other localised instability. 

A second major challenge is to determine how far the route to failure is selected by the bulk disordered structure and how far it is instead controlled by edges, defects, and free surfaces. Real soft materials are finite in size, processed, and handled, and are therefore rarely ideal bulk systems. Boundaries may concentrate stress, defects may act as preferred nucleation sites, and surface effects may compete with the intrinsic susceptibility of the bulk structures. Distinguishing between these possibilities is essential both for interpreting experiments and for material design: a strategy which may suppress a bulk instability may not be the same as that required to stabilise a boundary or defect.

\section{Concluding Remarks}
Finally, this thesis has brought together a number of different phenomena and shown how they can be interpreted within a single perspective. Although this picture is necessarily incomplete, as all research must finish somewhere, it provides a basis on which further work can build. We hope that this work, and the work it inspires, provide a small step towards understanding the origin of fracture and yielding in amorphous solids.
 \newpage 
                                                    %
% Bibliography                                      %
 \newpage % REFERENCES.
\begin{singlespace}
    %\bibliography{references/bib}
    \printbibliography[heading=bibintoc]
\end{singlespace} \newpage                             %
                                                    %
% % Appendix                                        %
%DIF <  \appendix                                         %
%DIF <  \include{appendices/a/appendix}                   %
\DIFaddbegin \appendix                                           %DIF > 
\DIFaddend % \include{appendices/b/appendix}                   %
                                                    %
\end{document}                                      %
                                                    %
% % % % % % % % % % % % % % % % % % % % % % % % % % %
